\input awebmac
% This program originally by D. E. Knuth is not copyrighted and can be used
% freely. Changes to ADA by U. Schweigert.
% ADA-Version 1.0 was released in May, 1988

% Here is TeX material that gets inserted after \input webmac
\def\hang{\hangindent 3em\indent\ignorespaces}
\font\ninerm=cmr9
\let\mc=\ninerm % medium caps for names like ADA
\def\ADA{{\mc ADA}}
\def\ab{$\.|\ldots\.|$} % ada brackets (|...|)
\def\v{\.{\char'174}} % vertical (|) in typewriter font
\def\dleft{[\![} \def\dright{]\!]} % double brackets
\mathchardef\RA="3221 % right arrow
\mathchardef\BA="3224 % double arrow
\def\({} % kludge for alphabetizing certain module names

\def\title{AWEAVE}
\def\contentspagenumber{13} % should be odd
\def\topofcontents{\null\vfill
  \titlefalse % include headline on the contents page
  \def\rheader{\mainfont Appendix D\hfil \contentspagenumber}
  \centerline{\titlefont The {\ttitlefont AWEAVE} processor}
  \vskip 15pt
  \centerline{(based on {\tt WEAVE} Version 2.8)}
  \centerline{\ninerm[ADA-Version 1.0]}
  \vfill}
\pageno=\contentspagenumber \advance\pageno by 1


\N1.  Introduction.
This program converts an \.{AWEB} file to a \TeX\ file. It was written
(initially in {\mc PASCAL})
by D. E. Knuth in October, 1981; a somewhat similar {\mc SAIL} program had
been developed in March, 1979, although the earlier program used a top-down
parsing method that is quite different from the present scheme.

The conversion from {\mc PASCAL} to {\mc ADA} has been done by
U. Schweigert in May, 1988.

The ``banner line'' defined here should be changed whenever \.{AWEAVE}
is modified.

\Y\P\4\D\\{banner}\?\S\?\.{"This\ is\ AWEAVE,\ Version\ 1.1\ (NRaD\ Version)"}%
\?\6
\?\par
\fi

\M2. The program consists of several packages that are declared right now;
each of these packages and both the specification and the body of
the packages are send to a seperate file. The main program itself
is declared later.

\fi

\M3\*. The first package we declare is called \\{int\_number\_io} and enables
us
to output integer numbers. Further we declare a package \\{data\_structures},
which will contain all the global constants, types, variables, and some
procedures used throughout the program; it will (nearly) contain code of
the sections ``The character set'' and ``Data structures.''

\Y\P\O{data\_structures.ads}\?\6
\?\&{with}\8\p{TEXT\_IO};\6
\&{with}\8\\{int\_number\_io};\6
\&{with}\8\p{TEXT\_IO};\6
\&{package}\8\1\\{data\_structures}\2\8\&{is}\8\1\6
\X14:Types and objects visible from \\{data\_structures}\X\6
\X48:Specification of procedures visible from \\{data\_structures}\X\6
\&{procedure}\8\1\\{initialize}\2;\2\6
\&{end}\8\\{data\_structures};\7
\O{data\_structures.adb}\?\6
\?\&{with}\8\\{input\_output};\8\&{use}\8\\{input\_output};\6
\&{package}\8\1\&{body}\8\\{data\_structures}\2\8\&{is}\8\1\6
\X49:Procedures in \\{data\_structures}\X\6
\&{procedure}\8\1\\{initialize}\2\8\&{is}\8\1\6
\4\&{begin}\6
\X19:Set initial values\X\2\6
\&{end}\8\\{initialize};\2\6
\&{end}\8\\{data\_structures};\par
\fi

\M4\*. The following package, called \\{input\_output}, is responsible for the
input/output features needed by \.{AWEAVE}, and for the error handling.
It will contain code of sections ``Input and output'' and ``Reporting
errors to the user.''

\Y\P\O{input\_output.ads}\?\6
\?\&{with}\8\p{TEXT\_IO};\6
\&{with}\8\\{data\_structures};\8\&{use}\8\\{data\_structures};\6
\&{package}\8\1\\{input\_output}\2\8\&{is}\8\1\6
\X27:Specification of procedures visible from \\{input\_output}\X\2\6
\&{end}\8\\{input\_output};\7
\O{input\_output.adb}\?\6
\?\&{with}\8\p{TEXT\_IO}\?,\39\?\\{int\_number\_io};\6
\&{package}\8\1\&{body}\8\\{input\_output}\2\8\&{is}\8\1\6
\X28:Procedures in \\{input\_output}\X\2\6
\&{end}\8\\{input\_output};\par
\fi

\M5\*. Additionally we have a package \\{searching}, which contains the code of
sections ``Searching for identifiers'' and ``Searching for module names.''

\Y\P\O{searching.ads}\?\6
\?\&{with}\8\\{data\_structures};\8\&{use}\8\\{data\_structures};\6
\&{package}\8\1\\{searching}\2\8\&{is}\8\1\6
\X60:Specification of procedures visible from \\{searching}\X\2\6
\&{end}\8\\{searching};\7
\O{searching.adb}\?\6
\?\&{with}\8\p{TEXT\_IO};\6
\&{with}\8\\{input\_output};\8\&{use}\8\\{input\_output};\6
\&{package}\8\1\&{body}\8\\{searching}\2\8\&{is}\8\1\6
\X61:Procedures in \\{searching}\X\2\6
\&{end}\8\\{searching};\par
\fi

\M6\*. The first of the major packages will manage the input process of
\.{AWEAVE} and contains code of sections ``Lexical Scanning''
and ``Inputting the next token.''

\Y\P\O{lexical\_scanning.ads}\?\6
\?\&{with}\8\\{data\_structures};\8\&{use}\8\\{data\_structures};\6
\&{package}\8\1\\{lexical\_scanning}\2\8\&{is}\8\1\6
\X92:Specification of procedures visible from \\{lexical\_scanning}\X\2\6
\&{end}\8\\{lexical\_scanning};\7
\O{lexical\_scanning.adb}\?\6
\?\&{with}\8\p{TEXT\_IO}\?,\39\?\\{int\_number\_io};\6
\&{with}\8\\{input\_output};\8\&{use}\8\\{input\_output};\6
\&{with}\8\\{searching};\8\&{use}\8\\{searching};\6
\&{package}\8\1\&{body}\8\\{lexical\_scanning}\2\8\&{is}\8\1\6
\X85:Procedures in \\{lexical\_scanning}\X\2\6
\&{end}\8\\{lexical\_scanning};\par
\fi

\M7\*. The next package is responsible for the first phase of \.{AWEAVE} and
contains code of section ``Phase one processing.''
\Y\P\O{first\_phase.ads}\?\6
\?\&{with}\8\\{data\_structures};\8\&{use}\8\\{data\_structures};\6
\&{package}\8\1\\{first\_phase}\2\8\&{is}\8\1\6
\X128:Specification of procedures visible from \\{first\_phase}\X\2\6
\&{end}\8\\{first\_phase};\7
\O{first\_phase.adb}\?\6
\?\&{with}\8\p{TEXT\_IO}\?,\39\?\\{int\_number\_io};\6
\&{with}\8\\{input\_output};\8\&{use}\8\\{input\_output};\6
\&{with}\8\\{searching};\8\&{use}\8\\{searching};\6
\&{with}\8\\{lexical\_scanning};\8\&{use}\8\\{lexical\_scanning};\6
\&{package}\8\1\&{body}\8\\{first\_phase}\2\8\&{is}\8\1\6
\X129:Procedures in \\{first\_phase}\X\2\6
\&{end}\8\\{first\_phase};\par
\fi

\M8\*. This package provides the output features needed in the second and
third phase of \.{AWEAVE} and contains code of sections ``Low-level
output routines'' and ``Routines that copy \TeX\ material.''
\Y\P\O{output.ads}\?\6
\?\&{with}\8\\{data\_structures};\8\&{use}\8\\{data\_structures};\6
\&{package}\8\1\\{output}\2\8\&{is}\8\1\6
\X140:Specification of procedures visible from \\{output}\X\2\6
\&{end}\8\\{output};\7
\O{output.adb}\?\6
\?\&{with}\8\p{TEXT\_IO}\?,\39\?\\{int\_number\_io};\6
\&{with}\8\\{input\_output};\8\&{use}\8\\{input\_output};\6
\&{with}\8\\{lexical\_scanning};\8\&{use}\8\\{lexical\_scanning};\6
\&{package}\8\1\&{body}\8\\{output}\2\8\&{is}\8\1\6
\X141:Procedures in \\{output}\X\2\6
\&{end}\8\\{output};\par
\fi

\M9\*. Here's a package that handles the parsing of \ADA-texts containing code
of sections ``Parsing,'' ``Implementing the productions,'' and
``Initializing the scraps.''
\Y\P\O{parsing.ads}\?\6
\?\&{with}\8\\{data\_structures};\8\&{use}\8\\{data\_structures};\6
\&{package}\8\1\\{parsing}\2\8\&{is}\8\1\6
\X217:Specification of procedures visible from \\{parsing}\X\2\6
\&{end}\8\\{parsing};\7
\O{parsing.adb}\?\6
\?\&{with}\8\p{TEXT\_IO}\?,\39\?\\{int\_number\_io};\6
\&{with}\8\\{input\_output};\8\&{use}\8\\{input\_output};\6
\&{with}\8\\{searching};\8\&{use}\8\\{searching};\6
\&{with}\8\\{lexical\_scanning};\8\&{use}\8\\{lexical\_scanning};\6
\&{with}\8\\{output};\8\&{use}\8\\{output};\6
\&{package}\8\1\&{body}\8\\{parsing}\2\8\&{is}\8\1\6
\X212:Procedures in \\{parsing}\X\2\6
\&{end}\8\\{parsing};\par
\fi

\M10\*. Now there's a package for handling phase two of \.{AWEAVE} containing
code of sections ``Output of tokens'' and ``Phase two processing.''
\Y\P\O{second\_phase.ads}\?\6
\?\&{with}\8\\{data\_structures};\8\&{use}\8\\{data\_structures};\6
\&{package}\8\1\\{second\_phase}\2\8\&{is}\8\1\6
\X244:Specification of procedures visible from \\{second\_phase}\X\2\6
\&{end}\8\\{second\_phase};\7
\O{second\_phase.adb}\?\6
\?\&{with}\8\p{TEXT\_IO}\?,\39\?\\{int\_number\_io};\6
\&{with}\8\\{input\_output};\8\&{use}\8\\{input\_output};\6
\&{with}\8\\{searching};\8\&{use}\8\\{searching};\6
\&{with}\8\\{lexical\_scanning};\8\&{use}\8\\{lexical\_scanning};\6
\&{with}\8\\{output};\8\&{use}\8\\{output};\6
\&{with}\8\\{parsing};\8\&{use}\8\\{parsing};\6
\&{package}\8\1\&{body}\8\\{second\_phase}\2\8\&{is}\8\1\6
\X245:Procedures in \\{second\_phase}\X\2\6
\&{end}\8\\{second\_phase};\par
\fi

\M11\*. Finally the last package handles the third phase of \.{AWEAVE} by
means of ``Phase three processing.''
\Y\P\O{third\_phase.ads}\?\6
\?\&{with}\8\\{data\_structures};\8\&{use}\8\\{data\_structures};\6
\&{package}\8\1\\{third\_phase}\2\8\&{is}\8\1\6
\X298:Specification of procedures visible from \\{third\_phase}\X\2\6
\&{end}\8\\{third\_phase};\7
\O{third\_phase.adb}\?\6
\?\&{with}\8\p{TEXT\_IO}\?,\39\?\\{int\_number\_io};\6
\&{with}\8\\{input\_output};\8\&{use}\8\\{input\_output};\6
\&{with}\8\\{output};\8\&{use}\8\\{output};\6
\&{with}\8\\{second\_phase};\8\&{use}\8\\{second\_phase};\6
\&{package}\8\1\&{body}\8\\{third\_phase}\2\8\&{is}\8\1\6
\X291:Procedures in \\{third\_phase}\X\2\6
\&{end}\8\\{third\_phase};\par
\fi

\M12. Some of this code is optional for use when
statistics about \.{AWEAVE}'s memory usage are desired;
such material is enclosed between the delimiters  \&{stat} and \&{tats}.
Other material is optional when a somewhat optimal code is wanted;
such is delimited by  \&{optim} and \&{mitpo}.

\Y\P\4\D\\{stat}\?\S\B\C{change this to `$\\{stat}\equiv\null$'   when
gathering usage statistics}\6
\?\par
\P\4\D\\{tats}\?\S\?\hbox{}\?\T\C{change this to `$\\{tats}\equiv\null$'   when
gathering usage statistics}\6
\?\par
\P\4\F\\{stat}\?\S\?\\{declare}\par
\P\4\F\\{tats}\?\S\?\\{end}\par
\P\4\D\\{optim}\?\S\B\C{change this to `$\\{optim}\equiv\null$' when optimal
 code is desired}\6
\?\par
\P\4\D\\{mitpo}\?\S\?\hbox{}\?\T\C{change this to `$\\{mitpo}\equiv\null$' when
optimal     code is desired}\6
\?\par
\P\4\F\\{optim}\?\S\?\\{declare}\par
\P\4\F\\{mitpo}\?\S\?\\{end}\par
\fi

\M13. Here are some macros for common programming idioms.

\Y\P\4\D\\{incr}\?\1(\?\#\?\2)\S\?\#\?\K\?\#\?+\?1\?\C{increase a variable by
unity}\6
\?\par
\P\4\D\\{decr}\?\1(\?\#\?\2)\S\?\#\?\K\?\#\?-\?1\?\C{decrease a variable by
unity}\6
\?\par
\fi

\M14. The following parameters are set big enough to handle \TeX, so they
should be sufficient for most applications of \.{AWEAVE}.

\Y\P\?\4\X14:Types and objects visible from \\{data\_structures}\X\S\?\6
\\{max\_bytes}\?:\?\&{constant}\?\8\K\?45\_000;\C{1\?/\?\\{ww} times the number
of bytes in identifiers,   index entries, and module names; must be less than
65\_536}\6
\\{max\_names}\?:\?\&{constant}\?\8\K\?5\_000;\C{number of identifiers, index
entries, and module names;   must be less than 10\_240}\6
\\{max\_modules}\?:\?\&{constant}\?\8\K\?2\_000;\C{greater than the total
number of modules}\6
\\{hash\_size}\?:\?\&{constant}\?\8\K\?353;\C{should be prime}\6
\\{buf\_size}\?:\?\&{constant}\?\8\K\?100;\C{maximum length of input line}\6
\\{longest\_name}\?:\?\&{constant}\?\8\K\?400;\C{module names shouldn't be
longer than this}\6
\\{long\_buf\_size}\?:\?\&{constant}\?\8\K\?500;\C{\\{buf\_size}\?+\?\\{longest%
\_name}}\6
\\{line\_length}\?:\?\&{constant}\?\8\K\?80;\C{lines of \TeX\ output have at
most this many characters,   should be less than 256}\6
\\{max\_refs}\?:\?\&{constant}\?\8\K\?60\_000;\C{number of cross references;
must be less than 65\_536}\6
\\{max\_toks}\?:\?\&{constant}\?\8\K\?60\_000;\C{number of symbols in \ADA\
texts being parsed;   must be less than 65\_536}\6
\\{max\_texts}\?:\?\&{constant}\?\8\K\?10\_000;\C{number of phrases in \ADA\
texts being parsed;   must be less than 10\_240}\6
\\{max\_scraps}\?:\?\&{constant}\?\8\K\?8\_000;\C{number of tokens in \ADA\
texts being parsed}\6
\\{stack\_size}\?:\?\&{constant}\?\8\K\?800;\C{number of simultaneous output
levels}\par
\A sections~15, 16, 17, 18, 20, 26, 29, 31, 34, 38, 41, 42, 43, 45, 46, 50, 51,
56, 57, 59, 66, 73, 74, 82, 84, 99, 109, 125, 131, 136, 139, 150, 167, 168,
171, 172, 243, 247, 263, 271, 276, 283, 285, 287, 288, 290, and~307.
\U section~3\*.\fi

\M15. A global variable called \\{history} will contain one of four values
at the end of every run: \\{spotless} means that no unusual messages were
printed; \\{harmless\_message} means that a message of possible interest
was printed but no serious errors were detected; \\{error\_message} means that
at least one error was found; \\{fatal\_message} means that the program
terminated abnormally. The value of \\{history} does not influence the
behavior of the program; it is simply computed for the convenience
of systems that might want to use such information.

\Y\P\4\D\\{mark\_harmless}\?\S\?\6
\&{if}\1\1\8\\{history}\?=\?\\{spotless}\2\8\&{then}\6
\\{history}\?\K\?\\{harmless\_message};\2\6
\&{end}\8\2\&{if}\1\par
\P\4\D\\{mark\_error}\?\S\?\\{history}\?\K\?\\{error\_message}\par
\P\4\D\\{mark\_fatal}\?\S\?\\{history}\?\K\?\\{fatal\_message}\par
\Y\P\?\4\X14:Types and objects visible from \\{data\_structures}\X\mathrel{+}\S%
\?\6
\&{type}\8\\{history\_type}\8\&{is}\?\8\1(\?\\{spotless}\?,\39\?\\{harmless%
\_message}\?,\39\?\\{error\_message}\?,\39\?\\{fatal\_message}\?\2)\?;\6
\\{history}\?:\?\\{history\_type}\?\K\?\\{spotless};\C{how bad was this run?}%
\par
\fi

\N16.  The character set.
One of the main goals in the design of \.{AWEB} has been to make it readily
portable between a wide variety of computers. Yet \.{AWEB} by its very
nature must use a greater variety of characters than most computer
programs deal with, and character encoding is one of the areas in which
existing machines differ most widely from each other.

To resolve this problem, all input to \.{AWEAVE} and \.{ATANGLE} is converted
to an internal seven-bit code that is essentially standard ASCII, the
``American Standard Code for Information Interchange.''  The conversion is
done immediately when each character is read in. Conversely, characters are
converted from ASCII to the user's external representation just before
they are output.

Such an internal code is relevant to users of \.{AWEB} only because it is
the code used for preprocessed constants like \.{@\'A\'}. If you are
writing a program in \.{AWEB} that makes use of such one-character constants,
you
should convert your input to ASCII form, like \.{AWEAVE} and \.{ATANGLE} do.
Otherwise \.{AWEB}'s internal coding scheme does not affect you.

Here is a table of the standard visible ASCII codes:
$$\def\:{\char\count255\global\advance\count255 by 1}
\count255='40
\vbox{
\hbox{\hbox to 40pt{\it\hfill0\/\hfill}%
\hbox to 40pt{\it\hfill1\/\hfill}%
\hbox to 40pt{\it\hfill2\/\hfill}%
\hbox to 40pt{\it\hfill3\/\hfill}%
\hbox to 40pt{\it\hfill4\/\hfill}%
\hbox to 40pt{\it\hfill5\/\hfill}%
\hbox to 40pt{\it\hfill6\/\hfill}%
\hbox to 40pt{\it\hfill7\/\hfill}}
\vskip 4pt
\hrule
\def\^{\vrule height 10.5pt depth 4.5pt}
\halign{\hbox to 0pt{\hskip -24pt\bn{8}{#0}\hfill}&\^
\hbox to 40pt{\tt\hfill#\hfill\^}&
&\hbox to 40pt{\tt\hfill#\hfill\^}\cr
04&\:&\:&\:&\:&\:&\:&\:&\:\cr\noalign{\hrule}
05&\:&\:&\:&\:&\:&\:&\:&\:\cr\noalign{\hrule}
06&\:&\:&\:&\:&\:&\:&\:&\:\cr\noalign{\hrule}
07&\:&\:&\:&\:&\:&\:&\:&\:\cr\noalign{\hrule}
10&\:&\:&\:&\:&\:&\:&\:&\:\cr\noalign{\hrule}
11&\:&\:&\:&\:&\:&\:&\:&\:\cr\noalign{\hrule}
12&\:&\:&\:&\:&\:&\:&\:&\:\cr\noalign{\hrule}
13&\:&\:&\:&\:&\:&\:&\:&\:\cr\noalign{\hrule}
14&\:&\:&\:&\:&\:&\:&\:&\:\cr\noalign{\hrule}
15&\:&\:&\:&\:&\:&\:&\:&\:\cr\noalign{\hrule}
16&\:&\:&\:&\:&\:&\:&\:&\:\cr\noalign{\hrule}
17&\:&\:&\:&\:&\:&\:&\:\cr}
\hrule width 280pt}$$
(Actually, of course, code \bn{8}{040} is an invisible blank space.)  Code
\bn{8}{136} was once an upward arrow (\.{\char'13}), and code \bn{8}{137} was
once a left arrow (\.^^X), in olden times when the first draft
of ASCII code was prepared; but \.{AWEB} works with today's standard
ASCII in which those codes represent circumflex and underline as shown.

\Y\P\?\4\X14:Types and objects visible from \\{data\_structures}\X\mathrel{+}\S%
\?\6
\&{subtype}\8\\{ASCII\_code}\8\&{is}\8\p{INTEGER}\8\&{range}\80\?\to\?127;%
\C{seven-bit numbers, a subrange           of the integers}\par
\fi

\M17. We shall use the name \\{text\_char} to stand for the data type of the
characters in the input and output files.  We shall also assume that
\\{text\_char} consists of the elements \\{text\_char}\?'\?\p{FIRST}
through  \\{text\_char}\?'\?\p{LAST}, inclusive. The following
definition should be adjusted if necessary.

\Y\P\?\4\X14:Types and objects visible from \\{data\_structures}\X\mathrel{+}\S%
\?\6
\&{subtype}\8\\{text\_char}\8\&{is}\8\p{CHARACTER};\C{the data type of
characters in text files}\6
\&{subtype}\8\\{text\_file}\8\&{is}\8\p{TEXT\_IO}\?.\?\p{FILE\_TYPE};\C{file
type consisting of \\{text\_char}}\par
\fi

\M18. The \.{AWEAVE} and \.{ATANGLE} processors convert between ASCII code and
the user's external character set by means of arrays \\{xord} and \\{xchr}
that are analogous to \ADA's \p{POS} and \p{VAL} attributes.

\Y\P\?\4\X14:Types and objects visible from \\{data\_structures}\X\mathrel{+}\S%
\?\6
\\{xord}\?:\1\&{array}\8\1(\?\\{text\_char}\?\2)\2\8\&{of}\?\8\\{ASCII\_code};%
\C{specifies conversion of input characters}\6
\\{xchr}\?:\1\&{array}\8\1(\?\\{ASCII\_code}\?\2)\2\8\&{of}\?\8\\{text\_char};%
\C{specifies conversion of output characters}\par
\fi

\M19. If we assume that every system using \.{AWEB} is able to read and write
the
visible characters of standard ASCII (although not necessarily using the
ASCII codes to represent them), the following assignment statements
initialize most of the \\{xchr} array properly, without needing any
system-dependent changes. For example, the statement
\.{xchr(8\#101\#):=\'A\';} that appears in the present \.{AWEB} file might be
encoded in, say, {\mc EBCDIC} code on the external medium on which it
resides, but \.{ATANGLE} will convert from this external code to ASCII and
back again. Therefore the assignment statement \.{XCHR(8\#101\#):=\'A\';}
will appear in the corresponding \ADA\ file, and \ADA\ will compile this
statement so that \\{xchr}\?(\bn{8}{101})\? receives  the character \.A in the
external
(\\{text\_char}) code. Note that it would be quite incorrect to say
\.{xchr(8\#101\#):=\AT!\'A\';}, because \H{A} is a constant of type
\p{INTEGER}, not \p{CHARACTER}, and because we have \H{A}\?=\?65 regardless of
the
external character set.

\Y\P\?\4\X19:Set initial values\X\S\?\6
\\{xchr}\?\1(\bn{8}{40}\2)\K\?\.{\'\ \'};\5
\\{xchr}\?\1(\bn{8}{41}\2)\K\?\.{\'!\'};\5
\\{xchr}\?\1(\bn{8}{42}\2)\K\?\.{\'"\'};\5
\\{xchr}\?\1(\bn{8}{43}\2)\K\?\.{\'\#\'};\5
\\{xchr}\?\1(\bn{8}{44}\2)\K\?\.{\'\$\'};\5
\\{xchr}\?\1(\bn{8}{45}\2)\K\?\.{\'\%\'};\5
\\{xchr}\?\1(\bn{8}{46}\2)\K\?\.{\'\&\'};\5
\\{xchr}\?\1(\bn{8}{47}\2)\K\?\.{\'\'\'};\6
\\{xchr}\?\1(\bn{8}{50}\2)\K\?\.{\'(\'};\5
\\{xchr}\?\1(\bn{8}{51}\2)\K\?\.{\')\'};\5
\\{xchr}\?\1(\bn{8}{52}\2)\K\?\.{\'*\'};\5
\\{xchr}\?\1(\bn{8}{53}\2)\K\?\.{\'+\'};\5
\\{xchr}\?\1(\bn{8}{54}\2)\K\?\.{\',\'};\5
\\{xchr}\?\1(\bn{8}{55}\2)\K\?\.{\'-\'};\5
\\{xchr}\?\1(\bn{8}{56}\2)\K\?\.{\'.\'};\5
\\{xchr}\?\1(\bn{8}{57}\2)\K\?\.{\'/\'};\6
\\{xchr}\?\1(\bn{8}{60}\2)\K\?\.{\'0\'};\5
\\{xchr}\?\1(\bn{8}{61}\2)\K\?\.{\'1\'};\5
\\{xchr}\?\1(\bn{8}{62}\2)\K\?\.{\'2\'};\5
\\{xchr}\?\1(\bn{8}{63}\2)\K\?\.{\'3\'};\5
\\{xchr}\?\1(\bn{8}{64}\2)\K\?\.{\'4\'};\5
\\{xchr}\?\1(\bn{8}{65}\2)\K\?\.{\'5\'};\5
\\{xchr}\?\1(\bn{8}{66}\2)\K\?\.{\'6\'};\5
\\{xchr}\?\1(\bn{8}{67}\2)\K\?\.{\'7\'};\6
\\{xchr}\?\1(\bn{8}{70}\2)\K\?\.{\'8\'};\5
\\{xchr}\?\1(\bn{8}{71}\2)\K\?\.{\'9\'};\5
\\{xchr}\?\1(\bn{8}{72}\2)\K\?\.{\':\'};\5
\\{xchr}\?\1(\bn{8}{73}\2)\K\?\.{\';\'};\5
\\{xchr}\?\1(\bn{8}{74}\2)\K\?\.{\'<\'};\5
\\{xchr}\?\1(\bn{8}{75}\2)\K\?\.{\'=\'};\5
\\{xchr}\?\1(\bn{8}{76}\2)\K\?\.{\'>\'};\5
\\{xchr}\?\1(\bn{8}{77}\2)\K\?\.{\'?\'};\6
\\{xchr}\?\1(\bn{8}{100}\2)\K\?\.{\'@\'};\5
\\{xchr}\?\1(\bn{8}{101}\2)\K\?\.{\'A\'};\5
\\{xchr}\?\1(\bn{8}{102}\2)\K\?\.{\'B\'};\5
\\{xchr}\?\1(\bn{8}{103}\2)\K\?\.{\'C\'};\5
\\{xchr}\?\1(\bn{8}{104}\2)\K\?\.{\'D\'};\5
\\{xchr}\?\1(\bn{8}{105}\2)\K\?\.{\'E\'};\5
\\{xchr}\?\1(\bn{8}{106}\2)\K\?\.{\'F\'};\5
\\{xchr}\?\1(\bn{8}{107}\2)\K\?\.{\'G\'};\6
\\{xchr}\?\1(\bn{8}{110}\2)\K\?\.{\'H\'};\5
\\{xchr}\?\1(\bn{8}{111}\2)\K\?\.{\'I\'};\5
\\{xchr}\?\1(\bn{8}{112}\2)\K\?\.{\'J\'};\5
\\{xchr}\?\1(\bn{8}{113}\2)\K\?\.{\'K\'};\5
\\{xchr}\?\1(\bn{8}{114}\2)\K\?\.{\'L\'};\5
\\{xchr}\?\1(\bn{8}{115}\2)\K\?\.{\'M\'};\5
\\{xchr}\?\1(\bn{8}{116}\2)\K\?\.{\'N\'};\5
\\{xchr}\?\1(\bn{8}{117}\2)\K\?\.{\'O\'};\6
\\{xchr}\?\1(\bn{8}{120}\2)\K\?\.{\'P\'};\5
\\{xchr}\?\1(\bn{8}{121}\2)\K\?\.{\'Q\'};\5
\\{xchr}\?\1(\bn{8}{122}\2)\K\?\.{\'R\'};\5
\\{xchr}\?\1(\bn{8}{123}\2)\K\?\.{\'S\'};\5
\\{xchr}\?\1(\bn{8}{124}\2)\K\?\.{\'T\'};\5
\\{xchr}\?\1(\bn{8}{125}\2)\K\?\.{\'U\'};\5
\\{xchr}\?\1(\bn{8}{126}\2)\K\?\.{\'V\'};\5
\\{xchr}\?\1(\bn{8}{127}\2)\K\?\.{\'W\'};\6
\\{xchr}\?\1(\bn{8}{130}\2)\K\?\.{\'X\'};\5
\\{xchr}\?\1(\bn{8}{131}\2)\K\?\.{\'Y\'};\5
\\{xchr}\?\1(\bn{8}{132}\2)\K\?\.{\'Z\'};\5
\\{xchr}\?\1(\bn{8}{133}\2)\K\?\.{\'[\'};\5
\\{xchr}\?\1(\bn{8}{134}\2)\K\?\.{\'\\\'};\5
\\{xchr}\?\1(\bn{8}{135}\2)\K\?\.{\']\'};\5
\\{xchr}\?\1(\bn{8}{136}\2)\K\?\.{\'\^\'};\5
\\{xchr}\?\1(\bn{8}{137}\2)\K\?\.{\'\_\'};\6
\\{xchr}\?\1(\bn{8}{140}\2)\K\?\.{\'\`\'};\5
\\{xchr}\?\1(\bn{8}{141}\2)\K\?\.{\'a\'};\5
\\{xchr}\?\1(\bn{8}{142}\2)\K\?\.{\'b\'};\5
\\{xchr}\?\1(\bn{8}{143}\2)\K\?\.{\'c\'};\5
\\{xchr}\?\1(\bn{8}{144}\2)\K\?\.{\'d\'};\5
\\{xchr}\?\1(\bn{8}{145}\2)\K\?\.{\'e\'};\5
\\{xchr}\?\1(\bn{8}{146}\2)\K\?\.{\'f\'};\5
\\{xchr}\?\1(\bn{8}{147}\2)\K\?\.{\'g\'};\6
\\{xchr}\?\1(\bn{8}{150}\2)\K\?\.{\'h\'};\5
\\{xchr}\?\1(\bn{8}{151}\2)\K\?\.{\'i\'};\5
\\{xchr}\?\1(\bn{8}{152}\2)\K\?\.{\'j\'};\5
\\{xchr}\?\1(\bn{8}{153}\2)\K\?\.{\'k\'};\5
\\{xchr}\?\1(\bn{8}{154}\2)\K\?\.{\'l\'};\5
\\{xchr}\?\1(\bn{8}{155}\2)\K\?\.{\'m\'};\5
\\{xchr}\?\1(\bn{8}{156}\2)\K\?\.{\'n\'};\5
\\{xchr}\?\1(\bn{8}{157}\2)\K\?\.{\'o\'};\6
\\{xchr}\?\1(\bn{8}{160}\2)\K\?\.{\'p\'};\5
\\{xchr}\?\1(\bn{8}{161}\2)\K\?\.{\'q\'};\5
\\{xchr}\?\1(\bn{8}{162}\2)\K\?\.{\'r\'};\5
\\{xchr}\?\1(\bn{8}{163}\2)\K\?\.{\'s\'};\5
\\{xchr}\?\1(\bn{8}{164}\2)\K\?\.{\'t\'};\5
\\{xchr}\?\1(\bn{8}{165}\2)\K\?\.{\'u\'};\5
\\{xchr}\?\1(\bn{8}{166}\2)\K\?\.{\'v\'};\5
\\{xchr}\?\1(\bn{8}{167}\2)\K\?\.{\'w\'};\6
\\{xchr}\?\1(\bn{8}{170}\2)\K\?\.{\'x\'};\5
\\{xchr}\?\1(\bn{8}{171}\2)\K\?\.{\'y\'};\5
\\{xchr}\?\1(\bn{8}{172}\2)\K\?\.{\'z\'};\5
\\{xchr}\?\1(\bn{8}{173}\2)\K\?\.{\'\{\'};\5
\\{xchr}\?\1(\bn{8}{174}\2)\K\?\.{\'|\'};\5
\\{xchr}\?\1(\bn{8}{175}\2)\K\?\.{\'\}\'};\5
\\{xchr}\?\1(\bn{8}{176}\2)\K\?\.{\'\~\'};\6
\\{xchr}\?\1(\?0\?\2)\K\?\.{\'\ \'};\5
\\{xchr}\?\1(\bn{8}{177}\2)\K\?\.{\'\ \'};\C{these ASCII codes are not used}\par
\A sections~21, 22, 44, 47, 52, 58, 117, 144, 146, and~289.
\U section~3\*.\fi

\M20. Some of the ASCII codes below \bn{8}{40} have been given symbolic names
in
\.{AWEAVE} and \.{ATANGLE} because they are used with a special meaning.

\Y\P\?\4\X14:Types and objects visible from \\{data\_structures}\X\mathrel{+}\S%
\?\6
\\{and\_sign}\?:\?\&{constant}\8\\{ASCII\_code}\?\K\bn{8}{4}\?;\C{equivalent to
`\.{and}'}\6
\\{not\_sign}\?:\?\&{constant}\8\\{ASCII\_code}\?\K\bn{8}{5}\?;\C{equivalent to
`\.{not}'}\6
\\{set\_element\_sign}\?:\?\&{constant}\8\\{ASCII\_code}\?\K\bn{8}{6}\?;%
\C{equivalent to `\.{in}'}\6
\\{tab\_mark}\?:\?\&{constant}\8\\{ASCII\_code}\?\K\bn{8}{11}\?;\C{ASCII code
used as tab-skip}\6
\\{line\_feed}\?:\?\&{constant}\8\\{ASCII\_code}\?\K\bn{8}{12}\?;\C{ASCII code
thrown away at end of line}\6
\\{form\_feed}\?:\?\&{constant}\8\\{ASCII\_code}\?\K\bn{8}{14}\?;\C{ASCII code
used at end of page}\6
\\{carriage\_return}\?:\?\&{constant}\8\\{ASCII\_code}\?\K\bn{8}{15}\?;\C{ASCII
code used at end of line}\6
\\{right\_arrow}\?:\?\&{constant}\8\\{ASCII\_code}\?\K\bn{8}{21}\?;%
\C{equivalent to `\.{=>}'}\6
\\{box\_sign}\?:\?\&{constant}\8\\{ASCII\_code}\?\K\bn{8}{22}\?;\C{equivalent
to `\.{<>}'}\6
\\{double\_star}\?:\?\&{constant}\8\\{ASCII\_code}\?\K\bn{8}{23}\?;%
\C{equivalent to `\.{**}'}\6
\\{left\_label\_bracket}\?:\?\&{constant}\8\\{ASCII\_code}\?\K\bn{8}{24}\?;%
\C{equivalent to `\.{<<}'}\6
\\{right\_label\_bracket}\?:\?\&{constant}\8\\{ASCII\_code}\?\K\bn{8}{25}\?;%
\C{equivalent to `\.{>>}'}\6
\\{left\_arrow}\?:\?\&{constant}\8\\{ASCII\_code}\?\K\bn{8}{30}\?;\C{equivalent
to `\.{:=}'}\6
\\{not\_equal}\?:\?\&{constant}\8\\{ASCII\_code}\?\K\bn{8}{32}\?;\C{equivalent
to `\.{/=}'}\6
\\{less\_or\_equal}\?:\?\&{constant}\8\\{ASCII\_code}\?\K\bn{8}{34}\?;%
\C{equivalent to `\.{<=}'}\6
\\{greater\_or\_equal}\?:\?\&{constant}\8\\{ASCII\_code}\?\K\bn{8}{35}\?;%
\C{equivalent to `\.{>=}'}\6
\\{equivalence\_sign}\?:\?\&{constant}\8\\{ASCII\_code}\?\K\bn{8}{36}\?;%
\C{equivalent to `\.{==}'}\6
\\{or\_sign}\?:\?\&{constant}\8\\{ASCII\_code}\?\K\bn{8}{37}\?;\C{equivalent to
`\.{or}'}\par
\fi

\M21. Here now is the system-dependent part of the character set.
If \.{AWEB} is being implemented on a garden-variety \ADA\ for which
only standard ASCII codes will appear in the input and output files, you
don't need to make any changes here.

Changes to the present module will make \.{AWEB} more friendly on computers
that have an extended character set, so that one can type things like
\.^^Z\ instead of \.{/=}. If you have an extended set of characters that
are easily incorporated into text files, you can assign codes arbitrarily
here, giving an \\{xchr} equivalent to whatever characters the users of
\.{AWEB} are allowed to have in their input files, provided that unsuitable
characters do not correspond to special codes like \\{carriage\_return}
that are listed above.

(The present file \.{AWEAVE.AWEB} does not contain any of the non-ASCII
characters, because it is intended to be used with all implementations of
\.{AWEB}.)

\Y\P\?\4\X19:Set initial values\X\mathrel{+}\S\?\6
\\{xchr}\?\1(\?1\?\to\bn{8}{37}\2)\K\1(\?\&{others}\?\8\KR\?\.{\'\ \'}\?\2)\?;%
\par
\fi

\M22. The following system-independent code makes the \\{xord} array contain a
suitable inverse to the information in \\{xchr}.

\Y\P\?\4\X19:Set initial values\X\mathrel{+}\S\?\6
\\{xord}\?\K\1(\?\&{others}\?\8\KR\bn{8}{40}\2)\?;\6
\1\&{for}\8\|i\?\in\?1\?\to\bn{8}{176}\?\8\&{loop}\6
\\{xord}\?\1(\?\\{xchr}\?\1(\?\|i\?\2)\2)\K\?\|i;\2\6
\&{end}\8\&{loop};\par
\fi

\N23.  Input and output.
The input conventions of this program are identical to those of \.{ATANGLE}.
Therefore people who need to make modifications to both systems
should be able to do so without too many headaches.

\fi

\M24. Terminal output is done by writing on file \p{STANDARD\_OUTPUT}, which is
assumed to consist of characters of type \\{text\_char}:

\Y\P\4\D\\{new\_ln}\?\S\?\p{TEXT\_IO}\?.\?\p{NEW\_LINE}\?\1(\?\p{TEXT\_IO}\?.\?%
\p{STANDARD\_OUTPUT}\?\2)\6
\?\par
\P\4\D\\{print}\?\1(\?\#\?\2)\S\?\p{TEXT\_IO}\?.\?\p{PUT}\?\1(\?\p{TEXT\_IO}\?.%
\?\p{STANDARD\_OUTPUT}\?,\39\?\#\?\2)\C{`\\{print}' means write on the
terminal}\6
\?\par
\P\4\D\\{print\_ln}\?\1(\?\#\?\2)\S\?\\{print}\?\1(\?\#\?\2)\?; \\{new\_ln}\?%
\C{`\\{print}' and then start new line}\6
\?\par
\P\4\D\\{print\_nl}\?\1(\?\#\?\2)\S\?\\{new\_ln}; \\{print}\?\1(\?\#\?\2)%
\C{print information starting on a
new line}\6
\?\par
\P\4\D\\{print\_int}\?\1(\?\#\?\2)\S\?\\{int\_number\_io}\?.\?\p{PUT}\?\1(\?%
\p{TEXT\_IO}\?.\?\p{STANDARD\_OUTPUT}\?,\39\?\#\?\2)\C{print out an integer
value}\6
\?\par
\fi

\M25. The \\{update\_terminal} procedure is called when we want
to make sure that everything we have output to the terminal so far has
actually left the computer's internal buffers and been sent.
(In the \ADA-system used to develop the initial program this is not
required.)

\Y\P\4\D\\{update\_terminal}\?\S\?\&{null}\?\C{empty the terminal output
buffer}\6
\?\par
\fi

\M26. The main input comes from \\{aweb\_file}; this input may be overridden
by changes in \\{change\_file}.
(If \\{change\_file} is empty, there are no changes.)

\Y\P\?\4\X14:Types and objects visible from \\{data\_structures}\X\mathrel{+}\S%
\?\6
\\{aweb\_file}\?:\?\\{text\_file};\C{primary input}\6
\\{change\_file}\?:\?\\{text\_file};\C{updates}\par
\fi

\M27. \P\?\X27:Specification of procedures visible from \\{input\_output}\X\S\?%
\6
\&{procedure}\8\1\\{open\_a\_file}\?(\?\|f\?:\&{in}\?\8\&{out}\8\\{text\_file};%
\?\,\35\?\\{mode\_of\_file}\?:\?\p{TEXT\_IO}\?.\?\p{FILE\_MODE};\?\,\35\?%
\\{default\_name}\?:\?\p{STRING};\?\,\35\1\?\\{success}\?:\?\&{out}\8%
\p{BOOLEAN}\?\2)\?\2;\6
\&{function}\8\1\\{create\_TeX\_file}\?(\1\?\\{TeX\_file\_name}\?:\?\p{STRING}%
\?\2)\&{return}\?\8\p{BOOLEAN}\2;\par
\A sections~32 and~35.
\U section~4\*.\fi

\M28. The following code opens a file and returns \p{TRUE} if succesfully.
If the file is already open the procedure has no effect; this gives
the ability to open the files by means of the change file, without
changing the \.{AWEB}-file of \.{AWEAVE}.


\Y\P\?\4\X28:Procedures in \\{input\_output}\X\S\?\6
\&{procedure}\8\1\\{open\_a\_file}\?(\?\|f\?:\&{in}\?\8\&{out}\8\\{text\_file};%
\?\,\35\?\\{mode\_of\_file}\?:\?\p{TEXT\_IO}\?.\?\p{FILE\_MODE};\?\,\35\?%
\\{default\_name}\?:\?\p{STRING};\?\,\35\1\?\\{success}\?:\?\&{out}\8%
\p{BOOLEAN}\?\2)\?\2\8\&{is}\8\1\6
\4\&{begin}\6
\&{if}\1\1\8\p{TEXT\_IO}\?.\?\p{IS\_OPEN}\?\1(\?\|f\?\2)\?\2\8\&{then}\6
\\{success}\?\K\?\p{TRUE};\6
\4\&{else}\6
\p{TEXT\_IO}\?.\?\p{OPEN}\?\1(\?\p{FILE}\?\KR\?\|f\?,\39\?\p{MODE}\?\KR\?%
\\{mode\_of\_file}\?,\39\?\p{NAME}\?\KR\?\\{default\_name}\?\2)\?;\5
\\{success}\?\K\?\p{TRUE};\2\6
\&{end}\8\2\&{if}\1;\6
\1\4\&{exception}\6
\4\&{when}\8\p{TEXT\_IO}\?.\?\p{STATUS\_ERROR}\?\mathrel{\.|}\?\p{TEXT\_IO}\?.%
\?\p{NAME\_ERROR}\?\mathrel{\.|}\?\p{TEXT\_IO}\?.\?\p{USE\_ERROR}\KR\6
\\{success}\?\K\?\p{FALSE};\5
\&{return};\5
\hbox{\2}\2\6
\&{end}\8\\{open\_a\_file};\par
\A sections~30, 33, and~36.
\U section~4\*.\fi

\M29. The output goes to \\{TeX\_file}.

\Y\P\?\4\X14:Types and objects visible from \\{data\_structures}\X\mathrel{+}\S%
\?\6
\\{TeX\_file}\?:\?\\{text\_file};\par
\fi

\M30. The following code creates \\{TeX\_file}.

\Y\P\?\4\X28:Procedures in \\{input\_output}\X\mathrel{+}\S\?\6
\&{function}\8\1\\{create\_TeX\_file}\?(\1\?\\{TeX\_file\_name}\?:\?\p{STRING}%
\?\2)\&{return}\?\8\p{BOOLEAN}\2\8\&{is}\8\1\6
\4\&{begin}\6
\&{if}\1\1\8\p{TEXT\_IO}\?.\?\p{IS\_OPEN}\?\1(\?\\{TeX\_file}\?\2)\?\2\8%
\&{then}\6
\&{if}\1\1\8\p{TEXT\_IO}\?.\?\p{NAME}\?\1(\?\\{TeX\_file}\?\2)=\?\\{TeX\_file%
\_name}\2\8\&{then}\6
\?\&{return}\?\8\p{TRUE};\6
\4\&{else}\6
\?\&{return}\?\8\p{FALSE};\2\6
\&{end}\8\2\&{if}\1;\6
\4\&{else}\6
\p{TEXT\_IO}\?.\?\p{CREATE}\?\1(\?\p{FILE}\?\KR\?\\{TeX\_file}\?,\39\?\p{MODE}%
\?\KR\?\p{TEXT\_IO}\?.\?\p{OUT\_FILE}\?,\39\?\p{NAME}\?\KR\?\\{TeX\_file\_name}%
\?\2)\?;\5
\?\&{return}\?\8\p{TRUE};\2\6
\&{end}\8\2\&{if}\1;\6
\1\4\&{exception}\6
\4\&{when}\8\p{TEXT\_IO}\?.\?\p{STATUS\_ERROR}\?\mathrel{\.|}\?\p{TEXT\_IO}\?.%
\?\p{NAME\_ERROR}\?\mathrel{\.|}\?\p{TEXT\_IO}\?.\?\p{USE\_ERROR}\KR\6
\?\&{return}\?\8\p{FALSE};\5
\hbox{\2}\2\6
\&{end}\8\\{create\_TeX\_file};\par
\fi

\M31. Input goes into an array called \\{buffer}.

\Y\P\?\4\X14:Types and objects visible from \\{data\_structures}\X\mathrel{+}\S%
\?\6
\&{type}\8\\{buffer\_type}\8\&{is}\?\8\1\&{array}\8\1(\?\p{INTEGER}\8\&{range}%
\?\8\BOX\2)\2\8\&{of}\?\8\\{ASCII\_code};\6
\&{subtype}\8\\{long\_buffer\_index}\8\&{is}\8\p{INTEGER}\8\&{range}\80\?\to\?%
\\{long\_buf\_size};\6
\&{subtype}\8\\{buffer\_index}\8\&{is}\8\\{long\_buffer\_index}\8\&{range}\80\?%
\to\?\\{buf\_size};\6
\\{buffer}\?:\?\\{buffer\_type}\?\1(\?\\{buffer\_index}\?\2)\?;\par
\fi

\M32. The \\{input\_ln} procedure brings the next line of input from the
specified
file into the \\{buffer} array and returns the value \p{TRUE}, unless the file
has already been entirely read, in which case it returns \p{FALSE}. The
conventions of \TeX\ are followed; i.e., \\{ASCII\_code} numbers representing
the next line of the file are input into \\{buffer}\?(\?0\?\to\?\\{limit}\?-\?1%
\?)\?;
trailing blanks are ignored;
and the global variable \\{limit} is set to the length of the
line. The value of \\{limit} must be strictly less than \\{buf\_size}.

We assume that none of the \\{ASCII\_code} values
of \\{buffer}\?(\?\|j\?)\? for 0\?\L\?\|j\?<\?\\{limit} is equal to 0, %
\bn{8}{177}, \\{line\_feed},
\\{form\_feed}, or \\{carriage\_return}.

\Y\P\?\4\X27:Specification of procedures visible from \\{input\_output}\X%
\mathrel{+}\S\?\6
\&{function}\8\1\\{input\_ln}\?(\1\?\|f\?:\?\\{text\_file}\?\2)\&{return}\?\8%
\p{BOOLEAN}\2;\par
\fi

\M33. \P\?\X28:Procedures in \\{input\_output}\X\mathrel{+}\S\?\6
\&{function}\8\1\\{input\_ln}\?(\1\?\|f\?:\?\\{text\_file}\?\2)\&{return}\?\8%
\p{BOOLEAN}\2\8\&{is}\8\C{inputs a line or returns \p{FALSE}}\1\6
\\{final\_limit}\?:\?\\{buffer\_index}\?\K\?0;\C{\\{limit} without trailing
blanks}\6
\|c\?:\?\\{text\_char};\6
\4\&{begin}\6
\\{limit}\?\K\?0;\6
\&{if}\1\1\?\8\R\?\p{TEXT\_IO}\?.\?\p{IS\_OPEN}\?\1(\?\|f\?\2)\ORE\?\p{TEXT%
\_IO}\?.\?\p{END\_OF\_FILE}\?\1(\?\|f\?\2)\?\2\8\&{then}\6
\?\&{return}\?\8\p{FALSE};\6
\4\&{else}\6
\1\&{while}\?\8\R\?\p{TEXT\_IO}\?.\?\p{END\_OF\_LINE}\?\1(\?\|f\?\2)\?\8%
\&{loop}\6
\p{TEXT\_IO}\?.\?\p{GET}\?\1(\?\|f\?,\39\?\|c\?\2)\?;\5
\\{buffer}\?\1(\?\\{limit}\?\2)\K\?\\{xord}\?\1(\?\|c\?\2)\?;\5
\\{incr}\?\1(\?\\{limit}\?\2)\?;\6
\&{if}\1\1\8\\{buffer}\?\1(\?\\{limit}\?-\?1\?\2)\I\?\H{\ }\2\8\&{then}\6
\\{final\_limit}\?\K\?\\{limit};\2\6
\&{end}\8\2\&{if}\1;\6
\&{if}\1\1\8\\{limit}\?=\?\\{buf\_size}\2\8\&{then}\6
\p{TEXT\_IO}\?.\?\p{SKIP\_LINE}\?\1(\?\|f\?\2)\?;\5
\\{decr}\?\1(\?\\{limit}\?\2)\?;\C{keep \\{buffer}\?(\?\\{buf\_size}\?)\?
empty}\6
\\{print\_nl}\?\1(\?\.{"!\ Input\ line\ too\ long"}\?\2)\?;\5
\\{loc}\?\K\?0;\5
\\{error};\5
\?\&{return}\?\8\p{TRUE};\2\6
\&{end}\8\2\&{if}\1;\2\6
\&{end}\8\&{loop};\6
\p{TEXT\_IO}\?.\?\p{SKIP\_LINE}\?\1(\?\|f\?\2)\?;\5
\\{limit}\?\K\?\\{final\_limit};\5
\?\&{return}\?\8\p{TRUE};\2\6
\&{end}\8\2\&{if}\1;\2\6
\&{end}\8\\{input\_ln};\par
\fi

\N34.  Reporting errors to the user.
The \.{AWEAVE} processor operates in three phases: first it inputs the source
file and stores cross-reference data, then it inputs the source once again and
produces the \TeX\ output file, and finally it sorts and outputs the index.

The global variables \\{phase\_one} and \\{phase\_three} tell which Phase we
are in.

\Y\P\?\4\X14:Types and objects visible from \\{data\_structures}\X\mathrel{+}\S%
\?\6
\\{phase\_one}\?:\?\p{BOOLEAN};\C{\p{TRUE} in Phase I, \p{FALSE} in Phases II
and III}\6
\\{phase\_three}\?:\?\p{BOOLEAN};\C{\p{TRUE} in Phase III, \p{FALSE} in Phases
I and II}\par
\fi

\M35. The command `\\{err\_print}\?(\?\.{"!\ Error\ message"}\?)\?' will report
a syntax error to
the user, by printing the error message at the beginning of a new line and
then giving an indication of where the error was spotted in the source file.
Note that no period follows the error message, since the error routine
will automatically supply a period.

The actual error indications are provided by a procedure called \\{error}.
However, error messages are not actually reported during phase one,
since errors detected on the first pass will be detected again
during the second.

\Y\P\4\D\\{err\_print}\?\1(\?\#\?\2)\S\?\6
\&{if}\1\1\?\8\R\?\\{phase\_one}\2\8\&{then}\6
\\{print\_nl}\?\1(\?\#\?\2)\?;\5
\\{error};\2\6
\&{end}\8\2\&{if}\1\par
\Y\P\?\4\X27:Specification of procedures visible from \\{input\_output}\X%
\mathrel{+}\S\?\6
\&{procedure}\8\1\\{error}\2;\par
\fi

\M36. \P\?\X28:Procedures in \\{input\_output}\X\mathrel{+}\S\?\6
\&{procedure}\8\1\\{error}\2\8\&{is}\8\C{prints "\.." and location of error
message}\1\6
\|l\?:\?\\{buffer\_index};\C{index into \\{buffer}}\6
\4\&{begin}\6
\X37:Print error location based on input buffer\X;\6
\\{update\_terminal};\5
\\{mark\_error};\2\6
\&{end}\8\\{error};\par
\fi

\M37. The error locations during Phase I can be indicated by using the global
variables \\{loc}, \\{line}, and \\{changing}, which tell respectively the
first
unlooked-at position in \\{buffer}, the current line number, and whether or not
the current line is from \\{change\_file} or \\{aweb\_file}.
This routine should be modified on systems whose standard text editor
has special line-numbering conventions.

\Y\P\?\4\X37:Print error location based on input buffer\X\S\?\6
\&{if}\1\1\8\\{changing}\2\8\&{then}\6
\\{print}\?\1(\?\.{".\ (change\ file\ "}\?\2)\?;\6
\4\&{else}\6
\\{print}\?\1(\?\.{".\ ("}\?\2)\?;\2\6
\&{end}\8\2\&{if}\1;\6
\\{print}\?\1(\?\.{"l."}\?\2)\?;\5
\\{print\_int}\?\1(\?\\{line}\?,\39\?1\?\2)\?;\5
\\{print\_ln}\?\1(\?\.{")"}\?\2)\?;\6
\&{if}\1\1\8\\{loc}\?\G\?\\{limit}\2\8\&{then}\6
\|l\?\K\?\\{limit};\6
\4\&{else}\6
\|l\?\K\?\\{loc};\2\6
\&{end}\8\2\&{if}\1;\6
\1\&{for}\8\|k\?\in\?1\?\to\?\|l\8\&{loop}\6
\&{if}\1\1\8\\{buffer}\?\1(\?\|k\?-\?1\?\2)=\?\\{tab\_mark}\2\8\&{then}\6
\\{print}\?\1(\?\.{"\ "}\?\2)\?;\6
\4\&{else}\6
\\{print}\?\1(\?\\{xchr}\?\1(\?\\{buffer}\?\1(\?\|k\?-\?1\?\2)\2)\2)\?;\C{print
the characters already read}\2\6
\&{end}\8\2\&{if}\1;\2\6
\&{end}\8\&{loop};\6
\\{new\_ln};\6
\1\&{for}\8\|k\?\in\?1\?\to\?\|l\8\&{loop}\6
\\{print}\?\1(\?\.{"\ "}\?\2)\?;\2\6
\&{end}\8\&{loop};\C{space out the next line}\6
\1\&{for}\8\|k\?\in\?\|l\?+\?1\?\to\?\\{limit}\8\&{loop}\6
\\{print}\?\1(\?\\{xchr}\?\1(\?\\{buffer}\?\1(\?\|k\?-\?1\?\2)\2)\2)\?;\2\6
\&{end}\8\&{loop};\C{print the part not yet read}\6
\&{if}\1\1\8\\{buffer}\?\1(\?\\{limit}\?\2)=\?\H{|}\2\8\&{then}\6
\\{print}\?\1(\?\\{xchr}\?\1(\?\H{|}\?\2)\2)\?;\2\6
\&{end}\8\2\&{if}\1;\C{end of \ADA\ text in module names}\6
\\{print}\?\1(\?\.{"\ "}\?\2)\?\C{this space separates the message from future
asterisks}\par
\U section~36.\fi

\M38. The \\{jump\_out} macro just raises the \\{end\_of\_AWEAVE} exception.
It is used when no recovery from a particular error has
been provided.

\Y\P\4\D\\{jump\_out}\?\S\?\&{raise}\8\\{end\_of\_AWEAVE}\par
\P\4\D\\{fatal\_error}\?\1(\?\#\?\2)\S\?\\{print\_nl}\?\1(\?\#\?\2)\?;\5
\\{error};\5
\\{mark\_fatal};\5
\\{jump\_out}\par
\Y\P\?\4\X14:Types and objects visible from \\{data\_structures}\X\mathrel{+}\S%
\?\6
\\{end\_of\_AWEAVE}\?:\&{exception}\?;\par
\fi

\M39. Sometimes the program's behavior is far different from what it should be,
and \.{AWEAVE} prints an error message that is really for the \.{AWEAVE}
maintenance person, not the user. In such cases the program says
\\{confusion}\?(\?\.{"indication\ of\ where\ we\ are"}\?)\?.

\Y\P\4\D\\{confusion}\?\1(\?\#\?\2)\S\?\\{fatal\_error}\?\1(\?\.{"!\ This\ can%
\'t\ happen\ ("}\?\mathbin{\.\&}\?\#\?\mathbin{\.\&}\?\.{")"}\?\2)\?\par
\fi

\M40. An overflow stop occurs if \.{AWEAVE}'s tables aren't large enough.

\Y\P\4\D\\{overflow}\?\1(\?\#\?\2)\S\?\\{fatal\_error}\?\1(\?\.{"!\ Sorry,\ "}%
\?\mathbin{\.\&}\?\#\?\mathbin{\.\&}\?\.{"\ capacity\ exceeded"}\?\2)\?\par
\fi

\N41.  Data structures.
During the first phase of its processing, \.{AWEAVE} puts identifier names,
index entries, and module names into the large \\{byte\_mem} array, which is
packed with seven-bit integers. Allocation is sequential, since names are
never deleted.

An auxiliary array \\{byte\_start} is used as a directory for \\{byte\_mem},
and the \\{link}, \\{ilk}, and \\{xref} arrays give further information about
names.
These auxiliary arrays consist of sixteen-bit items.

\Y\P\?\4\X14:Types and objects visible from \\{data\_structures}\X\mathrel{+}\S%
\?\6
\&{subtype}\8\\{eight\_bits}\8\&{is}\8\p{INTEGER}\8\&{range}\80\?\to\?255;%
\C{unsigned one-byte quantity}\6
\&{subtype}\8\\{sixteen\_bits}\8\&{is}\8\p{INTEGER}\8\&{range}\80\?\to\?65%
\_535;\C{unsigned two-byte quantity}\par
\fi

\M42. \.{AWEAVE} has been designed to avoid the need for indices that are more
than sixteen bits wide, so that it can be used on most computers. But
there are programs that need more than 65\_536 bytes; \TeX\ is one of these.
To get around this problem, a slight complication has been added to the
data structures:  \\{byte\_mem} is a two-dimensional array, whose first index
is either 0 or 1. (For generality, the first index is actually allowed to
run between 0 and \\{ww}\?-\?1, where \\{ww} is defined to be 2; the program
will
work for any positive value of \\{ww}, and it can be simplified in obvious
ways if \\{ww}\?=\?1.)

\Y\P\?\4\X14:Types and objects visible from \\{data\_structures}\X\mathrel{+}\S%
\?\6
\\{ww}\?:\?\&{constant}\8\p{INTEGER}\?\K\?2;\C{we multiply the byte capacity by
approximately this amount}\6
\&{subtype}\8\\{w\_range}\8\&{is}\8\p{INTEGER}\8\&{range}\80\?\to\?\\{ww}\?-%
\?1;\6
\&{subtype}\8\\{byte\_index}\8\&{is}\8\p{INTEGER}\8\&{range}\80\?\to\?\\{max%
\_bytes};\6
\&{subtype}\8\\{name\_index}\8\&{is}\8\p{INTEGER}\8\&{range}\80\?\to\?\\{max%
\_names};\7
\&{type}\8\\{byte\_memory}\8\&{is}\?\8\1\&{array}\8\1(\?\\{w\_range}\?,\39\?%
\\{byte\_index}\?\2)\2\8\&{of}\?\8\\{ASCII\_code};\6
\&{optim}\1\6
\&{pragma}\8\1\p{PACK}\?\1(\?\\{byte\_memory}\?\2)\?\2;\2\6
\&{mitpo}\7
\&{type}\8\\{name\_info}\8\&{is}\?\8\1\&{array}\8\1(\?\\{name\_index}\?\2)\2\8%
\&{of}\?\8\\{sixteen\_bits};\7
\\{byte\_mem}\?:\?\\{byte\_memory};\C{characters of names}\6
\\{byte\_start}\?:\?\\{name\_info};\C{directory into \\{byte\_mem}}\6
\\{link}\?:\?\\{name\_info};\C{hash table or tree links}\6
\\{ilk}\?:\?\\{name\_info};\C{type codes or tree links}\6
\\{xref}\?:\?\\{name\_info};\C{heads of cross-reference lists}\par
\fi

\M43. The names of identifiers are found by computing a hash address \|h and
then looking at strings of bytes signified by \\{hash}\?(\?\|h\?)\?, \\{link}%
\?(\?\\{hash}\?(\?\|h\?))\?,
\\{link}\?(\?\\{link}\?(\?\\{hash}\?(\?\|h\?)))\?, \dots, until either finding
the desired name
or encountering a zero.

A `\\{name\_pointer}' variable, which signifies a name, is an index into
\\{byte\_start}. The actual sequence of characters in the name pointed to by
\|p appears in positions \\{byte\_start}\?(\?\|p\?)\? to \\{byte\_start}\?(\?%
\|p\?+\?\\{ww}\?)-\?1, inclusive,
in the segment of \\{byte\_mem} whose first index is \|p\?\mathbin{\&{mod}}\?%
\\{ww}. Thus, when
\\{ww}\?=\?2 the even-numbered name bytes appear in \\{byte\_mem}\?(\?0\?,\?%
\hbox{$*$}\?)\?
and the odd-numbered ones appear in \\{byte\_mem}\?(\?1\?,\?\hbox{$*$}\?)\?.
The pointer 0 is used for undefined module names; we don't
want to use it for the names of identifiers, since 0 stands for a null
pointer in a linked list.

We usually have \\{byte\_start}\?(\?\\{name\_ptr}\?+\?\|w\?)=\?\\{byte\_ptr}%
\?((\?\\{name\_ptr}\?+\?\|w\?)\mathbin{\&{mod}}\?\\{ww}\?)\?
for 0\?\L\?\|w\?<\?\\{ww}, since these are the starting positions for the next %
\\{ww}
names to be stored in \\{byte\_mem}.

\Y\P\4\D\\{length}\?\1(\?\#\?\2)\S\?\\{byte\_start}\?\1(\?\#\?+\?\\{ww}\?\2)-\?%
\\{byte\_start}\?\1(\?\#\?\2)\C{the length of a name}\6
\?\par
\P\4\D\\{null\_array}\?\S\1(\?\&{others}\?\8\KR\?0\?\2)\C{all elements have
value 0}\6
\?\par
\Y\P\?\4\X14:Types and objects visible from \\{data\_structures}\X\mathrel{+}\S%
\?\6
\&{subtype}\8\\{name\_pointer}\8\&{is}\8\\{name\_index};\C{identifies a name}\7
\\{name\_ptr}\?:\?\\{name\_pointer}\?\K\?1;\C{first unused position in \\{byte%
\_start}}\6
\\{byte\_ptr}\?:\1\&{array}\8\1(\?\\{w\_range}\?\2)\2\8\&{of}\?\8\\{byte%
\_index}\?\K\?\\{null\_array};\C{first unused position in \\{byte\_mem}}\par
\fi

\M44. \P\?\X19:Set initial values\X\mathrel{+}\S\?\6
\\{byte\_start}\?\1(\?0\?\to\?\\{ww}\?\2)\K\?\\{null\_array};\par
\fi

\M45. Several types of identifiers are distinguished by their \\{ilk}:

\yskip\hang \\{normal} identifiers are part of the \ADA\ program and
will appear in italic type.

\yskip\hang \\{predefined} identifiers are part of the \ADA\ program and
will appear like `\p{INTEGER}'.

\yskip\hang \\{roman} identifiers are index entries that appear after
\.{@\^} in the \.{AWEB} file.

\yskip\hang \\{wildcard} identifiers are index entries that appear after
\.{@:} in the \.{AWEB} file.

\yskip\hang \\{italic} identifiers are index entries that appear after
\.{@i} in the \.{AWEB} file.

\yskip\hang \\{bold\_face} identifiers are index entries that appear after
\.{@b} in the \.{AWEB} file.

\yskip\hang \\{slanted} identifiers are index entries that appear after
\.{@s} in the \.{AWEB} file.

\yskip\hang \\{typewriter} identifiers are index entries that appear after
\.{@.} in the \.{AWEB} file.

\yskip\hang \\{accept\_like}, \\{and\_like}, \dots, \\{with\_like}
identifiers are \ADA\ reserved words whose \\{ilk} explains how they are
to be treated when \ADA\ code is being formatted.

\Y\P\4\D\\{reserved}\?\1(\?\#\?\2)\S\1(\?\\{ilk}\?\1(\?\#\?\2)>\?\\{predefined}%
\?\2)\C{tells if a name is a reserved word}\6
\?\par
\P\4\D\\{defined}\?\1(\?\#\?\2)\S\1(\?\\{ilk}\?\1(\?\#\?\2)>\?\\{typewriter}\?%
\2)\C{tells if a name is predefined or reserved}\6
\?\par
\Y\P\?\4\X14:Types and objects visible from \\{data\_structures}\X\mathrel{+}\S%
\?\6
\\{normal}\?:\?\&{constant}\8\p{INTEGER}\?\K\?0;\C{ordinary identifiers have %
\\{normal} ilk}\6
\\{roman}\?:\?\&{constant}\8\p{INTEGER}\?\K\?1;\C{normal index entries have %
\\{roman} ilk}\6
\\{wildcard}\?:\?\&{constant}\8\p{INTEGER}\?\K\?2;\C{user-formatted index
entries have \\{wildcard} ilk}\6
\\{italic}\?:\?\&{constant}\8\p{INTEGER}\?\K\?3;\C{`italic-type' entries have %
\\{italic} ilk}\6
\\{bold\_face}\?:\?\&{constant}\8\p{INTEGER}\?\K\?4;\C{`boldface-type' entries
have \\{bold\_face} ilk}\6
\\{slanted}\?:\?\&{constant}\8\p{INTEGER}\?\K\?5;\C{`slanted-type' entries have
\\{slanted} ilk}\6
\\{typewriter}\?:\?\&{constant}\8\p{INTEGER}\?\K\?6;\C{`typewriter type'
entries have \\{typewriter} ilk}\6
\\{predefined}\?:\?\&{constant}\8\p{INTEGER}\?\K\?7;\C{predefined identifiers
have \\{predefined} ilk}\6
\\{accept\_like}\?:\?\&{constant}\8\p{INTEGER}\?\K\?8;\C{\&{accept}}\6
\\{all\_like}\?:\?\&{constant}\8\p{INTEGER}\?\K\?9;\C{\&{all}, \&{null}, %
\&{terminate}}\6
\\{and\_like}\?:\?\&{constant}\8\p{INTEGER}\?\K\?10;\C{\&{and}}\6
\\{array\_like}\?:\?\&{constant}\8\p{INTEGER}\?\K\?11;\C{\&{array}}\6
\\{at\_like}\?:\?\&{constant}\8\p{INTEGER}\?\K\?12;\C{\&{at}}\6
\\{begin\_like}\?:\?\&{constant}\8\p{INTEGER}\?\K\?13;\C{\&{begin}, \&{select}}%
\6
\\{body\_like}\?:\?\&{constant}\8\p{INTEGER}\?\K\?14;\C{\&{body}}\6
\\{case\_like}\?:\?\&{constant}\8\p{INTEGER}\?\K\?15;\C{\&{case}}\6
\\{declare\_like}\?:\?\&{constant}\8\p{INTEGER}\?\K\?16;\C{\&{declare}}\6
\\{else\_like}\?:\?\&{constant}\8\p{INTEGER}\?\K\?17;\C{\&{else}}\6
\\{elsif\_like}\?:\?\&{constant}\8\p{INTEGER}\?\K\?18;\C{\&{elsif}}\6
\\{end\_like}\?:\?\&{constant}\8\p{INTEGER}\?\K\?19;\C{\&{end}}\6
\\{exception\_like}\?:\?\&{constant}\8\p{INTEGER}\?\K\?20;\C{\&{exception}}\6
\\{generic\_like}\?:\?\&{constant}\8\p{INTEGER}\?\K\?21;\C{\&{generic}}\6
\\{goto\_like}\?:\?\&{constant}\8\p{INTEGER}\?\K\?22;\C{\&{abort}, \dots, %
\&{while}, if not stated                                     otherwise}\6
\\{if\_like}\?:\?\&{constant}\8\p{INTEGER}\?\K\?23;\C{\&{if}}\6
\\{in\_like}\?:\?\&{constant}\8\p{INTEGER}\?\K\?24;\C{\&{in}}\6
\\{is\_like}\?:\?\&{constant}\8\p{INTEGER}\?\K\?25;\C{\&{do}, \&{is}}\6
\\{loop\_like}\?:\?\&{constant}\8\p{INTEGER}\?\K\?26;\C{\&{loop}}\6
\\{mod\_like}\?:\?\&{constant}\8\p{INTEGER}\?\K\?27;\C{\&{mod}, \&{rem}, %
\&{xor}}\6
\\{new\_like}\?:\?\&{constant}\8\p{INTEGER}\?\K\?28;\C{\&{new}}\6
\\{not\_like}\?:\?\&{constant}\8\p{INTEGER}\?\K\?29;\C{\&{not}}\6
\\{of\_like}\?:\?\&{constant}\8\p{INTEGER}\?\K\?30;\C{\&{of}}\6
\\{or\_like}\?:\?\&{constant}\8\p{INTEGER}\?\K\?31;\C{\&{or}}\6
\\{pragma\_like}\?:\?\&{constant}\8\p{INTEGER}\?\K\?32;\C{\&{pragma}}\6
\\{private\_like}\?:\?\&{constant}\8\p{INTEGER}\?\K\?33;\C{\&{private}}\6
\\{proc\_like}\?:\?\&{constant}\8\p{INTEGER}\?\K\?34;\C{\&{entry}, %
\&{function}, \&{package},                                    \&{procedure}, %
\&{task}}\6
\\{range\_like}\?:\?\&{constant}\8\p{INTEGER}\?\K\?35;\C{\&{range}}\6
\\{record\_like}\?:\?\&{constant}\8\p{INTEGER}\?\K\?36;\C{\&{record}}\6
\\{return\_like}\?:\?\&{constant}\8\p{INTEGER}\?\K\?37;\C{\&{return}, \&{exit}}%
\6
\\{separate\_like}\?:\?\&{constant}\8\p{INTEGER}\?\K\?38;\C{\&{separate}}\6
\\{then\_like}\?:\?\&{constant}\8\p{INTEGER}\?\K\?39;\C{\&{then}}\6
\\{type\_like}\?:\?\&{constant}\8\p{INTEGER}\?\K\?40;\C{\&{subtype}, \&{type}, %
\&{for}}\6
\\{when\_like}\?:\?\&{constant}\8\p{INTEGER}\?\K\?41;\C{\&{when}}\6
\\{with\_like}\?:\?\&{constant}\8\p{INTEGER}\?\K\?42;\C{\&{use}, \&{with}}\par
\fi

\M46. The names of modules are stored in \\{byte\_mem} together
with the identifier names, but a hash table is not used for them because
\.{AWEAVE} needs to be able to recognize a module name when given a prefix of
that name. A conventional binary seach tree is used to retrieve module names,
with fields called \\{llink} and \\{rlink} in place of \\{link} and \\{ilk}.
The
root of this tree is \\{rlink}\?(\?0\?)\?.

\Y\P\?\4\X14:Types and objects visible from \\{data\_structures}\X\mathrel{+}\S%
\?\6
\\{llink}\?:\?\\{name\_info}\8\&{renames}\8\\{link};\C{left link in binary
search tree for module names}\6
\\{rlink}\?:\?\\{name\_info}\8\&{renames}\8\\{ilk};\C{right link in binary
search tree for module names}\6
\\{root}\?:\?\\{sixteen\_bits}\8\&{renames}\8\\{rlink}\?\1(\?0\?\2)\?;\C{the
root of the binary search tree for module names}\par
\fi

\M47. \P\?\X19:Set initial values\X\mathrel{+}\S\?\6
\\{root}\?\K\?0;\C{the binary search tree starts out with nothing in it}\par
\fi

\M48. Here is a little procedure that prints the text of a given name.

\Y\P\?\4\X48:Specification of procedures visible from \\{data\_structures}\X\S%
\?\6
\&{procedure}\8\1\\{print\_id}\?(\1\?\|p\?:\?\\{name\_pointer}\?\2)\?\2;\par
\A section~53.
\U section~3\*.\fi

\M49. \P\?\X49:Procedures in \\{data\_structures}\X\S\?\6
\&{procedure}\8\1\\{print\_id}\?(\1\?\|p\?:\?\\{name\_pointer}\?\2)\?\2\8\&{is}%
\8\C{print identifier or module name}\1\6
\|w\?:\?\\{w\_range};\C{segment of \\{byte\_mem}}\6
\4\&{begin}\6
\&{if}\1\1\8\|p\?\G\?\\{name\_ptr}\2\8\&{then}\6
\\{print}\?\1(\?\.{"IMPOSSIBLE"}\?\2)\?;\6
\4\&{else}\6
\|w\?\K\?\|p\?\mathbin{\&{mod}}\?\\{ww};\6
\1\&{for}\8\|k\?\in\?\\{byte\_start}\?\1(\?\|p\?\2)\to\?\\{byte\_start}\?\1(\?%
\|p\?+\?\\{ww}\?\2)-\?1\8\&{loop}\6
\\{print}\?\1(\?\\{xchr}\?\1(\?\\{byte\_mem}\?\1(\?\|w\?,\39\?\|k\?\2)\2)\2)\?;%
\2\6
\&{end}\8\&{loop};\2\6
\&{end}\8\2\&{if}\1;\2\6
\&{end}\8\\{print\_id};\par
\A sections~54 and~55.
\U section~3\*.\fi

\M50. We keep track of the current module number in
\\{module\_count}, which is the total number of modules that have started.
Modules which have been altered by a change file entry
have their \\{changed\_module} flag turned on during the first phase.

\Y\P\?\4\X14:Types and objects visible from \\{data\_structures}\X\mathrel{+}\S%
\?\6
\&{subtype}\8\\{module\_index}\8\&{is}\8\p{INTEGER}\8\&{range}\80\?\to\?\\{max%
\_modules};\6
\\{module\_count}\?:\?\\{module\_index};\C{the current module number}\6
\\{changed\_module}\?:\1\&{array}\8\1(\?\\{module\_index}\?\2)\2\8\&{of}\?\8%
\p{BOOLEAN};\C{is it changed?}\6
\\{change\_exists}\?:\?\p{BOOLEAN};\C{has any module changed?}\par
\fi

\M51. The other large memory area in \.{AWEAVE} keeps the cross-reference data.
All uses of the name \|p are recorded in a linked list beginning at
\\{xref}\?(\?\|p\?)\?, which points into the \\{xmem} array. Entries in %
\\{xmem} consist
of two sixteen-bit items per word, called the \\{num} and \\{xlink} fields.
If \|x is an index into \\{xmem}, reached from name \|p, the value of \\{num}%
\?(\?\|x\?)\?
is either a module number where \|p is used, or it is \\{def\_flag} plus a
module number where \|p is defined, or it is \\{body\_flag} plus a module
number where the body of \|p is defined; and \\{xlink}\?(\?\|x\?)\? points to
the next
such cross reference for \|p, if any. This list of cross references is in
decreasing order by module number. The current number of cross references
is \\{xref\_ptr}.

The global variable \\{xref\_switch} is set either to \\{def\_flag}, \\{body%
\_flag} or
to zero, depending on whether the next cross reference to an identifier is to
be underlined, to be set in italic type, or neither of these in the index.
This  switch is incremented by \\{def\_flag} when
\.{@!} or \.{@d} or \.{@f} is scanned, and it is cleared to zero when
the next identifier or index entry cross reference has been made. Similarly,
the global variable \\{mod\_xref\_switch} is either \\{def\_flag} or zero,
depending
on whether a module name is being defined or used.

\Y\P\4\D\\{num}\?\1(\?\#\?\2)\S\?\\{xmem}\?\1(\?\#\?\2).\?\\{num\_field}\?\6
\?\par
\P\4\D\\{xlink}\?\1(\?\#\?\2)\S\?\\{xmem}\?\1(\?\#\?\2).\?\\{xlink\_field}\?\6
\?\par
\Y\P\?\4\X14:Types and objects visible from \\{data\_structures}\X\mathrel{+}\S%
\?\6
\\{def\_flag}\?:\?\&{constant}\8\p{INTEGER}\?\K\?10\_240;\C{must be strictly
larger than \\{max\_modules}}\6
\\{body\_flag}\?:\?\&{constant}\8\p{INTEGER}\?\K\?\\{def\_flag}\?+\?\\{def%
\_flag};\7
\&{subtype}\8\\{xref\_number}\8\&{is}\8\p{INTEGER}\8\&{range}\80\?\to\?\\{max%
\_refs};\6
\&{type}\8\\{xmem\_field}\8\&{is}\8\1\6
\&{record}\1\6
\\{num\_field}\?:\?\\{sixteen\_bits};\C{module number plus zero or \\{def%
\_flag}}\6
\\{xlink\_field}\?:\?\\{sixteen\_bits};\C{pointer to the previous cross
reference}\2\6
\2\&{end}\8\&{record};\7
\\{xmem}\?:\1\&{array}\8\1(\?\\{xref\_number}\?\2)\2\8\&{of}\?\8\\{xmem%
\_field};\6
\\{xref\_ptr}\?:\?\\{xref\_number}\?\K\?0;\C{the largest occupied position in %
\\{xmem}}\6
\\{xref\_switch}\?,\39\?\\{mod\_xref\_switch}\?:\?\p{INTEGER}\8\&{range}\80\?%
\to\?\\{body\_flag}\?\K\?0;\C{zero, \\{def\_flag} or \\{body\_flag}}\par
\fi

\M52. \P\?\X19:Set initial values\X\mathrel{+}\S\?\6
\\{num}\?\1(\?0\?\2)\K\?0;\5
\\{xref}\?\1(\?0\?\2)\K\?0;\C{cross references to undefined modules}\par
\fi

\M53. A new cross reference for an identifier is formed by calling \\{new%
\_xref},
which discards duplicate entries and ignores non-underlined references
to one-letter identifiers or \ADA's reserved words.

\Y\P\4\D\\{append\_xref}\?\1(\?\#\?\2)\S\?\6
\&{if}\1\1\8\\{xref\_ptr}\?=\?\\{max\_refs}\2\8\&{then}\6
\\{overflow}\?\1(\?\.{"cross\ reference"}\?\2)\?;\6
\4\&{else}\6
\\{incr}\?\1(\?\\{xref\_ptr}\?\2)\?;\5
\\{num}\?\1(\?\\{xref\_ptr}\?\2)\K\?\#;\2\6
\&{end}\8\2\&{if}\1\par
\Y\P\?\4\X48:Specification of procedures visible from \\{data\_structures}\X%
\mathrel{+}\S\?\6
\&{procedure}\8\1\\{new\_xref}\?(\1\?\|p\?:\?\\{name\_pointer}\?\2)\?\2;\6
\&{procedure}\8\1\\{new\_mod\_xref}\?(\1\?\|p\?:\?\\{name\_pointer}\?\2)\?\2;%
\par
\fi

\M54. \P\?\X49:Procedures in \\{data\_structures}\X\mathrel{+}\S\?\6
\&{procedure}\8\1\\{new\_xref}\?(\1\?\|p\?:\?\\{name\_pointer}\?\2)\?\2\8\&{is}%
\8\1\6
\|q\?:\?\\{xref\_number};\C{pointer to previous cross reference}\6
\|m\?,\39\?\|n\?:\?\\{sixteen\_bits};\C{new and previous cross-reference value}%
\6
\4\&{begin}\6
\&{if}\1\1\?\8\1(\?\\{reserved}\?\1(\?\|p\?\2)\V\1(\?\\{byte\_start}\?\1(\?\|p%
\?\2)+\?1\?=\?\\{byte\_start}\?\1(\?\|p\?+\?\\{ww}\?\2)\2)\2)\W\1(\?\\{xref%
\_switch}\?=\?0\?\2)\?\2\8\&{then}\6
\&{return};\2\6
\&{end}\8\2\&{if}\1;\6
\|m\?\K\?\\{module\_count}\?+\?\\{xref\_switch};\5
\\{xref\_switch}\?\K\?0;\5
\|q\?\K\?\\{xref}\?\1(\?\|p\?\2)\?;\6
\&{if}\1\1\8\|q\?>\?0\2\8\&{then}\6
\|n\?\K\?\\{num}\?\1(\?\|q\?\2)\?;\6
\&{if}\1\1\?\8\1(\?\|n\?=\?\|m\?\2)\V\1(\?\|n\?-\?\\{def\_flag}\?=\?\|m\?\2)\V%
\1(\?\|n\?-\?\\{body\_flag}\?=\?\|m\?\2)\?\2\8\&{then}\6
\&{return};\6
\4\&{elsif}\1\?\8\1(\?\|m\?-\?\\{def\_flag}\?=\?\|n\?\2)\V\1(\?\|m\?-\?\\{body%
\_flag}\?=\?\|n\?\2)\?\2\8\&{then}\6
\\{num}\?\1(\?\|q\?\2)\K\?\|m;\5
\&{return};\2\6
\&{end}\8\2\&{if}\1;\2\6
\&{end}\8\2\&{if}\1;\6
\\{append\_xref}\?\1(\?\|m\?\2)\?;\5
\\{xlink}\?\1(\?\\{xref\_ptr}\?\2)\K\?\|q;\5
\\{xref}\?\1(\?\|p\?\2)\K\?\\{xref\_ptr};\2\6
\&{end}\8\\{new\_xref};\par
\fi

\M55. The cross reference lists for module names are slightly different.
Suppose
that a module name is defined in modules $m_1$, \dots, $m_k$ and used in
modules $n_1$, \dots, $n_l$. Then its list will contain $m_1+\\{def\_flag}$,
$m_k+\\{def\_flag}$, \dots, $m_2+\\{def\_flag}$, $n_l$, \dots, $n_1$, in
this order.  After Phase II, however, the order will be
$m_1+\\{def\_flag}$, \dots, $m_k+\\{def\_flag}$, $n_1$, \dots, $n_l$.

\Y\P\?\4\X49:Procedures in \\{data\_structures}\X\mathrel{+}\S\?\6
\&{procedure}\8\1\\{new\_mod\_xref}\?(\1\?\|p\?:\?\\{name\_pointer}\?\2)\?\2\8%
\&{is}\8\1\6
\|q\?,\39\?\|r\?:\?\\{xref\_number};\C{pointers to previous cross references}\6
\4\&{begin}\6
\|q\?\K\?\\{xref}\?\1(\?\|p\?\2)\?;\5
\|r\?\K\?0;\6
\&{if}\1\1\8\|q\?>\?0\2\8\&{then}\6
\&{if}\1\1\8\\{mod\_xref\_switch}\?=\?0\2\8\&{then}\6
\1\&{while}\8\\{num}\?\1(\?\|q\?\2)\G\?\\{def\_flag}\8\&{loop}\6
\|r\?\K\?\|q;\5
\|q\?\K\?\\{xlink}\?\1(\?\|q\?\2)\?;\2\6
\&{end}\8\&{loop};\6
\4\&{elsif}\1\8\\{num}\?\1(\?\|q\?\2)\G\?\\{def\_flag}\2\8\&{then}\6
\|r\?\K\?\|q;\5
\|q\?\K\?\\{xlink}\?\1(\?\|q\?\2)\?;\2\6
\&{end}\8\2\&{if}\1;\2\6
\&{end}\8\2\&{if}\1;\6
\\{append\_xref}\?\1(\?\\{module\_count}\?+\?\\{mod\_xref\_switch}\?\2)\?;\5
\\{xlink}\?\1(\?\\{xref\_ptr}\?\2)\K\?\|q;\5
\\{mod\_xref\_switch}\?\K\?0;\6
\&{if}\1\1\8\|r\?=\?0\2\8\&{then}\6
\\{xref}\?\1(\?\|p\?\2)\K\?\\{xref\_ptr};\6
\4\&{else}\6
\\{xlink}\?\1(\?\|r\?\2)\K\?\\{xref\_ptr};\2\6
\&{end}\8\2\&{if}\1;\2\6
\&{end}\8\\{new\_mod\_xref};\par
\fi

\M56. A third large area of memory is used for sixteen-bit `tokens', which
appear
in short lists similar to the strings of characters in \\{byte\_mem}. Token
lists
are used to contain the result of \ADA\ code translated into \TeX\ form;
further details about them will be explained later. A \\{text\_pointer}
variable
is an index into \\{tok\_start}.

\Y\P\?\4\X14:Types and objects visible from \\{data\_structures}\X\mathrel{+}\S%
\?\6
\&{subtype}\8\\{text\_pointer}\8\&{is}\8\p{INTEGER}\8\&{range}\80\?\to\?\\{max%
\_texts};\C{identifies a token list}\par
\fi

\M57. The first position of \\{tok\_mem}
that is unoccupied by replacement text is called \\{tok\_ptr}, and the first
unused location of \\{tok\_start} is called \\{text\_ptr}.
Thus, we usually have \\{tok\_start}\?(\?\\{text\_ptr}\?)=\?\\{tok\_ptr}.

\Y\P\?\4\X14:Types and objects visible from \\{data\_structures}\X\mathrel{+}\S%
\?\6
\&{subtype}\8\\{token\_index}\8\&{is}\8\p{INTEGER}\8\&{range}\80\?\to\?\\{max%
\_toks};\6
\\{tok\_mem}\?:\1\&{array}\8\1(\?\\{token\_index}\?\2)\2\8\&{of}\?\8\\{sixteen%
\_bits};\C{tokens}\6
\\{tok\_start}\?:\1\&{array}\8\1(\?\\{text\_pointer}\?\2)\2\8\&{of}\?\8%
\\{sixteen\_bits};\C{directory into \\{tok\_mem}}\6
\\{text\_ptr}\?:\?\\{text\_pointer}\?\K\?1;\C{first unused position in \\{tok%
\_start}}\6
\\{tok\_ptr}\?:\?\\{token\_index}\?\K\?1;\C{first unused position in \\{tok%
\_mem}}\6
\&{stat}\1\6
\\{max\_tok\_ptr}\?,\39\?\\{max\_txt\_ptr}\?:\?\\{token\_index}\?\K\?1;%
\C{largest values occurring}\2\6
\&{tats}\par
\fi

\M58. \P\?\X19:Set initial values\X\mathrel{+}\S\?\6
\\{tok\_start}\?\1(\?0\?\to\?1\?\2)\K\1(\?1\?,\39\?1\?\2)\?;\par
\fi

\N59.  Searching for identifiers.
The hash table described above is updated by the \\{id\_lookup} procedure,
which finds a given identifier and returns a pointer to its index in
\\{byte\_start}. The identifier is supposed to match character by character
and it is also supposed to have a given \\{ilk} code; the same name may be
present more than once if it is supposed to appear in the index with
different typesetting conventions.
If the identifier was not already present, it is inserted into the table.

Because of the way \.{AWEAVE}'s scanning mechanism works, it is most convenient
to let \\{id\_lookup} search for an identifier that is present in the %
\\{buffer}
array. Two other global variables specify its position in the buffer: the
first character is \\{buffer}\?(\?\\{id\_first}\?)\?, and the last is %
\\{buffer}\?(\?\\{id\_loc}\?-\?1\?)\?.

\Y\P\?\4\X14:Types and objects visible from \\{data\_structures}\X\mathrel{+}\S%
\?\6
\\{id\_first}\?:\?\\{buffer\_index};\C{where the current identifier begins in
the buffer}\6
\\{id\_loc}\?:\?\\{buffer\_index};\C{just after the current identifier in the
buffer}\7
\&{subtype}\8\\{hash\_index}\8\&{is}\8\p{INTEGER}\8\&{range}\80\?\to\?\\{hash%
\_size};\6
\\{hash}\?:\1\&{array}\8\1(\?\\{hash\_index}\?\2)\2\8\&{of}\?\8\\{sixteen%
\_bits}\?\K\?\\{null\_array};\C{heads of hash lists---initially empty}\par
\fi

\M60. Here now is the main procedure for finding identifiers (and strings).
The parameter \|t is set to the desired \\{ilk} code.
The identifier must either have \\{ilk}\?=\?\|t, or we must have \|t\?=\?%
\\{normal} and the
identifier must be predefined or a reserved word.

\Y\P\?\4\X60:Specification of procedures visible from \\{searching}\X\S\?\6
\&{function}\8\1\\{id\_lookup}\?(\1\?\|t\?:\?\\{eight\_bits}\?\2)\&{return}\?\8%
\\{name\_pointer}\2;\par
\A sections~75 and~79.
\U section~5\*.\fi

\M61. \P\?\X61:Procedures in \\{searching}\X\S\?\6
\&{function}\8\1\\{id\_lookup}\?(\1\?\|t\?:\?\\{eight\_bits}\?\2)\&{return}\?\8%
\\{name\_pointer}\2\8\&{is}\8\C{finds current identifier}\1\6
\|c\?:\?\\{eight\_bits};\C{byte being chopped}\6
\|i\?:\?\\{buffer\_index};\C{index into \\{buffer}}\6
\|h\?:\?\\{hash\_index};\C{hash code}\6
\|k\?:\?\\{byte\_index};\C{index into \\{byte\_mem}}\6
\|w\?:\?\\{w\_range};\C{segment of \\{byte\_mem}}\6
\|l\?:\?\\{buffer\_index}\?\K\?\\{id\_loc}\?-\?\\{id\_first};\C{length of the
given identifier}\6
\|p\?,\39\?\|q\?:\?\\{name\_pointer};\C{where the identifier is being sought}\6
\4\&{begin}\6
\X62:Compute the hash code \|h\X;\6
\X63:Compute the name location \|p\X;\6
\4\BL\\{found}\EL\8\&{if}\1\1\8\|p\?=\?\\{name\_ptr}\2\8\&{then}\6
\X65:Enter a new name into the table at position \|p\X;\2\6
\&{end}\8\2\&{if}\1;\6
\?\&{return}\?\8\|p;\2\6
\&{end}\8\\{id\_lookup};\par
\A sections~76 and~80.
\U section~5\*.\fi

\M62. A simple hash code is used: If the sequence of
ASCII codes is $c_1c_2\ldots c_n$, its hash value will be
$$(2^{n-1}c_1+2^{n-2}c_2+\cdots+c_n)\,\bmod\,\\{hash\_size}.$$

\Y\P\?\4\X62:Compute the hash code \|h\X\S\?\6
\|h\?\K\?\\{buffer}\?\1(\?\\{id\_first}\?\2)\?;\5
\|i\?\K\?\\{id\_first}\?+\?1;\6
\1\&{while}\8\|i\?<\?\\{id\_loc}\8\&{loop}\6
\|h\?\K\1(\?\|h\?+\?\|h\?+\?\\{buffer}\?\1(\?\|i\?\2)\2)\mathbin{\&{mod}}\?%
\\{hash\_size};\5
\\{incr}\?\1(\?\|i\?\2)\?;\2\6
\&{end}\8\&{loop}\par
\U section~61.\fi

\M63. If the identifier is new, it will be placed in position \|p\?=\?\\{name%
\_ptr},
otherwise \|p will point to its existing location.

\Y\P\?\4\X63:Compute the name location \|p\X\S\?\6
\|p\?\K\?\\{hash}\?\1(\?\|h\?\2)\?;\6
\1\&{while}\8\|p\?\I\?0\8\&{loop}\6
\&{if}\1\1\?\8\1(\?\\{length}\?\1(\?\|p\?\2)=\?\|l\?\2)\W\1(\1(\?\\{ilk}\?\1(\?%
\|p\?\2)=\?\|t\?\2)\V\1(\1(\?\|t\?=\?\\{normal}\?\2)\W\?\\{defined}\?\1(\?\|p\?%
\2)\2)\2)\?\2\8\&{then}\6
\X64:Compare name \|p with current identifier, \&{goto}\8\\{found} if equal\X;%
\2\6
\&{end}\8\2\&{if}\1;\6
\|p\?\K\?\\{link}\?\1(\?\|p\?\2)\?;\2\6
\&{end}\8\&{loop};\6
\|p\?\K\?\\{name\_ptr};\C{the current identifier is new}\6
\\{link}\?\1(\?\|p\?\2)\K\?\\{hash}\?\1(\?\|h\?\2)\?;\5
\\{hash}\?\1(\?\|h\?\2)\K\?\|p\C{insert \|p at beginning of hash list}\par
\U section~61.\fi

\M64. \P\?\X64:Compare name \|p with current identifier, \&{goto}\8\\{found} if
equal\X\S\?\6
\|i\?\K\?\\{id\_first};\5
\|k\?\K\?\\{byte\_start}\?\1(\?\|p\?\2)\?;\5
\|w\?\K\?\|p\?\mathbin{\&{mod}}\?\\{ww};\6
\1\&{while}\?\8\1(\?\|i\?<\?\\{id\_loc}\?\2)\W\1(\?\\{buffer}\?\1(\?\|i\?\2)=\?%
\\{byte\_mem}\?\1(\?\|w\?,\39\?\|k\?\2)\2)\?\8\&{loop}\6
\\{incr}\?\1(\?\|i\?\2)\?;\5
\\{incr}\?\1(\?\|k\?\2)\?;\2\6
\&{end}\8\&{loop};\6
\&{if}\1\1\8\|i\?=\?\\{id\_loc}\2\8\&{then}\6
\&{goto}\8\\{found};\C{all characters agree}\2\6
\&{end}\8\2\&{if}\1\par
\U section~63.\fi

\M65. When we begin the following segment of the program, \|p\?=\?\\{name%
\_ptr}.

\Y\P\?\4\X65:Enter a new name into the table at position \|p\X\S\?\6
\|w\?\K\?\\{name\_ptr}\?\mathbin{\&{mod}}\?\\{ww};\5
\|k\?\K\?\\{byte\_ptr}\?\1(\?\|w\?\2)\?;\6
\&{if}\1\1\8\|k\?+\?\|l\?>\?\\{max\_bytes}\2\8\&{then}\6
\\{overflow}\?\1(\?\.{"byte\ memory"}\?\2)\?;\2\6
\&{end}\8\2\&{if}\1;\6
\&{if}\1\1\8\\{name\_ptr}\?>\?\\{max\_names}\?-\?\\{ww}\2\8\&{then}\6
\\{overflow}\?\1(\?\.{"name"}\?\2)\?;\2\6
\&{end}\8\2\&{if}\1;\6
\|i\?\K\?\\{id\_first};\C{get ready to move the identifier into \\{byte\_mem}}\6
\1\&{while}\8\|i\?<\?\\{id\_loc}\8\&{loop}\6
\\{byte\_mem}\?\1(\?\|w\?,\39\?\|k\?\2)\K\?\\{buffer}\?\1(\?\|i\?\2)\?;\5
\\{incr}\?\1(\?\|k\?\2)\?;\5
\\{incr}\?\1(\?\|i\?\2)\?;\2\6
\&{end}\8\&{loop};\6
\\{byte\_ptr}\?\1(\?\|w\?\2)\K\?\|k;\5
\\{byte\_start}\?\1(\?\\{name\_ptr}\?+\?\\{ww}\?\2)\K\?\|k;\5
\\{incr}\?\1(\?\\{name\_ptr}\?\2)\?;\5
\\{ilk}\?\1(\?\|p\?\2)\K\?\|t;\5
\\{xref}\?\1(\?\|p\?\2)\K\?0\par
\U section~61.\fi

\N66.  Initializing the table of reserved and predefined words.
We have to get \ADA's reserved words and predefined identifiers into the
hash table, and the
simplest way to do this is to insert them every time \.{AWEAVE} is run.
A few macros permit us to do the initialization with a compact program.

\Y\P\4\D\\{sid}\?\1(\?\#\?\2)\S\1(\?\#\?\2)\?; \\{cur\_name}\?\K\?\\{id%
\_lookup}\?\6
\?\par
\P\4\D\\{id}\?\1(\?\#\?\2)\S\?\\{id\_loc}\?\K\?\#\?+\?1; \\{buffer}\?\1(\?1\?%
\to\?\#\?\2)\K\?\\{sid}\?\6
\?\par
\P\4\D\\{psid}\?\1(\?\#\?\2)\S\1(\?\#\?\2)\?;\5
\\{cur\_name}\?\K\?\\{id\_lookup}\?\1(\?\\{predefined}\?\2)\?\par
\P\4\D\\{pid}\?\1(\?\#\?\2)\S\?\\{id\_loc}\?\K\?\#\?+\?1; \\{buffer}\?\1(\?1\?%
\to\?\#\?\2)\K\?\\{psid}\?\6
\?\par
\Y\P\?\4\X14:Types and objects visible from \\{data\_structures}\X\mathrel{+}\S%
\?\6
\\{cur\_name}\?:\?\\{name\_pointer};\C{points to the identifier just inserted}%
\par
\fi

\M67. First we define the reserved words of \ADA.

\Y\P\?\4\X67:Store all the reserved words\X\S\?\6
\\{id\_first}\?\K\?1;\6
\\{id}\?\1(\?5\?\2)\1(\?\H{a}\?,\39\?\H{b}\?,\39\?\H{o}\?,\39\?\H{r}\?,\39\?%
\H{t}\?\2)\1(\?\\{goto\_like}\?\2)\?;\6
\\{id}\?\1(\?3\?\2)\1(\?\H{a}\?,\39\?\H{b}\?,\39\?\H{s}\?\2)\1(\?\\{goto\_like}%
\?\2)\?;\6
\\{id}\?\1(\?6\?\2)\1(\?\H{a}\?,\39\?\H{c}\?,\39\?\H{c}\?,\39\?\H{e}\?,\39\?%
\H{p}\?,\39\?\H{t}\?\2)\1(\?\\{accept\_like}\?\2)\?;\6
\\{id}\?\1(\?6\?\2)\1(\?\H{a}\?,\39\?\H{c}\?,\39\?\H{c}\?,\39\?\H{e}\?,\39\?%
\H{s}\?,\39\?\H{s}\?\2)\1(\?\\{goto\_like}\?\2)\?;\6
\\{id}\?\1(\?3\?\2)\1(\?\H{a}\?,\39\?\H{l}\?,\39\?\H{l}\?\2)\1(\?\\{all\_like}%
\?\2)\?;\6
\\{id}\?\1(\?3\?\2)\1(\?\H{a}\?,\39\?\H{n}\?,\39\?\H{d}\?\2)\1(\?\\{and\_like}%
\?\2)\?;\6
\\{id}\?\1(\?5\?\2)\1(\?\H{a}\?,\39\?\H{r}\?,\39\?\H{r}\?,\39\?\H{a}\?,\39\?%
\H{y}\?\2)\1(\?\\{array\_like}\?\2)\?;\6
\\{id}\?\1(\?2\?\2)\1(\?\H{a}\?,\39\?\H{t}\?\2)\1(\?\\{at\_like}\?\2)\?;\6
\\{id}\?\1(\?5\?\2)\1(\?\H{b}\?,\39\?\H{e}\?,\39\?\H{g}\?,\39\?\H{i}\?,\39\?%
\H{n}\?\2)\1(\?\\{begin\_like}\?\2)\?;\6
\\{id}\?\1(\?4\?\2)\1(\?\H{b}\?,\39\?\H{o}\?,\39\?\H{d}\?,\39\?\H{y}\?\2)\1(\?%
\\{body\_like}\?\2)\?;\6
\\{id}\?\1(\?4\?\2)\1(\?\H{c}\?,\39\?\H{a}\?,\39\?\H{s}\?,\39\?\H{e}\?\2)\1(\?%
\\{case\_like}\?\2)\?;\6
\\{id}\?\1(\?8\?\2)\1(\?\H{c}\?,\39\?\H{o}\?,\39\?\H{n}\?,\39\?\H{s}\?,\39\?%
\H{t}\?,\39\?\H{a}\?,\39\?\H{n}\?,\39\?\H{t}\?\2)\1(\?\\{goto\_like}\?\2)\?;\6
\\{id}\?\1(\?7\?\2)\1(\?\H{d}\?,\39\?\H{e}\?,\39\?\H{c}\?,\39\?\H{l}\?,\39\?%
\H{a}\?,\39\?\H{r}\?,\39\?\H{e}\?\2)\1(\?\\{declare\_like}\?\2)\?;\6
\\{id}\?\1(\?5\?\2)\1(\?\H{d}\?,\39\?\H{e}\?,\39\?\H{l}\?,\39\?\H{a}\?,\39\?%
\H{y}\?\2)\1(\?\\{goto\_like}\?\2)\?;\6
\\{id}\?\1(\?5\?\2)\1(\?\H{d}\?,\39\?\H{e}\?,\39\?\H{l}\?,\39\?\H{t}\?,\39\?%
\H{a}\?\2)\1(\?\\{goto\_like}\?\2)\?;\6
\\{id}\?\1(\?6\?\2)\1(\?\H{d}\?,\39\?\H{i}\?,\39\?\H{g}\?,\39\?\H{i}\?,\39\?%
\H{t}\?,\39\?\H{s}\?\2)\1(\?\\{goto\_like}\?\2)\?;\6
\\{id}\?\1(\?2\?\2)\1(\?\H{d}\?,\39\?\H{o}\?\2)\1(\?\\{is\_like}\?\2)\?;\6
\\{id}\?\1(\?4\?\2)\1(\?\H{e}\?,\39\?\H{l}\?,\39\?\H{s}\?,\39\?\H{e}\?\2)\1(\?%
\\{else\_like}\?\2)\?;\6
\\{id}\?\1(\?5\?\2)\1(\?\H{e}\?,\39\?\H{l}\?,\39\?\H{s}\?,\39\?\H{i}\?,\39\?%
\H{f}\?\2)\1(\?\\{elsif\_like}\?\2)\?;\6
\\{id}\?\1(\?3\?\2)\1(\?\H{e}\?,\39\?\H{n}\?,\39\?\H{d}\?\2)\1(\?\\{end\_like}%
\?\2)\?;\6
\\{id}\?\1(\?5\?\2)\1(\?\H{e}\?,\39\?\H{n}\?,\39\?\H{t}\?,\39\?\H{r}\?,\39\?%
\H{y}\?\2)\1(\?\\{proc\_like}\?\2)\?;\6
\\{id}\?\1(\?9\?\2)\1(\?\H{e}\?,\39\?\H{x}\?,\39\?\H{c}\?,\39\?\H{e}\?,\39\?%
\H{p}\?,\39\?\H{t}\?,\39\?\H{i}\?,\39\?\H{o}\?,\39\?\H{n}\?\2)\1(\?\\{exception%
\_like}\?\2)\?;\6
\\{id}\?\1(\?4\?\2)\1(\?\H{e}\?,\39\?\H{x}\?,\39\?\H{i}\?,\39\?\H{t}\?\2)\1(\?%
\\{return\_like}\?\2)\?;\6
\\{id}\?\1(\?3\?\2)\1(\?\H{f}\?,\39\?\H{o}\?,\39\?\H{r}\?\2)\1(\?\\{type\_like}%
\?\2)\?;\6
\\{id}\?\1(\?8\?\2)\1(\?\H{f}\?,\39\?\H{u}\?,\39\?\H{n}\?,\39\?\H{c}\?,\39\?%
\H{t}\?,\39\?\H{i}\?,\39\?\H{o}\?,\39\?\H{n}\?\2)\1(\?\\{proc\_like}\?\2)\?;\6
\\{id}\?\1(\?7\?\2)\1(\?\H{g}\?,\39\?\H{e}\?,\39\?\H{n}\?,\39\?\H{e}\?,\39\?%
\H{r}\?,\39\?\H{i}\?,\39\?\H{c}\?\2)\1(\?\\{generic\_like}\?\2)\?;\6
\\{id}\?\1(\?4\?\2)\1(\?\H{g}\?,\39\?\H{o}\?,\39\?\H{t}\?,\39\?\H{o}\?\2)\1(\?%
\\{goto\_like}\?\2)\?;\6
\\{id}\?\1(\?2\?\2)\1(\?\H{i}\?,\39\?\H{f}\?\2)\1(\?\\{if\_like}\?\2)\?;\6
\\{id}\?\1(\?2\?\2)\1(\?\H{i}\?,\39\?\H{n}\?\2)\1(\?\\{in\_like}\?\2)\?;\6
\\{id}\?\1(\?2\?\2)\1(\?\H{i}\?,\39\?\H{s}\?\2)\1(\?\\{is\_like}\?\2)\?;\6
\\{id}\?\1(\?7\?\2)\1(\?\H{l}\?,\39\?\H{i}\?,\39\?\H{m}\?,\39\?\H{i}\?,\39\?%
\H{t}\?,\39\?\H{e}\?,\39\?\H{d}\?\2)\1(\?\\{goto\_like}\?\2)\?;\6
\\{id}\?\1(\?4\?\2)\1(\?\H{l}\?,\39\?\H{o}\?,\39\?\H{o}\?,\39\?\H{p}\?\2)\1(\?%
\\{loop\_like}\?\2)\?;\6
\\{id}\?\1(\?3\?\2)\1(\?\H{m}\?,\39\?\H{o}\?,\39\?\H{d}\?\2)\1(\?\\{mod\_like}%
\?\2)\?;\6
\\{id}\?\1(\?3\?\2)\1(\?\H{n}\?,\39\?\H{e}\?,\39\?\H{w}\?\2)\1(\?\\{new\_like}%
\?\2)\?;\6
\\{id}\?\1(\?3\?\2)\1(\?\H{n}\?,\39\?\H{o}\?,\39\?\H{t}\?\2)\1(\?\\{not\_like}%
\?\2)\?;\6
\\{id}\?\1(\?4\?\2)\1(\?\H{n}\?,\39\?\H{u}\?,\39\?\H{l}\?,\39\?\H{l}\?\2)\1(\?%
\\{all\_like}\?\2)\?;\6
\\{id}\?\1(\?2\?\2)\1(\?\H{o}\?,\39\?\H{f}\?\2)\1(\?\\{of\_like}\?\2)\?;\6
\\{id}\?\1(\?2\?\2)\1(\?\H{o}\?,\39\?\H{r}\?\2)\1(\?\\{or\_like}\?\2)\?;\6
\\{id}\?\1(\?6\?\2)\1(\?\H{o}\?,\39\?\H{t}\?,\39\?\H{h}\?,\39\?\H{e}\?,\39\?%
\H{r}\?,\39\?\H{s}\?\2)\1(\?\\{goto\_like}\?\2)\?;\par
\A section~68.
\U section~308\*.\fi

\M68. \P\?\X67:Store all the reserved words\X\mathrel{+}\S\?\6
\\{id}\?\1(\?3\?\2)\1(\?\H{o}\?,\39\?\H{u}\?,\39\?\H{t}\?\2)\1(\?\\{goto\_like}%
\?\2)\?;\6
\\{id}\?\1(\?7\?\2)\1(\?\H{p}\?,\39\?\H{a}\?,\39\?\H{c}\?,\39\?\H{k}\?,\39\?%
\H{a}\?,\39\?\H{g}\?,\39\?\H{e}\?\2)\1(\?\\{proc\_like}\?\2)\?;\6
\\{id}\?\1(\?6\?\2)\1(\?\H{p}\?,\39\?\H{r}\?,\39\?\H{a}\?,\39\?\H{g}\?,\39\?%
\H{m}\?,\39\?\H{a}\?\2)\1(\?\\{pragma\_like}\?\2)\?;\6
\\{id}\?\1(\?7\?\2)\1(\?\H{p}\?,\39\?\H{r}\?,\39\?\H{i}\?,\39\?\H{v}\?,\39\?%
\H{a}\?,\39\?\H{t}\?,\39\?\H{e}\?\2)\1(\?\\{private\_like}\?\2)\?;\6
\\{id}\?\1(\?9\?\2)\1(\?\H{p}\?,\39\?\H{r}\?,\39\?\H{o}\?,\39\?\H{c}\?,\39\?%
\H{e}\?,\39\?\H{d}\?,\39\?\H{u}\?,\39\?\H{r}\?,\39\?\H{e}\?\2)\1(\?\\{proc%
\_like}\?\2)\?;\6
\\{id}\?\1(\?5\?\2)\1(\?\H{r}\?,\39\?\H{a}\?,\39\?\H{i}\?,\39\?\H{s}\?,\39\?%
\H{e}\?\2)\1(\?\\{goto\_like}\?\2)\?;\6
\\{id}\?\1(\?5\?\2)\1(\?\H{r}\?,\39\?\H{a}\?,\39\?\H{n}\?,\39\?\H{g}\?,\39\?%
\H{e}\?\2)\1(\?\\{range\_like}\?\2)\?;\6
\\{id}\?\1(\?6\?\2)\1(\?\H{r}\?,\39\?\H{e}\?,\39\?\H{c}\?,\39\?\H{o}\?,\39\?%
\H{r}\?,\39\?\H{d}\?\2)\1(\?\\{record\_like}\?\2)\?;\6
\\{id}\?\1(\?3\?\2)\1(\?\H{r}\?,\39\?\H{e}\?,\39\?\H{m}\?\2)\1(\?\\{mod\_like}%
\?\2)\?;\6
\\{id}\?\1(\?7\?\2)\1(\?\H{r}\?,\39\?\H{e}\?,\39\?\H{n}\?,\39\?\H{a}\?,\39\?%
\H{m}\?,\39\?\H{e}\?,\39\?\H{s}\?\2)\1(\?\\{goto\_like}\?\2)\?;\6
\\{id}\?\1(\?6\?\2)\1(\?\H{r}\?,\39\?\H{e}\?,\39\?\H{t}\?,\39\?\H{u}\?,\39\?%
\H{r}\?,\39\?\H{n}\?\2)\1(\?\\{return\_like}\?\2)\?;\6
\\{id}\?\1(\?7\?\2)\1(\?\H{r}\?,\39\?\H{e}\?,\39\?\H{v}\?,\39\?\H{e}\?,\39\?%
\H{r}\?,\39\?\H{s}\?,\39\?\H{e}\?\2)\1(\?\\{goto\_like}\?\2)\?;\6
\\{id}\?\1(\?6\?\2)\1(\?\H{s}\?,\39\?\H{e}\?,\39\?\H{l}\?,\39\?\H{e}\?,\39\?%
\H{c}\?,\39\?\H{t}\?\2)\1(\?\\{begin\_like}\?\2)\?;\6
\\{id}\?\1(\?8\?\2)\1(\?\H{s}\?,\39\?\H{e}\?,\39\?\H{p}\?,\39\?\H{a}\?,\39\?%
\H{r}\?,\39\?\H{a}\?,\39\?\H{t}\?,\39\?\H{e}\?\2)\1(\?\\{separate\_like}\?\2)%
\?;\6
\\{id}\?\1(\?7\?\2)\1(\?\H{s}\?,\39\?\H{u}\?,\39\?\H{b}\?,\39\?\H{t}\?,\39\?%
\H{y}\?,\39\?\H{p}\?,\39\?\H{e}\?\2)\1(\?\\{type\_like}\?\2)\?;\6
\\{id}\?\1(\?4\?\2)\1(\?\H{t}\?,\39\?\H{a}\?,\39\?\H{s}\?,\39\?\H{k}\?\2)\1(\?%
\\{proc\_like}\?\2)\?;\6
\\{id}\?\1(\?9\?\2)\1(\?\H{t}\?,\39\?\H{e}\?,\39\?\H{r}\?,\39\?\H{m}\?,\39\?%
\H{i}\?,\39\?\H{n}\?,\39\?\H{a}\?,\39\?\H{t}\?,\39\?\H{e}\?\2)\1(\?\\{all%
\_like}\?\2)\?;\6
\\{id}\?\1(\?4\?\2)\1(\?\H{t}\?,\39\?\H{h}\?,\39\?\H{e}\?,\39\?\H{n}\?\2)\1(\?%
\\{then\_like}\?\2)\?;\6
\\{id}\?\1(\?4\?\2)\1(\?\H{t}\?,\39\?\H{y}\?,\39\?\H{p}\?,\39\?\H{e}\?\2)\1(\?%
\\{type\_like}\?\2)\?;\6
\\{id}\?\1(\?3\?\2)\1(\?\H{u}\?,\39\?\H{s}\?,\39\?\H{e}\?\2)\1(\?\\{with\_like}%
\?\2)\?;\6
\\{id}\?\1(\?4\?\2)\1(\?\H{w}\?,\39\?\H{h}\?,\39\?\H{e}\?,\39\?\H{n}\?\2)\1(\?%
\\{when\_like}\?\2)\?;\6
\\{id}\?\1(\?5\?\2)\1(\?\H{w}\?,\39\?\H{h}\?,\39\?\H{i}\?,\39\?\H{l}\?,\39\?%
\H{e}\?\2)\1(\?\\{goto\_like}\?\2)\?;\6
\\{id}\?\1(\?4\?\2)\1(\?\H{w}\?,\39\?\H{i}\?,\39\?\H{t}\?,\39\?\H{h}\?\2)\1(\?%
\\{with\_like}\?\2)\?;\6
\\{id}\?\1(\?3\?\2)\1(\?\H{x}\?,\39\?\H{o}\?,\39\?\H{r}\?\2)\1(\?\\{mod\_like}%
\?\2)\?\par
\fi

\M69. Here are the definitions for all the predefined identifiers of the
\ADA\ language; these are related to attributes, pragmas, exceptions,
input/output facilities, and the standard environment.

\Y\P\?\4\X69:Store all the predefined identifiers\X\S\?\6
\\{pid}\?\1(\?7\?\2)\1(\?\H{A}\?,\39\?\H{D}\?,\39\?\H{D}\?,\39\?\H{R}\?,\39\?%
\H{E}\?,\39\?\H{S}\?,\39\?\H{S}\?\2)\?;\6
\\{pid}\?\1(\?3\?\2)\1(\?\H{A}\?,\39\?\H{F}\?,\39\?\H{T}\?\2)\?;\6
\\{pid}\?\1(\?5\?\2)\1(\?\H{A}\?,\39\?\H{S}\?,\39\?\H{C}\?,\39\?\H{I}\?,\39\?%
\H{I}\?\2)\?;\6
\\{pid}\?\1(\?4\?\2)\1(\?\H{B}\?,\39\?\H{A}\?,\39\?\H{S}\?,\39\?\H{E}\?\2)\?;\6
\\{pid}\?\1(\?7\?\2)\1(\?\H{B}\?,\39\?\H{O}\?,\39\?\H{O}\?,\39\?\H{L}\?,\39\?%
\H{E}\?,\39\?\H{A}\?,\39\?\H{N}\?\2)\?;\6
\\{pid}\?\1(\?8\?\2)\1(\?\H{C}\?,\39\?\H{A}\?,\39\?\H{L}\?,\39\?\H{E}\?,\39\?%
\H{N}\?,\39\?\H{D}\?,\39\?\H{A}\?,\39\?\H{R}\?\2)\?;\6
\\{pid}\?\1(\?8\?\2)\1(\?\H{C}\?,\39\?\H{A}\?,\39\?\H{L}\?,\39\?\H{L}\?,\39\?%
\H{A}\?,\39\?\H{B}\?,\39\?\H{L}\?,\39\?\H{E}\?\2)\?;\6
\\{pid}\?\1(\?9\?\2)\1(\?\H{C}\?,\39\?\H{H}\?,\39\?\H{A}\?,\39\?\H{R}\?,\39\?%
\H{A}\?,\39\?\H{C}\?,\39\?\H{T}\?,\39\?\H{E}\?,\39\?\H{R}\?\2)\?;\6
\\{pid}\?\1(\?5\?\2)\1(\?\H{C}\?,\39\?\H{L}\?,\39\?\H{O}\?,\39\?\H{C}\?,\39\?%
\H{K}\?\2)\?;\6
\\{pid}\?\1(\?5\?\2)\1(\?\H{C}\?,\39\?\H{L}\?,\39\?\H{O}\?,\39\?\H{S}\?,\39\?%
\H{E}\?\2)\?;\6
\\{pid}\?\1(\?3\?\2)\1(\?\H{C}\?,\39\?\H{O}\?,\39\?\H{L}\?\2)\?;\6
\\{pid}\?\1(\?11\?\2)\1(\?\H{C}\?,\39\?\H{O}\?,\39\?\H{N}\?,\39\?\H{S}\?,\39\?%
\H{T}\?,\39\?\H{R}\?,\39\?\H{A}\?,\39\?\H{I}\?,\39\?\H{N}\?,\39\?\H{E}\?,\39\?%
\H{D}\?\2)\?;\6
\\{pid}\?\1(\?16\?\2)\1(\?\H{C}\?,\39\?\H{O}\?,\39\?\H{N}\?,\39\?\H{S}\?,\39\?%
\H{T}\?,\39\?\H{R}\?,\39\?\H{A}\?,\39\?\H{I}\?,\39\?\H{N}\?,\39\?\H{T}\?,\39\?%
\H{\_}\?,\39\?\H{E}\?,\39\?\H{R}\?,\39\?\H{R}\?,\39\?\H{O}\?,\39\?\H{R}\?\2)\?;%
\6
\\{pid}\?\1(\?10\?\2)\1(\?\H{C}\?,\39\?\H{O}\?,\39\?\H{N}\?,\39\?\H{T}\?,\39\?%
\H{R}\?,\39\?\H{O}\?,\39\?\H{L}\?,\39\?\H{L}\?,\39\?\H{E}\?,\39\?\H{D}\?\2)\?;\6
\\{pid}\?\1(\?5\?\2)\1(\?\H{C}\?,\39\?\H{O}\?,\39\?\H{U}\?,\39\?\H{N}\?,\39\?%
\H{T}\?\2)\?;\6
\\{pid}\?\1(\?6\?\2)\1(\?\H{C}\?,\39\?\H{R}\?,\39\?\H{E}\?,\39\?\H{A}\?,\39\?%
\H{T}\?,\39\?\H{E}\?\2)\?;\6
\\{pid}\?\1(\?13\?\2)\1(\?\H{C}\?,\39\?\H{U}\?,\39\?\H{R}\?,\39\?\H{R}\?,\39\?%
\H{E}\?,\39\?\H{N}\?,\39\?\H{T}\?,\39\?\H{\_}\?,\39\?\H{I}\?,\39\?\H{N}\?,\39\?%
\H{P}\?,\39\?\H{U}\?,\39\?\H{T}\?\2)\?;\6
\\{pid}\?\1(\?14\?\2)\1(\?\H{C}\?,\39\?\H{U}\?,\39\?\H{R}\?,\39\?\H{R}\?,\39\?%
\H{E}\?,\39\?\H{N}\?,\39\?\H{T}\?,\39\?\H{\_}\?,\39\?\H{O}\?,\39\?\H{U}\?,\39\?%
\H{T}\?,\39\?\H{P}\?,\39\?\H{U}\?,\39\?\H{T}\?\2)\?;\par
\A sections~70, 71, and~72.
\U section~308\*.\fi

\M70. \P\?\X69:Store all the predefined identifiers\X\mathrel{+}\S\?\6
\\{pid}\?\1(\?4\?\2)\1(\?\H{D}\?,\39\?\H{A}\?,\39\?\H{T}\?,\39\?\H{A}\?\2)\?;\6
\\{pid}\?\1(\?10\?\2)\1(\?\H{D}\?,\39\?\H{A}\?,\39\?\H{T}\?,\39\?\H{A}\?,\39\?%
\H{\_}\?,\39\?\H{E}\?,\39\?\H{R}\?,\39\?\H{R}\?,\39\?\H{O}\?,\39\?\H{R}\?\2)\?;%
\6
\\{pid}\?\1(\?4\?\2)\1(\?\H{D}\?,\39\?\H{A}\?,\39\?\H{T}\?,\39\?\H{E}\?\2)\?;\6
\\{pid}\?\1(\?3\?\2)\1(\?\H{D}\?,\39\?\H{A}\?,\39\?\H{Y}\?\2)\?;\6
\\{pid}\?\1(\?12\?\2)\1(\?\H{D}\?,\39\?\H{A}\?,\39\?\H{Y}\?,\39\?\H{\_}\?,\39\?%
\H{D}\?,\39\?\H{U}\?,\39\?\H{R}\?,\39\?\H{A}\?,\39\?\H{T}\?,\39\?\H{I}\?,\39\?%
\H{O}\?,\39\?\H{N}\?\2)\?;\6
\\{pid}\?\1(\?10\?\2)\1(\?\H{D}\?,\39\?\H{A}\?,\39\?\H{Y}\?,\39\?\H{\_}\?,\39\?%
\H{N}\?,\39\?\H{U}\?,\39\?\H{M}\?,\39\?\H{B}\?,\39\?\H{E}\?,\39\?\H{R}\?\2)\?;\6
\\{pid}\?\1(\?11\?\2)\1(\?\H{D}\?,\39\?\H{E}\?,\39\?\H{F}\?,\39\?\H{A}\?,\39\?%
\H{U}\?,\39\?\H{L}\?,\39\?\H{T}\?,\39\?\H{\_}\?,\39\?\H{A}\?,\39\?\H{F}\?,\39\?%
\H{T}\?\2)\?;\6
\\{pid}\?\1(\?12\?\2)\1(\?\H{D}\?,\39\?\H{E}\?,\39\?\H{F}\?,\39\?\H{A}\?,\39\?%
\H{U}\?,\39\?\H{L}\?,\39\?\H{T}\?,\39\?\H{\_}\?,\39\?\H{B}\?,\39\?\H{A}\?,\39\?%
\H{S}\?,\39\?\H{E}\?\2)\?;\6
\\{pid}\?\1(\?11\?\2)\1(\?\H{D}\?,\39\?\H{E}\?,\39\?\H{F}\?,\39\?\H{A}\?,\39\?%
\H{U}\?,\39\?\H{L}\?,\39\?\H{T}\?,\39\?\H{\_}\?,\39\?\H{E}\?,\39\?\H{X}\?,\39\?%
\H{P}\?\2)\?;\6
\\{pid}\?\1(\?12\?\2)\1(\?\H{D}\?,\39\?\H{E}\?,\39\?\H{F}\?,\39\?\H{A}\?,\39\?%
\H{U}\?,\39\?\H{L}\?,\39\?\H{T}\?,\39\?\H{\_}\?,\39\?\H{F}\?,\39\?\H{O}\?,\39\?%
\H{R}\?,\39\?\H{E}\?\2)\?;\6
\\{pid}\?\1(\?15\?\2)\1(\?\H{D}\?,\39\?\H{E}\?,\39\?\H{F}\?,\39\?\H{A}\?,\39\?%
\H{U}\?,\39\?\H{L}\?,\39\?\H{T}\?,\39\?\H{\_}\?,\39\?\H{S}\?,\39\?\H{E}\?,\39\?%
\H{T}\?,\39\?\H{T}\?,\39\?\H{I}\?,\39\?\H{N}\?,\39\?\H{G}\?\2)\?;\6
\\{pid}\?\1(\?13\?\2)\1(\?\H{D}\?,\39\?\H{E}\?,\39\?\H{F}\?,\39\?\H{A}\?,\39\?%
\H{U}\?,\39\?\H{L}\?,\39\?\H{T}\?,\39\?\H{\_}\?,\39\?\H{W}\?,\39\?\H{I}\?,\39\?%
\H{D}\?,\39\?\H{T}\?,\39\?\H{H}\?\2)\?;\6
\\{pid}\?\1(\?6\?\2)\1(\?\H{D}\?,\39\?\H{E}\?,\39\?\H{L}\?,\39\?\H{E}\?,\39\?%
\H{T}\?,\39\?\H{E}\?\2)\?;\6
\\{pid}\?\1(\?5\?\2)\1(\?\H{D}\?,\39\?\H{E}\?,\39\?\H{L}\?,\39\?\H{T}\?,\39\?%
\H{A}\?\2)\?;\6
\\{pid}\?\1(\?6\?\2)\1(\?\H{D}\?,\39\?\H{E}\?,\39\?\H{V}\?,\39\?\H{I}\?,\39\?%
\H{C}\?,\39\?\H{E}\?\2)\?;\6
\\{pid}\?\1(\?12\?\2)\1(\?\H{D}\?,\39\?\H{E}\?,\39\?\H{V}\?,\39\?\H{I}\?,\39\?%
\H{C}\?,\39\?\H{E}\?,\39\?\H{\_}\?,\39\?\H{E}\?,\39\?\H{R}\?,\39\?\H{R}\?,\39\?%
\H{O}\?,\39\?\H{R}\?\2)\?;\6
\\{pid}\?\1(\?6\?\2)\1(\?\H{D}\?,\39\?\H{I}\?,\39\?\H{G}\?,\39\?\H{I}\?,\39\?%
\H{T}\?,\39\?\H{S}\?\2)\?;\6
\\{pid}\?\1(\?9\?\2)\1(\?\H{D}\?,\39\?\H{I}\?,\39\?\H{R}\?,\39\?\H{E}\?,\39\?%
\H{C}\?,\39\?\H{T}\?,\39\?\H{\_}\?,\39\?\H{I}\?,\39\?\H{O}\?\2)\?;\6
\\{pid}\?\1(\?8\?\2)\1(\?\H{D}\?,\39\?\H{U}\?,\39\?\H{R}\?,\39\?\H{A}\?,\39\?%
\H{T}\?,\39\?\H{I}\?,\39\?\H{O}\?,\39\?\H{N}\?\2)\?;\6
\\{pid}\?\1(\?8\?\2)\1(\?\H{E}\?,\39\?\H{L}\?,\39\?\H{B}\?,\39\?\H{O}\?,\39\?%
\H{R}\?,\39\?\H{A}\?,\39\?\H{T}\?,\39\?\H{E}\?\2)\?;\6
\\{pid}\?\1(\?12\?\2)\1(\?\H{E}\?,\39\?\H{L}\?,\39\?\H{E}\?,\39\?\H{M}\?,\39\?%
\H{E}\?,\39\?\H{N}\?,\39\?\H{T}\?,\39\?\H{\_}\?,\39\?\H{T}\?,\39\?\H{Y}\?,\39\?%
\H{P}\?,\39\?\H{E}\?\2)\?;\6
\\{pid}\?\1(\?4\?\2)\1(\?\H{E}\?,\39\?\H{M}\?,\39\?\H{A}\?,\39\?\H{X}\?\2)\?;\6
\\{pid}\?\1(\?9\?\2)\1(\?\H{E}\?,\39\?\H{N}\?,\39\?\H{D}\?,\39\?\H{\_}\?,\39\?%
\H{E}\?,\39\?\H{R}\?,\39\?\H{R}\?,\39\?\H{O}\?,\39\?\H{R}\?\2)\?;\6
\\{pid}\?\1(\?11\?\2)\1(\?\H{E}\?,\39\?\H{N}\?,\39\?\H{D}\?,\39\?\H{\_}\?,\39\?%
\H{O}\?,\39\?\H{F}\?,\39\?\H{\_}\?,\39\?\H{F}\?,\39\?\H{I}\?,\39\?\H{L}\?,\39\?%
\H{E}\?\2)\?;\6
\\{pid}\?\1(\?11\?\2)\1(\?\H{E}\?,\39\?\H{N}\?,\39\?\H{D}\?,\39\?\H{\_}\?,\39\?%
\H{O}\?,\39\?\H{F}\?,\39\?\H{\_}\?,\39\?\H{L}\?,\39\?\H{I}\?,\39\?\H{N}\?,\39\?%
\H{E}\?\2)\?;\6
\\{pid}\?\1(\?11\?\2)\1(\?\H{E}\?,\39\?\H{N}\?,\39\?\H{D}\?,\39\?\H{\_}\?,\39\?%
\H{O}\?,\39\?\H{F}\?,\39\?\H{\_}\?,\39\?\H{P}\?,\39\?\H{A}\?,\39\?\H{G}\?,\39\?%
\H{E}\?\2)\?;\6
\\{pid}\?\1(\?4\?\2)\1(\?\H{E}\?,\39\?\H{N}\?,\39\?\H{U}\?,\39\?\H{M}\?\2)\?;\6
\\{pid}\?\1(\?14\?\2)\1(\?\H{E}\?,\39\?\H{N}\?,\39\?\H{U}\?,\39\?\H{M}\?,\39\?%
\H{E}\?,\39\?\H{R}\?,\39\?\H{A}\?,\39\?\H{T}\?,\39\?\H{I}\?,\39\?\H{O}\?,\39\?%
\H{N}\?,\39\?\H{\_}\?,\39\?\H{I}\?,\39\?\H{O}\?\2)\?;\6
\\{pid}\?\1(\?7\?\2)\1(\?\H{E}\?,\39\?\H{P}\?,\39\?\H{S}\?,\39\?\H{I}\?,\39\?%
\H{L}\?,\39\?\H{O}\?,\39\?\H{N}\?\2)\?;\6
\\{pid}\?\1(\?3\?\2)\1(\?\H{E}\?,\39\?\H{X}\?,\39\?\H{P}\?\2)\?;\6
\\{pid}\?\1(\?5\?\2)\1(\?\H{F}\?,\39\?\H{A}\?,\39\?\H{L}\?,\39\?\H{S}\?,\39\?%
\H{E}\?\2)\?;\6
\\{pid}\?\1(\?5\?\2)\1(\?\H{F}\?,\39\?\H{I}\?,\39\?\H{E}\?,\39\?\H{L}\?,\39\?%
\H{D}\?\2)\?;\6
\\{pid}\?\1(\?4\?\2)\1(\?\H{F}\?,\39\?\H{I}\?,\39\?\H{L}\?,\39\?\H{E}\?\2)\?;\6
\\{pid}\?\1(\?9\?\2)\1(\?\H{F}\?,\39\?\H{I}\?,\39\?\H{L}\?,\39\?\H{E}\?,\39\?%
\H{\_}\?,\39\?\H{M}\?,\39\?\H{O}\?,\39\?\H{D}\?,\39\?\H{E}\?\2)\?;\6
\\{pid}\?\1(\?9\?\2)\1(\?\H{F}\?,\39\?\H{I}\?,\39\?\H{L}\?,\39\?\H{E}\?,\39\?%
\H{\_}\?,\39\?\H{T}\?,\39\?\H{Y}\?,\39\?\H{P}\?,\39\?\H{E}\?\2)\?;\6
\\{pid}\?\1(\?10\?\2)\1(\?\H{F}\?,\39\?\H{I}\?,\39\?\H{N}\?,\39\?\H{E}\?,\39\?%
\H{\_}\?,\39\?\H{D}\?,\39\?\H{E}\?,\39\?\H{L}\?,\39\?\H{T}\?,\39\?\H{A}\?\2)\?;%
\6
\\{pid}\?\1(\?5\?\2)\1(\?\H{F}\?,\39\?\H{I}\?,\39\?\H{R}\?,\39\?\H{S}\?,\39\?%
\H{T}\?\2)\?;\6
\\{pid}\?\1(\?9\?\2)\1(\?\H{F}\?,\39\?\H{I}\?,\39\?\H{R}\?,\39\?\H{S}\?,\39\?%
\H{T}\?,\39\?\H{\_}\?,\39\?\H{B}\?,\39\?\H{I}\?,\39\?\H{T}\?\2)\?;\6
\\{pid}\?\1(\?8\?\2)\1(\?\H{F}\?,\39\?\H{I}\?,\39\?\H{X}\?,\39\?\H{E}\?,\39\?%
\H{D}\?,\39\?\H{\_}\?,\39\?\H{I}\?,\39\?\H{O}\?\2)\?;\6
\\{pid}\?\1(\?5\?\2)\1(\?\H{F}\?,\39\?\H{L}\?,\39\?\H{O}\?,\39\?\H{A}\?,\39\?%
\H{T}\?\2)\?;\6
\\{pid}\?\1(\?8\?\2)\1(\?\H{F}\?,\39\?\H{L}\?,\39\?\H{O}\?,\39\?\H{A}\?,\39\?%
\H{T}\?,\39\?\H{\_}\?,\39\?\H{I}\?,\39\?\H{O}\?\2)\?;\6
\\{pid}\?\1(\?4\?\2)\1(\?\H{F}\?,\39\?\H{O}\?,\39\?\H{R}\?,\39\?\H{E}\?\2)\?;\6
\\{pid}\?\1(\?4\?\2)\1(\?\H{F}\?,\39\?\H{O}\?,\39\?\H{R}\?,\39\?\H{M}\?\2)\?;\6
\\{pid}\?\1(\?4\?\2)\1(\?\H{F}\?,\39\?\H{R}\?,\39\?\H{O}\?,\39\?\H{M}\?\2)\?;\6
\\{pid}\?\1(\?3\?\2)\1(\?\H{G}\?,\39\?\H{E}\?,\39\?\H{T}\?\2)\?;\6
\\{pid}\?\1(\?8\?\2)\1(\?\H{G}\?,\39\?\H{E}\?,\39\?\H{T}\?,\39\?\H{\_}\?,\39\?%
\H{L}\?,\39\?\H{I}\?,\39\?\H{N}\?,\39\?\H{E}\?\2)\?;\6
\\{pid}\?\1(\?5\?\2)\1(\?\H{I}\?,\39\?\H{M}\?,\39\?\H{A}\?,\39\?\H{G}\?,\39\?%
\H{E}\?\2)\?;\6
\\{pid}\?\1(\?5\?\2)\1(\?\H{I}\?,\39\?\H{N}\?,\39\?\H{D}\?,\39\?\H{E}\?,\39\?%
\H{X}\?\2)\?;\6
\\{pid}\?\1(\?7\?\2)\1(\?\H{I}\?,\39\?\H{N}\?,\39\?\H{\_}\?,\39\?\H{F}\?,\39\?%
\H{I}\?,\39\?\H{L}\?,\39\?\H{E}\?\2)\?;\6
\\{pid}\?\1(\?6\?\2)\1(\?\H{I}\?,\39\?\H{N}\?,\39\?\H{L}\?,\39\?\H{I}\?,\39\?%
\H{N}\?,\39\?\H{E}\?\2)\?;\6
\\{pid}\?\1(\?10\?\2)\1(\?\H{I}\?,\39\?\H{N}\?,\39\?\H{O}\?,\39\?\H{U}\?,\39\?%
\H{T}\?,\39\?\H{\_}\?,\39\?\H{F}\?,\39\?\H{I}\?,\39\?\H{L}\?,\39\?\H{E}\?\2)\?;%
\6
\\{pid}\?\1(\?7\?\2)\1(\?\H{I}\?,\39\?\H{N}\?,\39\?\H{T}\?,\39\?\H{E}\?,\39\?%
\H{G}\?,\39\?\H{E}\?,\39\?\H{R}\?\2)\?;\6
\\{pid}\?\1(\?10\?\2)\1(\?\H{I}\?,\39\?\H{N}\?,\39\?\H{T}\?,\39\?\H{E}\?,\39\?%
\H{G}\?,\39\?\H{E}\?,\39\?\H{R}\?,\39\?\H{\_}\?,\39\?\H{I}\?,\39\?\H{O}\?\2)\?;%
\6
\\{pid}\?\1(\?13\?\2)\1(\?\H{I}\?,\39\?\H{O}\?,\39\?\H{\_}\?,\39\?\H{E}\?,\39\?%
\H{X}\?,\39\?\H{C}\?,\39\?\H{E}\?,\39\?\H{P}\?,\39\?\H{T}\?,\39\?\H{I}\?,\39\?%
\H{O}\?,\39\?\H{N}\?,\39\?\H{S}\?\2)\?;\6
\\{pid}\?\1(\?7\?\2)\1(\?\H{I}\?,\39\?\H{S}\?,\39\?\H{\_}\?,\39\?\H{O}\?,\39\?%
\H{P}\?,\39\?\H{E}\?,\39\?\H{N}\?\2)\?;\6
\\{pid}\?\1(\?4\?\2)\1(\?\H{I}\?,\39\?\H{T}\?,\39\?\H{E}\?,\39\?\H{M}\?\2)\?;%
\par
\fi

\M71. \P\?\X69:Store all the predefined identifiers\X\mathrel{+}\S\?\6
\\{pid}\?\1(\?5\?\2)\1(\?\H{L}\?,\39\?\H{A}\?,\39\?\H{R}\?,\39\?\H{G}\?,\39\?%
\H{E}\?\2)\?;\6
\\{pid}\?\1(\?4\?\2)\1(\?\H{L}\?,\39\?\H{A}\?,\39\?\H{S}\?,\39\?\H{T}\?\2)\?;\6
\\{pid}\?\1(\?8\?\2)\1(\?\H{L}\?,\39\?\H{A}\?,\39\?\H{S}\?,\39\?\H{T}\?,\39\?%
\H{\_}\?,\39\?\H{B}\?,\39\?\H{I}\?,\39\?\H{T}\?\2)\?;\6
\\{pid}\?\1(\?12\?\2)\1(\?\H{L}\?,\39\?\H{A}\?,\39\?\H{Y}\?,\39\?\H{O}\?,\39\?%
\H{U}\?,\39\?\H{T}\?,\39\?\H{\_}\?,\39\?\H{E}\?,\39\?\H{R}\?,\39\?\H{R}\?,\39\?%
\H{O}\?,\39\?\H{R}\?\2)\?;\6
\\{pid}\?\1(\?6\?\2)\1(\?\H{L}\?,\39\?\H{E}\?,\39\?\H{N}\?,\39\?\H{G}\?,\39\?%
\H{T}\?,\39\?\H{H}\?\2)\?;\6
\\{pid}\?\1(\?4\?\2)\1(\?\H{L}\?,\39\?\H{I}\?,\39\?\H{N}\?,\39\?\H{E}\?\2)\?;\6
\\{pid}\?\1(\?11\?\2)\1(\?\H{L}\?,\39\?\H{I}\?,\39\?\H{N}\?,\39\?\H{E}\?,\39\?%
\H{\_}\?,\39\?\H{L}\?,\39\?\H{E}\?,\39\?\H{N}\?,\39\?\H{G}\?,\39\?\H{T}\?,\39\?%
\H{H}\?\2)\?;\6
\\{pid}\?\1(\?4\?\2)\1(\?\H{L}\?,\39\?\H{I}\?,\39\?\H{S}\?,\39\?\H{T}\?\2)\?;\6
\\{pid}\?\1(\?10\?\2)\1(\?\H{L}\?,\39\?\H{O}\?,\39\?\H{W}\?,\39\?\H{E}\?,\39\?%
\H{R}\?,\39\?\H{\_}\?,\39\?\H{C}\?,\39\?\H{A}\?,\39\?\H{S}\?,\39\?\H{E}\?\2)\?;%
\6
\\{pid}\?\1(\?12\?\2)\1(\?\H{L}\?,\39\?\H{O}\?,\39\?\H{W}\?,\39\?\H{\_}\?,\39\?%
\H{L}\?,\39\?\H{E}\?,\39\?\H{V}\?,\39\?\H{E}\?,\39\?\H{L}\?,\39\?\H{\_}\?,\39\?%
\H{I}\?,\39\?\H{O}\?\2)\?;\6
\\{pid}\?\1(\?12\?\2)\1(\?\H{M}\?,\39\?\H{A}\?,\39\?\H{C}\?,\39\?\H{H}\?,\39\?%
\H{I}\?,\39\?\H{N}\?,\39\?\H{E}\?,\39\?\H{\_}\?,\39\?\H{E}\?,\39\?\H{M}\?,\39\?%
\H{A}\?,\39\?\H{X}\?\2)\?;\6
\\{pid}\?\1(\?12\?\2)\1(\?\H{M}\?,\39\?\H{A}\?,\39\?\H{C}\?,\39\?\H{H}\?,\39\?%
\H{I}\?,\39\?\H{N}\?,\39\?\H{E}\?,\39\?\H{\_}\?,\39\?\H{E}\?,\39\?\H{M}\?,\39\?%
\H{I}\?,\39\?\H{N}\?\2)\?;\6
\\{pid}\?\1(\?16\?\2)\1(\?\H{M}\?,\39\?\H{A}\?,\39\?\H{C}\?,\39\?\H{H}\?,\39\?%
\H{I}\?,\39\?\H{N}\?,\39\?\H{E}\?,\39\?\H{\_}\?,\39\?\H{M}\?,\39\?\H{A}\?,\39\?%
\H{N}\?,\39\?\H{T}\?,\39\?\H{I}\?,\39\?\H{S}\?,\39\?\H{S}\?,\39\?\H{A}\?\2)\?;\6
\\{pid}\?\1(\?17\?\2)\1(\?\H{M}\?,\39\?\H{A}\?,\39\?\H{C}\?,\39\?\H{H}\?,\39\?%
\H{I}\?,\39\?\H{N}\?,\39\?\H{E}\?,\39\?\H{\_}\?,\39\?\H{O}\?,\39\?\H{V}\?,\39\?%
\H{E}\?,\39\?\H{R}\?,\39\?\H{F}\?,\39\?\H{L}\?,\39\?\H{O}\?,\39\?\H{W}\?,\39\?%
\H{S}\?\2)\?;\6
\\{pid}\?\1(\?13\?\2)\1(\?\H{M}\?,\39\?\H{A}\?,\39\?\H{C}\?,\39\?\H{H}\?,\39\?%
\H{I}\?,\39\?\H{N}\?,\39\?\H{E}\?,\39\?\H{\_}\?,\39\?\H{R}\?,\39\?\H{A}\?,\39\?%
\H{D}\?,\39\?\H{I}\?,\39\?\H{X}\?\2)\?;\6
\\{pid}\?\1(\?14\?\2)\1(\?\H{M}\?,\39\?\H{A}\?,\39\?\H{C}\?,\39\?\H{H}\?,\39\?%
\H{I}\?,\39\?\H{N}\?,\39\?\H{E}\?,\39\?\H{\_}\?,\39\?\H{R}\?,\39\?\H{O}\?,\39\?%
\H{U}\?,\39\?\H{N}\?,\39\?\H{D}\?,\39\?\H{S}\?\2)\?;\6
\\{pid}\?\1(\?8\?\2)\1(\?\H{M}\?,\39\?\H{A}\?,\39\?\H{N}\?,\39\?\H{T}\?,\39\?%
\H{I}\?,\39\?\H{S}\?,\39\?\H{S}\?,\39\?\H{A}\?\2)\?;\6
\\{pid}\?\1(\?10\?\2)\1(\?\H{M}\?,\39\?\H{A}\?,\39\?\H{X}\?,\39\?\H{\_}\?,\39\?%
\H{D}\?,\39\?\H{I}\?,\39\?\H{G}\?,\39\?\H{I}\?,\39\?\H{T}\?,\39\?\H{S}\?\2)\?;\6
\\{pid}\?\1(\?7\?\2)\1(\?\H{M}\?,\39\?\H{A}\?,\39\?\H{X}\?,\39\?\H{\_}\?,\39\?%
\H{I}\?,\39\?\H{N}\?,\39\?\H{T}\?\2)\?;\6
\\{pid}\?\1(\?12\?\2)\1(\?\H{M}\?,\39\?\H{A}\?,\39\?\H{X}\?,\39\?\H{\_}\?,\39\?%
\H{M}\?,\39\?\H{A}\?,\39\?\H{N}\?,\39\?\H{T}\?,\39\?\H{I}\?,\39\?\H{S}\?,\39\?%
\H{S}\?,\39\?\H{A}\?\2)\?;\6
\\{pid}\?\1(\?11\?\2)\1(\?\H{M}\?,\39\?\H{E}\?,\39\?\H{M}\?,\39\?\H{O}\?,\39\?%
\H{R}\?,\39\?\H{Y}\?,\39\?\H{\_}\?,\39\?\H{S}\?,\39\?\H{I}\?,\39\?\H{Z}\?,\39\?%
\H{E}\?\2)\?;\6
\\{pid}\?\1(\?7\?\2)\1(\?\H{M}\?,\39\?\H{I}\?,\39\?\H{N}\?,\39\?\H{\_}\?,\39\?%
\H{I}\?,\39\?\H{N}\?,\39\?\H{T}\?\2)\?;\6
\\{pid}\?\1(\?4\?\2)\1(\?\H{M}\?,\39\?\H{O}\?,\39\?\H{D}\?,\39\?\H{E}\?\2)\?;\6
\\{pid}\?\1(\?10\?\2)\1(\?\H{M}\?,\39\?\H{O}\?,\39\?\H{D}\?,\39\?\H{E}\?,\39\?%
\H{\_}\?,\39\?\H{E}\?,\39\?\H{R}\?,\39\?\H{R}\?,\39\?\H{O}\?,\39\?\H{R}\?\2)\?;%
\6
\\{pid}\?\1(\?5\?\2)\1(\?\H{M}\?,\39\?\H{O}\?,\39\?\H{N}\?,\39\?\H{T}\?,\39\?%
\H{H}\?\2)\?;\6
\\{pid}\?\1(\?12\?\2)\1(\?\H{M}\?,\39\?\H{O}\?,\39\?\H{N}\?,\39\?\H{T}\?,\39\?%
\H{H}\?,\39\?\H{\_}\?,\39\?\H{N}\?,\39\?\H{U}\?,\39\?\H{M}\?,\39\?\H{B}\?,\39\?%
\H{E}\?,\39\?\H{R}\?\2)\?;\6
\\{pid}\?\1(\?4\?\2)\1(\?\H{N}\?,\39\?\H{A}\?,\39\?\H{M}\?,\39\?\H{E}\?\2)\?;\6
\\{pid}\?\1(\?10\?\2)\1(\?\H{N}\?,\39\?\H{A}\?,\39\?\H{M}\?,\39\?\H{E}\?,\39\?%
\H{\_}\?,\39\?\H{E}\?,\39\?\H{R}\?,\39\?\H{R}\?,\39\?\H{O}\?,\39\?\H{R}\?\2)\?;%
\6
\\{pid}\?\1(\?7\?\2)\1(\?\H{N}\?,\39\?\H{A}\?,\39\?\H{T}\?,\39\?\H{U}\?,\39\?%
\H{R}\?,\39\?\H{A}\?,\39\?\H{L}\?\2)\?;\6
\\{pid}\?\1(\?8\?\2)\1(\?\H{N}\?,\39\?\H{E}\?,\39\?\H{W}\?,\39\?\H{\_}\?,\39\?%
\H{L}\?,\39\?\H{I}\?,\39\?\H{N}\?,\39\?\H{E}\?\2)\?;\6
\\{pid}\?\1(\?8\?\2)\1(\?\H{N}\?,\39\?\H{E}\?,\39\?\H{W}\?,\39\?\H{\_}\?,\39\?%
\H{P}\?,\39\?\H{A}\?,\39\?\H{G}\?,\39\?\H{E}\?\2)\?;\6
\\{pid}\?\1(\?3\?\2)\1(\?\H{N}\?,\39\?\H{U}\?,\39\?\H{M}\?\2)\?;\6
\\{pid}\?\1(\?11\?\2)\1(\?\H{N}\?,\39\?\H{U}\?,\39\?\H{M}\?,\39\?\H{B}\?,\39\?%
\H{E}\?,\39\?\H{R}\?,\39\?\H{\_}\?,\39\?\H{B}\?,\39\?\H{A}\?,\39\?\H{S}\?,\39\?%
\H{E}\?\2)\?;\6
\\{pid}\?\1(\?13\?\2)\1(\?\H{N}\?,\39\?\H{U}\?,\39\?\H{M}\?,\39\?\H{E}\?,\39\?%
\H{R}\?,\39\?\H{I}\?,\39\?\H{C}\?,\39\?\H{\_}\?,\39\?\H{E}\?,\39\?\H{R}\?,\39\?%
\H{R}\?,\39\?\H{O}\?,\39\?\H{R}\?\2)\?;\6
\\{pid}\?\1(\?3\?\2)\1(\?\H{O}\?,\39\?\H{F}\?,\39\?\H{F}\?\2)\?;\6
\\{pid}\?\1(\?2\?\2)\1(\?\H{O}\?,\39\?\H{N}\?\2)\?;\6
\\{pid}\?\1(\?4\?\2)\1(\?\H{O}\?,\39\?\H{P}\?,\39\?\H{E}\?,\39\?\H{N}\?\2)\?;\6
\\{pid}\?\1(\?8\?\2)\1(\?\H{O}\?,\39\?\H{P}\?,\39\?\H{T}\?,\39\?\H{I}\?,\39\?%
\H{M}\?,\39\?\H{I}\?,\39\?\H{Z}\?,\39\?\H{E}\?\2)\?;\6
\\{pid}\?\1(\?8\?\2)\1(\?\H{O}\?,\39\?\H{U}\?,\39\?\H{T}\?,\39\?\H{\_}\?,\39\?%
\H{F}\?,\39\?\H{I}\?,\39\?\H{L}\?,\39\?\H{E}\?\2)\?;\6
\\{pid}\?\1(\?4\?\2)\1(\?\H{P}\?,\39\?\H{A}\?,\39\?\H{C}\?,\39\?\H{K}\?\2)\?;\6
\\{pid}\?\1(\?4\?\2)\1(\?\H{P}\?,\39\?\H{A}\?,\39\?\H{G}\?,\39\?\H{E}\?\2)\?;\6
\\{pid}\?\1(\?11\?\2)\1(\?\H{P}\?,\39\?\H{A}\?,\39\?\H{G}\?,\39\?\H{E}\?,\39\?%
\H{\_}\?,\39\?\H{L}\?,\39\?\H{E}\?,\39\?\H{N}\?,\39\?\H{G}\?,\39\?\H{T}\?,\39\?%
\H{H}\?\2)\?;\6
\\{pid}\?\1(\?3\?\2)\1(\?\H{P}\?,\39\?\H{O}\?,\39\?\H{S}\?\2)\?;\6
\\{pid}\?\1(\?8\?\2)\1(\?\H{P}\?,\39\?\H{O}\?,\39\?\H{S}\?,\39\?\H{I}\?,\39\?%
\H{T}\?,\39\?\H{I}\?,\39\?\H{O}\?,\39\?\H{N}\?\2)\?;\6
\\{pid}\?\1(\?8\?\2)\1(\?\H{P}\?,\39\?\H{O}\?,\39\?\H{S}\?,\39\?\H{I}\?,\39\?%
\H{T}\?,\39\?\H{I}\?,\39\?\H{V}\?,\39\?\H{E}\?\2)\?;\6
\\{pid}\?\1(\?14\?\2)\1(\?\H{P}\?,\39\?\H{O}\?,\39\?\H{S}\?,\39\?\H{I}\?,\39\?%
\H{T}\?,\39\?\H{I}\?,\39\?\H{V}\?,\39\?\H{E}\?,\39\?\H{\_}\?,\39\?\H{C}\?,\39\?%
\H{O}\?,\39\?\H{U}\?,\39\?\H{N}\?,\39\?\H{T}\?\2)\?;\6
\\{pid}\?\1(\?4\?\2)\1(\?\H{P}\?,\39\?\H{R}\?,\39\?\H{E}\?,\39\?\H{D}\?\2)\?;\6
\\{pid}\?\1(\?8\?\2)\1(\?\H{P}\?,\39\?\H{R}\?,\39\?\H{I}\?,\39\?\H{O}\?,\39\?%
\H{R}\?,\39\?\H{I}\?,\39\?\H{T}\?,\39\?\H{Y}\?\2)\?;\6
\\{pid}\?\1(\?13\?\2)\1(\?\H{P}\?,\39\?\H{R}\?,\39\?\H{O}\?,\39\?\H{G}\?,\39\?%
\H{R}\?,\39\?\H{A}\?,\39\?\H{M}\?,\39\?\H{\_}\?,\39\?\H{E}\?,\39\?\H{R}\?,\39\?%
\H{R}\?,\39\?\H{O}\?,\39\?\H{R}\?\2)\?;\6
\\{pid}\?\1(\?3\?\2)\1(\?\H{P}\?,\39\?\H{U}\?,\39\?\H{T}\?\2)\?;\6
\\{pid}\?\1(\?8\?\2)\1(\?\H{P}\?,\39\?\H{U}\?,\39\?\H{T}\?,\39\?\H{\_}\?,\39\?%
\H{L}\?,\39\?\H{I}\?,\39\?\H{N}\?,\39\?\H{E}\?\2)\?;\6
\\{pid}\?\1(\?5\?\2)\1(\?\H{R}\?,\39\?\H{A}\?,\39\?\H{N}\?,\39\?\H{G}\?,\39\?%
\H{E}\?\2)\?;\6
\\{pid}\?\1(\?4\?\2)\1(\?\H{R}\?,\39\?\H{E}\?,\39\?\H{A}\?,\39\?\H{D}\?\2)\?;\6
\\{pid}\?\1(\?15\?\2)\1(\?\H{R}\?,\39\?\H{E}\?,\39\?\H{C}\?,\39\?\H{E}\?,\39\?%
\H{I}\?,\39\?\H{V}\?,\39\?\H{E}\?,\39\?\H{\_}\?,\39\?\H{C}\?,\39\?\H{O}\?,\39\?%
\H{N}\?,\39\?\H{T}\?,\39\?\H{R}\?,\39\?\H{O}\?,\39\?\H{L}\?\2)\?;\6
\\{pid}\?\1(\?5\?\2)\1(\?\H{R}\?,\39\?\H{E}\?,\39\?\H{S}\?,\39\?\H{E}\?,\39\?%
\H{T}\?\2)\?;\par
\fi

\M72. \P\?\X69:Store all the predefined identifiers\X\mathrel{+}\S\?\6
\\{pid}\?\1(\?9\?\2)\1(\?\H{S}\?,\39\?\H{A}\?,\39\?\H{F}\?,\39\?\H{E}\?,\39\?%
\H{\_}\?,\39\?\H{E}\?,\39\?\H{M}\?,\39\?\H{A}\?,\39\?\H{X}\?\2)\?;\6
\\{pid}\?\1(\?10\?\2)\1(\?\H{S}\?,\39\?\H{A}\?,\39\?\H{F}\?,\39\?\H{E}\?,\39\?%
\H{\_}\?,\39\?\H{L}\?,\39\?\H{A}\?,\39\?\H{R}\?,\39\?\H{G}\?,\39\?\H{E}\?\2)\?;%
\6
\\{pid}\?\1(\?10\?\2)\1(\?\H{S}\?,\39\?\H{A}\?,\39\?\H{F}\?,\39\?\H{E}\?,\39\?%
\H{\_}\?,\39\?\H{S}\?,\39\?\H{M}\?,\39\?\H{A}\?,\39\?\H{L}\?,\39\?\H{L}\?\2)\?;%
\6
\\{pid}\?\1(\?7\?\2)\1(\?\H{S}\?,\39\?\H{E}\?,\39\?\H{C}\?,\39\?\H{O}\?,\39\?%
\H{N}\?,\39\?\H{D}\?,\39\?\H{S}\?\2)\?;\6
\\{pid}\?\1(\?12\?\2)\1(\?\H{S}\?,\39\?\H{E}\?,\39\?\H{N}\?,\39\?\H{D}\?,\39\?%
\H{\_}\?,\39\?\H{C}\?,\39\?\H{O}\?,\39\?\H{N}\?,\39\?\H{T}\?,\39\?\H{R}\?,\39\?%
\H{O}\?,\39\?\H{L}\?\2)\?;\6
\\{pid}\?\1(\?13\?\2)\1(\?\H{S}\?,\39\?\H{E}\?,\39\?\H{Q}\?,\39\?\H{U}\?,\39\?%
\H{E}\?,\39\?\H{N}\?,\39\?\H{T}\?,\39\?\H{I}\?,\39\?\H{A}\?,\39\?\H{L}\?,\39\?%
\H{\_}\?,\39\?\H{I}\?,\39\?\H{O}\?\2)\?;\6
\\{pid}\?\1(\?3\?\2)\1(\?\H{S}\?,\39\?\H{E}\?,\39\?\H{T}\?\2)\?;\6
\\{pid}\?\1(\?7\?\2)\1(\?\H{S}\?,\39\?\H{E}\?,\39\?\H{T}\?,\39\?\H{\_}\?,\39\?%
\H{C}\?,\39\?\H{O}\?,\39\?\H{L}\?\2)\?;\6
\\{pid}\?\1(\?9\?\2)\1(\?\H{S}\?,\39\?\H{E}\?,\39\?\H{T}\?,\39\?\H{\_}\?,\39\?%
\H{I}\?,\39\?\H{N}\?,\39\?\H{D}\?,\39\?\H{E}\?,\39\?\H{X}\?\2)\?;\6
\\{pid}\?\1(\?9\?\2)\1(\?\H{S}\?,\39\?\H{E}\?,\39\?\H{T}\?,\39\?\H{\_}\?,\39\?%
\H{I}\?,\39\?\H{N}\?,\39\?\H{P}\?,\39\?\H{U}\?,\39\?\H{T}\?\2)\?;\6
\\{pid}\?\1(\?8\?\2)\1(\?\H{S}\?,\39\?\H{E}\?,\39\?\H{T}\?,\39\?\H{\_}\?,\39\?%
\H{L}\?,\39\?\H{I}\?,\39\?\H{N}\?,\39\?\H{E}\?\2)\?;\6
\\{pid}\?\1(\?15\?\2)\1(\?\H{S}\?,\39\?\H{E}\?,\39\?\H{T}\?,\39\?\H{\_}\?,\39\?%
\H{L}\?,\39\?\H{I}\?,\39\?\H{N}\?,\39\?\H{E}\?,\39\?\H{\_}\?,\39\?\H{L}\?,\39\?%
\H{E}\?,\39\?\H{N}\?,\39\?\H{G}\?,\39\?\H{T}\?,\39\?\H{H}\?\2)\?;\6
\\{pid}\?\1(\?10\?\2)\1(\?\H{S}\?,\39\?\H{E}\?,\39\?\H{T}\?,\39\?\H{\_}\?,\39\?%
\H{O}\?,\39\?\H{U}\?,\39\?\H{T}\?,\39\?\H{P}\?,\39\?\H{U}\?,\39\?\H{T}\?\2)\?;\6
\\{pid}\?\1(\?15\?\2)\1(\?\H{S}\?,\39\?\H{E}\?,\39\?\H{T}\?,\39\?\H{\_}\?,\39\?%
\H{P}\?,\39\?\H{A}\?,\39\?\H{G}\?,\39\?\H{E}\?,\39\?\H{\_}\?,\39\?\H{L}\?,\39\?%
\H{E}\?,\39\?\H{N}\?,\39\?\H{G}\?,\39\?\H{T}\?,\39\?\H{H}\?\2)\?;\6
\\{pid}\?\1(\?6\?\2)\1(\?\H{S}\?,\39\?\H{H}\?,\39\?\H{A}\?,\39\?\H{R}\?,\39\?%
\H{E}\?,\39\?\H{D}\?\2)\?;\6
\\{pid}\?\1(\?4\?\2)\1(\?\H{S}\?,\39\?\H{I}\?,\39\?\H{Z}\?,\39\?\H{E}\?\2)\?;\6
\\{pid}\?\1(\?9\?\2)\1(\?\H{S}\?,\39\?\H{K}\?,\39\?\H{I}\?,\39\?\H{P}\?,\39\?%
\H{\_}\?,\39\?\H{L}\?,\39\?\H{I}\?,\39\?\H{N}\?,\39\?\H{E}\?\2)\?;\6
\\{pid}\?\1(\?9\?\2)\1(\?\H{S}\?,\39\?\H{K}\?,\39\?\H{I}\?,\39\?\H{P}\?,\39\?%
\H{\_}\?,\39\?\H{P}\?,\39\?\H{A}\?,\39\?\H{G}\?,\39\?\H{E}\?\2)\?;\6
\\{pid}\?\1(\?5\?\2)\1(\?\H{S}\?,\39\?\H{M}\?,\39\?\H{A}\?,\39\?\H{L}\?,\39\?%
\H{L}\?\2)\?;\6
\\{pid}\?\1(\?5\?\2)\1(\?\H{S}\?,\39\?\H{P}\?,\39\?\H{A}\?,\39\?\H{C}\?,\39\?%
\H{E}\?\2)\?;\6
\\{pid}\?\1(\?7\?\2)\1(\?\H{S}\?,\39\?\H{P}\?,\39\?\H{A}\?,\39\?\H{C}\?,\39\?%
\H{I}\?,\39\?\H{N}\?,\39\?\H{G}\?\2)\?;\6
\\{pid}\?\1(\?5\?\2)\1(\?\H{S}\?,\39\?\H{P}\?,\39\?\H{L}\?,\39\?\H{I}\?,\39\?%
\H{T}\?\2)\?;\6
\\{pid}\?\1(\?8\?\2)\1(\?\H{S}\?,\39\?\H{T}\?,\39\?\H{A}\?,\39\?\H{N}\?,\39\?%
\H{D}\?,\39\?\H{A}\?,\39\?\H{R}\?,\39\?\H{D}\?\2)\?;\6
\\{pid}\?\1(\?14\?\2)\1(\?\H{S}\?,\39\?\H{T}\?,\39\?\H{A}\?,\39\?\H{N}\?,\39\?%
\H{D}\?,\39\?\H{A}\?,\39\?\H{R}\?,\39\?\H{D}\?,\39\?\H{\_}\?,\39\?\H{I}\?,\39\?%
\H{N}\?,\39\?\H{P}\?,\39\?\H{U}\?,\39\?\H{T}\?\2)\?;\6
\\{pid}\?\1(\?15\?\2)\1(\?\H{S}\?,\39\?\H{T}\?,\39\?\H{A}\?,\39\?\H{N}\?,\39\?%
\H{D}\?,\39\?\H{A}\?,\39\?\H{R}\?,\39\?\H{D}\?,\39\?\H{\_}\?,\39\?\H{O}\?,\39\?%
\H{U}\?,\39\?\H{T}\?,\39\?\H{P}\?,\39\?\H{U}\?,\39\?\H{T}\?\2)\?;\6
\\{pid}\?\1(\?12\?\2)\1(\?\H{S}\?,\39\?\H{T}\?,\39\?\H{A}\?,\39\?\H{T}\?,\39\?%
\H{U}\?,\39\?\H{S}\?,\39\?\H{\_}\?,\39\?\H{E}\?,\39\?\H{R}\?,\39\?\H{R}\?,\39\?%
\H{O}\?,\39\?\H{R}\?\2)\?;\6
\\{pid}\?\1(\?13\?\2)\1(\?\H{S}\?,\39\?\H{T}\?,\39\?\H{O}\?,\39\?\H{R}\?,\39\?%
\H{A}\?,\39\?\H{G}\?,\39\?\H{E}\?,\39\?\H{\_}\?,\39\?\H{E}\?,\39\?\H{R}\?,\39\?%
\H{R}\?,\39\?\H{O}\?,\39\?\H{R}\?\2)\?;\6
\\{pid}\?\1(\?12\?\2)\1(\?\H{S}\?,\39\?\H{T}\?,\39\?\H{O}\?,\39\?\H{R}\?,\39\?%
\H{A}\?,\39\?\H{G}\?,\39\?\H{E}\?,\39\?\H{\_}\?,\39\?\H{S}\?,\39\?\H{I}\?,\39\?%
\H{Z}\?,\39\?\H{E}\?\2)\?;\6
\\{pid}\?\1(\?12\?\2)\1(\?\H{S}\?,\39\?\H{T}\?,\39\?\H{O}\?,\39\?\H{R}\?,\39\?%
\H{A}\?,\39\?\H{G}\?,\39\?\H{E}\?,\39\?\H{\_}\?,\39\?\H{U}\?,\39\?\H{N}\?,\39\?%
\H{I}\?,\39\?\H{T}\?\2)\?;\6
\\{pid}\?\1(\?6\?\2)\1(\?\H{S}\?,\39\?\H{T}\?,\39\?\H{R}\?,\39\?\H{I}\?,\39\?%
\H{N}\?,\39\?\H{G}\?\2)\?;\6
\\{pid}\?\1(\?4\?\2)\1(\?\H{S}\?,\39\?\H{U}\?,\39\?\H{C}\?,\39\?\H{C}\?\2)\?;\6
\\{pid}\?\1(\?8\?\2)\1(\?\H{S}\?,\39\?\H{U}\?,\39\?\H{P}\?,\39\?\H{P}\?,\39\?%
\H{R}\?,\39\?\H{E}\?,\39\?\H{S}\?,\39\?\H{S}\?\2)\?;\6
\\{pid}\?\1(\?6\?\2)\1(\?\H{S}\?,\39\?\H{Y}\?,\39\?\H{S}\?,\39\?\H{T}\?,\39\?%
\H{E}\?,\39\?\H{M}\?\2)\?;\6
\\{pid}\?\1(\?11\?\2)\1(\?\H{S}\?,\39\?\H{Y}\?,\39\?\H{S}\?,\39\?\H{T}\?,\39\?%
\H{E}\?,\39\?\H{M}\?,\39\?\H{\_}\?,\39\?\H{N}\?,\39\?\H{A}\?,\39\?\H{M}\?,\39\?%
\H{E}\?\2)\?;\6
\\{pid}\?\1(\?13\?\2)\1(\?\H{T}\?,\39\?\H{A}\?,\39\?\H{S}\?,\39\?\H{K}\?,\39\?%
\H{I}\?,\39\?\H{N}\?,\39\?\H{G}\?,\39\?\H{\_}\?,\39\?\H{E}\?,\39\?\H{R}\?,\39\?%
\H{R}\?,\39\?\H{O}\?,\39\?\H{R}\?\2)\?;\6
\\{pid}\?\1(\?10\?\2)\1(\?\H{T}\?,\39\?\H{E}\?,\39\?\H{R}\?,\39\?\H{M}\?,\39\?%
\H{I}\?,\39\?\H{N}\?,\39\?\H{A}\?,\39\?\H{T}\?,\39\?\H{E}\?,\39\?\H{D}\?\2)\?;\6
\\{pid}\?\1(\?7\?\2)\1(\?\H{T}\?,\39\?\H{E}\?,\39\?\H{X}\?,\39\?\H{T}\?,\39\?%
\H{\_}\?,\39\?\H{I}\?,\39\?\H{O}\?\2)\?;\6
\\{pid}\?\1(\?4\?\2)\1(\?\H{T}\?,\39\?\H{I}\?,\39\?\H{C}\?,\39\?\H{K}\?\2)\?;\6
\\{pid}\?\1(\?4\?\2)\1(\?\H{T}\?,\39\?\H{I}\?,\39\?\H{M}\?,\39\?\H{E}\?\2)\?;\6
\\{pid}\?\1(\?10\?\2)\1(\?\H{T}\?,\39\?\H{I}\?,\39\?\H{M}\?,\39\?\H{E}\?,\39\?%
\H{\_}\?,\39\?\H{E}\?,\39\?\H{R}\?,\39\?\H{R}\?,\39\?\H{O}\?,\39\?\H{R}\?\2)\?;%
\6
\\{pid}\?\1(\?7\?\2)\1(\?\H{T}\?,\39\?\H{I}\?,\39\?\H{M}\?,\39\?\H{E}\?,\39\?%
\H{\_}\?,\39\?\H{O}\?,\39\?\H{F}\?\2)\?;\6
\\{pid}\?\1(\?2\?\2)\1(\?\H{T}\?,\39\?\H{O}\?\2)\?;\6
\\{pid}\?\1(\?4\?\2)\1(\?\H{T}\?,\39\?\H{R}\?,\39\?\H{U}\?,\39\?\H{E}\?\2)\?;\6
\\{pid}\?\1(\?8\?\2)\1(\?\H{T}\?,\39\?\H{Y}\?,\39\?\H{P}\?,\39\?\H{E}\?,\39\?%
\H{\_}\?,\39\?\H{S}\?,\39\?\H{E}\?,\39\?\H{T}\?\2)\?;\6
\\{pid}\?\1(\?9\?\2)\1(\?\H{U}\?,\39\?\H{N}\?,\39\?\H{B}\?,\39\?\H{O}\?,\39\?%
\H{U}\?,\39\?\H{N}\?,\39\?\H{D}\?,\39\?\H{E}\?,\39\?\H{D}\?\2)\?;\6
\\{pid}\?\1(\?20\?\2)\1(\?\H{U}\?,\39\?\H{N}\?,\39\?\H{C}\?,\39\?\H{H}\?,\39\?%
\H{E}\?,\39\?\H{C}\?,\39\?\H{K}\?,\39\?\H{E}\?,\39\?\H{D}\?,\39\?\H{\_}\?,\39\?%
\H{C}\?,\39\?\H{O}\?,\39\?\H{N}\?,\39\?\H{V}\?,\39\?\H{E}\?,\39\?\H{R}\?,\39\?%
\H{S}\?,\39\?\H{I}\?,\39\?\H{O}\?,\39\?\H{N}\?\2)\?;\6
\\{pid}\?\1(\?22\?\2)\1(\?\H{U}\?,\39\?\H{N}\?,\39\?\H{C}\?,\39\?\H{H}\?,\39\?%
\H{E}\?,\39\?\H{C}\?,\39\?\H{K}\?,\39\?\H{E}\?,\39\?\H{D}\?,\39\?\H{\_}\?,\39\?%
\H{D}\?,\39\?\H{E}\?,\39\?\H{A}\?,\39\?\H{L}\?,\39\?\H{L}\?,\39\?\H{O}\?,\39\?%
\H{C}\?,\39\?\H{A}\?,\39\?\H{T}\?,\39\?\H{I}\?,\39\?\H{O}\?,\39\?\H{N}\?\2)\?;\6
\\{pid}\?\1(\?10\?\2)\1(\?\H{U}\?,\39\?\H{P}\?,\39\?\H{P}\?,\39\?\H{E}\?,\39\?%
\H{R}\?,\39\?\H{\_}\?,\39\?\H{C}\?,\39\?\H{A}\?,\39\?\H{S}\?,\39\?\H{E}\?\2)\?;%
\6
\\{pid}\?\1(\?9\?\2)\1(\?\H{U}\?,\39\?\H{S}\?,\39\?\H{E}\?,\39\?\H{\_}\?,\39\?%
\H{E}\?,\39\?\H{R}\?,\39\?\H{R}\?,\39\?\H{O}\?,\39\?\H{R}\?\2)\?;\6
\\{pid}\?\1(\?3\?\2)\1(\?\H{V}\?,\39\?\H{A}\?,\39\?\H{L}\?\2)\?;\6
\\{pid}\?\1(\?5\?\2)\1(\?\H{V}\?,\39\?\H{A}\?,\39\?\H{L}\?,\39\?\H{U}\?,\39\?%
\H{E}\?\2)\?;\6
\\{pid}\?\1(\?5\?\2)\1(\?\H{W}\?,\39\?\H{I}\?,\39\?\H{D}\?,\39\?\H{T}\?,\39\?%
\H{H}\?\2)\?;\6
\\{pid}\?\1(\?5\?\2)\1(\?\H{W}\?,\39\?\H{R}\?,\39\?\H{I}\?,\39\?\H{T}\?,\39\?%
\H{E}\?\2)\?;\6
\\{pid}\?\1(\?4\?\2)\1(\?\H{Y}\?,\39\?\H{E}\?,\39\?\H{A}\?,\39\?\H{R}\?\2)\?;\6
\\{pid}\?\1(\?11\?\2)\1(\?\H{Y}\?,\39\?\H{E}\?,\39\?\H{A}\?,\39\?\H{R}\?,\39\?%
\H{\_}\?,\39\?\H{N}\?,\39\?\H{U}\?,\39\?\H{M}\?,\39\?\H{B}\?,\39\?\H{E}\?,\39\?%
\H{R}\?\2)\?\par
\fi

\N73.  Searching for module names.
The \\{mod\_lookup} procedure finds the module name \\{mod\_text}\?(\?1\?\to\?%
\|l\?)\? in the
search tree, after inserting it if necessary, and returns a pointer to
where it was found.

\Y\P\?\4\X14:Types and objects visible from \\{data\_structures}\X\mathrel{+}\S%
\?\6
\&{subtype}\8\\{long\_name\_index}\8\&{is}\8\p{INTEGER}\8\&{range}\80\?\to\?%
\\{longest\_name};\6
\\{mod\_text}\?:\1\&{array}\8\1(\?\\{long\_name\_index}\?\2)\2\8\&{of}\?\8%
\\{ASCII\_code};\C{name being sought for}\par
\fi

\M74. According to the rules of \.{AWEB}, no module name
should be a proper prefix of another, so a ``clean'' comparison should
occur between any two names. The result of \\{mod\_lookup} is 0 if this
prefix condition is violated. An error message is printed when such violations
are detected during phase two of \.{AWEAVE}.

\Y\P\?\4\X14:Types and objects visible from \\{data\_structures}\X\mathrel{+}\S%
\?\6
\&{type}\8\\{compare}\8\&{is}\?\8\1(\?\\{less}\?,\39\?\\{equal}\?,\39\?%
\\{greater}\?,\39\?\\{prefix}\?,\39\?\\{extension}\?\2)\?;\par
\fi

\M75. \P\?\X60:Specification of procedures visible from \\{searching}\X%
\mathrel{+}\S\?\6
\&{function}\8\1\\{mod\_lookup}\?(\1\?\|l\?:\?\\{sixteen\_bits}\?\2)\&{return}%
\?\8\\{name\_pointer}\2;\par
\fi

\M76. \P\?\X61:Procedures in \\{searching}\X\mathrel{+}\S\?\6
\&{function}\8\1\\{mod\_lookup}\?(\1\?\|l\?:\?\\{sixteen\_bits}\?\2)\&{return}%
\?\8\\{name\_pointer}\2\8\&{is}\8\C{finds module name}\1\6
\|c\?:\?\\{compare}\?\K\?\\{greater};\C{comparison between two names}\6
\|j\?:\?\\{long\_name\_index};\C{index into \\{mod\_text}}\6
\|k\?:\?\\{byte\_index};\C{index into \\{byte\_mem}}\6
\|w\?:\?\\{w\_range};\C{segment of \\{byte\_mem}}\6
\|p\?:\?\\{name\_pointer}\?\K\?\\{root};\C{current node of the search tree}\6
\|q\?:\?\\{name\_pointer}\?\K\?0;\C{father of node \|p}\6
\4\&{begin}\6
\1\&{while}\8\|p\?\I\?0\8\&{loop}\6
\X78:Set variable \|c to the result of comparing the given name to name \|p\X;\6
\|q\?\K\?\|p;\6
\&{if}\1\1\8\|c\?=\?\\{less}\2\8\&{then}\6
\|p\?\K\?\\{llink}\?\1(\?\|q\?\2)\?;\6
\4\&{elsif}\1\8\|c\?=\?\\{greater}\2\8\&{then}\6
\|p\?\K\?\\{rlink}\?\1(\?\|q\?\2)\?;\6
\4\&{else}\6
\&{goto}\8\\{found};\2\6
\&{end}\8\2\&{if}\1;\2\6
\&{end}\8\&{loop};\6
\X77:Enter a new module name into the tree\X;\6
\4\BL\\{found}\EL\8\&{if}\1\1\8\|c\?\I\?\\{equal}\2\8\&{then}\6
\\{err\_print}\?\1(\?\.{"!\ Incompatible\ section\ names"}\?\2)\?;\5
\|p\?\K\?0;\2\6
\&{end}\8\2\&{if}\1;\6
\?\&{return}\?\8\|p;\2\6
\&{end}\8\\{mod\_lookup};\par
\fi

\M77. \P\?\X77:Enter a new module name into the tree\X\S\?\6
\|w\?\K\?\\{name\_ptr}\?\mathbin{\&{mod}}\?\\{ww};\5
\|k\?\K\?\\{byte\_ptr}\?\1(\?\|w\?\2)\?;\6
\&{if}\1\1\8\|k\?+\?\|l\?>\?\\{max\_bytes}\2\8\&{then}\6
\\{overflow}\?\1(\?\.{"byte\ memory"}\?\2)\?;\2\6
\&{end}\8\2\&{if}\1;\6
\&{if}\1\1\8\\{name\_ptr}\?>\?\\{max\_names}\?-\?\\{ww}\2\8\&{then}\6
\\{overflow}\?\1(\?\.{"name"}\?\2)\?;\2\6
\&{end}\8\2\&{if}\1;\6
\|p\?\K\?\\{name\_ptr};\6
\&{if}\1\1\8\|c\?=\?\\{less}\2\8\&{then}\6
\\{llink}\?\1(\?\|q\?\2)\K\?\|p;\6
\4\&{else}\6
\\{rlink}\?\1(\?\|q\?\2)\K\?\|p;\2\6
\&{end}\8\2\&{if}\1;\6
\\{llink}\?\1(\?\|p\?\2)\K\?0;\5
\\{rlink}\?\1(\?\|p\?\2)\K\?0;\5
\|c\?\K\?\\{equal};\6
\1\&{for}\8\|m\?\in\?1\?\to\?\|l\8\&{loop}\6
\\{byte\_mem}\?\1(\?\|w\?,\39\?\|k\?+\?\|m\?-\?1\?\2)\K\?\\{mod\_text}\?\1(\?%
\|m\?\2)\?;\2\6
\&{end}\8\&{loop};\6
\\{byte\_ptr}\?\1(\?\|w\?\2)\K\?\|k\?+\?\|l;\5
\\{byte\_start}\?\1(\?\\{name\_ptr}\?+\?\\{ww}\?\2)\K\?\|k\?+\?\|l;\5
\\{incr}\?\1(\?\\{name\_ptr}\?\2)\?\par
\U section~76.\fi

\M78. \P\?\X78:Set variable \|c to the result of comparing the given name to
name \|p\X\S\?\6
\|k\?\K\?\\{byte\_start}\?\1(\?\|p\?\2)\?;\5
\|w\?\K\?\|p\?\mathbin{\&{mod}}\?\\{ww};\5
\|c\?\K\?\\{equal};\5
\|j\?\K\?1;\6
\1\&{while}\?\8\1(\?\|k\?<\?\\{byte\_start}\?\1(\?\|p\?+\?\\{ww}\?\2)\2)\W\1(\?%
\|j\?\L\?\|l\?\2)\W\1(\?\\{mod\_text}\?\1(\?\|j\?\2)=\?\\{byte\_mem}\?\1(\?\|w%
\?,\39\?\|k\?\2)\2)\?\8\&{loop}\6
\\{incr}\?\1(\?\|k\?\2)\?;\5
\\{incr}\?\1(\?\|j\?\2)\?;\2\6
\&{end}\8\&{loop};\6
\&{if}\1\1\8\|k\?=\?\\{byte\_start}\?\1(\?\|p\?+\?\\{ww}\?\2)\?\2\8\&{then}\6
\&{if}\1\1\8\|j\?>\?\|l\2\8\&{then}\6
\|c\?\K\?\\{equal};\6
\4\&{else}\6
\|c\?\K\?\\{extension};\2\6
\&{end}\8\2\&{if}\1;\6
\4\&{elsif}\1\8\|j\?>\?\|l\2\8\&{then}\6
\|c\?\K\?\\{prefix};\6
\4\&{elsif}\1\8\\{mod\_text}\?\1(\?\|j\?\2)<\?\\{byte\_mem}\?\1(\?\|w\?,\39\?%
\|k\?\2)\?\2\8\&{then}\6
\|c\?\K\?\\{less};\6
\4\&{else}\6
\|c\?\K\?\\{greater};\2\6
\&{end}\8\2\&{if}\1\par
\U sections~76 and~80.\fi

\M79. The \\{prefix\_lookup} procedure is supposed to find exactly one module
name that has \\{mod\_text}\?(\?1\?\to\?\|l\?)\? as a prefix. Actually the
algorithm silently
accepts also the situation that some module name is a prefix of
\\{mod\_text}\?(\?1\?\to\?\|l\?)\?, because the user who painstakingly typed in
more than
necessary probably doesn't want to be told about the wasted effort.

\Y\P\?\4\X60:Specification of procedures visible from \\{searching}\X%
\mathrel{+}\S\?\6
\&{function}\8\1\\{prefix\_lookup}\?(\1\?\|l\?:\?\\{sixteen\_bits}\?\2)%
\&{return}\?\8\\{name\_pointer}\2;\par
\fi

\M80. \P\?\X61:Procedures in \\{searching}\X\mathrel{+}\S\?\6
\&{function}\8\1\\{prefix\_lookup}\?(\1\?\|l\?:\?\\{sixteen\_bits}\?\2)%
\&{return}\?\8\\{name\_pointer}\2\8\&{is}\8\C{finds name extension}\1\6
\|c\?:\?\\{compare};\C{comparison between two names}\6
\\{count}\?:\?\\{name\_index}\?\K\?0;\C{the number of hits}\6
\|j\?:\?\\{long\_name\_index};\C{index into \\{mod\_text}}\6
\|k\?:\?\\{byte\_index};\C{index into \\{byte\_mem}}\6
\|w\?:\?\\{w\_range};\C{segment of \\{byte\_mem}}\6
\|p\?:\?\\{name\_pointer}\?\K\?\\{root};\C{current node of the search tree}\6
\|q\?:\?\\{name\_pointer}\?\K\?0;\C{another place to resume the search after
one branch is done}\6
\|r\?:\?\\{name\_pointer}\?\K\?0;\C{extension found}\6
\4\&{begin}\6
\1\&{while}\8\|p\?\I\?0\8\&{loop}\6
\X78:Set variable \|c to the result of comparing the given name to name \|p\X;\6
\&{if}\1\1\8\|c\?=\?\\{less}\2\8\&{then}\6
\|p\?\K\?\\{llink}\?\1(\?\|p\?\2)\?;\6
\4\&{elsif}\1\8\|c\?=\?\\{greater}\2\8\&{then}\6
\|p\?\K\?\\{rlink}\?\1(\?\|p\?\2)\?;\6
\4\&{else}\6
\|r\?\K\?\|p;\5
\\{incr}\?\1(\?\\{count}\?\2)\?;\5
\|q\?\K\?\\{rlink}\?\1(\?\|p\?\2)\?;\5
\|p\?\K\?\\{llink}\?\1(\?\|p\?\2)\?;\2\6
\&{end}\8\2\&{if}\1;\6
\&{if}\1\1\8\|p\?=\?0\2\8\&{then}\6
\|p\?\K\?\|q;\5
\|q\?\K\?0;\2\6
\&{end}\8\2\&{if}\1;\2\6
\&{end}\8\&{loop};\6
\&{if}\1\1\8\\{count}\?\I\?1\2\8\&{then}\6
\&{if}\1\1\8\\{count}\?=\?0\2\8\&{then}\6
\\{err\_print}\?\1(\?\.{"!\ Name\ does\ not\ match"}\?\2)\?;\6
\4\&{else}\6
\\{err\_print}\?\1(\?\.{"!\ Ambiguous\ prefix"}\?\2)\?;\2\6
\&{end}\8\2\&{if}\1;\2\6
\&{end}\8\2\&{if}\1;\6
\?\&{return}\?\8\|r;\C{the result will be 0 if there was no match}\2\6
\&{end}\8\\{prefix\_lookup};\par
\fi

\N81.  Lexical scanning.
Let us now consider the subroutines that read the \.{AWEB} source file
and break it into meaningful units. There are four such procedures:
One simply skips to the next `\.{@\ }' or `\.{@*}' that begins a
module; another passes over the \TeX\ text at the beginning of a
module; the third passes over the \TeX\ text in a \ADA\ comment;
and the last, which is the most interesting, gets the next token of
a \ADA\ text.

\fi

\M82. But first we need to consider the low-level routine \\{get\_line}
that takes care of merging \\{change\_file} into \\{aweb\_file}. The \\{get%
\_line}
procedure also updates the line numbers for error messages.

\Y\P\?\4\X14:Types and objects visible from \\{data\_structures}\X\mathrel{+}\S%
\?\6
\\{line}\?:\?\p{INTEGER};\C{the number of the current line in the current file}%
\6
\\{other\_line}\?:\?\p{INTEGER};\C{the number of the current line in the input
file that   is not currently being read}\6
\\{temp\_line}\?:\?\p{INTEGER};\C{used when interchanging \\{line} with %
\\{other\_line}}\6
\\{limit}\?:\?\\{buffer\_index};\C{the last character position occupied in
                               the buffer}\6
\\{loc}\?:\?\\{buffer\_index};\C{the next character position to be read from
                               the buffer}\6
\\{input\_has\_ended}\?:\?\p{BOOLEAN};\C{if \p{TRUE}, there is no more input}\6
\\{changing}\?:\?\p{BOOLEAN};\C{if \p{TRUE}, the current line is from \\{change%
\_file}}\par
\fi

\M83. As we change \\{changing} from \p{TRUE} to \p{FALSE} and back again, we
must
remember to swap the values of \\{line} and \\{other\_line} so that the
\\{err\_print} routine will be sure to report the correct line number.

\Y\P\4\D\\{change\_changing}\?\S\?\\{changing}\?\K\R\?\\{changing};\5
\\{temp\_line}\?\K\?\\{other\_line};\5
\\{other\_line}\?\K\?\\{line};\5
\\{line}\?\K\?\\{temp\_line}\C{\\{line}\hbox{$\null\BA\null$}\\{other\_line}}%
\par
\fi

\M84. When \\{changing} is \p{FALSE}, the next line of \\{change\_file} is kept
in
\\{change\_buffer}\?(\?0\?\to\?\\{change\_limit}\?)\?, for purposes of
comparison with the next
line of \\{aweb\_file}. After the change file has been completely input, we
set \\{change\_limit}\?\K\?0, so that no further matches will be made.

\Y\P\?\4\X14:Types and objects visible from \\{data\_structures}\X\mathrel{+}\S%
\?\6
\\{change\_buffer}\?:\?\\{buffer\_type}\?\1(\?\\{buffer\_index}\?\2)\?;\6
\\{change\_limit}\?:\?\\{buffer\_index};\C{the last position occupied in %
\\{change\_buffer}}\par
\fi

\M85. Here's a simple function that checks if the two buffers are different.

\Y\P\?\4\X85:Procedures in \\{lexical\_scanning}\X\S\?\6
\&{function}\8\1\\{lines\_dont\_match}\?\8\&{return}\?\8\p{BOOLEAN}\2\8\&{is}\8%
\1\6
\4\&{begin}\6
\&{if}\1\1\8\\{change\_limit}\?\I\?\\{limit}\2\8\&{then}\6
\?\&{return}\?\8\p{TRUE};\2\6
\&{end}\8\2\&{if}\1;\6
\&{if}\1\1\?\8\1(\?\\{limit}\?>\?0\?\2)\ANT\?\\{change\_buffer}\?\1(\?0\?\to\?%
\\{limit}\?-\?1\?\2)\I\?\\{buffer}\?\1(\?0\?\to\?\\{limit}\?-\?1\?\2)\?\2\8%
\&{then}\6
\?\&{return}\?\8\p{TRUE};\2\6
\&{end}\8\2\&{if}\1;\6
\?\&{return}\?\8\p{FALSE};\2\6
\&{end}\8\\{lines\_dont\_match};\par
\A sections~86, 90, 93, 95, 101, 103, 105, 107, and~111.
\U section~6\*.\fi

\M86. Procedure \\{prime\_the\_change\_buffer} sets \\{change\_buffer} in
preparation
for the next matching operation. Since blank lines in the change file are
not used for matching, we have \?(\?\\{change\_limit}\?=\?0\?)\W\R\?%
\\{changing} if and
only if the change file is exhausted. This procedure is called only
when \\{changing} is true; hence error messages will be reported correctly.

\Y\P\?\4\X85:Procedures in \\{lexical\_scanning}\X\mathrel{+}\S\?\6
\&{procedure}\8\1\\{prime\_the\_change\_buffer}\2\8\&{is}\8\1\6
\4\&{begin}\6
\\{change\_limit}\?\K\?0;\C{this value will be used if the change file ends}\6
\X87:Skip over comment lines in the change file; \&{return} if end of file\X;\6
\X88:Skip to the next nonblank line; \&{return} if end of file\X;\6
\X89:Move \\{buffer} and \\{limit} to \\{change\_buffer} and \\{change\_limit}%
\X;\2\6
\&{end}\8\\{prime\_the\_change\_buffer};\par
\fi

\M87. While looking for a line that begins with \.{@x} in the change file,
we allow lines that begin with \.{@}, as long as they don't begin with
\.{@y} or \.{@z} (which would probably indicate that the change file is
fouled up).

\Y\P\?\4\X87:Skip over comment lines in the change file; \&{return} if end of
file\X\S\?\6
\1\&{loop}\6
\\{incr}\?\1(\?\\{line}\?\2)\?;\6
\&{if}\1\1\?\8\R\?\\{input\_ln}\?\1(\?\\{change\_file}\?\2)\?\2\8\&{then}\6
\&{return};\6
\4\&{elsif}\1\8\\{limit}\?<\?2\2\8\&{then}\6
\&{goto}\8\\{continue};\6
\4\&{elsif}\1\8\\{buffer}\?\1(\?0\?\2)\I\?\H{@}\2\8\&{then}\6
\&{goto}\8\\{continue};\6
\4\&{elsif}\1\8\\{buffer}\?\1(\?1\?\2)\in\?\H{X}\?\to\?\H{Z}\2\8\&{then}\6
\\{buffer}\?\1(\?1\?\2)\K\?\\{buffer}\?\1(\?1\?\2)+\bn{8}{40}\?;\C{lowercasify}%
\2\6
\&{end}\8\2\&{if}\1;\6
\?\&{exit}\8\&{when}\?\8\\{buffer}\?\1(\?1\?\2)=\?\H{x};\6
\&{if}\1\1\8\\{buffer}\?\1(\?1\?\2)\in\?\H{y}\?\to\?\H{z}\2\8\&{then}\6
\\{loc}\?\K\?2;\5
\\{err\_print}\?\1(\?\.{"!\ Where\ is\ the\ matching\ @x?"}\?\2)\?;\2\6
\&{end}\8\2\&{if}\1;\6
\4\BL\\{continue}\EL\8\&{null};\2\6
\&{end}\8\&{loop}\par
\U section~86.\fi

\M88. Here we are looking at lines following the \.{@x}.

\Y\P\?\4\X88:Skip to the next nonblank line; \&{return} if end of file\X\S\?\6
\1\&{loop}\6
\\{incr}\?\1(\?\\{line}\?\2)\?;\6
\&{if}\1\1\?\8\R\?\\{input\_ln}\?\1(\?\\{change\_file}\?\2)\?\2\8\&{then}\6
\\{err\_print}\?\1(\?\.{"!\ Change\ file\ ended\ after\ @x"}\?\2)\?;\6
\&{return};\2\6
\&{end}\8\2\&{if}\1;\6
\?\&{exit}\8\&{when}\?\8\\{limit}\?>\?0;\2\6
\&{end}\8\&{loop}\par
\U section~86.\fi

\M89. \P\?\X89:Move \\{buffer} and \\{limit} to \\{change\_buffer} and %
\\{change\_limit}\X\S\?\6
\\{change\_limit}\?\K\?\\{limit};\6
\&{if}\1\1\8\\{limit}\?>\?0\2\8\&{then}\6
\\{change\_buffer}\?\1(\?0\?\to\?\\{limit}\?-\?1\?\2)\K\?\\{buffer}\?\1(\?0\?%
\to\?\\{limit}\?-\?1\?\2)\?;\2\6
\&{end}\8\2\&{if}\1\par
\U sections~86 and~90.\fi

\M90. The following procedure is used to see if the next change entry should
go into effect; it is called only when \\{changing} is false.
The idea is to test whether or not the current
contents of \\{buffer} matches the current contents of \\{change\_buffer}.
If not, there's nothing more to do; but if so, a change is called for:
All of the text down to the \.{@y} is supposed to match. An error
message is issued if any discrepancy is found. Then the procedure
prepares to read the next line from \\{change\_file}.

\Y\P\?\4\X85:Procedures in \\{lexical\_scanning}\X\mathrel{+}\S\?\6
\&{procedure}\8\1\\{check\_change}\2\8\&{is}\8\C{switches to \\{change\_file}
if the buffers match}\1\6
\|n\?:\?\p{INTEGER}\?\K\?0;\C{the number of discrepancies found}\6
\4\&{begin}\6
\&{if}\1\1\8\\{lines\_dont\_match}\2\8\&{then}\6
\&{return};\2\6
\&{end}\8\2\&{if}\1;\6
\1\&{loop}\6
\\{change\_changing};\C{now it's \p{TRUE}}\6
\\{incr}\?\1(\?\\{line}\?\2)\?;\6
\&{if}\1\1\?\8\R\?\\{input\_ln}\?\1(\?\\{change\_file}\?\2)\?\2\8\&{then}\6
\\{err\_print}\?\1(\?\.{"!\ Change\ file\ ended\ before\ @y"}\?\2)\?;\6
\\{change\_limit}\?\K\?0;\5
\\{change\_changing};\C{\p{FALSE} again}\6
\&{return};\2\6
\&{end}\8\2\&{if}\1;\6
\X91:If the current line starts with \.{@y}, report any discrepancies and %
\&{return}\X;\6
\X89:Move \\{buffer} and \\{limit} to \\{change\_buffer} and \\{change\_limit}%
\X;\6
\\{change\_changing};\C{now it's \p{FALSE}}\6
\\{incr}\?\1(\?\\{line}\?\2)\?;\6
\&{if}\1\1\?\8\R\?\\{input\_ln}\?\1(\?\\{aweb\_file}\?\2)\?\2\8\&{then}\6
\\{err\_print}\?\1(\?\.{"!\ AWEB\ file\ ended\ during\ a\ change"}\?\2)\?;\6
\\{input\_has\_ended}\?\K\?\p{TRUE};\5
\&{return};\2\6
\&{end}\8\2\&{if}\1;\6
\&{if}\1\1\8\\{lines\_dont\_match}\2\8\&{then}\6
\\{incr}\?\1(\?\|n\?\2)\?;\2\6
\&{end}\8\2\&{if}\1;\2\6
\&{end}\8\&{loop};\2\6
\&{end}\8\\{check\_change};\par
\fi

\M91. \P\?\X91:If the current line starts with \.{@y}, report any discrepancies
and \&{return}\X\S\?\6
\&{if}\1\1\?\8\1(\?\\{limit}\?>\?1\?\2)\ANT\?\\{buffer}\?\1(\?0\?\2)=\?\H{@}\2%
\8\&{then}\6
\&{if}\1\1\8\\{buffer}\?\1(\?1\?\2)\in\?\H{X}\?\to\?\H{Z}\2\8\&{then}\6
\\{buffer}\?\1(\?1\?\2)\K\?\\{buffer}\?\1(\?1\?\2)+\bn{8}{40}\?;\C{lowercasify}%
\2\6
\&{end}\8\2\&{if}\1;\6
\&{if}\1\1\?\8\1(\?\\{buffer}\?\1(\?1\?\2)=\?\H{x}\?\2)\V\1(\?\\{buffer}\?\1(%
\?1\?\2)=\?\H{z}\?\2)\?\2\8\&{then}\6
\\{loc}\?\K\?2;\5
\\{err\_print}\?\1(\?\.{"!\ Where\ is\ the\ matching\ @y?"}\?\2)\?;\6
\4\&{elsif}\1\8\\{buffer}\?\1(\?1\?\2)=\?\H{y}\2\8\&{then}\6
\&{if}\1\1\8\|n\?>\?0\2\8\&{then}\6
\\{loc}\?\K\?2;\5
\\{new\_ln};\5
\\{print}\?\1(\?\.{"!\ Hmm...\ "}\?\2)\?;\5
\\{print\_int}\?\1(\?\|n\?,\39\?1\?\2)\?;\5
\\{print}\?\1(\?\.{"\ of\ the\ preceding\ lines\ failed\ to\ match"}\?\2)\?;\5
\\{error};\2\6
\&{end}\8\2\&{if}\1;\6
\&{return};\2\6
\&{end}\8\2\&{if}\1;\2\6
\&{end}\8\2\&{if}\1\par
\U section~90.\fi

\M92. The \\{reset\_input} procedure, which gets \.{AWEAVE} ready to read the
user's \.{AWEB} input, is used at the beginning of phases one and two.

\Y\P\?\4\X92:Specification of procedures visible from \\{lexical\_scanning}\X\S%
\?\6
\&{procedure}\8\1\\{reset\_input}\2;\par
\A sections~94, 100, 102, 104, 106, and~110.
\U section~6\*.\fi

\M93. \P\?\X85:Procedures in \\{lexical\_scanning}\X\mathrel{+}\S\?\6
\&{procedure}\8\1\\{reset\_input}\2\8\&{is}\8\1\6
\4\&{begin}\6
\p{TEXT\_IO}\?.\?\p{RESET}\?\1(\?\\{aweb\_file}\?\2)\?;\6
\&{if}\1\1\8\p{TEXT\_IO}\?.\?\p{IS\_OPEN}\?\1(\?\\{change\_file}\?\2)\?\2\8%
\&{then}\6
\p{TEXT\_IO}\?.\?\p{RESET}\?\1(\?\\{change\_file}\?\2)\?;\2\6
\&{end}\8\2\&{if}\1;\6
\\{line}\?\K\?0;\5
\\{other\_line}\?\K\?0;\6
\\{changing}\?\K\?\p{TRUE};\5
\\{prime\_the\_change\_buffer};\5
\\{change\_changing};\6
\\{limit}\?\K\?0;\5
\\{loc}\?\K\?1;\5
\\{buffer}\?\1(\?0\?\2)\K\?\H{\ };\5
\\{input\_has\_ended}\?\K\?\p{FALSE};\2\6
\&{end}\8\\{reset\_input};\par
\fi

\M94. The \\{get\_line} procedure is called when \\{loc}\?>\?\\{limit}; it puts
the next
line of merged input into the buffer and updates the other variables
appropriately. A space is placed at the right end of the line.

\Y\P\?\4\X92:Specification of procedures visible from \\{lexical\_scanning}\X%
\mathrel{+}\S\?\6
\&{procedure}\8\1\\{get\_line}\2;\par
\fi

\M95. \P\?\X85:Procedures in \\{lexical\_scanning}\X\mathrel{+}\S\?\6
\&{procedure}\8\1\\{get\_line}\2\8\&{is}\8\C{inputs the next line}\1\6
\4\&{begin}\6
\4\BL\\{restart}\EL\8\&{if}\1\1\8\\{changing}\2\8\&{then}\6
\\{changed\_module}\?\1(\?\\{module\_count}\?\2)\K\?\p{TRUE};\6
\4\&{else}\6
\X96:Read from \\{aweb\_file} and maybe turn on \\{changing}\X;\2\6
\&{end}\8\2\&{if}\1;\6
\&{if}\1\1\8\\{changing}\2\8\&{then}\6
\X97:Read from \\{change\_file} and maybe turn off \\{changing}\X;\6
\&{if}\1\1\?\8\R\?\\{changing}\2\8\&{then}\6
\\{changed\_module}\?\1(\?\\{module\_count}\?\2)\K\?\p{TRUE};\5
\&{goto}\8\\{restart};\2\6
\&{end}\8\2\&{if}\1;\2\6
\&{end}\8\2\&{if}\1;\6
\\{loc}\?\K\?0;\5
\\{buffer}\?\1(\?\\{limit}\?\2)\K\?\H{\ };\2\6
\&{end}\8\\{get\_line};\par
\fi

\M96. \P\?\X96:Read from \\{aweb\_file} and maybe turn on \\{changing}\X\S\?\6
\\{incr}\?\1(\?\\{line}\?\2)\?;\6
\&{if}\1\1\?\8\R\?\\{input\_ln}\?\1(\?\\{aweb\_file}\?\2)\?\2\8\&{then}\6
\\{input\_has\_ended}\?\K\?\p{TRUE};\6
\4\&{elsif}\1\?\8\1(\?\\{limit}\?=\?\\{change\_limit}\?\2)\W\1(\1(\?\\{buffer}%
\?\1(\?0\?\2)=\?\\{change\_buffer}\?\1(\?0\?\2)\2)\W\1(\?\\{change\_limit}\?>%
\?0\?\2)\2)\?\2\8\&{then}\6
\\{check\_change};\2\6
\&{end}\8\2\&{if}\1\par
\U section~95.\fi

\M97. \P\?\X97:Read from \\{change\_file} and maybe turn off \\{changing}\X\S\?%
\6
\\{incr}\?\1(\?\\{line}\?\2)\?;\6
\&{if}\1\1\?\8\R\?\\{input\_ln}\?\1(\?\\{change\_file}\?\2)\?\2\8\&{then}\6
\\{err\_print}\?\1(\?\.{"!\ Change\ file\ ended\ without\ @z"}\?\2)\?;\6
\\{buffer}\?\1(\?0\?\to\?1\?\2)\K\1(\?\H{@}\?,\39\?\H{z}\?\2)\?;\5
\\{limit}\?\K\?2;\2\6
\&{end}\8\2\&{if}\1;\6
\&{if}\1\1\?\8\1(\?\\{limit}\?>\?1\?\2)\ANT\1(\?\\{buffer}\?\1(\?0\?\2)=\?\H{@}%
\?\2)\?\2\8\&{then}\C{check if the change has ended}\6
\&{if}\1\1\8\\{buffer}\?\1(\?1\?\2)\in\?\H{X}\?\to\?\H{Z}\2\8\&{then}\6
\\{buffer}\?\1(\?1\?\2)\K\?\\{buffer}\?\1(\?1\?\2)+\bn{8}{40}\?;\C{lowercasify}%
\2\6
\&{end}\8\2\&{if}\1;\6
\&{if}\1\1\8\\{buffer}\?\1(\?1\?\2)\in\?\H{x}\?\to\?\H{y}\2\8\&{then}\6
\\{loc}\?\K\?2;\5
\\{err\_print}\?\1(\?\.{"!\ Where\ is\ the\ matching\ @z?"}\?\2)\?;\6
\4\&{elsif}\1\8\\{buffer}\?\1(\?1\?\2)=\?\H{z}\2\8\&{then}\6
\\{prime\_the\_change\_buffer};\5
\\{change\_changing};\2\6
\&{end}\8\2\&{if}\1;\2\6
\&{end}\8\2\&{if}\1\par
\U section~95.\fi

\M98. At the end of the program, we will tell the user if the change file
had a line that didn't match any relevant line in \\{aweb\_file}.

\Y\P\?\4\X98:Check that all changes have been read\X\S\?\6
\&{if}\1\1\8\\{change\_limit}\?\I\?0\2\8\&{then}\C{\\{changing} is \p{FALSE}}\6
\\{buffer}\?\1(\?0\?\to\?\\{change\_limit}\?\2)\K\?\\{change\_buffer}\?\1(\?0\?%
\to\?\\{change\_limit}\?\2)\?;\5
\\{limit}\?\K\?\\{change\_limit};\5
\\{changing}\?\K\?\p{TRUE};\5
\\{line}\?\K\?\\{other\_line};\5
\\{loc}\?\K\?\\{change\_limit};\5
\\{err\_print}\?\1(\?\.{"!\ Change\ file\ entry\ did\ not\ match"}\?\2)\?;\2\6
\&{end}\8\2\&{if}\1\par
\U section~308\*.\fi

\M99. Control codes in \.{AWEB}, which begin with `\.{@}', are converted
into a numeric code designed to simplify \.{AWEAVE}'s logic; for example,
larger numbers are given to the control codes that denote more significant
milestones, and the code of \\{new\_module} should be the largest of
all. Some of these numeric control codes take the place of ASCII
control codes that will not otherwise appear in the output of the
scanning routines.

\Y\P\?\4\X14:Types and objects visible from \\{data\_structures}\X\mathrel{+}\S%
\?\6
\\{ignore}\?:\?\&{constant}\8\\{eight\_bits}\?\K\?0;\C{control code of no
interest to \.{AWEAVE}}\6
\\{verbatim}\?:\?\&{constant}\8\\{eight\_bits}\?\K\bn{8}{2}\?;\C{extended ASCII
alpha will not appear}\6
\\{force\_line}\?:\?\&{constant}\8\\{eight\_bits}\?\K\bn{8}{3}\?;\C{extended
ASCII beta will not appear}\6
\\{begin\_comment}\?:\?\&{constant}\8\\{eight\_bits}\?\K\bn{8}{11}\?;\C{ASCII
tab mark will not appear}\6
\\{end\_comment}\?:\?\&{constant}\8\\{eight\_bits}\?\K\bn{8}{12}\?;\C{ASCII
line feed will not appear}\6
\\{character\_tok}\?:\?\&{constant}\8\\{eight\_bits}\?\K\bn{8}{14}\?;\C{ASCII
form feed will not appear}\6
\\{ascii\_tok}\?:\?\&{constant}\8\\{eight\_bits}\?\K\bn{8}{15}\?;\C{ASCII
carriage return will not appear}\6
\\{output\_tok}\?:\?\&{constant}\8\\{eight\_bits}\?\K\bn{8}{16}\?;\C{ASCII %
\.{DLE} will not appear}\6
\\{double\_dot}\?:\?\&{constant}\8\\{eight\_bits}\?\K\bn{8}{40}\?;\C{ASCII
space will not appear except in strings}\6
\\{no\_underline}\?:\?\&{constant}\8\\{eight\_bits}\?\K\bn{8}{175}\?;\C{this
code will be intercepted without confusion}\6
\\{underline}\?:\?\&{constant}\8\\{eight\_bits}\?\K\bn{8}{176}\?;\C{this code
will be intercepted without confusion}\6
\\{param}\?:\?\&{constant}\8\\{eight\_bits}\?\K\bn{8}{177}\?;\C{ASCII delete
will not appear}\6
\\{xref\_roman}\?:\?\&{constant}\8\\{eight\_bits}\?\K\bn{8}{204}\?;\C{control
code for `\.{@\^}'}\6
\\{xref\_wildcard}\?:\?\&{constant}\8\\{eight\_bits}\?\K\bn{8}{205}\?;%
\C{control code for `\.{@:}'}\6
\\{xref\_italic}\?:\?\&{constant}\8\\{eight\_bits}\?\K\bn{8}{206}\?;\C{control
code for `\.{@i}'}\6
\\{xref\_bold\_face}\?:\?\&{constant}\8\\{eight\_bits}\?\K\bn{8}{207}\?;%
\C{control code for `\.{@b}'}\6
\\{xref\_slanted}\?:\?\&{constant}\8\\{eight\_bits}\?\K\bn{8}{210}\?;\C{control
code for `\.{@s}'}\6
\\{xref\_typewriter}\?:\?\&{constant}\8\\{eight\_bits}\?\K\bn{8}{211}\?;%
\C{control code for `\.{@.}'}\6
\\{xref\_predefined}\?:\?\&{constant}\8\\{eight\_bits}\?\K\bn{8}{212}\?;%
\C{control code for `\.{@p}'}\6
\\{xref\_reserved}\?:\?\&{constant}\8\\{eight\_bits}\?\K\bn{8}{213}\?;%
\C{control code for `\.{@r}'}\6
\\{TeX\_string}\?:\?\&{constant}\8\\{eight\_bits}\?\K\bn{8}{214}\?;\C{control
code for `\.{@t}'}\6
\\{join}\?:\?\&{constant}\8\\{eight\_bits}\?\K\bn{8}{215}\?;\C{control code for
`\.{@\&}'}\6
\\{thin\_space}\?:\?\&{constant}\8\\{eight\_bits}\?\K\bn{8}{216}\?;\C{control
code for `\.{@,}'}\6
\\{math\_break}\?:\?\&{constant}\8\\{eight\_bits}\?\K\bn{8}{217}\?;\C{control
code for `\.{@\char'174}'}\6
\\{line\_break}\?:\?\&{constant}\8\\{eight\_bits}\?\K\bn{8}{220}\?;\C{control
code for `\.{@/}'}\6
\\{big\_line\_break}\?:\?\&{constant}\8\\{eight\_bits}\?\K\bn{8}{221}\?;%
\C{control code for `\.{@\#}'}\6
\\{no\_line\_break}\?:\?\&{constant}\8\\{eight\_bits}\?\K\bn{8}{222}\?;%
\C{control code for `\.{@+}'}\6
\\{pseudo\_semi}\?:\?\&{constant}\8\\{eight\_bits}\?\K\bn{8}{223}\?;\C{control
code for `\.{@;}'}\6
\\{format}\?:\?\&{constant}\8\\{eight\_bits}\?\K\bn{8}{224}\?;\C{control code
for `\.{@f}'}\6
\\{definition}\?:\?\&{constant}\8\\{eight\_bits}\?\K\bn{8}{225}\?;\C{control
code for `\.{@d}'}\6
\\{begin\_ada}\?:\?\&{constant}\8\\{eight\_bits}\?\K\bn{8}{226}\?;\C{control
code for `\.{@a}'}\6
\\{module\_name}\?:\?\&{constant}\8\\{eight\_bits}\?\K\bn{8}{227}\?;\C{control
code for `\.{@<}'}\6
\\{new\_module}\?:\?\&{constant}\8\\{eight\_bits}\?\K\bn{8}{230}\?;\C{control
code for `\.{@\ }' and `\.{@*}'}\par
\fi

\M100. Control codes are converted from ASCII to \.{AWEAVE}'s internal
representation by the \\{control\_code} routine.

\Y\P\?\4\X92:Specification of procedures visible from \\{lexical\_scanning}\X%
\mathrel{+}\S\?\6
\&{function}\8\1\\{control\_code}\?(\1\?\|c\?:\?\\{ASCII\_code}\?\2)\&{return}%
\?\8\\{eight\_bits}\2;\par
\fi

\M101. \P\?\X85:Procedures in \\{lexical\_scanning}\X\mathrel{+}\S\?\6
\&{function}\8\1\\{control\_code}\?(\1\?\|c\?:\?\\{ASCII\_code}\?\2)\&{return}%
\?\8\\{eight\_bits}\2\8\&{is}\8\C{convert \|c   after \.{@}}\1\6
\4\&{begin}\6
\&{case}\1\8\|c\1\8\&{is}\8\6
\4\&{when}\8\H{@}\KR\8\?\&{return}\?\8\H{@};\C{`quoted' at sign}\6
\4\&{when}\8\H{\'}\KR\8\?\&{return}\?\8\\{ascii\_tok};\C{precedes an ASCII
constant}\6
\4\&{when}\8\H{\~}\KR\8\?\&{return}\?\8\\{output\_tok};\C{output file name}\6
\4\&{when}\8\H{\ }\?\mathrel{\.|}\?\\{tab\_mark}\?\mathrel{\.|}\?\H{*}\KR\8\?%
\&{return}\?\8\\{new\_module};\C{beginning of a new module}\6
\4\&{when}\8\H{=}\KR\8\?\&{return}\?\8\\{verbatim};\C{verbatim string}\6
\4\&{when}\8\H{\\}\KR\8\?\&{return}\?\8\\{force\_line};\C{force line in \ADA\
output of \.{ATANGLE}}\6
\4\&{when}\8\H{D}\?\mathrel{\.|}\?\H{d}\KR\8\?\&{return}\?\8\\{definition};%
\C{macro definition}\6
\4\&{when}\8\H{F}\?\mathrel{\.|}\?\H{f}\KR\8\?\&{return}\?\8\\{format};%
\C{format definition}\6
\4\&{when}\8\H{\{}\KR\8\?\&{return}\?\8\\{begin\_comment};\C{begin-comment
delimiter}\6
\4\&{when}\8\H{\}}\KR\8\?\&{return}\?\8\\{end\_comment};\C{end-comment
delimiter}\6
\4\&{when}\8\H{A}\?\mathrel{\.|}\?\H{a}\KR\8\?\&{return}\?\8\\{begin\_ada};\C{%
\ADA\ text in unnamed module}\6
\4\&{when}\8\H{\&}\KR\8\?\&{return}\?\8\\{join};\C{concatenate two tokens}\6
\4\&{when}\8\H{<}\KR\8\?\&{return}\?\8\\{module\_name};\C{beginning of a module
name}\6
\4\&{when}\8\H{>}\KR\6
\\{err\_print}\?\1(\?\.{"!\ Extra\ @>"}\?\2)\?;\6
\?\&{return}\?\8\\{ignore};\C{end of module name should not be discovered in
this way}\6
\4\&{when}\8\H{T}\?\mathrel{\.|}\?\H{t}\KR\8\?\&{return}\?\8\\{TeX\_string};\C{%
\TeX\ box within \ADA}\6
\4\&{when}\8\H{!}\KR\8\?\&{return}\?\8\\{underline};\C{increment definition
flag}\6
\4\&{when}\8\H{?}\KR\8\?\&{return}\?\8\\{no\_underline};\C{decrement definition
flag}\6
\4\&{when}\8\H{\^}\KR\8\?\&{return}\?\8\\{xref\_roman};\C{index entry to be
typeset normally}\6
\4\&{when}\8\H{:}\KR\8\?\&{return}\?\8\\{xref\_wildcard};\C{index entry to be
in user format}\6
\4\&{when}\8\H{i}\?\mathrel{\.|}\?\H{I}\KR\8\?\&{return}\?\8\\{xref\_italic};%
\C{index entry to be in italic type}\6
\4\&{when}\8\H{b}\?\mathrel{\.|}\?\H{B}\KR\8\?\&{return}\?\8\\{xref\_bold%
\_face};\C{index entry to be in boldface type}\6
\4\&{when}\8\H{s}\?\mathrel{\.|}\?\H{S}\KR\8\?\&{return}\?\8\\{xref\_slanted};%
\C{index entry to be in slanted type}\6
\4\&{when}\8\H{p}\?\mathrel{\.|}\?\H{P}\KR\8\?\&{return}\?\8\\{xref%
\_predefined};\C{index entry like predefined identifier}\6
\4\&{when}\8\H{r}\?\mathrel{\.|}\?\H{R}\KR\8\?\&{return}\?\8\\{xref\_reserved};%
\C{reserved word to be indexed}\6
\4\&{when}\8\H{.}\KR\8\?\&{return}\?\8\\{xref\_typewriter};\C{index entry to be
in typewriter type}\6
\4\&{when}\8\H{,}\KR\8\?\&{return}\?\8\\{thin\_space};\C{puts extra space in %
\ADA\ format}\6
\4\&{when}\8\H{|}\KR\8\?\&{return}\?\8\\{math\_break};\C{allows a break in a
formula}\6
\4\&{when}\8\H{/}\KR\8\?\&{return}\?\8\\{line\_break};\C{forces end-of-line in %
\ADA\ format}\6
\4\&{when}\8\H{\#}\KR\8\?\&{return}\?\8\\{big\_line\_break};\C{forces
end-of-line and some space besides}\6
\4\&{when}\8\H{+}\KR\8\?\&{return}\?\8\\{no\_line\_break};\C{cancels
end-of-line down to single space}\6
\4\&{when}\8\H{;}\KR\8\?\&{return}\?\8\\{pseudo\_semi};\C{acts like a
semicolon, but is invisible}\6
\4\&{when}\8\&{others}\8\KR\6
\\{err\_print}\?\1(\?\.{"!\ Unknown\ control\ code"}\?\2)\?;\6
\?\&{return}\?\8\\{ignore};\2\6
\4\2\&{end}\8\&{case};\2\6
\&{end}\8\\{control\_code};\par
\fi

\M102. The \\{skip\_limbo} routine is used on the first pass to skip through
portions of the input that are not in any modules, i.e., that precede
the first module. After this procedure has been called, the value of
\\{input\_has\_ended} will tell whether or not a new module has
actually been found.

\Y\P\?\4\X92:Specification of procedures visible from \\{lexical\_scanning}\X%
\mathrel{+}\S\?\6
\&{procedure}\8\1\\{skip\_limbo}\2;\par
\fi

\M103. \P\?\X85:Procedures in \\{lexical\_scanning}\X\mathrel{+}\S\?\6
\&{procedure}\8\1\\{skip\_limbo}\2\8\&{is}\8\C{skip to next module}\1\6
\|c\?:\?\\{ASCII\_code};\C{character following \.{@}}\6
\4\&{begin}\6
\1\&{loop}\6
\&{if}\1\1\8\\{loc}\?>\?\\{limit}\2\8\&{then}\6
\\{get\_line};\6
\&{if}\1\1\8\\{input\_has\_ended}\2\8\&{then}\6
\&{return};\2\6
\&{end}\8\2\&{if}\1;\6
\4\&{else}\6
\\{buffer}\?\1(\?\\{limit}\?+\?1\?\2)\K\?\H{@};\6
\1\&{while}\8\\{buffer}\?\1(\?\\{loc}\?\2)\I\?\H{@}\8\&{loop}\6
\\{incr}\?\1(\?\\{loc}\?\2)\?;\2\6
\&{end}\8\&{loop};\6
\&{if}\1\1\8\\{loc}\?\L\?\\{limit}\2\8\&{then}\6
\\{loc}\?\K\?\\{loc}\?+\?2;\5
\|c\?\K\?\\{buffer}\?\1(\?\\{loc}\?-\?1\?\2)\?;\6
\&{if}\1\1\?\8\1(\?\|c\?=\?\H{\ }\?\2)\V\1(\?\|c\?=\?\\{tab\_mark}\?\2)\V\1(\?%
\|c\?=\?\H{*}\?\2)\?\2\8\&{then}\6
\&{return};\2\6
\&{end}\8\2\&{if}\1;\2\6
\&{end}\8\2\&{if}\1;\2\6
\&{end}\8\2\&{if}\1;\2\6
\&{end}\8\&{loop};\2\6
\&{end}\8\\{skip\_limbo};\par
\fi

\M104. The \\{skip\_TeX} routine is used on the first pass to skip through
the \TeX\ code at the beginning of a module. It returns the next
control code or `\v' found in the input. A \\{new\_module} is
assumed to exist at the very end of the file.

\Y\P\?\4\X92:Specification of procedures visible from \\{lexical\_scanning}\X%
\mathrel{+}\S\?\6
\&{function}\8\1\\{skip\_TeX}\?\8\&{return}\?\8\\{eight\_bits}\2;\par
\fi

\M105. \P\?\X85:Procedures in \\{lexical\_scanning}\X\mathrel{+}\S\?\6
\&{function}\8\1\\{skip\_TeX}\?\8\&{return}\?\8\\{eight\_bits}\2\8\&{is}\8%
\C{skip past pure \TeX\ code}\1\6
\|c\?:\?\\{eight\_bits};\C{control code found}\6
\4\&{begin}\6
\1\&{loop}\6
\&{if}\1\1\8\\{loc}\?>\?\\{limit}\2\8\&{then}\6
\\{get\_line};\6
\&{if}\1\1\8\\{input\_has\_ended}\2\8\&{then}\6
\?\&{return}\?\8\\{new\_module};\2\6
\&{end}\8\2\&{if}\1;\2\6
\&{end}\8\2\&{if}\1;\6
\\{buffer}\?\1(\?\\{limit}\?+\?1\?\2)\K\?\H{@};\6
\1\&{loop}\6
\|c\?\K\?\\{buffer}\?\1(\?\\{loc}\?\2)\?;\5
\\{incr}\?\1(\?\\{loc}\?\2)\?;\6
\&{if}\1\1\8\|c\?=\?\H{|}\2\8\&{then}\6
\?\&{return}\?\8\H{|};\2\6
\&{end}\8\2\&{if}\1;\6
\?\&{exit}\8\&{when}\?\8\|c\?=\?\H{@};\2\6
\&{end}\8\&{loop};\6
\&{if}\1\1\8\\{loc}\?\L\?\\{limit}\2\8\&{then}\6
\|c\?\K\?\\{control\_code}\?\1(\?\\{buffer}\?\1(\?\\{loc}\?\2)\2)\?;\5
\\{incr}\?\1(\?\\{loc}\?\2)\?;\5
\?\&{return}\?\8\|c;\2\6
\&{end}\8\2\&{if}\1;\2\6
\&{end}\8\&{loop};\2\6
\&{end}\8\\{skip\_TeX};\par
\fi

\M106. The \\{skip\_comment} routine is used on the first pass to skip
through \TeX\ code in \ADA\ comments. The \\{bal} parameter
tells how many left braces are assumed to have been scanned when
this routine is called, and the procedure returns a corresponding
value of \\{bal} at the point that scanning has stopped. Scanning
stops either at a `\v' that introduces \ADA\ text,
in which case the returned value is positive, or it stops at the
end of the comment, in which case the returned value is zero.
The scanning also stops in anomalous situations when the comment
doesn't end or when it contains an illegal use of \.{@}.
One should call \\{skip\_comment}\?(\?1\?)\? when beginning to scan a comment.

\Y\P\?\4\X92:Specification of procedures visible from \\{lexical\_scanning}\X%
\mathrel{+}\S\?\6
\&{function}\8\1\\{skip\_comment}\?(\1\?\\{bal}\?:\?\\{eight\_bits}\?\2)%
\&{return}\?\8\\{eight\_bits}\2;\par
\fi

\M107. \P\?\X85:Procedures in \\{lexical\_scanning}\X\mathrel{+}\S\?\6
\&{function}\8\1\\{skip\_comment}\?(\1\?\\{bal}\?:\?\\{eight\_bits}\?\2)%
\&{return}\?\8\\{eight\_bits}\2\8\&{is}\8\C{skips \TeX\   code in comments}\1\6
\|b\?:\?\\{eight\_bits}\?\K\?\\{bal};\C{depth of brace nesting}\6
\|c\?:\?\\{ASCII\_code};\C{the current character}\6
\4\&{begin}\6
\1\&{loop}\6
\&{if}\1\1\8\\{loc}\?>\?\\{limit}\2\8\&{then}\6
\\{get\_line};\6
\&{if}\1\1\8\\{input\_has\_ended}\2\8\&{then}\6
\?\&{return}\?\80;\C{an error message will occur in phase two}\2\6
\&{end}\8\2\&{if}\1;\2\6
\&{end}\8\2\&{if}\1;\6
\|c\?\K\?\\{buffer}\?\1(\?\\{loc}\?\2)\?;\5
\\{incr}\?\1(\?\\{loc}\?\2)\?;\6
\&{if}\1\1\8\|c\?=\?\H{|}\2\8\&{then}\6
\?\&{return}\?\8\|b;\2\6
\&{end}\8\2\&{if}\1;\6
\X108:Do special things when \|c\?=\?\H{@}\?,\?\H{\\}\?,\?\H{\{}\?,\?\H{\}}; %
\&{return} at end\X;\2\6
\&{end}\8\&{loop};\2\6
\&{end}\8\\{skip\_comment};\par
\fi

\M108. \P\?\X108:Do special things when \|c\?=\?\H{@}\?,\?\H{\\}\?,\?\H{\{}\?,%
\?\H{\}}; \&{return} at end\X\S\?\6
\&{if}\1\1\8\|c\?=\?\H{@}\2\8\&{then}\6
\|c\?\K\?\\{buffer}\?\1(\?\\{loc}\?\2)\?;\6
\&{if}\1\1\?\8\1(\?\|c\?\I\?\H{\ }\?\2)\W\1(\?\|c\?\I\?\\{tab\_mark}\?\2)\W\1(%
\?\|c\?\I\?\H{*}\?\2)\?\2\8\&{then}\6
\\{incr}\?\1(\?\\{loc}\?\2)\?;\6
\4\&{else}\6
\\{decr}\?\1(\?\\{loc}\?\2)\?;\5
\?\&{return}\?\80;\C{an error message will occur in phase two}\2\6
\&{end}\8\2\&{if}\1;\6
\4\&{elsif}\1\?\8\1(\?\|c\?=\?\H{\\}\?\2)\W\1(\?\\{buffer}\?\1(\?\\{loc}\?\2)\I%
\?\H{@}\?\2)\?\2\8\&{then}\6
\\{incr}\?\1(\?\\{loc}\?\2)\?;\6
\4\&{elsif}\1\8\|c\?=\?\H{\{}\2\8\&{then}\6
\\{incr}\?\1(\?\|b\?\2)\?;\6
\4\&{elsif}\1\8\|c\?=\?\H{\}}\2\8\&{then}\6
\\{decr}\?\1(\?\|b\?\2)\?;\6
\&{if}\1\1\8\|b\?=\?0\2\8\&{then}\6
\?\&{return}\?\80;\2\6
\&{end}\8\2\&{if}\1;\2\6
\&{end}\8\2\&{if}\1\par
\U section~107.\fi

\N109.  Inputting the next token.
As stated above, \.{AWEAVE}'s most interesting lexical scanning routine is the
\\{get\_next} function that inputs the next token of \ADA\ input. However,
\\{get\_next} is not especially complicated.

The result of \\{get\_next} is either an ASCII code for some special character,
or it is a special code representing a pair of characters (e.g., `\.{:=}'
or `\.{..}'), or it is the numeric value computed by the \\{control\_code}
procedure, or it is one of the following special codes:

\yskip\hang \\{exponent}: The `\.E' in a real constant.

\yskip\hang \\{identifier}: In this case the global variables \\{id\_first}
and \\{id\_loc} will have been set to the appropriate values needed by the
\\{id\_lookup} routine.

\yskip\hang \\{str}: In this case the global variables \\{id\_first} and
\\{id\_loc} will have been set to the beginning and ending-plus-one locations
in the buffer.  The string ends with the first reappearance of its initial
delimiter; thus, for example, $$\.{"This is ""not"" a single string"}$$
will be treated as three consecutive strings, the first being \.{"This
is"}.

\yskip\noindent There are some special cases that need to be considered:

\yskip\hang \\{character\_tok} is returned if a character literal, e.g.\
\.{\'a\'}, is sensed and \\{id\_first} and \\{id\_loc} are set appropriate
as described for \\{str}.

\yskip\hang based literals: Between two \# signs characters are treated
as digits; this is done be means of the global variable
\\{scanning\_based\_number}.

\yskip\noindent Furthermore, some of the control codes cause
\\{get\_next} to take additional actions:

\yskip\hang $\\{xref\_roman},\ldots,\\{TeX\_string}$: The values of
\\{id\_first} and \\{id\_loc} will be set so that the string in question
appears
in \\{buffer}\?(\?\\{id\_first}\?\to(\?\\{id\_loc}\?-\?1\?))\?.

\yskip\hang \\{module\_name}: In this case the global variable \\{cur\_module}
will
point to the \\{byte\_start} entry for the module name that has just been
scanned.

\yskip\hang \\{output\_tok}: The file name is treated the same way
as decribed above for \\{str}.

\yskip\hang \\{ascii\_tok}: The ASCII code constant is treated the same way
as described above for \\{character\_tok}.

\yskip\noindent If \\{get\_next} sees `\.{@!}' or `\.{@?}',
it increments respectively decrements \\{xref\_switch} by \\{def\_flag},
supposed that 0\?\L\?\\{xref\_switch}\?\L\?\\{body\_flag} afterwards,
and goes on to the next token.

\Y\P\?\4\X14:Types and objects visible from \\{data\_structures}\X\mathrel{+}\S%
\?\6
\\{exponent}\?:\?\&{constant}\8\\{eight\_bits}\?\K\bn{8}{201}\?;\C{\.E or \.e
following a digit}\6
\\{str}\?:\?\&{constant}\8\\{eight\_bits}\?\K\bn{8}{202}\?;\C{\ADA\ string}\6
\\{identifier}\?:\?\&{constant}\8\\{eight\_bits}\?\K\bn{8}{203}\?;\C{\ADA\
identifier or reserved word}\7
\\{cur\_module}\?:\?\\{name\_pointer};\C{name of module just scanned}\6
\\{scanning\_based\_number}\?:\?\p{BOOLEAN}\?\K\?\p{FALSE};\C{are we scanning a
based number?}\par
\fi

\M110. As one might expect, \\{get\_next} consists mostly of a big switch
that branches to the various special cases that can arise. Some care is
needed for the handling of \ADA's based literals.

\Y\P\?\4\X92:Specification of procedures visible from \\{lexical\_scanning}\X%
\mathrel{+}\S\?\6
\&{function}\8\1\\{get\_next}\?\8\&{return}\?\8\\{eight\_bits}\2;\par
\fi

\M111. \P\?\X85:Procedures in \\{lexical\_scanning}\X\mathrel{+}\S\?\6
\&{function}\8\1\\{get\_next}\?\8\&{return}\?\8\\{eight\_bits}\2\8\&{is}\8%
\C{produces the next input token}\1\6
\|c\?:\?\\{eight\_bits};\C{the current character}\6
\|d\?:\?\\{eight\_bits};\C{the next character}\6
\|j\?,\39\?\|k\?:\?\\{long\_name\_index};\C{indices into \\{mod\_text}}\6
\4\&{begin}\6
\4\BL\\{restart}\EL\8\&{if}\1\1\8\\{loc}\?>\?\\{limit}\2\8\&{then}\6
\\{get\_line};\6
\&{if}\1\1\8\\{input\_has\_ended}\2\8\&{then}\6
\?\&{return}\?\8\\{new\_module};\6
\4\&{elsif}\1\8\\{scanning\_based\_number}\2\8\&{then}\6
\\{err\_print}\?\1(\?\.{"!\ Improper\ based\ number,\ \#\ sign\ missing"}\?\2)%
\?;\C{a based literal should fit in one line}\2\6
\&{end}\8\2\&{if}\1;\2\6
\&{end}\8\2\&{if}\1;\6
\|c\?\K\?\\{buffer}\?\1(\?\\{loc}\?\2)\?;\5
\\{incr}\?\1(\?\\{loc}\?\2)\?;\6
\&{if}\1\1\8\\{scanning\_based\_number}\2\8\&{then}\6
\&{if}\1\1\8\|c\?=\?\H{\#}\2\8\&{then}\6
\\{scanning\_based\_number}\?\K\?\p{FALSE};\2\6
\&{end}\8\2\&{if}\1;\6
\?\&{return}\?\8\|c;\2\6
\&{end}\8\2\&{if}\1;\6
\&{case}\1\8\|c\1\8\&{is}\8\6
\4\&{when}\8\H{\#}\KR\6
\&{if}\1\1\?\8\1(\?\\{loc}\?>\?1\?\2)\ANT\1(\?\\{buffer}\?\1(\?\\{loc}\?-\?2\?%
\2)\in\?\H{0}\?\to\?\H{9}\?\2)\?\2\8\&{then}\6
\\{scanning\_based\_number}\?\K\?\p{TRUE};\2\6
\&{end}\8\2\&{if}\1;\6
\?\&{return}\?\8\H{\#};\6
\4\&{when}\8\H{A}\?\to\?\H{Z}\?\mathrel{\.|}\?\H{a}\?\to\?\H{z}\KR\6
\X113:Get an identifier\X;\6
\4\&{when}\8\H{\'}\KR\6
\X123:Get a character token\X;\6
\4\&{when}\8\H{"}\KR\6
\X114:Get a string\X;\6
\4\&{when}\8\H{@}\KR\6
\X115:Get control code and possible module name\X;\6
\hbox{\4}\X112:Compress two-symbol combinations like `\.{:=}'\X\6
\4\&{when}\8\H{\ }\?\mathrel{\.|}\?\\{tab\_mark}\KR\6
\&{goto}\8\\{restart};\C{ignore spaces and tabs}\6
\4\&{when}\8\&{others}\8\KR\6
\&{null};\2\6
\4\2\&{end}\8\&{case};\6
\?\&{return}\?\8\|c;\2\6
\&{end}\8\\{get\_next};\par
\fi

\M112.
\Y\P\4\D\\{compress}\?\1(\?\#\?\2)\S\?\6
\&{if}\1\1\8\\{loc}\?\L\?\\{limit}\2\8\&{then}\6
\|c\?\K\?\#;\5
\\{incr}\?\1(\?\\{loc}\?\2)\?;\2\6
\&{end}\8\2\&{if}\1\par
\Y\P\?\4\X112:Compress two-symbol combinations like `\.{:=}'\X\S\?\6
\4\&{when}\8\H{.}\KR\6
\&{if}\1\1\8\\{buffer}\?\1(\?\\{loc}\?\2)=\?\H{.}\2\8\&{then}\6
\\{compress}\?\1(\?\\{double\_dot}\?\2)\?;\2\6
\&{end}\8\2\&{if}\1;\6
\4\&{when}\8\H{:}\KR\6
\&{if}\1\1\8\\{buffer}\?\1(\?\\{loc}\?\2)=\?\H{=}\2\8\&{then}\6
\\{compress}\?\1(\?\\{left\_arrow}\?\2)\?;\2\6
\&{end}\8\2\&{if}\1;\6
\4\&{when}\8\H{=}\KR\6
\&{if}\1\1\8\\{buffer}\?\1(\?\\{loc}\?\2)=\?\H{=}\2\8\&{then}\6
\\{compress}\?\1(\?\\{equivalence\_sign}\?\2)\?;\6
\4\&{elsif}\1\8\\{buffer}\?\1(\?\\{loc}\?\2)=\?\H{>}\2\8\&{then}\6
\\{compress}\?\1(\?\\{right\_arrow}\?\2)\?;\2\6
\&{end}\8\2\&{if}\1;\6
\4\&{when}\8\H{>}\KR\6
\&{if}\1\1\8\\{buffer}\?\1(\?\\{loc}\?\2)=\?\H{=}\2\8\&{then}\6
\\{compress}\?\1(\?\\{greater\_or\_equal}\?\2)\?;\6
\4\&{elsif}\1\8\\{buffer}\?\1(\?\\{loc}\?\2)=\?\H{>}\2\8\&{then}\6
\\{compress}\?\1(\?\\{right\_label\_bracket}\?\2)\?;\2\6
\&{end}\8\2\&{if}\1;\6
\4\&{when}\8\H{<}\KR\6
\&{if}\1\1\8\\{buffer}\?\1(\?\\{loc}\?\2)=\?\H{=}\2\8\&{then}\6
\\{compress}\?\1(\?\\{less\_or\_equal}\?\2)\?;\6
\4\&{elsif}\1\8\\{buffer}\?\1(\?\\{loc}\?\2)=\?\H{<}\2\8\&{then}\6
\\{compress}\?\1(\?\\{left\_label\_bracket}\?\2)\?;\6
\4\&{elsif}\1\8\\{buffer}\?\1(\?\\{loc}\?\2)=\?\H{>}\2\8\&{then}\6
\\{compress}\?\1(\?\\{box\_sign}\?\2)\?;\2\6
\&{end}\8\2\&{if}\1;\6
\4\&{when}\8\H{/}\KR\6
\&{if}\1\1\8\\{buffer}\?\1(\?\\{loc}\?\2)=\?\H{=}\2\8\&{then}\6
\\{compress}\?\1(\?\\{not\_equal}\?\2)\?;\2\6
\&{end}\8\2\&{if}\1;\6
\4\&{when}\8\H{*}\KR\6
\&{if}\1\1\8\\{buffer}\?\1(\?\\{loc}\?\2)=\?\H{*}\2\8\&{then}\6
\\{compress}\?\1(\?\\{double\_star}\?\2)\?;\2\6
\&{end}\8\2\&{if}\1;\par
\U section~111.\fi

\M113. \P\?\X113:Get an identifier\X\S\?\6
\&{if}\1\1\?\8\1(\1(\1(\?\|c\?=\?\H{E}\?\2)\V\1(\?\|c\?=\?\H{e}\?\2)\2)\W\1(\?%
\\{loc}\?>\?1\?\2)\2)\ANT\1(\1(\?\\{buffer}\?\1(\?\\{loc}\?-\?2\?\2)\in\?\H{0}%
\?\to\?\H{9}\?\2)\V\?\\{buffer}\?\1(\?\\{loc}\?-\?2\?\2)=\?\H{\#}\?\2)\?\2\8%
\&{then}\6
\|c\?\K\?\\{exponent};\2\6
\&{end}\8\2\&{if}\1;\6
\&{if}\1\1\8\|c\?\I\?\\{exponent}\2\8\&{then}\6
\\{decr}\?\1(\?\\{loc}\?\2)\?;\5
\\{id\_first}\?\K\?\\{loc};\6
\1\&{loop}\6
\\{incr}\?\1(\?\\{loc}\?\2)\?;\5
\|d\?\K\?\\{buffer}\?\1(\?\\{loc}\?\2)\?;\6
\?\&{exit}\8\&{when}\8\1(\?\|d\?\notin\?\H{0}\?\to\?\H{9}\?\2)\W\1(\?\|d\?%
\notin\?\H{A}\?\to\?\H{Z}\?\2)\W\1(\?\|d\?\notin\?\H{a}\?\to\?\H{z}\?\2)\W\1(\?%
\|d\?\I\?\H{\_}\?\2)\?;\2\6
\&{end}\8\&{loop};\6
\|c\?\K\?\\{identifier};\5
\\{id\_loc}\?\K\?\\{loc};\2\6
\&{end}\8\2\&{if}\1\par
\U section~111.\fi

\M114. A string that starts and ends with double quote marks is
scanned by the following piece of the program.

\Y\P\?\4\X114:Get a string\X\S\?\6
\\{id\_first}\?\K\?\\{loc}\?-\?1;\6
\1\&{loop}\6
\|d\?\K\?\\{buffer}\?\1(\?\\{loc}\?\2)\?;\5
\\{incr}\?\1(\?\\{loc}\?\2)\?;\6
\&{if}\1\1\8\\{loc}\?>\?\\{limit}\2\8\&{then}\6
\\{err\_print}\?\1(\?\.{"!\ String\ constant\ didn\'t\ end"}\?\2)\?;\6
\\{loc}\?\K\?\\{limit};\5
\|d\?\K\?\|c;\2\6
\&{end}\8\2\&{if}\1;\6
\?\&{exit}\8\&{when}\?\8\|d\?=\?\|c;\2\6
\&{end}\8\&{loop};\6
\\{id\_loc}\?\K\?\\{loc};\5
\|c\?\K\?\\{str}\par
\U section~111.\fi

\M115. After an \.{@} sign has been scanned, the next character tells us
whether there is more work to do.

\Y\P\?\4\X115:Get control code and possible module name\X\S\?\6
\|c\?\K\?\\{control\_code}\?\1(\?\\{buffer}\?\1(\?\\{loc}\?\2)\2)\?;\5
\\{incr}\?\1(\?\\{loc}\?\2)\?;\6
\&{case}\1\8\|c\1\8\&{is}\8\6
\4\&{when}\8\\{underline}\KR\6
\&{if}\1\1\8\\{xref\_switch}\?<\?\\{body\_flag}\2\8\&{then}\6
\\{xref\_switch}\?\K\?\\{xref\_switch}\?+\?\\{def\_flag};\2\6
\&{end}\8\2\&{if}\1;\6
\&{goto}\8\\{restart};\6
\4\&{when}\8\\{no\_underline}\KR\6
\&{if}\1\1\8\\{xref\_switch}\?>\?0\2\8\&{then}\6
\\{xref\_switch}\?\K\?\\{xref\_switch}\?-\?\\{def\_flag};\2\6
\&{end}\8\2\&{if}\1;\6
\&{goto}\8\\{restart};\6
\4\&{when}\8\\{xref\_roman}\?\to\?\\{TeX\_string}\KR\6
\X121:Scan to the next \.{@>}\X;\6
\4\&{when}\8\\{module\_name}\KR\6
\X116:Scan the module name and make \\{cur\_module} point to it\X;\6
\4\&{when}\8\\{verbatim}\?\mathrel{\.|}\?\\{output\_tok}\KR\6
\X122:Scan a verbatim or output string\X;\6
\4\&{when}\8\\{ascii\_tok}\KR\6
\X124:Scan an ASCII token\X;\6
\4\&{when}\8\&{others}\8\KR\6
\&{null};\2\6
\4\2\&{end}\8\&{case}\par
\U section~111.\fi

\M116. The occurrence of a module name sets \\{xref\_switch} to zero,
because the module name might (for example) follow \&{type}.

\Y\P\?\4\X116:Scan the module name and make \\{cur\_module} point to it\X\S\?\6
\X118:Put module name into \\{mod\_text}\?(\?1\?\to\?\|k\?)\?\X;\6
\&{if}\1\1\?\8\1(\?\|k\?>\?3\?\2)\ANT\1(\?\\{mod\_text}\?\1(\?\|k\?-\?2\?\to\?%
\|k\?\2)=\1(\?\H{.}\?,\39\?\H{.}\?,\39\?\H{.}\?\2)\2)\?\2\8\&{then}\6
\\{cur\_module}\?\K\?\\{prefix\_lookup}\?\1(\?\|k\?-\?3\?\2)\?;\6
\4\&{else}\6
\\{cur\_module}\?\K\?\\{mod\_lookup}\?\1(\?\|k\?\2)\?;\2\6
\&{end}\8\2\&{if}\1;\6
\\{xref\_switch}\?\K\?0\par
\U section~115.\fi

\M117. Module names are placed into the \\{mod\_text} array with consecutive
spaces,
tabs, and carriage-returns replaced by single spaces. There will be no
spaces at the beginning or the end. (We set \\{mod\_text}\?(\?0\?)\K\?\H{\ } to
facilitate
this, since the \\{mod\_lookup} routine uses \\{mod\_text}\?(\?1\?)\? as the
first
character of the name.)

\Y\P\?\4\X19:Set initial values\X\mathrel{+}\S\?\6
\\{mod\_text}\?\1(\?0\?\2)\K\?\H{\ };\par
\fi

\M118. \P\?\X118:Put module name into \\{mod\_text}\?(\?1\?\to\?\|k\?)\?\X\S\?\6
\|k\?\K\?0;\6
\1\&{loop}\6
\&{if}\1\1\8\\{loc}\?>\?\\{limit}\2\8\&{then}\6
\\{get\_line};\6
\&{if}\1\1\8\\{input\_has\_ended}\2\8\&{then}\6
\\{err\_print}\?\1(\?\.{"!\ Input\ ended\ in\ section\ name"}\?\2)\?;\6
\\{loc}\?\K\?1;\6
\&{exit};\2\6
\&{end}\8\2\&{if}\1;\2\6
\&{end}\8\2\&{if}\1;\6
\|d\?\K\?\\{buffer}\?\1(\?\\{loc}\?\2)\?;\5
\X119:If end of name, \&{exit}\X;\6
\\{incr}\?\1(\?\\{loc}\?\2)\?;\6
\&{if}\1\1\8\|k\?<\?\\{longest\_name}\?-\?1\2\8\&{then}\6
\\{incr}\?\1(\?\|k\?\2)\?;\2\6
\&{end}\8\2\&{if}\1;\6
\&{if}\1\1\?\8\1(\?\|d\?=\?\H{\ }\?\2)\V\1(\?\|d\?=\?\\{tab\_mark}\?\2)\?\2\8%
\&{then}\6
\|d\?\K\?\H{\ };\6
\&{if}\1\1\8\\{mod\_text}\?\1(\?\|k\?-\?1\?\2)=\?\H{\ }\2\8\&{then}\6
\\{decr}\?\1(\?\|k\?\2)\?;\2\6
\&{end}\8\2\&{if}\1;\2\6
\&{end}\8\2\&{if}\1;\6
\\{mod\_text}\?\1(\?\|k\?\2)\K\?\|d;\2\6
\&{end}\8\&{loop};\6
\X120:Check for overlong name\X;\6
\&{if}\1\1\?\8\1(\?\\{mod\_text}\?\1(\?\|k\?\2)=\?\H{\ }\?\2)\W\1(\?\|k\?>\?0\?%
\2)\?\2\8\&{then}\6
\\{decr}\?\1(\?\|k\?\2)\?;\2\6
\&{end}\8\2\&{if}\1\par
\U section~116.\fi

\M119. \P\?\X119:If end of name, \&{exit}\X\S\?\6
\&{if}\1\1\8\|d\?=\?\H{@}\2\8\&{then}\6
\|d\?\K\?\\{buffer}\?\1(\?\\{loc}\?+\?1\?\2)\?;\6
\&{if}\1\1\8\|d\?=\?\H{>}\2\8\&{then}\6
\\{loc}\?\K\?\\{loc}\?+\?2;\6
\&{exit};\2\6
\&{end}\8\2\&{if}\1;\6
\&{if}\1\1\?\8\1(\?\|d\?=\?\H{\ }\?\2)\V\1(\?\|d\?=\?\\{tab\_mark}\?\2)\V\1(\?%
\|d\?=\?\H{*}\?\2)\?\2\8\&{then}\6
\\{err\_print}\?\1(\?\.{"!\ Section\ name\ didn\'t\ end"}\?\2)\?;\6
\&{exit};\2\6
\&{end}\8\2\&{if}\1;\6
\\{incr}\?\1(\?\|k\?\2)\?;\5
\\{mod\_text}\?\1(\?\|k\?\2)\K\?\H{@};\5
\\{incr}\?\1(\?\\{loc}\?\2)\?;\C{now \|d\?=\?\\{buffer}\?(\?\\{loc}\?)\? again}%
\2\6
\&{end}\8\2\&{if}\1\par
\U section~118.\fi

\M120. \P\?\X120:Check for overlong name\X\S\?\6
\&{if}\1\1\8\|k\?\G\?\\{longest\_name}\?-\?2\2\8\&{then}\6
\\{print\_nl}\?\1(\?\.{"!\ Section\ name\ too\ long:\ "}\?\2)\?;\6
\1\&{for}\8\|i\?\in\?1\?\to\?25\8\&{loop}\6
\\{print}\?\1(\?\\{xchr}\?\1(\?\\{mod\_text}\?\1(\?\|i\?\2)\2)\2)\?;\2\6
\&{end}\8\&{loop};\6
\\{print}\?\1(\?\.{"..."}\?\2)\?;\5
\\{mark\_harmless};\2\6
\&{end}\8\2\&{if}\1\par
\U section~118.\fi

\M121. \P\?\X121:Scan to the next \.{@>}\X\S\?\6
\\{id\_first}\?\K\?\\{loc};\5
\\{buffer}\?\1(\?\\{limit}\?+\?1\?\2)\K\?\H{@};\6
\1\&{while}\8\\{buffer}\?\1(\?\\{loc}\?\2)\I\?\H{@}\8\&{loop}\6
\\{incr}\?\1(\?\\{loc}\?\2)\?;\2\6
\&{end}\8\&{loop};\6
\\{id\_loc}\?\K\?\\{loc};\6
\&{if}\1\1\8\\{loc}\?>\?\\{limit}\2\8\&{then}\6
\\{err\_print}\?\1(\?\.{"!\ Control\ text\ didn\'t\ end"}\?\2)\?;\6
\\{loc}\?\K\?\\{limit};\6
\4\&{else}\6
\\{loc}\?\K\?\\{loc}\?+\?2;\6
\&{if}\1\1\8\\{buffer}\?\1(\?\\{loc}\?-\?1\?\2)\I\?\H{>}\2\8\&{then}\6
\\{err\_print}\?\1(\?\.{"!\ Control\ codes\ are\ forbidden\ in\ control\ text"}%
\?\2)\?;\2\6
\&{end}\8\2\&{if}\1;\2\6
\&{end}\8\2\&{if}\1\par
\U section~115.\fi

\M122. A verbatim \ADA\ string respectively an output string will be treated
like
ordinary strings, but with no surrounding delimiters.  At the present point
in the program we have \\{buffer}\?(\?\\{loc}\?-\?1\?)=\?\\{verbatim} or
\\{buffer}\?(\?\\{loc}\?-\?1\?)=\?\\{output\_tok}; we must set \\{id\_first} to
the beginning
of the string itself, and \\{id\_loc} to its ending-plus-one location in the
buffer.  We also set \\{loc} to the position just after the ending delimiter.

\Y\P\?\4\X122:Scan a verbatim or output string\X\S\?\6
\\{id\_first}\?\K\?\\{loc};\5
\\{incr}\?\1(\?\\{loc}\?\2)\?;\5
\\{buffer}\?\1(\?\\{limit}\?+\?1\?\2)\K\?\H{@};\5
\\{buffer}\?\1(\?\\{limit}\?+\?2\?\2)\K\?\H{>};\6
\1\&{while}\?\8\1(\?\\{buffer}\?\1(\?\\{loc}\?\2)\I\?\H{@}\?\2)\V\1(\?%
\\{buffer}\?\1(\?\\{loc}\?+\?1\?\2)\I\?\H{>}\?\2)\?\8\&{loop}\6
\\{incr}\?\1(\?\\{loc}\?\2)\?;\2\6
\&{end}\8\&{loop};\6
\&{if}\1\1\8\\{loc}\?\G\?\\{limit}\2\8\&{then}\6
\\{err\_print}\?\1(\?\.{"!\ Verbatim\ string\ didn\'t\ end"}\?\2)\?;\2\6
\&{end}\8\2\&{if}\1;\6
\\{id\_loc}\?\K\?\\{loc};\5
\\{loc}\?\K\?\\{loc}\?+\?2\par
\U section~115.\fi

\M123. Getting a character token and getting an ASCII token is essentially
the same: \\{id\_first} is set to the position of the first `\.\'' and
\\{id\_loc} to the position after the final `\.\''.

\Y\P\?\4\X123:Get a character token\X\S\?\6
\\{id\_first}\?\K\?\\{loc}\?-\?1;\6
\&{if}\1\1\8\\{buffer}\?\1(\?\\{loc}\?+\?1\?\2)=\?\H{\'}\2\8\&{then}\6
\\{loc}\?\K\?\\{loc}\?+\?2;\5
\\{id\_loc}\?\K\?\\{loc};\5
\|c\?\K\?\\{character\_tok};\6
\4\&{elsif}\1\8\\{buffer}\?\1(\?\\{loc}\?\to\?\\{loc}\?+\?2\?\2)=\1(\?\H{@}\?,%
\39\?\H{@}\?,\39\?\H{\'}\?\2)\?\2\8\&{then}\6
\\{loc}\?\K\?\\{loc}\?+\?3;\5
\\{id\_loc}\?\K\?\\{loc};\5
\|c\?\K\?\\{character\_tok};\2\6
\&{end}\8\2\&{if}\1\par
\U section~111.\fi

\M124. \P\?\X124:Scan an ASCII token\X\S\?\6
\\{id\_first}\?\K\?\\{loc}\?-\?1;\5
\\{incr}\?\1(\?\\{loc}\?\2)\?;\6
\&{if}\1\1\8\\{buffer}\?\1(\?\\{loc}\?\2)\I\?\H{\'}\2\8\&{then}\6
\&{if}\1\1\8\\{buffer}\?\1(\?\\{loc}\?-\?1\?\to\?\\{loc}\?\2)\I\1(\?\H{@}\?,\39%
\?\H{@}\?\2)\?\2\8\&{then}\6
\\{err\_print}\?\1(\?\.{"!\ Improper\ ASCII-code\ sequence,\ \'\ sign\
missing"}\?\2)\?;\6
\4\&{else}\6
\\{incr}\?\1(\?\\{loc}\?\2)\?;\2\6
\&{end}\8\2\&{if}\1;\2\6
\&{end}\8\2\&{if}\1;\6
\\{incr}\?\1(\?\\{loc}\?\2)\?;\5
\\{id\_loc}\?\K\?\\{loc}\par
\U section~115.\fi

\N125.  Phase one processing.
We now have accumulated enough subroutines to make it possible to carry out
\.{AWEAVE}'s first pass over the source file. If everything works right,
both phase one and phase two of \.{AWEAVE} will assign the same numbers to
modules, and these numbers will agree with what \.{ATANGLE} does.

The global variable \\{next\_control} often contains the most recent output of
\\{get\_next}; in interesting cases, this will be the control code that
ended a module or part of a module.

\Y\P\?\4\X14:Types and objects visible from \\{data\_structures}\X\mathrel{+}\S%
\?\6
\\{next\_control}\?:\?\\{eight\_bits};\C{control code waiting to be acting
upon}\par
\fi

\M126. The overall processing strategy in phase one has the following
straightforward outline.

\Y\P\?\4\X126:Phase I: Read all the user's text and store the cross references%
\X\S\?\6
\\{phase\_one}\?\K\?\p{TRUE};\5
\\{phase\_three}\?\K\?\p{FALSE};\5
\\{reset\_input};\5
\\{module\_count}\?\K\?0;\5
\\{skip\_limbo};\5
\\{change\_exists}\?\K\?\p{FALSE};\6
\1\&{while}\?\8\R\?\\{input\_has\_ended}\8\&{loop}\6
\X127:Store cross reference data for the current module\X;\2\6
\&{end}\8\&{loop};\6
\\{changed\_module}\?\1(\?\\{module\_count}\?\2)\K\?\\{change\_exists};\C{the
index changes if anything does}\6
\\{phase\_one}\?\K\?\p{FALSE};\C{prepare for second phase}\6
\X138:Print error messages about unused or undefined module names\X\par
\U section~302.\fi

\M127. \P\?\X127:Store cross reference data for the current module\X\S\?\6
\\{incr}\?\1(\?\\{module\_count}\?\2)\?;\6
\&{if}\1\1\8\\{module\_count}\?=\?\\{max\_modules}\2\8\&{then}\6
\\{overflow}\?\1(\?\.{"section\ number"}\?\2)\?;\2\6
\&{end}\8\2\&{if}\1;\6
\\{changed\_module}\?\1(\?\\{module\_count}\?\2)\K\?\p{FALSE};\C{it will become
\p{TRUE} if any line changes}\6
\&{if}\1\1\8\\{buffer}\?\1(\?\\{loc}\?-\?1\?\2)=\?\H{*}\2\8\&{then}\6
\\{print}\?\1(\?\.{"*"}\?\2)\?;\5
\\{print\_int}\?\1(\?\\{module\_count}\?,\39\?1\?\2)\?;\5
\\{update\_terminal};\C{print a progress report}\2\6
\&{end}\8\2\&{if}\1;\6
\X132:Store cross references in the \TeX\ part of a module\X;\6
\X133:Store cross references in the \(definition part of a module\X;\6
\X135:Store cross references in the \(\ADA\ part of a module\X;\6
\&{if}\1\1\8\\{changed\_module}\?\1(\?\\{module\_count}\?\2)\?\2\8\&{then}\6
\\{change\_exists}\?\K\?\p{TRUE};\2\6
\&{end}\8\2\&{if}\1\par
\U section~126.\fi

\M128. The \\{ADA\_xref} subroutine stores references to identifiers in
\ADA\ text material beginning with the current value of \\{next\_control}
and continuing until \\{next\_control} is `\.\{' or `\v', or until the next
``milestone'' is passed (i.e., \\{next\_control}\?\G\?\\{format}). If
\\{next\_control}\?\G\?\\{format} when \\{ADA\_xref} is called, nothing will
happen;
but if \\{next\_control}\?=\?\H{|} upon entry, the procedure assumes that this
is
the `\v' preceding \ADA\ text that is to be processed.

The program uses the fact that our internal code numbers satisfy
the relations \\{xref\_roman}\?=\?\\{identifier}\?+\?\\{roman}, \\{xref%
\_wildcard}\?=\?\\{identifier}\?+\?\\{wildcard}, $\ldots$, \\{xref\_typewriter}%
\?=\?\\{identifier}\?+\?\\{typewriter},
and \\{normal}\?=\?0.

After \&{entry}, \&{for}, \&{function}, \&{package}, \&{procedure},
\&{subtype}, \&{task}, and \&{type} \\{xref\_switch} is set to \\{def\_flag};
after \&{accept} and \&{body} it is set to \\{body\_flag}.

\Y\P\?\4\X128:Specification of procedures visible from \\{first\_phase}\X\S\?\6
\&{procedure}\8\1\\{ADA\_xref}\2;\par
\A section~301.
\U section~7\*.\fi

\M129. \P\?\X129:Procedures in \\{first\_phase}\X\S\?\6
\&{procedure}\8\1\\{ADA\_xref}\2\8\&{is}\8\C{makes cross references for \ADA\
identifiers}\1\6
\|p\?:\?\\{name\_pointer};\C{a referenced name}\6
\|i\?:\?\p{INTEGER};\C{temporary ilk code of \|p}\6
\4\&{begin}\6
\1\&{while}\8\\{next\_control}\?<\?\\{format}\8\&{loop}\6
\&{if}\1\1\8\\{next\_control}\?\in\?\\{identifier}\?\to\?\\{xref\_predefined}\2%
\8\&{then}\6
\|p\?\K\?\\{id\_lookup}\?\1(\?\\{next\_control}\?-\?\\{identifier}\?\2)\?;\6
\&{if}\1\1\?\8\1(\?\\{ilk}\?\1(\?\|p\?\2)=\?\\{body\_like}\?\2)\V\1(\?\\{ilk}\?%
\1(\?\|p\?\2)=\?\\{type\_like}\?\2)\?\2\8\&{then}\6
\\{xref\_switch}\?\K\?0;\C{no index entry for \&{body} and \&{type}}\2\6
\&{end}\8\2\&{if}\1;\6
\\{new\_xref}\?\1(\?\|p\?\2)\?;\6
\&{if}\1\1\?\8\1(\?\\{ilk}\?\1(\?\|p\?\2)=\?\\{proc\_like}\?\2)\V\1(\?\\{ilk}\?%
\1(\?\|p\?\2)=\?\\{type\_like}\?\2)\?\2\8\&{then}\6
\\{xref\_switch}\?\K\?\\{def\_flag};\C{implied \\{def\_flag}}\6
\4\&{elsif}\1\?\8\1(\?\\{ilk}\?\1(\?\|p\?\2)=\?\\{accept\_like}\?\2)\V\1(\?%
\\{ilk}\?\1(\?\|p\?\2)=\?\\{body\_like}\?\2)\?\2\8\&{then}\6
\\{xref\_switch}\?\K\?\\{body\_flag};\C{implied \\{body\_flag}}\2\6
\&{end}\8\2\&{if}\1;\6
\4\&{elsif}\1\8\\{next\_control}\?=\?\\{xref\_reserved}\2\8\&{then}\6
\|p\?\K\?\\{id\_lookup}\?\1(\?\\{normal}\?\2)\?;\5
\|i\?\K\?\\{ilk}\?\1(\?\|p\?\2)\?;\5
\\{ilk}\?\1(\?\|p\?\2)\K\?\\{normal};\5
\\{new\_xref}\?\1(\?\|p\?\2)\?;\5
\\{ilk}\?\1(\?\|p\?\2)\K\?\|i;\2\6
\&{end}\8\2\&{if}\1;\6
\\{next\_control}\?\K\?\\{get\_next};\6
\&{if}\1\1\?\8\1(\?\\{next\_control}\?=\?\H{|}\?\2)\V\1(\?\\{next\_control}\?=%
\?\H{\{}\?\2)\?\2\8\&{then}\6
\&{return};\2\6
\&{end}\8\2\&{if}\1;\2\6
\&{end}\8\&{loop};\2\6
\&{end}\8\\{ADA\_xref};\par
\A sections~130, 137, and~302.
\U section~7\*.\fi

\M130. The \\{outer\_xref} subroutine is like \\{ADA\_xref} but it begins
with \\{next\_control}\?\I\?\H{|} and ends with \\{next\_control}\?\G\?%
\\{format}. Thus, it
handles \ADA\ text with embedded comments.

\Y\P\?\4\X129:Procedures in \\{first\_phase}\X\mathrel{+}\S\?\6
\&{procedure}\8\1\\{outer\_xref}\2\8\&{is}\8\C{extension of \\{ADA\_xref}}\1\6
\\{bal}\?:\?\\{eight\_bits};\C{brace level in comment}\6
\4\&{begin}\6
\1\&{while}\8\\{next\_control}\?<\?\\{format}\8\&{loop}\6
\&{if}\1\1\8\\{next\_control}\?\I\?\H{\{}\2\8\&{then}\6
\\{ADA\_xref};\6
\4\&{else}\6
\\{bal}\?\K\?\\{skip\_comment}\?\1(\?1\?\2)\?;\5
\\{next\_control}\?\K\?\H{|};\6
\1\&{while}\8\\{bal}\?>\?0\8\&{loop}\6
\\{ADA\_xref};\6
\&{if}\1\1\8\\{next\_control}\?=\?\H{|}\2\8\&{then}\6
\\{bal}\?\K\?\\{skip\_comment}\?\1(\?\\{bal}\?\2)\?;\6
\4\&{else}\6
\\{bal}\?\K\?0;\C{an error will be reported in phase two}\2\6
\&{end}\8\2\&{if}\1;\2\6
\&{end}\8\&{loop};\2\6
\&{end}\8\2\&{if}\1;\2\6
\&{end}\8\&{loop};\2\6
\&{end}\8\\{outer\_xref};\par
\fi

\M131. During the definition and \ADA\ parts of a module, cross references
are made for all identifiers except reserved words; however, the
identifiers in a format definition are referenced even if they are
reserved. The \TeX\ code in comments is, of course, ignored, except for
\ADA\ portions enclosed in \ab; the text of a module name is skipped
entirely, even if it contains \ab\ constructions.

The variables \\{lhs} and \\{rhs} point to the respective identifiers involved
in a format definition.

Additionally to the points just stated a reserved identifier is
referenced by means of the `\.{@r}' control code; the variable
\\{temp\_ilk} contains the \\{ilk} of an identifier during the \\{new\_xref}
subroutine is executing, because we have to set the \\{ilk} to
\\{normal} before the call.

\Y\P\?\4\X14:Types and objects visible from \\{data\_structures}\X\mathrel{+}\S%
\?\6
\\{lhs}\?,\39\?\\{rhs}\?:\?\\{name\_pointer};\C{indices into \\{byte\_start}
for format identifiers}\6
\\{temp\_ilk}\?:\?\\{sixteen\_bits};\C{temporarily contains the \\{ilk} of %
\\{lhs}}\par
\fi

\M132. In the \TeX\ part of a module, cross reference entries are made only for
the identifiers in \ADA\ texts enclosed in \ab, or for control texts
enclosed in \.{@\^}$\,\ldots\,$\.{@>} or \.{@.}$\,\ldots\,$\.{@>}
or \.{@:}$\,\ldots\,$\.{@>}
or \.{@i}$\,\ldots\,$\.{@>}
or \.{@p}$\,\ldots\,$\.{@>}
or \.{@r}$\,\ldots\,$\.{@>}
or \.{@b}$\,\ldots\,$\.{@>}
or \.{@s}$\,\ldots\,$\.{@>}.


\Y\P\?\4\X132:Store cross references in the \TeX\ part of a module\X\S\?\6
\1\&{loop}\6
\\{next\_control}\?\K\?\\{skip\_TeX};\6
\&{case}\1\8\\{next\_control}\1\8\&{is}\8\6
\4\&{when}\8\\{underline}\KR\6
\&{if}\1\1\8\\{xref\_switch}\?<\?\\{body\_flag}\2\8\&{then}\6
\\{xref\_switch}\?\K\?\\{xref\_switch}\?+\?\\{def\_flag};\2\6
\&{end}\8\2\&{if}\1;\6
\4\&{when}\8\\{no\_underline}\KR\6
\&{if}\1\1\8\\{xref\_switch}\?>\?0\2\8\&{then}\6
\\{xref\_switch}\?\K\?\\{xref\_switch}\?-\?\\{def\_flag};\2\6
\&{end}\8\2\&{if}\1;\6
\4\&{when}\8\H{|}\KR\6
\\{ADA\_xref};\6
\4\&{when}\8\\{xref\_roman}\?\to\?\\{xref\_predefined}\?\mathrel{\.|}\?%
\\{module\_name}\KR\6
\\{loc}\?\K\?\\{loc}\?-\?2;\5
\\{next\_control}\?\K\?\\{get\_next};\C{scan to \.{@>}}\6
\&{if}\1\1\8\\{next\_control}\?\I\?\\{module\_name}\2\8\&{then}\6
\\{new\_xref}\?\1(\?\\{id\_lookup}\?\1(\?\\{next\_control}\?-\?\\{identifier}\?%
\2)\2)\?;\2\6
\&{end}\8\2\&{if}\1;\6
\4\&{when}\8\\{xref\_reserved}\KR\6
\\{loc}\?\K\?\\{loc}\?-\?2;\5
\\{next\_control}\?\K\?\\{get\_next};\C{scan to \.{@>}}\6
\\{lhs}\?\K\?\\{id\_lookup}\?\1(\?\\{normal}\?\2)\?;\5
\\{temp\_ilk}\?\K\?\\{ilk}\?\1(\?\\{lhs}\?\2)\?;\C{get the \\{ilk}}\6
\\{ilk}\?\1(\?\\{lhs}\?\2)\K\?\\{normal};\5
\\{new\_xref}\?\1(\?\\{lhs}\?\2)\?;\C{make the cross reference}\6
\\{ilk}\?\1(\?\\{lhs}\?\2)\K\?\\{temp\_ilk};\C{reset the \\{ilk}}\6
\4\&{when}\8\&{others}\8\KR\6
\&{null};\2\6
\4\2\&{end}\8\&{case};\6
\?\&{exit}\8\&{when}\?\8\\{next\_control}\?\G\?\\{format};\2\6
\&{end}\8\&{loop}\par
\U section~127.\fi

\M133. When we get to the following code we have \\{next\_control}\?\G\?%
\\{format}.

\Y\P\?\4\X133:Store cross references in the \(definition part of a module\X\S\?%
\6
\1\&{while}\8\\{next\_control}\?\L\?\\{definition}\8\&{loop}\C{\\{format} or %
\\{definition}}\6
\\{xref\_switch}\?\K\?\\{def\_flag};\C{implied \.{@!}}\6
\&{if}\1\1\8\\{next\_control}\?=\?\\{definition}\2\8\&{then}\6
\\{next\_control}\?\K\?\\{get\_next};\6
\4\&{else}\6
\X134:Process a format definition\X;\2\6
\&{end}\8\2\&{if}\1;\6
\\{outer\_xref};\2\6
\&{end}\8\&{loop}\par
\U section~127.\fi

\M134. Error messages for improper format definitions will be issued in phase
two. Our job in phase one is to define the \\{ilk} of a properly formatted
identifier, and to fool the \\{new\_xref} routine into thinking that the
identifier on the right-hand side of the format definition is not a
reserved word.

\Y\P\?\4\X134:Process a format definition\X\S\?\6
\\{next\_control}\?\K\?\\{get\_next};\6
\&{if}\1\1\8\\{next\_control}\?=\?\\{identifier}\2\8\&{then}\6
\\{lhs}\?\K\?\\{id\_lookup}\?\1(\?\\{normal}\?\2)\?;\5
\\{ilk}\?\1(\?\\{lhs}\?\2)\K\?\\{normal};\5
\\{new\_xref}\?\1(\?\\{lhs}\?\2)\?;\5
\\{next\_control}\?\K\?\\{get\_next};\6
\&{if}\1\1\8\\{next\_control}\?=\?\\{equivalence\_sign}\2\8\&{then}\6
\\{next\_control}\?\K\?\\{get\_next};\6
\&{if}\1\1\8\\{next\_control}\?=\?\\{identifier}\2\8\&{then}\6
\\{rhs}\?\K\?\\{id\_lookup}\?\1(\?\\{normal}\?\2)\?;\5
\\{ilk}\?\1(\?\\{lhs}\?\2)\K\?\\{ilk}\?\1(\?\\{rhs}\?\2)\?;\5
\\{ilk}\?\1(\?\\{rhs}\?\2)\K\?\\{normal};\5
\\{new\_xref}\?\1(\?\\{rhs}\?\2)\?;\5
\\{ilk}\?\1(\?\\{rhs}\?\2)\K\?\\{ilk}\?\1(\?\\{lhs}\?\2)\?;\5
\\{next\_control}\?\K\?\\{get\_next};\2\6
\&{end}\8\2\&{if}\1;\2\6
\&{end}\8\2\&{if}\1;\2\6
\&{end}\8\2\&{if}\1\par
\U section~133.\fi

\M135. Finally, when the \TeX\ and definition parts have been treated, we have
\\{next\_control}\?\G\?\\{begin\_ada}.

\Y\P\?\4\X135:Store cross references in the \(\ADA\ part of a module\X\S\?\6
\&{if}\1\1\8\\{next\_control}\?\L\?\\{module\_name}\2\8\&{then}\C{\\{begin%
\_ada} or \\{module\_name}}\6
\&{if}\1\1\8\\{next\_control}\?=\?\\{begin\_ada}\2\8\&{then}\6
\\{mod\_xref\_switch}\?\K\?0;\6
\4\&{else}\6
\\{mod\_xref\_switch}\?\K\?\\{def\_flag};\2\6
\&{end}\8\2\&{if}\1;\6
\1\&{loop}\6
\&{if}\1\1\8\\{next\_control}\?=\?\\{module\_name}\2\8\&{then}\6
\\{new\_mod\_xref}\?\1(\?\\{cur\_module}\?\2)\?;\2\6
\&{end}\8\2\&{if}\1;\6
\\{next\_control}\?\K\?\\{get\_next};\6
\\{outer\_xref};\6
\?\&{exit}\8\&{when}\?\8\\{next\_control}\?>\?\\{module\_name};\2\6
\&{end}\8\&{loop};\2\6
\&{end}\8\2\&{if}\1\par
\U section~127.\fi

\M136. After phase one has looked at everything, we want to check that each
module name was both defined and used.
The variable \\{cur\_xref} will point to cross references for the
current module name of interest.

\Y\P\?\4\X14:Types and objects visible from \\{data\_structures}\X\mathrel{+}\S%
\?\6
\\{cur\_xref}\?:\?\\{xref\_number};\C{temporary cross reference pointer}\par
\fi

\M137. The following recursive procedure
walks through the tree of module names and prints out anomalies.

\Y\P\?\4\X129:Procedures in \\{first\_phase}\X\mathrel{+}\S\?\6
\&{procedure}\8\1\\{mod\_check}\?(\1\?\|p\?:\?\\{name\_pointer}\?\2)\?\2\8%
\&{is}\8\C{print anomalies in subtree \|p}\1\6
\4\&{begin}\6
\&{if}\1\1\8\|p\?>\?0\2\8\&{then}\6
\\{mod\_check}\?\1(\?\\{llink}\?\1(\?\|p\?\2)\2)\?;\6
\\{cur\_xref}\?\K\?\\{xref}\?\1(\?\|p\?\2)\?;\6
\&{if}\1\1\8\\{num}\?\1(\?\\{cur\_xref}\?\2)<\?\\{def\_flag}\2\8\&{then}\6
\\{print\_nl}\?\1(\?\.{"!\ Never\ defined:\ <"}\?\2)\?;\5
\\{print\_id}\?\1(\?\|p\?\2)\?;\5
\\{print}\?\1(\?\.{">"}\?\2)\?;\5
\\{mark\_harmless};\2\6
\&{end}\8\2\&{if}\1;\6
\1\&{while}\8\\{num}\?\1(\?\\{cur\_xref}\?\2)\G\?\\{def\_flag}\8\&{loop}\6
\\{cur\_xref}\?\K\?\\{xlink}\?\1(\?\\{cur\_xref}\?\2)\?;\2\6
\&{end}\8\&{loop};\6
\&{if}\1\1\8\\{cur\_xref}\?=\?0\2\8\&{then}\6
\\{print\_nl}\?\1(\?\.{"!\ Never\ used:\ <"}\?\2)\?;\5
\\{print\_id}\?\1(\?\|p\?\2)\?;\5
\\{print}\?\1(\?\.{">"}\?\2)\?;\5
\\{mark\_harmless};\2\6
\&{end}\8\2\&{if}\1;\6
\\{mod\_check}\?\1(\?\\{rlink}\?\1(\?\|p\?\2)\2)\?;\2\6
\&{end}\8\2\&{if}\1;\2\6
\&{end}\8\\{mod\_check};\par
\fi

\M138. \P\?\X138:Print error messages about unused or undefined module names\X%
\S\?\8\\{mod\_check}\?\1(\?\\{root}\?\2)\?\par
\U section~126.\fi

\N139.  Low-level output routines.
The \TeX\ output is supposed to appear in lines at most \\{line\_length}
characters long, so we place it into an output buffer. During the output
process, \\{out\_line} will hold the current line number of the line about to
be output.

\Y\P\?\4\X14:Types and objects visible from \\{data\_structures}\X\mathrel{+}\S%
\?\6
\&{subtype}\8\\{length\_index}\8\&{is}\8\p{INTEGER}\8\&{range}\80\?\to\?\\{line%
\_length};\6
\\{out\_buf}\?:\1\&{array}\8\1(\?\\{length\_index}\?\2)\2\8\&{of}\?\8\\{ASCII%
\_code};\C{assembled characters}\6
\\{out\_ptr}\?:\?\\{length\_index};\C{number of characters in \\{out\_buf}}\6
\\{out\_line}\?:\?\p{INTEGER};\C{coordinates of next line to be output}\par
\fi

\M140. The \\{flush\_buffer} routine empties the buffer up to a given
breakpoint,
and moves any remaining characters to the beginning of the next line.
If the \\{per\_cent} parameter is \p{TRUE}, a \H{\%} is appended to the line
that is being output; in this case the breakpoint \|b should be strictly
less than \\{line\_length}. If the \\{per\_cent} parameter is \p{FALSE},
trailing blanks are suppressed.
The characters emptied from the buffer form a new line of output.

\Y\P\?\4\X140:Specification of procedures visible from \\{output}\X\S\?\6
\&{procedure}\8\1\\{flush\_buffer}\?(\?\|b\?:\?\\{eight\_bits};\?\,\35\1\?%
\\{per\_cent}\?:\?\p{BOOLEAN}\?\2)\?\2;\par
\A sections~142, 147, 151, 154, 156, 159, and~162.
\U section~8\*.\fi

\M141. \P\?\X141:Procedures in \\{output}\X\S\?\6
\&{procedure}\8\1\\{flush\_buffer}\?(\?\|b\?:\?\\{eight\_bits};\?\,\35\1\?%
\\{per\_cent}\?:\?\p{BOOLEAN}\?\2)\?\2\8\&{is}\8\C{outputs \\{out\_buf}\?(\?1\?%
\to\?\|b\?)\?, where \|b\?\L\?\\{out\_ptr}}\1\6
\|j\?:\?\\{length\_index}\?\K\?\|b;\6
\4\&{begin}\6
\&{if}\1\1\?\8\R\?\\{per\_cent}\2\8\&{then}\C{remove trailing blanks}\6
\1\&{loop}\6
\?\&{exit}\8\&{when}\8\1(\?\|j\?=\?0\?\2)\ORE\1(\?\\{out\_buf}\?\1(\?\|j\?\2)\I%
\?\H{\ }\?\2)\?;\6
\\{decr}\?\1(\?\|j\?\2)\?;\2\6
\&{end}\8\&{loop};\2\6
\&{end}\8\2\&{if}\1;\6
\1\&{for}\8\|k\?\in\?1\?\to\?\|j\8\&{loop}\6
\p{TEXT\_IO}\?.\?\p{PUT}\?\1(\?\\{TeX\_file}\?,\39\?\\{xchr}\?\1(\?\\{out\_buf}%
\?\1(\?\|k\?\2)\2)\2)\?;\2\6
\&{end}\8\&{loop};\6
\&{if}\1\1\8\\{per\_cent}\2\8\&{then}\6
\p{TEXT\_IO}\?.\?\p{PUT}\?\1(\?\\{TeX\_file}\?,\39\?\\{xchr}\?\1(\?\H{\%}\?\2)%
\2)\?;\2\6
\&{end}\8\2\&{if}\1;\6
\p{TEXT\_IO}\?.\?\p{NEW\_LINE}\?\1(\?\\{TeX\_file}\?\2)\?;\5
\\{incr}\?\1(\?\\{out\_line}\?\2)\?;\6
\&{if}\1\1\8\|b\?<\?\\{out\_ptr}\2\8\&{then}\6
\\{out\_buf}\?\1(\?1\?\to\?\\{out\_ptr}\?-\?\|b\?\2)\K\?\\{out\_buf}\?\1(\?\|b%
\?+\?1\?\to\?\\{out\_ptr}\?\2)\?;\2\6
\&{end}\8\2\&{if}\1;\6
\\{out\_ptr}\?\K\?\\{out\_ptr}\?-\?\|b;\2\6
\&{end}\8\\{flush\_buffer};\par
\A sections~143, 148, 152, 155, 157, 160, and~163.
\U section~8\*.\fi

\M142. When we are copying \TeX\ source material, we retain line breaks
that occur in the input, except that an empty line is not
output when the \TeX\ source line was nonempty. For example, a line
of the \TeX\ file that contains only an index cross-reference entry
will not be copied. The \\{finish\_line} routine is called just before
\\{get\_line} inputs a new line, and just after a line break token has
been emitted during the output of translated \ADA\ text.

\Y\P\?\4\X140:Specification of procedures visible from \\{output}\X\mathrel{+}%
\S\?\6
\&{procedure}\8\1\\{finish\_line}\2;\par
\fi

\M143. \P\?\X141:Procedures in \\{output}\X\mathrel{+}\S\?\6
\&{procedure}\8\1\\{finish\_line}\2\8\&{is}\8\C{do this at the end of a line}\1%
\6
\4\&{begin}\6
\&{if}\1\1\8\\{out\_ptr}\?>\?0\2\8\&{then}\6
\\{flush\_buffer}\?\1(\?\\{out\_ptr}\?,\39\?\p{FALSE}\?\2)\?;\6
\4\&{else}\6
\1\&{for}\8\|k\?\in\?0\?\to\?\\{limit}\8\&{loop}\6
\&{if}\1\1\?\8\1(\?\\{buffer}\?\1(\?\|k\?\2)\I\?\H{\ }\?\2)\W\1(\?\\{buffer}\?%
\1(\?\|k\?\2)\I\?\\{tab\_mark}\?\2)\?\2\8\&{then}\6
\&{return};\2\6
\&{end}\8\2\&{if}\1;\2\6
\&{end}\8\&{loop};\6
\\{flush\_buffer}\?\1(\?0\?,\39\?\p{FALSE}\?\2)\?;\2\6
\&{end}\8\2\&{if}\1;\2\6
\&{end}\8\\{finish\_line};\par
\fi

\M144. In particular, the \\{finish\_line} procedure is called near the very
beginning of phase two. We initialize the output variables in a slightly
tricky way so that the first line of the output file will be
`\.{\\input awebmac}'.

\Y\P\?\4\X19:Set initial values\X\mathrel{+}\S\?\6
\\{out\_ptr}\?\K\?1;\5
\\{out\_line}\?\K\?1;\5
\\{out\_buf}\?\1(\?1\?\2)\K\?\H{c};\5
\p{TEXT\_IO}\?.\?\p{PUT}\?\1(\?\\{TeX\_file}\?,\39\?\.{"\\input\ awebma"}\?\2)%
\?;\par
\fi

\M145. When we wish to append the character \|c to the output buffer, we write
`$\\{out1}(c)$'; this will cause the buffer to be emptied if it was already
full. Similarly, `$\\{out2}(c_1)(c_2)$' appends a pair of characters.
A line break will occur at a space or after a single-nonletter
\TeX\ control sequence.

\Y\P\4\D\\{out1}\?\1(\?\#\?\2)\S\?\6
\&{if}\1\1\8\\{out\_ptr}\?=\?\\{line\_length}\2\8\&{then}\6
\\{break\_out};\2\6
\&{end}\8\2\&{if}\1;\6
\\{incr}\?\1(\?\\{out\_ptr}\?\2)\?;\5
\\{out\_buf}\?\1(\?\\{out\_ptr}\?\2)\K\?\#\par
\P\4\D\\{out2}\?\1(\?\#\?\2)\S\?\\{out1}\?\1(\?\#\?\2)\?;\5
\\{out1}\par
\P\4\D\\{out3}\?\1(\?\#\?\2)\S\?\\{out1}\?\1(\?\#\?\2)\?;\5
\\{out2}\par
\P\4\D\\{out4}\?\1(\?\#\?\2)\S\?\\{out1}\?\1(\?\#\?\2)\?;\5
\\{out3}\par
\P\4\D\\{out5}\?\1(\?\#\?\2)\S\?\\{out1}\?\1(\?\#\?\2)\?;\5
\\{out4}\par
\fi

\M146. The \\{break\_out} routine is called just before the output buffer is
about
to overflow. To make this routine a little faster, we initialize position
0 of the output buffer to `\.\\'; this character isn't really output.

\Y\P\?\4\X19:Set initial values\X\mathrel{+}\S\?\6
\\{out\_buf}\?\1(\?0\?\2)\K\?\H{\\};\par
\fi

\M147. A long line is broken at a blank space or just before a backslash that
isn't
preceded by another backslash. In the latter case, a \H{\%} is output at
the break.

\Y\P\?\4\X140:Specification of procedures visible from \\{output}\X\mathrel{+}%
\S\?\6
\&{procedure}\8\1\\{break\_out}\2;\par
\fi

\M148. \P\?\X141:Procedures in \\{output}\X\mathrel{+}\S\?\6
\&{procedure}\8\1\\{break\_out}\2\8\&{is}\8\C{finds a way to break the output
line}\1\6
\|k\?:\?\\{length\_index}\?\K\?\\{out\_ptr};\C{index into \\{out\_buf}}\6
\|c\?,\39\?\|d\?:\?\\{ASCII\_code};\C{characters from the buffer}\6
\4\&{begin}\6
\1\&{loop}\6
\&{if}\1\1\8\|k\?=\?0\2\8\&{then}\6
\X149:Print warning message, break the line, \&{return}\X;\2\6
\&{end}\8\2\&{if}\1;\6
\|d\?\K\?\\{out\_buf}\?\1(\?\|k\?\2)\?;\6
\&{if}\1\1\8\|d\?=\?\H{\ }\2\8\&{then}\6
\\{flush\_buffer}\?\1(\?\|k\?,\39\?\p{FALSE}\?\2)\?;\5
\&{return};\2\6
\&{end}\8\2\&{if}\1;\6
\&{if}\1\1\?\8\1(\?\|d\?=\?\H{\\}\?\2)\W\1(\?\\{out\_buf}\?\1(\?\|k\?-\?1\?\2)%
\I\?\H{\\}\?\2)\?\2\8\&{then}\C{in this case \|k\?>\?1}\6
\\{flush\_buffer}\?\1(\?\|k\?-\?1\?,\39\?\p{TRUE}\?\2)\?;\5
\&{return};\2\6
\&{end}\8\2\&{if}\1;\6
\\{decr}\?\1(\?\|k\?\2)\?;\2\6
\&{end}\8\&{loop};\2\6
\&{end}\8\\{break\_out};\par
\fi

\M149. We get to this module only in unusual cases that the entire output line
consists of a string of backslashes followed by a string of nonblank
non-backslashes. In such cases it is almost always safe to break the
line by putting a \H{\%} just before the last character.

\Y\P\?\4\X149:Print warning message, break the line, \&{return}\X\S\?\6
\\{print\_nl}\?\1(\?\.{"!\ Line\ had\ to\ be\ broken\ (output\ l."}\?\2)\?;\5
\\{print\_int}\?\1(\?\\{out\_line}\?,\39\?1\?\2)\?;\5
\\{print\_ln}\?\1(\?\.{"):"}\?\2)\?;\6
\1\&{for}\8\|j\?\in\?1\?\to\?\\{out\_ptr}\?-\?1\8\&{loop}\6
\\{print}\?\1(\?\\{xchr}\?\1(\?\\{out\_buf}\?\1(\?\|j\?\2)\2)\2)\?;\2\6
\&{end}\8\&{loop};\6
\\{new\_ln};\5
\\{mark\_harmless};\5
\\{flush\_buffer}\?\1(\?\\{out\_ptr}\?-\?1\?,\39\?\p{TRUE}\?\2)\?;\5
\&{return}\par
\U section~148.\fi

\M150. Here is a procedure that outputs a module number in decimal notation.

\Y\P\?\4\X14:Types and objects visible from \\{data\_structures}\X\mathrel{+}\S%
\?\6
\\{dig}\?:\1\&{array}\8\1(\?0\?\to\?4\?\2)\2\8\&{of}\?\8\p{INTEGER}\8\&{range}%
\80\?\to\?9;\C{digits to output}\par
\fi

\M151. The number to be converted by \\{out\_mod} is known to be less than
\\{def\_flag}, so it cannot have more than five decimal digits.  If
the module is changed, we output `\.{\\*}' just after the number and if
\|c\?>\?1 we output `$\.(\,c\,\.)$'.


\Y\P\?\4\X140:Specification of procedures visible from \\{output}\X\mathrel{+}%
\S\?\6
\&{procedure}\8\1\\{out\_mod}\?(\?\|m\?:\?\p{INTEGER};\?\,\35\1\?\|c\?:\?%
\p{INTEGER}\?\K\?1\?\2)\?\2;\par
\fi

\M152. \P\?\X141:Procedures in \\{output}\X\mathrel{+}\S\?\6
\&{procedure}\8\1\\{out\_mod}\?(\?\|m\?:\?\p{INTEGER};\?\,\35\1\?\|c\?:\?%
\p{INTEGER}\?\K\?1\?\2)\?\2\8\&{is}\8\C{output a module number}\1\6
\|k\?:\?\p{INTEGER}\8\&{range}\80\?\to\?5;\C{index into \\{dig}}\6
\|a\?:\?\p{INTEGER};\C{accumulator}\6
\4\&{begin}\6
\|a\?\K\?\|m;\5
\X153:Output the digits of \|a\X;\6
\&{if}\1\1\8\\{changed\_module}\?\1(\?\|m\?\2)\?\2\8\&{then}\6
\\{out2}\?\1(\?\H{\\}\?\2)\1(\?\H{*}\?\2)\?;\2\6
\&{end}\8\2\&{if}\1;\6
\&{if}\1\1\8\|c\?>\?1\2\8\&{then}\6
\\{out1}\?\1(\?\H{(}\?\2)\?;\5
\|a\?\K\?\|c;\5
\X153:Output the digits of \|a\X;\6
\\{out1}\?\1(\?\H{)}\?\2)\?;\2\6
\&{end}\8\2\&{if}\1;\2\6
\&{end}\8\\{out\_mod};\par
\fi

\M153. Here we output the digits of \|a using the \\{dig} array.
\Y\P\?\4\X153:Output the digits of \|a\X\S\?\6
\|k\?\K\?0;\6
\1\&{loop}\6
\\{dig}\?\1(\?\|k\?\2)\K\?\|a\?\mathbin{\&{mod}}\?10;\5
\|a\?\K\?\|a\?/\?10;\5
\\{incr}\?\1(\?\|k\?\2)\?;\6
\?\&{exit}\8\&{when}\?\8\|a\?=\?0;\2\6
\&{end}\8\&{loop};\6
\1\&{loop}\6
\\{decr}\?\1(\?\|k\?\2)\?;\5
\\{out1}\?\1(\?\\{dig}\?\1(\?\|k\?\2)+\?\H{0}\?\2)\?;\6
\?\&{exit}\8\&{when}\?\8\|k\?=\?0;\2\6
\&{end}\8\&{loop}\par
\U section~152(2).\fi

\M154. The \\{out\_name} subroutine is used to output an identifier or index
entry, enclosing it in braces.

\Y\P\?\4\X140:Specification of procedures visible from \\{output}\X\mathrel{+}%
\S\?\6
\&{procedure}\8\1\\{out\_name}\?(\1\?\|p\?:\?\\{name\_pointer}\?\2)\?\2;\par
\fi

\M155. \P\?\X141:Procedures in \\{output}\X\mathrel{+}\S\?\6
\&{procedure}\8\1\\{out\_name}\?(\1\?\|p\?:\?\\{name\_pointer}\?\2)\?\2\8\&{is}%
\8\C{outputs a name}\1\6
\|w\?:\?\\{w\_range};\C{row of \\{byte\_mem}}\6
\4\&{begin}\6
\\{out1}\?\1(\?\H{\{}\?\2)\?;\5
\|w\?\K\?\|p\?\mathbin{\&{mod}}\?\\{ww};\6
\1\&{for}\8\|k\?\in\?\\{byte\_start}\?\1(\?\|p\?\2)\to\?\\{byte\_start}\?\1(\?%
\|p\?+\?\\{ww}\?\2)-\?1\8\&{loop}\6
\&{if}\1\1\8\\{byte\_mem}\?\1(\?\|w\?,\39\?\|k\?\2)=\?\H{\_}\2\8\&{then}\6
\\{out1}\?\1(\?\H{\\}\?\2)\?;\2\6
\&{end}\8\2\&{if}\1;\6
\\{out1}\?\1(\?\\{byte\_mem}\?\1(\?\|w\?,\39\?\|k\?\2)\2)\?;\2\6
\&{end}\8\&{loop};\6
\\{out1}\?\1(\?\H{\}}\?\2)\?;\2\6
\&{end}\8\\{out\_name};\par
\fi

\N156.  Routines that copy \TeX\ material.
During phase two, we use the subroutines \\{copy\_limbo}, \\{copy\_TeX}, and
\\{copy\_comment} in place of the analogous \\{skip\_limbo}, \\{skip\_TeX}, and
\\{skip\_comment} that were used in phase one.

The \\{copy\_limbo} routine, for example, takes \TeX\ material that is not
part of any module and transcribes it almost verbatim to the output file.
No `\.{@}' signs should occur in such material except in `\.{@@}'
pairs; such pairs are replaced by singletons.

\Y\P\?\4\X140:Specification of procedures visible from \\{output}\X\mathrel{+}%
\S\?\6
\&{procedure}\8\1\\{copy\_limbo}\2;\par
\fi

\M157. \P\?\X141:Procedures in \\{output}\X\mathrel{+}\S\?\6
\&{procedure}\8\1\\{copy\_limbo}\2\8\&{is}\8\C{copy \TeX\ code until the next
module begins}\1\6
\|c\?:\?\\{ASCII\_code};\C{character following \.{@} sign}\6
\4\&{begin}\6
\1\&{loop}\6
\&{if}\1\1\8\\{loc}\?>\?\\{limit}\2\8\&{then}\6
\\{finish\_line};\5
\\{get\_line};\6
\&{if}\1\1\8\\{input\_has\_ended}\2\8\&{then}\6
\&{return};\2\6
\&{end}\8\2\&{if}\1;\6
\4\&{else}\6
\\{buffer}\?\1(\?\\{limit}\?+\?1\?\2)\K\?\H{@};\6
\X158:Copy up to control code, \&{return} if finished\X;\2\6
\&{end}\8\2\&{if}\1;\2\6
\&{end}\8\&{loop};\2\6
\&{end}\8\\{copy\_limbo};\par
\fi

\M158. \P\?\X158:Copy up to control code, \&{return} if finished\X\S\?\6
\1\&{while}\8\\{buffer}\?\1(\?\\{loc}\?\2)\I\?\H{@}\8\&{loop}\6
\\{out1}\?\1(\?\\{buffer}\?\1(\?\\{loc}\?\2)\2)\?;\5
\\{incr}\?\1(\?\\{loc}\?\2)\?;\2\6
\&{end}\8\&{loop};\6
\&{if}\1\1\8\\{loc}\?\L\?\\{limit}\2\8\&{then}\6
\\{loc}\?\K\?\\{loc}\?+\?2;\5
\|c\?\K\?\\{buffer}\?\1(\?\\{loc}\?-\?1\?\2)\?;\6
\&{if}\1\1\?\8\1(\?\|c\?=\?\H{\ }\?\2)\V\1(\?\|c\?=\?\\{tab\_mark}\?\2)\V\1(\?%
\|c\?=\?\H{*}\?\2)\?\2\8\&{then}\6
\&{return};\2\6
\&{end}\8\2\&{if}\1;\6
\&{if}\1\1\?\8\1(\?\|c\?\I\?\H{z}\?\2)\W\1(\?\|c\?\I\?\H{Z}\?\2)\?\2\8\&{then}\6
\\{out1}\?\1(\?\H{@}\?\2)\?;\6
\&{if}\1\1\8\|c\?\I\?\H{@}\2\8\&{then}\6
\\{err\_print}\?\1(\?\.{"!\ Double\ @\ required\ outside\ of\ sections"}\?\2)%
\?;\2\6
\&{end}\8\2\&{if}\1;\2\6
\&{end}\8\2\&{if}\1;\2\6
\&{end}\8\2\&{if}\1\par
\U section~157.\fi

\M159. The \\{copy\_TeX} routine processes the \TeX\ code at the beginning of a
module; for example, the words you are now reading were copied in this
way. It returns the next control code or `\v' found in the input.

\Y\P\?\4\X140:Specification of procedures visible from \\{output}\X\mathrel{+}%
\S\?\6
\&{function}\8\1\\{copy\_TeX}\?\8\&{return}\?\8\\{eight\_bits}\2;\par
\fi

\M160. \P\?\X141:Procedures in \\{output}\X\mathrel{+}\S\?\6
\&{function}\8\1\\{copy\_TeX}\?\8\&{return}\?\8\\{eight\_bits}\2\8\&{is}\8%
\C{copy pure \TeX\ material}\1\6
\|c\?:\?\\{eight\_bits};\C{control code found}\6
\4\&{begin}\6
\1\&{loop}\6
\&{if}\1\1\8\\{loc}\?>\?\\{limit}\2\8\&{then}\6
\\{finish\_line};\5
\\{get\_line};\6
\&{if}\1\1\8\\{input\_has\_ended}\2\8\&{then}\6
\?\&{return}\?\8\\{new\_module};\2\6
\&{end}\8\2\&{if}\1;\2\6
\&{end}\8\2\&{if}\1;\6
\\{buffer}\?\1(\?\\{limit}\?+\?1\?\2)\K\?\H{@};\6
\X161:Copy up to `\v' or control code, \?\&{return}\?\8\|c if finished\X;\2\6
\&{end}\8\&{loop};\2\6
\&{end}\8\\{copy\_TeX};\par
\fi

\M161. We don't copy spaces or tab marks into the beginning of a line. This
makes the test for empty lines in \\{finish\_line} work.

\Y\P\?\4\X161:Copy up to `\v' or control code, \?\&{return}\?\8\|c if finished%
\X\S\?\6
\1\&{loop}\6
\|c\?\K\?\\{buffer}\?\1(\?\\{loc}\?\2)\?;\5
\\{incr}\?\1(\?\\{loc}\?\2)\?;\6
\&{if}\1\1\8\|c\?=\?\H{|}\2\8\&{then}\6
\?\&{return}\?\8\H{|};\2\6
\&{end}\8\2\&{if}\1;\6
\&{if}\1\1\8\|c\?\I\?\H{@}\2\8\&{then}\6
\\{out1}\?\1(\?\|c\?\2)\?;\6
\&{if}\1\1\?\8\1(\?\\{out\_ptr}\?=\?1\?\2)\W\1(\1(\?\|c\?=\?\H{\ }\?\2)\V\1(\?%
\|c\?=\?\\{tab\_mark}\?\2)\2)\?\2\8\&{then}\6
\\{decr}\?\1(\?\\{out\_ptr}\?\2)\?;\2\6
\&{end}\8\2\&{if}\1;\2\6
\&{end}\8\2\&{if}\1;\6
\?\&{exit}\8\&{when}\?\8\|c\?=\?\H{@};\2\6
\&{end}\8\&{loop};\6
\&{if}\1\1\8\\{loc}\?\L\?\\{limit}\2\8\&{then}\6
\|c\?\K\?\\{control\_code}\?\1(\?\\{buffer}\?\1(\?\\{loc}\?\2)\2)\?;\5
\\{incr}\?\1(\?\\{loc}\?\2)\?;\5
\?\&{return}\?\8\|c;\2\6
\&{end}\8\2\&{if}\1\par
\U section~160.\fi

\M162. The \\{copy\_comment} uses and returns a brace-balance value, following
the
conventions of \\{skip\_comment} above. Instead of copying the \TeX\ material
into the output buffer, this procedure copies it into the token memory.
The abbreviation \\{app\_tok}\?(\?\|t\?)\? is used to append token \|t to the
current
token list, and it also makes sure that it is possible to append at least
one further token without overflow.

\Y\P\4\D\\{app\_tok}\?\1(\?\#\?\2)\S\?\6
\&{if}\1\1\8\\{tok\_ptr}\?+\?2\?>\?\\{max\_toks}\2\8\&{then}\6
\\{overflow}\?\1(\?\.{"token"}\?\2)\?;\2\6
\&{end}\8\2\&{if}\1;\6
\\{tok\_mem}\?\1(\?\\{tok\_ptr}\?\2)\K\?\#;\5
\\{incr}\?\1(\?\\{tok\_ptr}\?\2)\?\par
\Y\P\?\4\X140:Specification of procedures visible from \\{output}\X\mathrel{+}%
\S\?\6
\&{function}\8\1\\{copy\_comment}\?(\1\?\\{bal}\?:\?\\{eight\_bits}\?\2)%
\&{return}\?\8\\{eight\_bits}\2;\par
\fi

\M163. \P\?\X141:Procedures in \\{output}\X\mathrel{+}\S\?\6
\&{function}\8\1\\{copy\_comment}\?(\1\?\\{bal}\?:\?\\{eight\_bits}\?\2)%
\&{return}\?\8\\{eight\_bits}\2\8\&{is}\8\C{copies \TeX\ code in comments}\1\6
\|c\?:\?\\{ASCII\_code};\C{current character being copied}\6
\|b\?:\?\\{eight\_bits}\?\K\?\\{bal};\C{brace-balance value}\6
\4\&{begin}\6
\1\&{loop}\6
\&{if}\1\1\8\\{loc}\?>\?\\{limit}\2\8\&{then}\6
\\{get\_line};\6
\&{if}\1\1\8\\{input\_has\_ended}\2\8\&{then}\6
\\{err\_print}\?\1(\?\.{"!\ Input\ ended\ in\ mid-comment"}\?\2)\?;\6
\\{loc}\?\K\?1;\5
\X165:Clear \|b and \?\&{return}\?\80\X;\2\6
\&{end}\8\2\&{if}\1;\2\6
\&{end}\8\2\&{if}\1;\6
\|c\?\K\?\\{buffer}\?\1(\?\\{loc}\?\2)\?;\5
\\{incr}\?\1(\?\\{loc}\?\2)\?;\6
\&{if}\1\1\8\|c\?=\?\H{|}\2\8\&{then}\6
\?\&{return}\?\8\|b;\2\6
\&{end}\8\2\&{if}\1;\6
\\{app\_tok}\?\1(\?\|c\?\2)\?;\6
\X164:Copy special things when \|c\?=\?\H{@}\?,\?\H{\\}\?,\?\H{\{}\?,\?\H{\}}; %
\&{return} at end\X;\2\6
\&{end}\8\&{loop};\2\6
\&{end}\8\\{copy\_comment};\par
\fi

\M164. \P\?\X164:Copy special things when \|c\?=\?\H{@}\?,\?\H{\\}\?,\?\H{\{}%
\?,\?\H{\}}; \&{return} at end\X\S\?\6
\&{if}\1\1\8\|c\?=\?\H{@}\2\8\&{then}\6
\\{incr}\?\1(\?\\{loc}\?\2)\?;\6
\&{if}\1\1\8\\{buffer}\?\1(\?\\{loc}\?-\?1\?\2)\I\?\H{@}\2\8\&{then}\6
\\{err\_print}\?\1(\?\.{"!\ Illegal\ use\ of\ @\ in\ comment"}\?\2)\?;\6
\\{loc}\?\K\?\\{loc}\?-\?2;\5
\\{decr}\?\1(\?\\{tok\_ptr}\?\2)\?;\5
\X165:Clear \|b and \?\&{return}\?\80\X;\2\6
\&{end}\8\2\&{if}\1;\6
\4\&{elsif}\1\?\8\1(\?\|c\?=\?\H{\\}\?\2)\W\1(\?\\{buffer}\?\1(\?\\{loc}\?\2)\I%
\?\H{@}\?\2)\?\2\8\&{then}\6
\\{app\_tok}\?\1(\?\\{buffer}\?\1(\?\\{loc}\?\2)\2)\?;\5
\\{incr}\?\1(\?\\{loc}\?\2)\?;\6
\4\&{elsif}\1\8\|c\?=\?\H{\{}\2\8\&{then}\6
\\{incr}\?\1(\?\|b\?\2)\?;\6
\4\&{elsif}\1\8\|c\?=\?\H{\}}\2\8\&{then}\6
\\{decr}\?\1(\?\|b\?\2)\?;\6
\&{if}\1\1\8\|b\?=\?0\2\8\&{then}\6
\?\&{return}\?\80;\2\6
\&{end}\8\2\&{if}\1;\2\6
\&{end}\8\2\&{if}\1\par
\U section~163.\fi

\M165. When the comment has terminated abruptly due to an error, we output
enough right braces to keep \TeX\ happy.

\Y\P\?\4\X165:Clear \|b and \?\&{return}\?\80\X\S\?\6
\\{app\_tok}\?\1(\?\H{\ }\?\2)\?;\C{this is done in case the previous character
was `\.\\'}\6
\1\&{while}\8\|b\?>\?0\8\&{loop}\6
\\{app\_tok}\?\1(\?\H{\}}\?\2)\?;\5
\\{decr}\?\1(\?\|b\?\2)\?;\2\6
\&{end}\8\&{loop};\6
\?\&{return}\?\80\par
\U sections~163 and~164.\fi

\N166.  Parsing.
The most intricate part of \.{AWEAVE} is its mechanism for converting
\ADA-like code into \TeX\ code, and we might as well plunge into this
aspect of the program now. A ``bottom up'' approach is used to parse the
\ADA-like material, since \.{AWEAVE} must deal with fragmentary
constructions whose overall ``part of speech'' is not known.

At the lowest level, the input is represented as a sequence of entities
that we shall call {\it scraps}, where each scrap of information consists
of two parts, its {\it category} and its {\it translation}. The category
is essentially a syntactic class, and the translation is a token list that
represents \TeX\ code. Rules of syntax and semantics tell us how to
combine adjacent scraps into larger ones, and if we are lucky an entire
\ADA\ text that starts out as hundreds of small scraps will join
together into one gigantic scrap whose translation is the desired \TeX\
code. If we are unlucky, we will be left with several scraps that don't
combine; their translations will simply be output, one by one.

The combination rules are given as context-sensitive productions that are
applied from left to right. Suppose that we are currently working on the
sequence of scraps $s_1\,s_2\ldots s_n$. We try first to find the longest
production that applies to an initial substring $s_1\,s_2\ldots\,$; but if
no such productions exist, we find to find the longest production
applicable to the next substring $s_2\,s_3\ldots\,$; and if that fails, we
try to match $s_3\,s_4\ldots\,$, etc.

A production applies if the category codes have a given pattern. For
example, one of the productions is
$$\\{open\_scrap}\;\\{decl}\;\RA\;\\{open\_scrap}$$
and it means that two consecutive scraps whose respective categories are
\\{open\_scrap} and \\{decl} are con\-verted to one scrap whose category
is \\{open\_scrap}. This production also has an associated rule that
tells how to combine the translation parts:
$$O_2=O_1\,\\{pdollar}\,D\,\\{pdollar}\,\.{\\,}\,\hbox{\\{opt}\thinspace\tt5}$$
This means that the new \\{open\_scrap} scrap has a translation $O_2$ composed
of
the translation $O_1$ of the original
\\{open\_scrap} scrap followed by \\{pdollar} followed by the translation \|D
of the
\\{decl} scrap followed by \\{pdollar} followed
by `\.{\\,}' followed by `\\{opt}' followed by `\.5'. (In the \TeX\ file,
this will specify an additional thin space after the \\{pdollar}, followed
by an optional line break with penalty 50.) Translation rules use subscripts
to distinguish between translations of scraps whose categories have the
same initial letter; these subscripts are assigned from left to right.

There is also the production rule
$$\\{loop\_scrap}\;\RA\;\\{beginning}$$
(meaning that a \&{loop} begins a block-like \ADA\ statement). Since
productions are applied from left to right, this rule will be activated
only if the \\{loop\_scrap} is not preceded by scraps that match other
productions.

The translation rule corresponding to $\\{loop\_scrap}\;\RA\;\\{beginning}$ is
$$B=L$$
but we shall not mention translation rules in the common case that the
translation of the new scrap on the right-hand side is simply the
concatenation of the disappearing scraps on the left-hand side.

\fi

\M167. Here is a list of the category codes that scraps can have.

\Y\P\?\4\X14:Types and objects visible from \\{data\_structures}\X\mathrel{+}\S%
\?\6
\\{accept\_scrap}\?:\?\&{constant}\8\\{eight\_bits}\?\K\?1;\C{denotes %
\&{accept}}\6
\\{and\_scrap}\?:\?\&{constant}\8\\{eight\_bits}\?\K\?2;\C{denotes \&{and}}\6
\\{arrow}\?:\?\&{constant}\8\\{eight\_bits}\?\K\?3;\C{denotes `$\KR$'}\6
\\{base}\?:\?\&{constant}\8\\{eight\_bits}\?\K\?4;\C{denotes `\#' in based
literals}\6
\\{begin\_label\_scrap}\?:\?\&{constant}\8\\{eight\_bits}\?\K\?5;\C{denotes `%
\.{<<}', i.e.\ the beginning of  a label}\6
\\{beginning}\?:\?\&{constant}\8\\{eight\_bits}\?\K\?6;\C{denotes beginning of
block, loop, etc.}\6
\\{case\_scrap}\?:\?\&{constant}\8\\{eight\_bits}\?\K\?7;\C{denotes \&{case}}\6
\\{clause}\?:\?\&{constant}\8\\{eight\_bits}\?\K\?8;\C{denotes a clause in %
\&{case} statements or   variant parts, i.e.\ $\&{when}\ldots\KR\null$}\6
\\{close\_scrap}\?:\?\&{constant}\8\\{eight\_bits}\?\K\?9;\C{denotes a right
parenthesis after which   indentation ends}\6
\\{colon}\?:\?\&{constant}\8\\{eight\_bits}\?\K\?10;\C{denotes a colon}\6
\\{decl}\?:\?\&{constant}\8\\{eight\_bits}\?\K\?11;\C{denotes a declaration}\6
\\{digit\_scrap}\?:\?\&{constant}\8\\{eight\_bits}\?\K\?12;\C{denotes a
digit---possibly (in based  literals) a character}\6
\\{else\_scrap}\?:\?\&{constant}\8\\{eight\_bits}\?\K\?13;\C{denotes a quantity
starting an else-part}\6
\\{end\_label\_scrap}\?:\?\&{constant}\8\\{eight\_bits}\?\K\?14;\C{denotes `%
\.{>>}', i.e.\ the ending of  a label}\6
\\{ending}\?:\?\&{constant}\8\\{eight\_bits}\?\K\?15;\C{denotes \&{end}}\6
\\{exception\_scrap}\?:\?\&{constant}\8\\{eight\_bits}\?\K\?16;\C{denotes %
\&{exception}}\6
\\{exp}\?:\?\&{constant}\8\\{eight\_bits}\?\K\?17;\C{denotes the `\.E' in a
real literal}\6
\\{follows}\?:\?\&{constant}\8\\{eight\_bits}\?\K\?18;\C{denotes \&{is} and %
\&{do}}\6
\\{generic\_scrap}\?:\?\&{constant}\8\\{eight\_bits}\?\K\?19;\C{denotes the
beginning of a generic formal  part}\6
\\{gproc}\?:\?\&{constant}\8\\{eight\_bits}\?\K\?20;\C{denotes \&{procedure}
and \&{function}  preceded by \&{with} in a generic formal part}\6
\\{if\_scrap}\?:\?\&{constant}\8\\{eight\_bits}\?\K\?21;\C{denotes \&{if} and %
\&{elsif}}\6
\\{in\_scrap}\?:\?\&{constant}\8\\{eight\_bits}\?\K\?22;\C{denotes \&{in}}\6
\\{label}\?:\?\&{constant}\8\\{eight\_bits}\?\K\?23;\C{denotes a complete
label}\6
\\{loop\_scrap}\?:\?\&{constant}\8\\{eight\_bits}\?\K\?24;\C{denotes \&{loop}}\6
\\{math}\?:\?\&{constant}\8\\{eight\_bits}\?\K\?25;\C{denotes a quantity to be
used in \TeX's  math mode only}\6
\\{mod\_scrap}\?:\?\&{constant}\8\\{eight\_bits}\?\K\?26;\C{denotes a module
name}\6
\\{new\_scrap}\?:\?\&{constant}\8\\{eight\_bits}\?\K\?27;\C{denotes \&{new}}\6
\\{not\_scrap}\?:\?\&{constant}\8\\{eight\_bits}\?\K\?28;\C{denotes \&{not}}\6
\\{open\_scrap}\?:\?\&{constant}\8\\{eight\_bits}\?\K\?29;\C{denotes a left
parenthesis after which  indentation begins}\6
\\{period}\?:\?\&{constant}\8\\{eight\_bits}\?\K\?30;\C{denotes a `\..'}\6
\\{private\_scrap}\?:\?\&{constant}\8\\{eight\_bits}\?\K\?31;\C{denotes %
\&{private}}\6
\\{proc}\?:\?\&{constant}\8\\{eight\_bits}\?\K\?32;\C{denotes \&{entry}, %
\&{function}, etc.}\6
\\{record\_scrap}\?:\?\&{constant}\8\\{eight\_bits}\?\K\?33;\C{denotes %
\&{record}}\6
\\{return\_scrap}\?:\?\&{constant}\8\\{eight\_bits}\?\K\?34;\C{denotes %
\&{return} and \&{exit}}\6
\\{semi}\?:\?\&{constant}\8\\{eight\_bits}\?\K\?35;\C{denotes a semicolon}\6
\\{separate\_scrap}\?:\?\&{constant}\8\\{eight\_bits}\?\K\?36;\C{denotes %
\&{separate}}\6
\\{stmt}\?:\?\&{constant}\8\\{eight\_bits}\?\K\?37;\C{denotes a statement}\6
\\{terminator}\?:\?\&{constant}\8\\{eight\_bits}\?\K\?38;\C{denotes the ending
of statements,  blocks, declaration parts, etc.}\6
\\{then\_scrap}\?:\?\&{constant}\8\\{eight\_bits}\?\K\?39;\C{denotes \&{then}}\6
\\{when\_scrap}\?:\?\&{constant}\8\\{eight\_bits}\?\K\?40;\C{denotes \&{when}}\6
\\{with\_scrap}\?:\?\&{constant}\8\\{eight\_bits}\?\K\?41;\C{denotes \&{use}
and \&{with}}\par
\fi

\M168. The token lists for translated \TeX\ output contain some special control
symbols as well as ordinary characters. These control symbols are
interpreted by \.{AWEAVE} before they are written to the output file.

\yskip\hang \\{pdollar} denotes a possible \.\$ sign in \TeX\ output;

\yskip\hang \\{pspace} denotes a possible blank space; two or more
consecutive \\{pspace} symbols combine to a \\{sspace};

\yskip\hang \\{sspace} denotes a blank space in \TeX\ output;

\yskip\hang \\{break\_space} denotes an optional line break or an en space;

\yskip\hang \\{force} denotes a line break;

\yskip\hang \\{big\_force} denotes a line break with additional vertical space;

\yskip\hang \\{opt} denotes an optional line break (with the continuation
line indented two ems with respect to the normal starting position)---this
code is followed by an integer \|n, and the break will occur with penalty
$10n$;

\yskip\hang \\{backup} denotes a backspace of one em;

\yskip\hang \\{cancel} obliterates any \\{break\_space} or \\{force} or \\{big%
\_force}
tokens that immediately precede or follow it and also cancels any
\\{backup} tokens that follow it;

\yskip\hang \\{indent} causes future lines to be indented one more em;

\yskip\hang \\{outdent} causes future lines to be indented one less em.

\yskip\noindent All of these tokens, except \\{pdollar} and \\{sspace},  are
removed from the \TeX\ output that
comes from \ADA\ text between \ab\ signs; \\{break\_space} and \\{force} and
\\{big\_force} become single spaces in this mode. The translation of other
\ADA\ texts results in \TeX\ control sequences \.{\\1}, \.{\\2},
\.{\\3}, \.{\\4}, \.{\\5}, \.{\\6}, \.{\\7} corresponding respectively to
\\{indent}, \\{outdent}, \\{opt}, \\{backup}, \\{break\_space}, \\{force}, and
\\{big\_force}. However, a sequence of consecutive `\.\ ', \\{break\_space},
\\{force}, and/or \\{big\_force} tokens is first replaced by a single token
(the maximum of the given ones). The \\{pdollar} token results in \.{\\?},
supposed there are $n$ consecutive \\{pdollar} tokens and $n\bmod2=1$, i.e.\
$n$ is odd. The \\{sspace} token results in \.{\\8}.

The tokens \\{math\_rel}, \\{math\_bin}, \\{math\_op} will be translated into
\.{\\mathrel\{}, \.{\\mathbin\{}, and \.{\\mathop\{}, respectively.
Other control sequences in the \TeX\ output will be
`\.{\\\\\{}$\,\ldots\,$\.\}' surrounding normal identifiers,
`\.{\\p\{}$\,\ldots\,$\.\}' surrounding predefined identifiers,
`\.{\\\&\{}$\,\ldots\,$\.\}' surrounding reserved words,
`\.{\\.\{}$\,\ldots\,$\.\}' surrounding strings,
`\.{\\C\{}$\,\ldots\,$\.\}$\,$\\{force}' surrounding comments, and
`\.{\\X$n$:}$\,\ldots\,$\.{\\X}' surrounding module names, where
\|n is the module number.

\Y\P\?\4\X14:Types and objects visible from \\{data\_structures}\X\mathrel{+}\S%
\?\6
\\{math\_bin}\?:\?\&{constant}\8\\{eight\_bits}\?\K\bn{8}{204}\?;\6
\\{math\_rel}\?:\?\&{constant}\8\\{eight\_bits}\?\K\bn{8}{205}\?;\6
\\{math\_op}\?:\?\&{constant}\8\\{eight\_bits}\?\K\bn{8}{206}\?;\6
\\{pdollar}\?:\?\&{constant}\8\\{eight\_bits}\?\K\bn{8}{207}\?;\C{possible
dollar sign (\.{\\?})}\6
\\{pspace}\?:\?\&{constant}\8\\{eight\_bits}\?\K\bn{8}{210}\?;\C{possible
space}\6
\\{sspace}\?:\?\&{constant}\8\\{eight\_bits}\?\K\bn{8}{211}\?;\C{sure space (%
\.{\\8})}\6
\\{big\_cancel}\?:\?\&{constant}\8\\{eight\_bits}\?\K\bn{8}{212}\?;\C{like %
\\{cancel}, also overrides spaces}\6
\\{cancel}\?:\?\&{constant}\8\\{eight\_bits}\?\K\bn{8}{213}\?;\C{overrides %
\\{backup}, \\{break\_space}, \\{force}, \\{big\_force}}\6
\\{indent}\?:\?\&{constant}\8\\{eight\_bits}\?\K\bn{8}{214}\?;\C{one more tab (%
\.{\\1})}\6
\\{outdent}\?:\?\&{constant}\8\\{eight\_bits}\?\K\bn{8}{215}\?;\C{one less tab
(\.{\\2})}\6
\\{opt}\?:\?\&{constant}\8\\{eight\_bits}\?\K\bn{8}{216}\?;\C{optional break in
mid-statement (\.{\\3})}\6
\\{backup}\?:\?\&{constant}\8\\{eight\_bits}\?\K\bn{8}{217}\?;\C{stick out one
unit to the left (\.{\\4})}\6
\\{break\_space}\?:\?\&{constant}\8\\{eight\_bits}\?\K\bn{8}{220}\?;\C{optional
break between statements (\.{\\5})}\6
\\{force}\?:\?\&{constant}\8\\{eight\_bits}\?\K\bn{8}{221}\?;\C{forced break
between statements (\.{\\6})}\6
\\{big\_force}\?:\?\&{constant}\8\\{eight\_bits}\?\K\bn{8}{222}\?;\C{forced
break with additional space (\.{\\7})}\6
\\{end\_translation}\?:\?\&{constant}\8\\{eight\_bits}\?\K\bn{8}{223}\?;%
\C{special sentinel token at end of list}\par
\fi

\M169. The raw input is converted into scraps according to the following table,
which gives category codes followed by the translations.
\def\stars {\.{***}}%
(The symbol `\stars' stands for `\.{\\\&\{{\rm identifier}\}}',
i.e., the identifier itself treated as a reserved word. In a few cases the
category is given as `\\{comment}'; this is not an actual category code, it
means that the translation will be treated as a comment, as explained
below.)

\yskip\halign{\quad#\hfil&\quad#\hfil\cr
\.{/=}&\\{math} : \.{\\I}\cr
\.{<=}&\\{math} : \.{\\L}\cr
\.{>=}&\\{math} : \.{\\G}\cr
\.{:=}&\\{math} : \.{\\K}\cr
\.{==}&\\{math} : \.{\\S}\cr
\.{=>}&\\{math} : \.{\\KR}\cr
\.{**}&\\{math} : \.{\\dst}\cr
\.{<<}&\\{math} : \.{\\BL}\cr
\.{>>}&\\{math} : \.{\\EL}\cr
\.{<>}&\\{math} : \.{\\BOX}\cr
\."$\,$string$\,$\."&\\{math} : $\\{pdollar}\,\.{\\.\{"{\rm$\,$modified string$%
\,$}"\}}
\,\\{pdollar}$\cr
\.\'$\,$character$\,$\.\'&\\{math} : $\\{pdollar}\,\.{\\.\{\\\'{\rm$\,$modified
character$\,$}\\\'\}}\,\\{pdollar}$\cr
\.{@=}$\,$string$\,$\.{@>}&\\{math} : $\\{pdollar}\,\.{\\=\{{\rm$\,$modified
string$\,$}\}}\,\\{pdollar}$\cr
\.{@\~}$\,$string$\,$\.{@>}&\\{math} : $\\{pdollar}\,\.{\\O\{{\rm$\,$modified
string$\,$}\}}\,\\{pdollar}\,\\{force}$\cr
\# in based literal&\\{base} :\cr
\# otherwise&\\{math} : $\\{pdollar}\,\.{\\\#}\,\\{pdollar}$\cr
\.\$&\\{math} : $\\{pdollar}\,\.{\\\$}\,\\{pdollar}$\cr
\.\_&\\{math} : $\\{pdollar}\,\.{\\\_}\,\\{pdollar}$\cr
\.\%&\\{math} : $\\{pdollar}\,\.{\\\%}\,\\{pdollar}$\cr
\.\^&\\{math} : $\\{pdollar}\,\.{\\\^}\,\\{pdollar}$\cr
\.\&&\\{math} : $\\{math\_bin}\,\.{\\.\\\&\}}$\cr
\.!&\\{math} : $\\{math\_rel}\,\.{\\.\char'174\}}$\cr
\.(&\\{open\_scrap} : \.(\cr
\.)&\\{close\_scrap} : \.)\cr
\.*&\\{math} : \.{\\ast}\cr
\.,&\\{math} : \.,\?\,\?\\{opt}\?\,\?\.9\cr
\.{..}&\\{math} : \.{\\to}\cr
\..&\\{period} : \..\cr
\.:&\\{colon} : \.:\cr
\.;&\\{semi} : \.;\cr
predefined identifier&\\{math} : $\\{pspace}\,\\{pdollar}\,\.{\\p\{{\rm$%
\,$identifier$\,$}\}}\,
\\{pdollar}\,\\{pspace}$\cr
normal identifier&\\{math} : $\\{pspace}\,\\{pdollar}\,\.{\\\\\{{\rm$%
\,$identifier$\,$}\}}\,
\\{pdollar}\,\\{pspace}$\cr
\.E in constant&\\{exp} : \.{\\E\{}\cr
digit $d$, not in based number&\\{digit\_scrap} : $\\{pdollar}\,d\,\\{pdollar}$%
\cr
digit $d$, in based number&\\{digit\_scrap} : $d$\cr
other character $c$&\\{math} : $c$\cr
\.{abort}&\\{math} : $\\{pspace}\,\\{pdollar}\,\stars\,\\{pdollar}\,\\{sspace}$%
\cr
\.{abs}&\\{math} : $\\{pspace}\,\\{pdollar}\,\stars\,\\{pdollar}\,\\{sspace}$%
\cr
\.{accept}&\\{accept\_scrap} : $\stars\,\\{sspace}$\cr
\.{access}&\\{math} : $\\{pspace}\,\\{pdollar}\,\stars\,\\{pdollar}\,%
\\{sspace}$\cr
\.{all}&\\{math} : $\\{pspace}\,\\{pdollar}\,\stars\,\\{pdollar}\,\\{pspace}$%
\cr
\.{and}&\\{and\_scrap} : \.{\\W}\cr
\.{array}&\\{open\_scrap} : $\stars\,\\{sspace}$\cr
\.{at}&\\{colon} : $\\{sspace}\,\stars\,\\{sspace}$\cr
\.{begin}&\\{beginning} : \stars\cr
\.{body}&\\{math} : $\\{pdollar}\,\stars\,\\{pdollar}\,\\{pspace}$\cr
\.{case}&\\{case\_scrap} : \stars\cr
\.{constant}&\\{math} : $\\{pspace}\,\\{pdollar}\,\stars\,\\{pdollar}\,%
\\{sspace}$\cr
\.{declare}&\\{decl} : $\\{force}\,\stars\,\\{indent}\,\\{force}$\cr
\.{delay}&\\{math} : $\\{pspace}\,\\{pdollar}\,\stars\,\\{pdollar}\,\\{sspace}$%
\cr
\.{delta}&\\{math} : $\\{pspace}\,\\{pdollar}\,\stars\,\\{pdollar}\,\\{sspace}$%
\cr
\.{digits}&\\{math} : $\\{pspace}\,\\{pdollar}\,\stars\,\\{pdollar}\,%
\\{sspace}$\cr
\.{do}&\\{follows} : $\\{sspace}\,\stars\,\\{sspace}$\cr
\.{else}&\\{else\_scrap} : \stars\cr
\.{elsif}&\\{if\_scrap} : $\\{backup}\,\stars$\cr
\.{end}&\\{ending} : \stars\cr
\.{entry}&\\{proc} : $\stars\,\\{sspace}$\cr
\.{exception}&\\{exception\_scrap} : \stars\cr
\.{exit}&\\{return\_scrap} : \stars\cr
\.{for}&\\{math} : $\\{pspace}\,\\{pdollar}\,\stars\,\\{pdollar}\,\\{sspace}$%
\cr
\.{function}&\\{proc} : $\stars\,\\{sspace}$\cr
\.{generic}&\\{generic\_scrap} : \stars\cr
\.{goto}&\\{math} : $\\{pspace}\,\\{pdollar}\,\stars\,\\{pdollar}\,\\{sspace}$%
\cr
\.{if}&\\{if\_scrap} : $\stars\,\\{indent}$\cr
\.{in}&\\{in\_scrap} : \.{\\in}\cr
\.{is}&\\{follows} : $\\{sspace}\,\stars\,\\{sspace}$\cr
\.{limited}&\\{math} : $\\{pspace}\,\\{pdollar}\,\stars\,\\{pdollar}\,%
\\{sspace}$\cr
\.{loop}&\\{loop\_scrap} : \stars\cr
\.{mod}&\\{math} : $\\{math\_bin}\,\stars\,\.\}$\cr
\.{new}&\\{new\_scrap} : $\\{pspace}\,\stars\,\\{pspace}$\cr
\.{not}&\\{not\_scrap} : \.{\\R}\cr
\.{null}&\\{math} : $\\{pspace}\,\\{pdollar}\,\stars\,\\{pdollar}\,\\{pspace}$%
\cr
\.{of}&\\{close\_scrap} : $\\{sspace}\,\stars\,\\{sspace}$\cr
\.{or}&\\{else\_scrap} : \stars\cr
\.{others}&\\{math} : $\\{pspace}\,\\{pdollar}\,\stars\,\\{pdollar}\,%
\\{sspace}$\cr
\.{out}&\\{math} : $\\{pspace}\,\\{pdollar}\,\stars\,\\{pdollar}\,\\{sspace}$%
\cr
\.{package}&\\{proc} : $\stars\,\\{sspace}$\cr
\.{pragma}&\\{proc} : $\stars\,\\{sspace}$\cr
\.{private}&\\{private\_scrap} : \stars\cr
\.{procedure}&\\{proc} : $\stars\,\\{sspace}$\cr
\.{raise}&\\{math} : $\\{pspace}\,\\{pdollar}\,\stars\,\\{pdollar}\,\\{sspace}$%
\cr
\.{range}&\\{math} : $\\{sspace}\,\\{pdollar}\,\stars\,\\{pdollar}\,\\{sspace}$%
\cr
\.{record}&\\{record\_scrap} : \stars\cr
\.{rem}&\\{math} : $\\{math\_bin}\,\stars\,\.\}$\cr
\.{renames}&\\{math} : $\\{pspace}\,\\{pdollar}\,\stars\,\\{pdollar}\,%
\\{sspace}$\cr
\.{return}&\\{return\_scrap} : \stars\cr
\.{reverse}&\\{math} : $\\{pspace}\,\\{pdollar}\,\stars\,\\{pdollar}\,%
\\{sspace}$\cr
\.{select}&\\{beginning} : \stars\cr
\.{separate}&\\{separate\_scrap} : \stars\cr
\.{subtype}&\\{math} : $\\{pdollar}\,\stars\,\\{pdollar}\,\\{sspace}$\cr
\.{task}&\\{proc} : $\stars\,\\{sspace}$\cr
\.{terminate}&\\{math} : $\\{pspace}\,\\{pdollar}\,\stars\,\\{pdollar}\,%
\\{pspace}$\cr
\.{then}&\\{then\_scrap} : $\\{sspace}\,\stars$\cr
\.{type}&\\{math} : $\\{pdollar}\,\stars\,\\{pdollar}\,\\{sspace}$\cr
\.{use}&\\{with\_scrap} : $\\{pspace}\stars\,\\{sspace}$\cr
\.{when}&\\{when\_scrap} : $\stars\,\\{sspace}$\cr
\.{while}&\\{math} : $\\{pspace}\,\\{pdollar}\,\stars\,\\{pdollar}\,\\{sspace}$%
\cr
\.{with}&\\{with\_scrap} : $\\{pspace}\stars\,\\{sspace}$\cr
\.{xor}&\\{math} : $\\{math\_bin}\,\stars\,\.\}$\cr


\.{@\'$c$\'}&\\{math} : $\\{pdollar}\,\.{\\H\{$c$\}}\,\\{pdollar}$\cr
\.{@\\}&\\{math} : \.{\\)}\cr
\.{@,}&\\{math} : \.{\\,}\cr
\.{@t}$\,$stuff$\,$\.{@>}&\\{math} : $\\{pdollar}\,\.{\\hbox\{{\rm$\,$stuff$%
\,$}\}}
\,\\{pdollar}$\cr
\.{@<}$\,$module$\,$\.{@>}&\\{mod\_scrap} : \.{\\X$n$:{\rm$\,$module$\,$}\\X}%
\cr
\.{@\#}&\\{comment} : \\{big\_force}\cr
\.{@/}&\\{comment} : \\{force}\cr
\.{@\char'174}&\\{math} : $\\{opt}\,\.0$\cr
\.{@+}&\\{comment} : $\\{big\_cancel}\,\.{\\8}\,\\{big\_cancel}$\cr
\.{@;}&\\{semi} :\cr
\.{@\&}&\\{math} : \.{\\J}\cr
\.{@\{}&\\{math} : \.{\\B}\cr
\.{@\}}&\\{math} : \.{\\T}\cr}
\yskip\noindent When a string is output, certain characters are preceded by
`\.\\' signs so that they will print properly.

A comment in the input will be combined with the preceding
scrap.
An additional ``comment'' is effectively appended at the end of the
\ADA\ text, just before translation begins; this consists of a \\{cancel}
token in the case of \ADA\ text in \ab, otherwise it consists of a
\\{force} token.

From this table it is evident that \.{AWEAVE} will parse a lot of non-\ADA\
programs. For example, the reserved words `\.{or}' and `\.{else}' are
treated in an identical way by \.{AWEAVE};
\ADA\ programmers may well be surprised at this
similarity. The idea is to keep \.{AWEAVE}'s rules as simple as possible,
consistent with doing a reasonable job on syntactically correct \ADA\
programs. The production rules below have been formulated in the same
spirit of ``almost anything goes.''

\fi

\M170. Here is a table of all the productions. The reader can best get a feel
for
how they work by trying them out by hand on small examples; no amount of
explanation will be as effective as watching the rules in action.

\def\[#1]{\quad$\dleft#1\dright$}
\def\sp{\\{sspace}}
\def\ps{\\{pspace}}
\def\pd{\\{pdollar}}
\yskip
\halign to\the\hsize{\hfil\it# &
#\hfil\hskip-200pt\tabskip 0pt plus 100pt&
\hbox to 128pt{#\hfil}\tabskip0pt\cr
&Production categories\[\hbox{translations}]&Remarks\cr
\noalign{\yskip}
1&\\{accept\_scrap}\?\,\?\\{math}\?\,\?\\{follows} $\RA$ \\{beginning}&%
\&{accept} \\{name} \&{do}\cr
&\qquad\[B=A\,\\{indent}\,\pd\,M\,\pd\,\\{outdent}\,F]\cr
2&\\{accept\_scrap}\?\,\?\\{math}\?\,\?\\{semi} $\RA$ \\{stmt}\[S_2=A\,\pd\,M\,%
\pd\,S_1]
&\&{accept} \\{name};\cr
3&\\{and\_scrap}\?\,\?\\{then\_scrap} $\RA$ \\{math}\[M=\.{\\ANT}]&\&{and} %
\&{then}\cr
4&\\{arrow} $\RA$ \\{math}& `$\KR$' in aggregate etc.\cr
5&\\{base}\?\,\?\\{base} $\RA$ \\{math}\[M=\.{\\bn\{}\,B_1\,\.{\}\{}\,B_2\,\.%
\}]
& complete based literal\cr
6&\\{begin\_label\_scrap}\?\,\?\\{math}\?\,\?\\{end\_label\_scrap} $\RA$ %
\\{label}&label\cr
&\qquad\[L=B\,\pd\,M\,\pd\,E]\cr
7&\\{beginning}\?\,\?\\{stmt} $\RA$ \\{stmt}\[S_2=\\{force}\,\\{indent}\,B\,%
\\{force}\,S_1]
&beginning of block\cr
8&\\{case\_scrap}\?\,\?\\{math}\?\,\?\\{follows}\?\,\?\\{decl} $\RA$ %
\\{decl}&beginning of variant part\cr
&\qquad\[D_2=C\,\\{indent}\,\pd\,\sp\,M\,\pd\,\\{indent}\,F\,\\{force}\,D_1]\cr
9&\\{case\_scrap}\?\,\?\\{math}\?\,\?\\{follows}\?\,\?\\{stmt} $\RA$ %
\\{stmt}&beginning of \&{case} statement\cr
&\qquad\[S_2=\\{force}\,C\,\\{indent}\,\pd\,\,\sp
\,M\,\pd\,\\{indent}\,F\,\\{force}\,S_1]\cr
10&\\{clause}\?\,\?\\{decl} $\RA$ \\{decl}\[D_2=\\{force}\,\\{backup}\,C\,%
\\{force}\,D_1]
&element of variant part\cr
11&\\{clause}\?\,\?\\{stmt} $\RA$ \\{stmt}\[S_2=\\{force}\,\\{backup}\,C\,%
\\{force}\,S_1]
&element of \&{case} statement\cr
12&\\{colon}\?\,\?\\{exception\_scrap} $\RA$ \\{colon}\?\,\?\\{math}&exception
declaration\cr
13&\\{colon}\?\,\?\\{in\_scrap} $\RA$ \\{colon}\[C_2=C_1\,\.{\\\&\{in\}}\,\sp]
&\&{in} as parameter mode\cr
14&\\{decl}\?\,\?\\{beginning} $\RA$ \\{stmt}\[S=D\,\\{force}\,\\{backup}\,B\,%
\\{force}]
&beginning of statements\cr
15&\\{decl}\?\,\?\\{decl} $\RA$ \\{decl}\[D_3=D_1\,\\{force}\,D_2]
&declaration part grows\cr
16&\\{decl}\?\,\?\\{mod\_scrap}\?\,\?\\{semi} $\RA$ \\{decl}\[D_2=D_1\,%
\\{force}\,M\,S]
&module as declaration\cr
17&\\{decl}\?\,\?\\{terminator} $\RA$ \\{decl}%
\[D_2=D_1\,\\{cancel}\,\\{outdent}\,\\{force}\,T\,\\{force}]
&end of declarative part\cr
18&\\{decl}\?\,\?\\{with\_scrap} $\RA$ \\{with\_scrap}\[W_2=D\,\\{force}%
\,W_1]&context clause grows\cr
19&\\{digit\_scrap} $\RA$ \\{math}&digit not in based number\cr
20&\\{digit\_scrap}\?\,\?\\{base} $\RA$ \\{base}\[B_2=D]&part of based number
grows\cr
21&\\{digit\_scrap}\?\,\?\\{digit\_scrap} $\RA$ \\{digit\_scrap}&concatenation%
\cr
22&\\{digit\_scrap}\?\,\?\\{period} $\RA$ \\{digit\_scrap}&concatenation\cr
23&\\{else\_scrap}\?\,\?\\{else\_scrap} $\RA$ \\{math}\[M=\.{\\ORE}]&\&{or
else}\cr
24&\\{else\_scrap}\?\,\?\\{when\_scrap} $\RA$ \\{when\_scrap}\[W_2=\.{\\ORW}]
&\&{or when} in \&{select} statement\cr
25&\\{ending}\?\,\?\\{beginning}\?\,\?\\{semi} $\RA$ \\{terminator}\[T=E\,\sp%
\,B\,S]
&\&{end} \&{select};\cr
26&\\{ending}\?\,\?\\{case\_scrap}\?\,\?\\{semi} $\RA$ \\{terminator}%
\[T=\\{backup}\,\\{outdent}\,E\,\sp\,C\,S]
&\&{end} \&{case};\cr
27&\\{ending}\?\,\?\\{if\_scrap}\?\,\?\\{semi} $\RA$ \\{terminator}\[T=E\,\sp\,%
\\{outdent}\,I\,S]
&\&{end} \&{if};\cr
28&\\{ending}\?\,\?\\{loop\_scrap}\?\,\?\\{semi} $\RA$ \\{terminator}\[T=E\,\sp%
\,L\,S]
&\&{end} \&{loop};\cr
29&\\{ending}\?\,\?\\{loop\_scrap}\?\,\?\\{math}\?\,\?\\{semi} $\RA$ %
\\{terminator}
&\&{end} \&{loop} \\{name};\cr
&\qquad\[T=E\,\sp\,L\,\pd\,\sp\,M\,\pd\,S]\cr
30&\\{ending}\?\,\?\\{math}\?\,\?\\{semi} $\RA$ \\{terminator}%
\[T=E\,\pd\,\sp\,M\,\pd\,S]&\&{end} \\{name};\cr
31&\\{ending}\?\,\?\\{record\_scrap}\?\,\?\\{semi} $\RA$ \\{terminator}\[T=%
\\{outdent}\,E\,\sp\,R\,S]
&\&{end} \&{record};\cr
32&\\{ending}\?\,\?\\{semi} $\RA$ \\{terminator}&\&{end};\cr
33&\\{exception\_scrap} $\RA$ \\{beginning}\[B=\\{backup}\,E]
&beginning of \&{exception} part\cr
34&$\\{exp}\\{digit\_scrap}^*$ $\RA$ \\{math}\[M=E\,D\,\.\}]&unsigned exponent%
\cr
35&\\{exp}\?\,\?\\{math}$\,\\{digit\_scrap}^*$ $\RA$ \\{math}\[M_2=E\,D\,M_1\,%
\.\}]
&signed exponent\cr
36&\\{generic\_scrap}\?\,\?\\{decl} $\RA$ \\{generic\_scrap}\[G_2=G_1\,%
\\{force}\,D]
&generic formal part grows\cr
37&\\{generic\_scrap}\?\,\?\\{proc} $\RA$ \\{proc}\[P_2=\\{indent}\,G\,%
\\{outdent}\,\\{force}\,P_1]
&generic formal part ends\cr
38&\\{gproc}\?\,\?\\{math}\?\,\?\\{semi} $\RA$ \\{decl}%
\[D=G\,\\{indent}\,\pd\,M\,\pd\,\\{outdent}\,S]
&generic formal subprogram\cr
39&\\{gproc}\?\,\?\\{math}\?\,\?\\{follows}\?\,\?\\{math}\?\,\?\\{semi} $\RA$ %
\\{decl}
&{\it 38\/} with default\cr
&\qquad\[D=G\,\\{indent}\,\pd\,M_1\,F\,M_2\,\pd\,\\{outdent}\,S]\cr
40&\\{if\_scrap}\?\,\?\\{math}\?\,\?\\{then\_scrap}\?\,\?\\{stmt} $\RA$ %
\\{stmt}
&part of \&{if} statement\cr
&\qquad\[S_2=\\{force}\,I\,\\{indent}\,\pd\,\sp\,M\,\pd
\,\\{outdent}\,T\,\\{force}\,S_1]\cr
41&\\{in\_scrap} $\RA$ \\{math}&\&{in} as `$\in$'\cr
42&\\{label}\?\,\?\\{label} $\RA$ \\{label}\[L_3=L_1\,\sp\,\\{opt}\,\.7\,L_2]
&multiple label\cr
43&\\{label}\?\,\?\\{stmt} $\RA$ \\{stmt}\[S_2=\\{force}\,\\{backup}\,L\,\sp\,%
\\{cancel}\,S_1]
&labeled statement\cr
44&\\{loop\_scrap} $\RA$ \\{beginning}&beginning of \&{loop} statement\cr
45&\\{math}\?\,\?\\{and\_scrap}\?\,\?\\{math} $\RA$ \\{math}&\&{and} as `$\W$'%
\cr
46&\\{math}\?\,\?\\{colon}\?\,\?\\{beginning} $\RA$ \\{beginning}&block or loop
with name\cr
&\qquad\[B_2=\\{backup}\,\pd\,M\,\sp\,\pd\,C\,\sp\,B]\cr
47&\\{math}\?\,\?\\{colon}\?\,\?\\{decl} $\RA$ \\{decl}&block with name and %
\&{declare}\cr
&\qquad\[D_2=\\{force}\,\\{backup}\,\pd\,M\,\sp\,\pd\,C\,\sp\,\\{cancel}\,D_1]%
\cr
48&\\{math}\?\,\?\\{colon}\?\,\?\\{math}\?\,\?\\{semi} $\RA$ \\{decl}%
\[D=\pd\,M_1\,C\,M_2\,\pd\,S]&declaration\cr
49&\\{math}\?\,\?\\{else\_scrap}\?\,\?\\{math} $\RA$ \\{math}\[M_3=M_1\,\.{\\V}%
\,M_2]
&\&{or} as `$\V$'\cr
50&\\{math}\?\,\?\\{follows}\?\,\?\\{math}\?\,\?\\{semi} $\RA$ \\{decl}&%
\&{type} declaration etc.\cr
&\qquad\[D=\pd\,M_1\,\pd\,F\,\pd\,M_2\,\pd\,S]\cr
51&\\{math}\?\,\?\\{follows}\?\,\?\\{record\_scrap} $\RA$ \\{decl}&\&{record}
type declaration\cr
&\qquad\[D=\pd\,M\,\pd\,F\,\\{indent}\,\\{force}\,R\,\\{indent}\,\\{force}]\cr
52&\\{math}\?\,\?\\{loop\_scrap} $\RA$ \\{beginning}\[B=\pd\,M\,\pd\,\sp\,L]
&\&{loop} with \&{for} or \&{while}\cr
53&\\{math}\?\,\?\\{math} $\RA$ \\{math}&simple concatenation\cr
54&\\{math}\?\,\?\\{semi} $\RA$ \\{stmt}\[S_2=\pd\,M\,\pd\,S_1]&simple
statement\cr
55&\\{math}\?\,\?\\{stmt} $\RA$ \\{stmt}\[S_2=\pd\,M\,\pd\,S_1]&simple
concatenation\cr
56&\\{math}\?\,\?\\{separate\_scrap} $\RA$ \\{separate\_scrap}\[S_2=\pd\,M\,\pd%
\,S_1]&in emergencies\cr
57&\\{math}\?\,\?\\{when\_scrap}\?\,\?\\{math} $\RA$ \\{math}&\&{exit} \&{when}
$x$;\cr
58&\\{math}\?\,\?\\{with\_scrap} $\RA$ \\{with\_scrap}\[W_2=\pd\,M\,\ps\,\pd%
\,W_1]
&\&{for} $x$ \&{use} \dots\cr
59&\\{mod\_scrap} $\RA$ \\{math}&module unlike a statement\cr
60&\\{mod\_scrap}\?\,\?\\{semi} $\RA$ \\{stmt}\[S_2=M\,S_1\,\\{force}]
&module like a statement\cr
61&\\{new\_scrap} $\RA$ \\{math}&\&{new} in allocator\cr
62&\\{not\_scrap} $\RA$ \\{math}&\&{not} as `$\R$'\cr
63&\\{not\_scrap}\?\,\?\\{in\_scrap} $\RA$ \\{math}\[M=\.{\\notin}]&\&{not} %
\&{in} as `$\notin$'\cr
64&\\{open\_scrap}\?\,\?\\{decl} $\RA$ \\{open\_scrap}\[O_2=O_1\,\pd\,D\,\pd\,%
\.{\\,}\,\\{opt}\,\.5]
&declaration in parentheses\cr
65&\\{open\_scrap}\?\,\?\\{math}\?\,\?\\{close\_scrap} $\RA$ \\{math}%
\[M_2=\\{indent}\,O\,M_1\,C\,\\{outdent}]&parenthesized quantity\cr
66&\\{open\_scrap}\?\,\?\\{math}\?\,\?\\{colon}\?\,\?\\{math}\?\,\?\\{close%
\_scrap} $\RA$ \\{math} &formal part ends\cr
&\qquad\[M_3=O\,\\{indent}\,M_1\,C_1\,M_2\,\\{outdent}\,C_2]\cr
67&\\{period} $\RA$ \\{math}&period in selected component\cr
68&\\{private\_scrap} $\RA$ \\{decl}\[D=\\{force}\,\\{backup}\,P\,\\{force}]
&beginning of private part\cr
69&\\{private\_scrap}\?\,\?\\{semi} $\RA$ \\{math}&private type declaration\cr
70&\\{proc}\?\,\?\\{math}\?\,\?\\{follows} $\RA$ \\{decl}&package declaration
etc.\cr
&\qquad\[D=P\,\\{indent}\,\pd\,M\,\pd\,\\{outdent}\,F\,\\{cancel}\,\\{indent}]%
\cr
71&\\{proc}\?\,\?\\{math}\?\,\?\\{follows}\?\,\?\\{new\_scrap} $\RA$ \\{proc}\?%
\,\?\\{math}\?\,\?\\{math}&instantiation\cr
72&\\{proc}\?\,\?\\{math}\?\,\?\\{follows}\?\,\?\\{separate\_scrap} $\RA$ %
\\{proc}\?\,\?\\{math}\?\,\?\\{math}&body stub\cr
73&\\{proc}\?\,\?\\{math}\?\,\?\\{semi} $\RA$ \\{decl}%
\[D=P\,\\{indent}\,\pd\,M\,\pd\,\\{outdent}\,S]&subprogram specification\cr
74&\\{return\_scrap} $\RA$ \\{math}\[M=\ps\,R\,\sp]&\&{return} with expression%
\cr
75&\\{return\_scrap}\?\,\?\\{semi} $\RA$ \\{stmt}&\&{return};\cr
76&\\{separate\_scrap}\?\,\?\\{math} $\RA$ \\{decl}\[D=S\,\pd\,\sp\,M\,\pd]&%
\&{separate}(\dots)\cr
77&\\{stmt}\?\,\?\\{else\_scrap}\?\,\?\\{stmt} $\RA$ \\{stmt}\[S_3=S_1\,%
\\{force}\,\\{backup}\,E\,
\\{force}\,S_2]&\&{else} part\cr
78&\\{stmt}\?\,\?\\{stmt} $\RA$ \\{stmt}\[S_3=S_1\,\\{break\_space}%
\,S_2]&concatenation\cr
79&\\{stmt}\?\,\?\\{terminator} $\RA$ \\{stmt}%
\[S_2=S_1\,\\{cancel}\,\\{outdent}\,\\{force}\,T\,\\{force}] &end of block etc.%
\cr
80&\\{when\_scrap}\?\,\?\\{math}\?\,\?\\{arrow} $\RA$ \\{clause}\[C=W\,\pd\,M\,%
\pd\,A]
&clause in \&{case}\cr
81&\\{with\_scrap}\?\,\?\\{colon} $\RA$ \\{with\_scrap}&\&{use} \&{at}\cr
82&\\{with\_scrap}\?\,\?\\{math}\?\,\?\\{semi} $\RA$ \\{decl}\[D=W\,\pd\,M\,\pd%
\,S]&context clause\cr
83&\\{with\_scrap}\?\,\?\\{proc} $\RA$ \\{gproc}&generic formal subprogram\cr
84&\\{with\_scrap}\?\,\?\\{record\_scrap} $\RA$ \\{decl}\[D=W\,\\{indent}\,%
\\{force}\,\\{indent}\,R]
&\&{record} representation\cr
85&\\{with\_scrap}\?\,\?\\{record\_scrap}\?\,\?\\{colon} $\RA$ \\{math}\?\,\?%
\\{colon}
&{\it 84\/} with alignment\cr
&\qquad\[M=\pd\,W\,\\{indent}\,\\{force}\,\\{indent}\,R\,\pd]\cr
}
\yskip\noindent
Translations are not specified here when they are simple concatenations
of the scraps that change. For example, the full translation of
`\\{colon}\?\,\?\\{exception\_scrap} $\RA$ \\{colon}\?\,\?\\{math}' is
$C_2=C_1$, $M=E$. The notation
`$\\{digit\_scrap}^*$' which occurs in the productions related to \\{exp}
(numbers 34 and~35) means a \\{digit\_scrap} scrap that isn't followed by
another
\\{digit\_scrap} scrap.



\fi

\N171.  Implementing the productions.
When \ADA\ text is to be processed with the grammar above, we put its
initial scraps $s_1\ldots s_n$ into two arrays \\{cat}\?(\?1\?\to\?\|n\?)\? and
\\{trans}\?(\?1\?\to\?\|n\?)\?.
The value of \\{cat}\?(\?\|k\?)\? is simply a category code from the list
above; the
value of \\{trans}\?(\?\|k\?)\? is a text pointer, i.e., an index into \\{tok%
\_start}.
Our production rules have the nice property that the right-hand side is never
longer than the left-hand side. Therefore it is convenient to use sequential
allocation for the current sequence of scraps. Five pointers are used to
manage the parsing:

\yskip\hang \\{pp} (the parsing pointer) is such that we are trying to match
the category codes \\{cat}\?(\?\\{pp}\?)\,\?\\{cat}\?(\?\\{pp}\?+\?1\?)\?$\,%
\ldots\,$ to the left-hand sides
of productions.

\yskip\hang \\{scrap\_base}, \\{lo\_ptr}, \\{hi\_ptr}, and \\{scrap\_ptr} are
such that
the current sequence of scraps appears in positions \\{scrap\_base} through
\\{lo\_ptr} and \\{hi\_ptr} through \\{scrap\_ptr}, inclusive, in the \\{cat}
and
\\{trans} arrays. Scraps located between \\{scrap\_base} and \\{lo\_ptr} have
been examined, while those in positions \?\G\?\\{hi\_ptr} have not yet been
looked at by the parsing process.

\yskip\noindent Initially \\{scrap\_ptr} is set to the position of the final
scrap to be parsed, and it doesn't change its value. The parsing process
makes sure that \\{lo\_ptr}\?\G\?\\{pp}\?+\?4, since productions have as many
as five terms,
by moving scraps from \\{hi\_ptr} to \\{lo\_ptr}. If there are
fewer than \\{pp}\?+\?4 scraps left, the positions up to \\{pp}\?+\?4 are
filled with
blanks that will not match in any productions. Parsing stops when
\\{pp}\?=\?\\{lo\_ptr}\?+\?1 and \\{hi\_ptr}\?=\?\\{scrap\_ptr}\?+\?1.

The \\{trans} array elements are declared to be of type 0\?\to\?10\_239 instead
of type \\{text\_pointer}, because the final sorting phase of \.{AWEAVE}
uses this array to contain elements of type \\{name\_pointer}. Both
of these types are subranges of 0\?\to\?10\_239.

\Y\P\?\4\X14:Types and objects visible from \\{data\_structures}\X\mathrel{+}\S%
\?\6
\&{subtype}\8\\{scrap\_index}\8\&{is}\8\p{INTEGER}\8\&{range}\80\?\to\?\\{max%
\_scraps};\6
\&{type}\8\\{cat\_type}\8\&{is}\?\8\1\&{array}\8\1(\?\\{scrap\_index}\?\2)\2\8%
\&{of}\?\8\\{eight\_bits};\6
\&{type}\8\\{trans\_type}\8\&{is}\?\8\1\&{array}\8\1(\?\\{scrap\_index}\?\2)\2%
\8\&{of}\?\8\p{INTEGER}\8\&{range}\80\?\to\?10\_239;\7
\\{cat}\?:\?\\{cat\_type};\C{category codes of scraps}\6
\\{trans}\?:\?\\{trans\_type};\C{translation texts of scraps}\6
\\{pp}\?:\?\\{scrap\_index};\C{current position for reducing productions}\6
\\{scrap\_base}\?:\?\\{scrap\_index}\?\K\?1;\C{beginning of the current scrap
sequence}\6
\\{scrap\_ptr}\?:\?\\{scrap\_index}\?\K\?0;\C{ending of the current scrap
sequence}\6
\\{lo\_ptr}\?:\?\\{scrap\_index};\C{last scrap that has been examined}\6
\\{hi\_ptr}\?:\?\\{scrap\_index};\C{first scrap that has not been examined}\6
\&{stat}\1\6
\\{max\_scr\_ptr}\?:\?\\{scrap\_index}\?\K\?0;\C{largest value assumed by %
\\{scrap\_ptr}}\2\6
\&{tats}\par
\fi

\M172. Token lists in \\{tok\_mem} are composed of the following kinds of
items for \TeX\ output.

\yskip\item{$\bullet$}ASCII codes and special codes like \\{force} and
\\{math\_rel} represent themselves;

\item{$\bullet$}\\{id\_flag}\?+\?\|p represents \.{\\\\\{{\rm identifier $p$}%
\}};

\item{$\bullet$}\\{res\_flag}\?+\?\|p represents \.{\\\&\{{\rm identifier $p$}%
\}};

\item{$\bullet$}\\{mod\_flag}\?+\?\|p represents module name \|p;

\item{$\bullet$}\\{tok\_flag}\?+\?\|p represents token list number \|p;

\item{$\bullet$}\\{inner\_tok\_flag}\?+\?\|p represents token list number \|p,
to be
translated without line-break controls.

\Y\P\?\4\X14:Types and objects visible from \\{data\_structures}\X\mathrel{+}\S%
\?\6
\\{id\_flag}\?:\?\&{constant}\8\\{sixteen\_bits}\?\K\?10\_240;\C{signifies an
identifier}\6
\\{res\_flag}\?:\?\&{constant}\8\\{sixteen\_bits}\?\K\?20\_480;\C{signifies a
reserved word}\6
\\{mod\_flag}\?:\?\&{constant}\8\\{sixteen\_bits}\?\K\?30\_720;\C{signifies a
module name}\6
\\{tok\_flag}\?:\?\&{constant}\8\\{sixteen\_bits}\?\K\?40\_960;\C{signifies a
token list}\6
\\{inner\_tok\_flag}\?:\?\&{constant}\8\\{sixteen\_bits}\?\K\?51\_200;%
\C{signifies a token list in `\ab'}\par
\fi

\M173. The production rules listed above are embedded directly into the %
\.{AWEAVE}
program, since it is easier to do this than to write an interpretive system
that would handle production systems in general. Several macros are defined
here so that the program for each production is fairly short.

All of our productions conform to the general notion that some \|k
consecutive scraps starting at some position \|j are to be replaced by a
single scrap of some category \|c whose translations is composed from the
translations of the disappearing scraps. After this production has been
applied, the production pointer \\{pp} should change by an amount \|d. Such
a production can be represented by the quadruple $(j,k,c,d)$. For example,
the production `\\{math}\?\,\?\\{math} $\RA$ \\{math}' would be represented by
`$(\\{pp},2,\\{math},-3)$'; in this case the pointer $pp$ should decrease by 3
after the production has been applied, because some productions with
\\{math} in their fourth positions might now match, but no productions have
\\{math} in the fifth position of their left-hand sides. Note that
the value of \|d is determined by the whole collection of productions, not
by an individual one.
The determination of \|d has been
done by hand in each case, based on the full set of productions but not on
the grammar of \ADA\ or on the rules for constructing the initial
scraps.

Before calling \\{reduce}, the program should have appended the tokens of
the new translation to the \\{tok\_mem} array. We commonly want to append
copies of several existing translations, and macros are defined to
simplify these common cases. For example, \\{app2}\?(\?\\{pp}\?)\? will append
the
translations of two consecutive scraps, \\{trans}\?(\?\\{pp}\?)\? and \\{trans}%
\?(\?\\{pp}\?+\?1\?)\?, to
the current token list. If the entire new translation is formed in this
way, we write `$\\{squash}(j,k,c,d)$' instead of `$\\{reduce}(j,k,c,d)$'. For
example, `\\{squash}\?(\?\\{pp}\?,\?2\?,\?\\{math}\?,-\?3\?)\?' is an
abbreviation for `\\{app2}\?(\?\\{pp}\?)\?; \\{reduce}\?(\?\\{pp}\?,\?2\?,\?%
\\{math}\?,-\?3\?)\?'.

The code below is an exact translation of the production rules into
\ADA, using such macros, and the reader should have no difficulty
understanding the format by comparing the code with the symbolic
productions as they were listed earlier.

\Y\P\4\D\\{reduce}\?\1(\?\#\?\2)\S\?\\{red}\?\1(\?\#\?\2)\?;\5
\&{goto}\8\\{found}\par
\P\4\D\\{squash}\?\1(\?\#\?\2)\S\?\\{sq}\?\1(\?\#\?\2)\?;\5
\&{goto}\8\\{found}\par
\P\4\D\\{app}\?\1(\?\#\?\2)\S\?\\{tok\_mem}\?\1(\?\\{tok\_ptr}\?\2)\K\?\#;\5
\\{incr}\?\1(\?\\{tok\_ptr}\?\2)\?\C{this is like \\{app\_tok},   but it
doesn't test for overflow}\par
\P\4\D\\{app1}\?\1(\?\#\?\2)\S\?\\{tok\_mem}\?\1(\?\\{tok\_ptr}\?\2)\K\?\\{tok%
\_flag}\?+\?\\{trans}\?\1(\?\#\?\2)\?;\5
\\{incr}\?\1(\?\\{tok\_ptr}\?\2)\?\par
\P\4\D\\{app2}\?\1(\?\#\?\2)\S\?\\{app1}\?\1(\?\#\?\2)\?;\5
\\{app1}\?\1(\?\#\?+\?1\?\2)\?\par
\P\4\D\\{app3}\?\1(\?\#\?\2)\S\?\\{app2}\?\1(\?\#\?\2)\?;\5
\\{app1}\?\1(\?\#\?+\?2\?\2)\?\par
\fi

\M174. Let us consider the big case statement for productions now, before
looking
at its context. We want to design the program so that this case statement
works, so we might as well not keep ourselves in suspense about exactly what
code needs to be provided with a proper environment.

The code here is more complicated than it need be, since some
compilers are unable to deal with procedures that contain a lot
of program text. The \\{translate} procedure, which incorporates the \&{case}
statement here, would become too long for those compilers if we did
not do something to split the cases into parts. Therefore
a separate procedure called \\{do\_cases} has been introduced.

\Y\P\?\4\X174:Match a production at \\{pp}, or increase \\{pp} if there is no
match\X\S\?\6
\&{if}\1\1\8\\{cat}\?\1(\?\\{pp}\?\2)<\?\\{if\_scrap}\2\8\&{then}\6
\\{do\_cases};\6
\4\&{else}\6
\&{case}\1\8\\{cat}\?\1(\?\\{pp}\?\2)\?\1\8\&{is}\8\6
\4\&{when}\8\\{if\_scrap}\KR\8\X194:Cases for \\{if\_scrap}\X;\6
\4\&{when}\8\\{in\_scrap}\KR\8\X195:Cases for \\{in\_scrap}\X;\6
\4\&{when}\8\\{label}\KR\8\X196:Cases for \\{label}\X;\6
\4\&{when}\8\\{loop\_scrap}\KR\8\X197:Cases for \\{loop\_scrap}\X;\6
\4\&{when}\8\\{math}\KR\8\X198:Cases for \\{math}\X;\6
\4\&{when}\8\\{mod\_scrap}\KR\8\X199:Cases for \\{mod\_scrap}\X;\6
\4\&{when}\8\\{new\_scrap}\KR\8\X200:Cases for \\{new\_scrap}\X;\6
\4\&{when}\8\\{not\_scrap}\KR\8\X201:Cases for \\{not\_scrap}\X;\6
\4\&{when}\8\\{open\_scrap}\KR\8\X202:Cases for \\{open\_scrap}\X;\6
\4\&{when}\8\\{period}\KR\8\X203:Cases for \\{period}\X;\6
\4\&{when}\8\\{private\_scrap}\KR\8\X204:Cases for \\{private\_scrap}\X;\6
\4\&{when}\8\\{proc}\KR\8\X205:Cases for \\{proc}\X;\6
\4\&{when}\8\\{return\_scrap}\KR\8\X206:Cases for \\{return\_scrap}\X;\6
\4\&{when}\8\\{separate\_scrap}\KR\8\X207:Cases for \\{separate\_scrap}\X;\6
\4\&{when}\8\\{stmt}\KR\8\X208:Cases for \\{stmt}\X;\6
\4\&{when}\8\\{when\_scrap}\KR\8\X209:Cases for \\{when\_scrap}\X;\6
\4\&{when}\8\\{with\_scrap}\KR\8\X210:Cases for \\{with\_scrap}\X;\6
\4\&{when}\8\&{others}\8\KR\8\&{null};\2\6
\4\2\&{end}\8\&{case};\6
\\{incr}\?\1(\?\\{pp}\?\2)\?;\C{if no match was found, we move to the right}\2\6
\&{end}\8\2\&{if}\1;\6
\4\BL\\{found}\EL\8\&{null}\par
\U section~215.\fi

\M175. Here is the procedure that need to be present for the reason just
explained.

\Y\P\?\4\X175:Declaration of subprocedure for \\{translate}\X\S\?\6
\&{procedure}\8\1\\{do\_cases}\?\8\&{is}\8\&{separate}\?\2;\par
\U section~218.\fi

\M176\*.
\Y\P\O{parsing-do\_cases.adb}\?\6
\?\&{separate}\?\8\1(\?\\{parsing}\?\2)\?\6
\&{procedure}\8\1\\{do\_cases}\2\8\&{is}\8\1\6
\4\&{begin}\6
\&{case}\1\8\\{cat}\?\1(\?\\{pp}\?\2)\?\1\8\&{is}\8\6
\4\&{when}\8\\{accept\_scrap}\KR\8\X177:Cases for \\{accept\_scrap}\X;\6
\4\&{when}\8\\{and\_scrap}\KR\8\X178:Cases for \\{and\_scrap}\X;\6
\4\&{when}\8\\{arrow}\KR\8\X179:Cases for \\{arrow}\X;\6
\4\&{when}\8\\{base}\KR\8\X180:Cases for \\{base}\X;\6
\4\&{when}\8\\{begin\_label\_scrap}\KR\8\X181:Cases for \\{begin\_label\_scrap}%
\X;\6
\4\&{when}\8\\{beginning}\KR\8\X182:Cases for \\{beginning}\X;\6
\4\&{when}\8\\{case\_scrap}\KR\8\X183:Cases for \\{case\_scrap}\X;\6
\4\&{when}\8\\{clause}\KR\8\X184:Cases for \\{clause}\X;\6
\4\&{when}\8\\{colon}\KR\8\X185:Cases for \\{colon}\X;\6
\4\&{when}\8\\{decl}\KR\8\X186:Cases for \\{decl}\X;\6
\4\&{when}\8\\{digit\_scrap}\KR\8\X187:Cases for \\{digit\_scrap}\X;\6
\4\&{when}\8\\{else\_scrap}\KR\8\X188:Cases for \\{else\_scrap}\X;\6
\4\&{when}\8\\{ending}\KR\8\X189:Cases for \\{ending}\X;\6
\4\&{when}\8\\{exception\_scrap}\KR\8\X190:Cases for \\{exception\_scrap}\X;\6
\4\&{when}\8\\{exp}\KR\8\X191:Cases for \\{exp}\X;\6
\4\&{when}\8\\{generic\_scrap}\KR\8\X192:Cases for \\{generic\_scrap}\X;\6
\4\&{when}\8\\{gproc}\KR\8\X193:Cases for \\{gproc}\X;\6
\4\&{when}\8\&{others}\8\KR\8\&{null};\2\6
\4\2\&{end}\8\&{case};\6
\\{incr}\?\1(\?\\{pp}\?\2)\?;\6
\4\BL\\{found}\EL\8\&{null};\2\6
\&{end}\8\\{do\_cases};\par
\fi

\M177. Now comes the code that tries to match each production that starts
with a particular type of scrap. Whenever a match is discovered,
the \\{squash} or \\{reduce} macro will cause the appropriate action
to be performed, followed by \&{goto}\8\\{found}.

\Y\P\?\4\X177:Cases for \\{accept\_scrap}\X\S\?\6
\&{if}\1\1\8\\{cat}\?\1(\?\\{pp}\?+\?1\?\to\?\\{pp}\?+\?2\?\2)=\1(\?\\{math}\?,%
\39\?\\{follows}\?\2)\?\2\8\&{then}\6
\\{app1}\?\1(\?\\{pp}\?\2)\?;\5
\\{app}\?\1(\?\\{indent}\?\2)\?;\5
\\{app}\?\1(\?\\{pdollar}\?\2)\?;\5
\\{app1}\?\1(\?\\{pp}\?+\?1\?\2)\?;\5
\\{app}\?\1(\?\\{pdollar}\?\2)\?;\5
\\{app}\?\1(\?\\{outdent}\?\2)\?;\5
\\{app1}\?\1(\?\\{pp}\?+\?2\?\2)\?;\6
\\{reduce}\?\1(\?\\{pp}\?,\39\?3\?,\39\?\\{beginning}\?,\39-\?2\?\2)\?;\6
\4\&{elsif}\1\8\\{cat}\?\1(\?\\{pp}\?+\?1\?\to\?\\{pp}\?+\?2\?\2)=\1(\?\\{math}%
\?,\39\?\\{semi}\?\2)\?\2\8\&{then}\6
\\{app1}\?\1(\?\\{pp}\?\2)\?;\5
\\{app}\?\1(\?\\{pdollar}\?\2)\?;\5
\\{app1}\?\1(\?\\{pp}\?+\?1\?\2)\?;\5
\\{app}\?\1(\?\\{pdollar}\?\2)\?;\5
\\{app1}\?\1(\?\\{pp}\?+\?2\?\2)\?;\6
\\{reduce}\?\1(\?\\{pp}\?,\39\?3\?,\39\?\\{stmt}\?,\39-\?3\?\2)\?;\2\6
\&{end}\8\2\&{if}\1\par
\U section~176\*.\fi

\M178. \P\?\X178:Cases for \\{and\_scrap}\X\S\?\6
\&{if}\1\1\8\\{cat}\?\1(\?\\{pp}\?+\?1\?\2)=\?\\{then\_scrap}\2\8\&{then}\6
\\{app}\?\1(\?\H{\\}\?\2)\?;\5
\\{app}\?\1(\?\H{A}\?\2)\?;\5
\\{app}\?\1(\?\H{N}\?\2)\?;\5
\\{app}\?\1(\?\H{T}\?\2)\?;\6
\\{reduce}\?\1(\?\\{pp}\?,\39\?2\?,\39\?\\{math}\?,\39-\?3\?\2)\?;\2\6
\&{end}\8\2\&{if}\1\par
\U section~176\*.\fi

\M179. \P\?\X179:Cases for \\{arrow}\X\S\?\6
\\{squash}\?\1(\?\\{pp}\?,\39\?1\?,\39\?\\{math}\?,\39-\?3\?\2)\?\par
\U section~176\*.\fi

\M180. \P\?\X180:Cases for \\{base}\X\S\?\6
\&{if}\1\1\8\\{cat}\?\1(\?\\{pp}\?+\?1\?\2)=\?\\{base}\2\8\&{then}\6
\\{app}\?\1(\?\H{\\}\?\2)\?;\5
\\{app}\?\1(\?\H{b}\?\2)\?;\5
\\{app}\?\1(\?\H{n}\?\2)\?;\5
\\{app}\?\1(\?\H{\{}\?\2)\?;\5
\\{app1}\?\1(\?\\{pp}\?\2)\?;\5
\\{app}\?\1(\?\H{\}}\?\2)\?;\5
\\{app}\?\1(\?\H{\{}\?\2)\?;\5
\\{app1}\?\1(\?\\{pp}\?+\?1\?\2)\?;\5
\\{app}\?\1(\?\H{\}}\?\2)\?;\6
\\{reduce}\?\1(\?\\{pp}\?,\39\?2\?,\39\?\\{math}\?,\39-\?3\?\2)\?;\2\6
\&{end}\8\2\&{if}\1\par
\U section~176\*.\fi

\M181. \P\?\X181:Cases for \\{begin\_label\_scrap}\X\S\?\6
\&{if}\1\1\8\\{cat}\?\1(\?\\{pp}\?+\?1\?\to\?\\{pp}\?+\?2\?\2)=\1(\?\\{math}\?,%
\39\?\\{end\_label\_scrap}\?\2)\?\2\8\&{then}\6
\\{app1}\?\1(\?\\{pp}\?\2)\?;\5
\\{app}\?\1(\?\\{pdollar}\?\2)\?;\5
\\{app1}\?\1(\?\\{pp}\?+\?1\?\2)\?;\5
\\{app}\?\1(\?\\{pdollar}\?\2)\?;\5
\\{app1}\?\1(\?\\{pp}\?+\?2\?\2)\?;\6
\\{reduce}\?\1(\?\\{pp}\?,\39\?3\?,\39\?\\{label}\?,\39-\?1\?\2)\?;\2\6
\&{end}\8\2\&{if}\1\par
\U section~176\*.\fi

\M182. \P\?\X182:Cases for \\{beginning}\X\S\?\6
\&{if}\1\1\8\\{cat}\?\1(\?\\{pp}\?+\?1\?\2)=\?\\{stmt}\2\8\&{then}\6
\\{app}\?\1(\?\\{force}\?\2)\?;\5
\\{app}\?\1(\?\\{indent}\?\2)\?;\5
\\{app1}\?\1(\?\\{pp}\?\2)\?;\5
\\{app}\?\1(\?\\{force}\?\2)\?;\5
\\{app1}\?\1(\?\\{pp}\?+\?1\?\2)\?;\6
\\{reduce}\?\1(\?\\{pp}\?,\39\?2\?,\39\?\\{stmt}\?,\39-\?3\?\2)\?;\2\6
\&{end}\8\2\&{if}\1\par
\U section~176\*.\fi

\M183. \P\?\X183:Cases for \\{case\_scrap}\X\S\?\6
\&{if}\1\1\8\\{cat}\?\1(\?\\{pp}\?+\?1\?\to\?\\{pp}\?+\?3\?\2)=\1(\?\\{math}\?,%
\39\?\\{follows}\?,\39\?\\{decl}\?\2)\?\2\8\&{then}\6
\\{app1}\?\1(\?\\{pp}\?\2)\?;\5
\\{app}\?\1(\?\\{indent}\?\2)\?;\5
\\{app}\?\1(\?\\{pdollar}\?\2)\?;\5
\\{app}\?\1(\?\\{sspace}\?\2)\?;\5
\\{app1}\?\1(\?\\{pp}\?+\?1\?\2)\?;\5
\\{app}\?\1(\?\\{pdollar}\?\2)\?;\5
\\{app}\?\1(\?\\{indent}\?\2)\?;\5
\\{app1}\?\1(\?\\{pp}\?+\?2\?\2)\?;\5
\\{app}\?\1(\?\\{force}\?\2)\?;\5
\\{app1}\?\1(\?\\{pp}\?+\?3\?\2)\?;\6
\\{reduce}\?\1(\?\\{pp}\?,\39\?4\?,\39\?\\{decl}\?,\39-\?3\?\2)\?;\6
\4\&{elsif}\1\8\\{cat}\?\1(\?\\{pp}\?+\?1\?\to\?\\{pp}\?+\?3\?\2)=\1(\?\\{math}%
\?,\39\?\\{follows}\?,\39\?\\{stmt}\?\2)\?\2\8\&{then}\6
\\{app}\?\1(\?\\{force}\?\2)\?;\5
\\{app1}\?\1(\?\\{pp}\?\2)\?;\5
\\{app}\?\1(\?\\{indent}\?\2)\?;\5
\\{app}\?\1(\?\\{pdollar}\?\2)\?;\5
\\{app}\?\1(\?\\{sspace}\?\2)\?;\5
\\{app1}\?\1(\?\\{pp}\?+\?1\?\2)\?;\5
\\{app}\?\1(\?\\{pdollar}\?\2)\?;\5
\\{app}\?\1(\?\\{indent}\?\2)\?;\5
\\{app1}\?\1(\?\\{pp}\?+\?2\?\2)\?;\5
\\{app}\?\1(\?\\{force}\?\2)\?;\5
\\{app1}\?\1(\?\\{pp}\?+\?3\?\2)\?;\6
\\{reduce}\?\1(\?\\{pp}\?,\39\?4\?,\39\?\\{stmt}\?,\39-\?3\?\2)\?;\2\6
\&{end}\8\2\&{if}\1\par
\U section~176\*.\fi

\M184. \P\?\X184:Cases for \\{clause}\X\S\?\6
\&{if}\1\1\8\\{cat}\?\1(\?\\{pp}\?+\?1\?\2)=\?\\{decl}\2\8\&{then}\6
\\{app}\?\1(\?\\{force}\?\2)\?;\5
\\{app}\?\1(\?\\{backup}\?\2)\?;\5
\\{app1}\?\1(\?\\{pp}\?\2)\?;\5
\\{app}\?\1(\?\\{force}\?\2)\?;\5
\\{app1}\?\1(\?\\{pp}\?+\?1\?\2)\?;\6
\\{reduce}\?\1(\?\\{pp}\?,\39\?2\?,\39\?\\{decl}\?,\39-\?3\?\2)\?;\6
\4\&{elsif}\1\8\\{cat}\?\1(\?\\{pp}\?+\?1\?\2)=\?\\{stmt}\2\8\&{then}\6
\\{app}\?\1(\?\\{force}\?\2)\?;\5
\\{app}\?\1(\?\\{backup}\?\2)\?;\5
\\{app1}\?\1(\?\\{pp}\?\2)\?;\5
\\{app}\?\1(\?\\{force}\?\2)\?;\5
\\{app1}\?\1(\?\\{pp}\?+\?1\?\2)\?;\6
\\{reduce}\?\1(\?\\{pp}\?,\39\?2\?,\39\?\\{stmt}\?,\39-\?3\?\2)\?;\2\6
\&{end}\8\2\&{if}\1\par
\U section~176\*.\fi

\M185. \P\?\X185:Cases for \\{colon}\X\S\?\6
\&{if}\1\1\8\\{cat}\?\1(\?\\{pp}\?+\?1\?\2)=\?\\{exception\_scrap}\2\8\&{then}\6
\\{squash}\?\1(\?\\{pp}\?+\?1\?,\39\?1\?,\39\?\\{math}\?,\39-\?1\?\2)\?;\6
\4\&{elsif}\1\8\\{cat}\?\1(\?\\{pp}\?+\?1\?\2)=\?\\{in\_scrap}\2\8\&{then}\6
\\{app1}\?\1(\?\\{pp}\?\2)\?;\5
\\{app}\?\1(\?\H{\\}\?\2)\?;\5
\\{app}\?\1(\?\H{\&}\?\2)\?;\5
\\{app}\?\1(\?\H{\{}\?\2)\?;\5
\\{app}\?\1(\?\H{i}\?\2)\?;\5
\\{app}\?\1(\?\H{n}\?\2)\?;\5
\\{app}\?\1(\?\H{\}}\?\2)\?;\5
\\{app}\?\1(\?\\{sspace}\?\2)\?;\6
\\{reduce}\?\1(\?\\{pp}\?,\39\?2\?,\39\?\\{colon}\?,\39-\?2\?\2)\?;\2\6
\&{end}\8\2\&{if}\1\par
\U section~176\*.\fi

\M186. \P\?\X186:Cases for \\{decl}\X\S\?\6
\&{if}\1\1\8\\{cat}\?\1(\?\\{pp}\?+\?1\?\2)=\?\\{beginning}\2\8\&{then}\6
\\{app1}\?\1(\?\\{pp}\?\2)\?;\5
\\{app}\?\1(\?\\{force}\?\2)\?;\5
\\{app}\?\1(\?\\{backup}\?\2)\?;\5
\\{app1}\?\1(\?\\{pp}\?+\?1\?\2)\?;\5
\\{app}\?\1(\?\\{force}\?\2)\?;\6
\\{reduce}\?\1(\?\\{pp}\?,\39\?2\?,\39\?\\{stmt}\?,\39-\?3\?\2)\?;\6
\4\&{elsif}\1\8\\{cat}\?\1(\?\\{pp}\?+\?1\?\2)=\?\\{decl}\2\8\&{then}\6
\\{app1}\?\1(\?\\{pp}\?\2)\?;\5
\\{app}\?\1(\?\\{force}\?\2)\?;\5
\\{app1}\?\1(\?\\{pp}\?+\?1\?\2)\?;\6
\\{reduce}\?\1(\?\\{pp}\?,\39\?2\?,\39\?\\{decl}\?,\39-\?3\?\2)\?;\6
\4\&{elsif}\1\8\\{cat}\?\1(\?\\{pp}\?+\?1\?\to\?\\{pp}\?+\?2\?\2)=\1(\?\\{mod%
\_scrap}\?,\39\?\\{semi}\?\2)\?\2\8\&{then}\6
\\{app1}\?\1(\?\\{pp}\?\2)\?;\5
\\{app}\?\1(\?\\{force}\?\2)\?;\5
\\{app2}\?\1(\?\\{pp}\?+\?1\?\2)\?;\6
\\{reduce}\?\1(\?\\{pp}\?,\39\?3\?,\39\?\\{decl}\?,\39-\?3\?\2)\?;\6
\4\&{elsif}\1\8\\{cat}\?\1(\?\\{pp}\?+\?1\?\2)=\?\\{terminator}\2\8\&{then}\6
\\{app1}\?\1(\?\\{pp}\?\2)\?;\5
\\{app}\?\1(\?\\{cancel}\?\2)\?;\5
\\{app}\?\1(\?\\{outdent}\?\2)\?;\5
\\{app}\?\1(\?\\{force}\?\2)\?;\5
\\{app1}\?\1(\?\\{pp}\?+\?1\?\2)\?;\5
\\{app}\?\1(\?\\{force}\?\2)\?;\6
\\{reduce}\?\1(\?\\{pp}\?,\39\?2\?,\39\?\\{decl}\?,\39-\?3\?\2)\?;\6
\4\&{elsif}\1\8\\{cat}\?\1(\?\\{pp}\?+\?1\?\2)=\?\\{with\_scrap}\2\8\&{then}\6
\\{app1}\?\1(\?\\{pp}\?\2)\?;\5
\\{app}\?\1(\?\\{force}\?\2)\?;\5
\\{app1}\?\1(\?\\{pp}\?+\?1\?\2)\?;\6
\\{reduce}\?\1(\?\\{pp}\?,\39\?2\?,\39\?\\{with\_scrap}\?,\39-\?1\?\2)\?;\2\6
\&{end}\8\2\&{if}\1\par
\U section~176\*.\fi

\M187. \P\?\X187:Cases for \\{digit\_scrap}\X\S\?\6
\&{if}\1\1\8\\{cat}\?\1(\?\\{pp}\?+\?1\?\2)=\?\\{base}\2\8\&{then}\6
\\{app1}\?\1(\?\\{pp}\?\2)\?;\6
\\{reduce}\?\1(\?\\{pp}\?,\39\?2\?,\39\?\\{base}\?,\39-\?1\?\2)\?;\6
\4\&{elsif}\1\8\\{cat}\?\1(\?\\{pp}\?+\?1\?\2)=\?\\{digit\_scrap}\2\8\&{then}\6
\\{squash}\?\1(\?\\{pp}\?,\39\?2\?,\39\?\\{digit\_scrap}\?,\39-\?2\?\2)\?;\6
\4\&{elsif}\1\8\\{cat}\?\1(\?\\{pp}\?+\?1\?\2)=\?\\{period}\2\8\&{then}\6
\\{squash}\?\1(\?\\{pp}\?,\39\?2\?,\39\?\\{digit\_scrap}\?,\39-\?2\?\2)\?;\6
\4\&{else}\6
\\{squash}\?\1(\?\\{pp}\?,\39\?1\?,\39\?\\{math}\?,\39-\?3\?\2)\?;\2\6
\&{end}\8\2\&{if}\1\par
\U section~176\*.\fi

\M188. \P\?\X188:Cases for \\{else\_scrap}\X\S\?\6
\&{if}\1\1\8\\{cat}\?\1(\?\\{pp}\?+\?1\?\2)=\?\\{else\_scrap}\2\8\&{then}\6
\\{app}\?\1(\?\H{\\}\?\2)\?;\5
\\{app}\?\1(\?\H{O}\?\2)\?;\5
\\{app}\?\1(\?\H{R}\?\2)\?;\5
\\{app}\?\1(\?\H{E}\?\2)\?;\6
\\{reduce}\?\1(\?\\{pp}\?,\39\?2\?,\39\?\\{math}\?,\39-\?3\?\2)\?;\6
\4\&{elsif}\1\8\\{cat}\?\1(\?\\{pp}\?+\?1\?\2)=\?\\{when\_scrap}\2\8\&{then}\6
\\{app}\?\1(\?\H{\\}\?\2)\?;\5
\\{app}\?\1(\?\H{O}\?\2)\?;\5
\\{app}\?\1(\?\H{R}\?\2)\?;\5
\\{app}\?\1(\?\H{W}\?\2)\?;\6
\\{reduce}\?\1(\?\\{pp}\?,\39\?2\?,\39\?\\{when\_scrap}\?,\39-\?1\?\2)\?;\2\6
\&{end}\8\2\&{if}\1\par
\U section~176\*.\fi

\M189.
\Y\P\4\D\\{compress\_end}\?\S\?\\{app1}\?\1(\?\\{pp}\?\2)\?;\5
\\{app}\?\1(\?\\{sspace}\?\2)\?;\5
\\{app2}\?\1(\?\\{pp}\?+\?1\?\2)\?\par
\Y\P\?\4\X189:Cases for \\{ending}\X\S\?\6
\&{if}\1\1\8\\{cat}\?\1(\?\\{pp}\?+\?1\?\to\?\\{pp}\?+\?2\?\2)=\1(\?%
\\{beginning}\?,\39\?\\{semi}\?\2)\?\2\8\&{then}\6
\\{compress\_end};\6
\\{reduce}\?\1(\?\\{pp}\?,\39\?3\?,\39\?\\{terminator}\?,\39-\?1\?\2)\?;\6
\4\&{elsif}\1\8\\{cat}\?\1(\?\\{pp}\?+\?1\?\to\?\\{pp}\?+\?2\?\2)=\1(\?\\{case%
\_scrap}\?,\39\?\\{semi}\?\2)\?\2\8\&{then}\6
\\{app}\?\1(\?\\{backup}\?\2)\?;\5
\\{app}\?\1(\?\\{outdent}\?\2)\?;\5
\\{compress\_end};\6
\\{reduce}\?\1(\?\\{pp}\?,\39\?3\?,\39\?\\{terminator}\?,\39-\?1\?\2)\?;\6
\4\&{elsif}\1\8\\{cat}\?\1(\?\\{pp}\?+\?1\?\to\?\\{pp}\?+\?2\?\2)=\1(\?\\{if%
\_scrap}\?,\39\?\\{semi}\?\2)\?\2\8\&{then}\6
\\{app1}\?\1(\?\\{pp}\?\2)\?;\5
\\{app}\?\1(\?\\{sspace}\?\2)\?;\5
\\{app}\?\1(\?\\{outdent}\?\2)\?;\5
\\{app2}\?\1(\?\\{pp}\?+\?1\?\2)\?;\6
\\{reduce}\?\1(\?\\{pp}\?,\39\?3\?,\39\?\\{terminator}\?,\39-\?1\?\2)\?;\6
\4\&{elsif}\1\8\\{cat}\?\1(\?\\{pp}\?+\?1\?\to\?\\{pp}\?+\?2\?\2)=\1(\?\\{loop%
\_scrap}\?,\39\?\\{semi}\?\2)\?\2\8\&{then}\6
\\{compress\_end};\6
\\{reduce}\?\1(\?\\{pp}\?,\39\?3\?,\39\?\\{terminator}\?,\39-\?1\?\2)\?;\6
\4\&{elsif}\1\8\\{cat}\?\1(\?\\{pp}\?+\?1\?\to\?\\{pp}\?+\?3\?\2)=\1(\?\\{loop%
\_scrap}\?,\39\?\\{math}\?,\39\?\\{semi}\?\2)\?\2\8\&{then}\6
\\{app1}\?\1(\?\\{pp}\?\2)\?;\5
\\{app}\?\1(\?\\{sspace}\?\2)\?;\5
\\{app1}\?\1(\?\\{pp}\?+\?1\?\2)\?;\5
\\{app}\?\1(\?\\{pdollar}\?\2)\?;\5
\\{app}\?\1(\?\\{sspace}\?\2)\?;\5
\\{app1}\?\1(\?\\{pp}\?+\?2\?\2)\?;\5
\\{app}\?\1(\?\\{pdollar}\?\2)\?;\5
\\{app1}\?\1(\?\\{pp}\?+\?3\?\2)\?;\6
\\{reduce}\?\1(\?\\{pp}\?,\39\?4\?,\39\?\\{terminator}\?,\39-\?1\?\2)\?;\6
\4\&{elsif}\1\8\\{cat}\?\1(\?\\{pp}\?+\?1\?\to\?\\{pp}\?+\?2\?\2)=\1(\?\\{math}%
\?,\39\?\\{semi}\?\2)\?\2\8\&{then}\6
\\{app1}\?\1(\?\\{pp}\?\2)\?;\5
\\{app}\?\1(\?\\{pdollar}\?\2)\?;\5
\\{app}\?\1(\?\\{sspace}\?\2)\?;\5
\\{app1}\?\1(\?\\{pp}\?+\?1\?\2)\?;\5
\\{app}\?\1(\?\\{pdollar}\?\2)\?;\5
\\{app1}\?\1(\?\\{pp}\?+\?2\?\2)\?;\6
\\{reduce}\?\1(\?\\{pp}\?,\39\?3\?,\39\?\\{terminator}\?,\39-\?1\?\2)\?;\6
\4\&{elsif}\1\8\\{cat}\?\1(\?\\{pp}\?+\?1\?\to\?\\{pp}\?+\?2\?\2)=\1(\?%
\\{record\_scrap}\?,\39\?\\{semi}\?\2)\?\2\8\&{then}\6
\\{app}\?\1(\?\\{outdent}\?\2)\?;\5
\\{compress\_end};\6
\\{reduce}\?\1(\?\\{pp}\?,\39\?3\?,\39\?\\{terminator}\?,\39-\?1\?\2)\?;\6
\4\&{elsif}\1\8\\{cat}\?\1(\?\\{pp}\?+\?1\?\2)=\?\\{semi}\2\8\&{then}\6
\\{squash}\?\1(\?\\{pp}\?,\39\?2\?,\39\?\\{terminator}\?,\39-\?1\?\2)\?;\2\6
\&{end}\8\2\&{if}\1\par
\U section~176\*.\fi

\M190. \P\?\X190:Cases for \\{exception\_scrap}\X\S\?\6
\\{app}\?\1(\?\\{backup}\?\2)\?;\5
\\{app1}\?\1(\?\\{pp}\?\2)\?;\6
\\{reduce}\?\1(\?\\{pp}\?,\39\?1\?,\39\?\\{beginning}\?,\39-\?2\?\2)\?\par
\U section~176\*.\fi

\M191. \P\?\X191:Cases for \\{exp}\X\S\?\6
\&{if}\1\1\?\8\1(\?\\{cat}\?\1(\?\\{pp}\?+\?1\?\2)=\?\\{digit\_scrap}\?\2)\W\1(%
\?\\{cat}\?\1(\?\\{pp}\?+\?2\?\2)\I\?\\{digit\_scrap}\?\2)\?\2\8\&{then}\6
\\{app2}\?\1(\?\\{pp}\?\2)\?;\5
\\{app}\?\1(\?\H{\}}\?\2)\?;\6
\\{reduce}\?\1(\?\\{pp}\?,\39\?2\?,\39\?\\{math}\?,\39-\?3\?\2)\?;\6
\4\&{elsif}\1\?\8\1(\?\\{cat}\?\1(\?\\{pp}\?+\?1\?\2)=\?\\{math}\?\2)\W\1(\1(\?%
\\{cat}\?\1(\?\\{pp}\?+\?2\?\2)=\?\\{digit\_scrap}\?\2)\W\1(\?\\{cat}\?\1(\?%
\\{pp}\?+\?3\?\2)\I\?\\{digit\_scrap}\?\2)\2)\?\2\8\&{then}\6
\\{app3}\?\1(\?\\{pp}\?\2)\?;\5
\\{app}\?\1(\?\H{\}}\?\2)\?;\6
\\{reduce}\?\1(\?\\{pp}\?,\39\?3\?,\39\?\\{math}\?,\39-\?3\?\2)\?;\2\6
\&{end}\8\2\&{if}\1\par
\U section~176\*.\fi

\M192. \P\?\X192:Cases for \\{generic\_scrap}\X\S\?\6
\&{if}\1\1\8\\{cat}\?\1(\?\\{pp}\?+\?1\?\2)=\?\\{decl}\2\8\&{then}\6
\\{app1}\?\1(\?\\{pp}\?\2)\?;\5
\\{app}\?\1(\?\\{force}\?\2)\?;\5
\\{app1}\?\1(\?\\{pp}\?+\?1\?\2)\?;\6
\\{reduce}\?\1(\?\\{pp}\?,\39\?2\?,\39\?\\{generic\_scrap}\?,\39\?0\?\2)\?;\6
\4\&{elsif}\1\8\\{cat}\?\1(\?\\{pp}\?+\?1\?\2)=\?\\{proc}\2\8\&{then}\6
\\{app}\?\1(\?\\{indent}\?\2)\?;\5
\\{app1}\?\1(\?\\{pp}\?\2)\?;\5
\\{app}\?\1(\?\\{outdent}\?\2)\?;\5
\\{app}\?\1(\?\\{force}\?\2)\?;\5
\\{app1}\?\1(\?\\{pp}\?+\?1\?\2)\?;\6
\\{reduce}\?\1(\?\\{pp}\?,\39\?2\?,\39\?\\{proc}\?,\39-\?1\?\2)\?;\2\6
\&{end}\8\2\&{if}\1\par
\U section~176\*.\fi

\M193. \P\?\X193:Cases for \\{gproc}\X\S\?\6
\&{if}\1\1\8\\{cat}\?\1(\?\\{pp}\?+\?1\?\to\?\\{pp}\?+\?2\?\2)=\1(\?\\{math}\?,%
\39\?\\{semi}\?\2)\?\2\8\&{then}\6
\\{app1}\?\1(\?\\{pp}\?\2)\?;\5
\\{app}\?\1(\?\\{indent}\?\2)\?;\5
\\{app}\?\1(\?\\{pdollar}\?\2)\?;\5
\\{app1}\?\1(\?\\{pp}\?+\?1\?\2)\?;\5
\\{app}\?\1(\?\\{pdollar}\?\2)\?;\5
\\{app}\?\1(\?\\{outdent}\?\2)\?;\5
\\{app1}\?\1(\?\\{pp}\?+\?2\?\2)\?;\6
\\{reduce}\?\1(\?\\{pp}\?,\39\?3\?,\39\?\\{decl}\?,\39-\?3\?\2)\?;\6
\4\&{elsif}\1\8\\{cat}\?\1(\?\\{pp}\?+\?1\?\to\?\\{pp}\?+\?4\?\2)=\1(\?\\{math}%
\?,\39\?\\{follows}\?,\39\?\\{math}\?,\39\?\\{semi}\?\2)\?\2\8\&{then}\6
\\{app1}\?\1(\?\\{pp}\?\2)\?;\5
\\{app}\?\1(\?\\{indent}\?\2)\?;\5
\\{app}\?\1(\?\\{pdollar}\?\2)\?;\5
\\{app3}\?\1(\?\\{pp}\?+\?1\?\2)\?;\5
\\{app}\?\1(\?\\{pdollar}\?\2)\?;\5
\\{app}\?\1(\?\\{outdent}\?\2)\?;\5
\\{app1}\?\1(\?\\{pp}\?+\?4\?\2)\?;\6
\\{reduce}\?\1(\?\\{pp}\?,\39\?5\?,\39\?\\{decl}\?,\39-\?3\?\2)\?;\2\6
\&{end}\8\2\&{if}\1\par
\U section~176\*.\fi

\M194. \P\?\X194:Cases for \\{if\_scrap}\X\S\?\6
\&{if}\1\1\8\\{cat}\?\1(\?\\{pp}\?+\?1\?\to\?\\{pp}\?+\?3\?\2)=\1(\?\\{math}\?,%
\39\?\\{then\_scrap}\?,\39\?\\{stmt}\?\2)\?\2\8\&{then}\6
\\{app}\?\1(\?\\{force}\?\2)\?;\5
\\{app1}\?\1(\?\\{pp}\?\2)\?;\5
\\{app}\?\1(\?\\{indent}\?\2)\?;\5
\\{app}\?\1(\?\\{pdollar}\?\2)\?;\5
\\{app}\?\1(\?\\{sspace}\?\2)\?;\5
\\{app1}\?\1(\?\\{pp}\?+\?1\?\2)\?;\5
\\{app}\?\1(\?\\{pdollar}\?\2)\?;\5
\\{app}\?\1(\?\\{outdent}\?\2)\?;\5
\\{app1}\?\1(\?\\{pp}\?+\?2\?\2)\?;\5
\\{app}\?\1(\?\\{force}\?\2)\?;\5
\\{app1}\?\1(\?\\{pp}\?+\?3\?\2)\?;\6
\\{reduce}\?\1(\?\\{pp}\?,\39\?4\?,\39\?\\{stmt}\?,\39-\?3\?\2)\?;\2\6
\&{end}\8\2\&{if}\1\par
\U section~174.\fi

\M195. \P\?\X195:Cases for \\{in\_scrap}\X\S\?\6
\\{squash}\?\1(\?\\{pp}\?,\39\?1\?,\39\?\\{math}\?,\39-\?3\?\2)\?\par
\U section~174.\fi

\M196. \P\?\X196:Cases for \\{label}\X\S\?\6
\&{if}\1\1\8\\{cat}\?\1(\?\\{pp}\?+\?1\?\2)=\?\\{label}\2\8\&{then}\6
\\{app1}\?\1(\?\\{pp}\?\2)\?;\5
\\{app}\?\1(\?\\{sspace}\?\2)\?;\5
\\{app}\?\1(\?\\{opt}\?\2)\?;\5
\\{app}\?\1(\?\H{7}\?\2)\?;\5
\\{app}\?\1(\?\\{pp}\?+\?1\?\2)\?;\6
\\{reduce}\?\1(\?\\{pp}\?,\39\?2\?,\39\?\\{label}\?,\39-\?1\?\2)\?;\6
\4\&{elsif}\1\8\\{cat}\?\1(\?\\{pp}\?+\?1\?\2)=\?\\{stmt}\2\8\&{then}\6
\\{app}\?\1(\?\\{force}\?\2)\?;\5
\\{app}\?\1(\?\\{backup}\?\2)\?;\5
\\{app1}\?\1(\?\\{pp}\?\2)\?;\5
\\{app}\?\1(\?\\{sspace}\?\2)\?;\5
\\{app}\?\1(\?\\{cancel}\?\2)\?;\5
\\{app1}\?\1(\?\\{pp}\?+\?1\?\2)\?;\6
\\{reduce}\?\1(\?\\{pp}\?,\39\?2\?,\39\?\\{stmt}\?,\39-\?3\?\2)\?;\2\6
\&{end}\8\2\&{if}\1\par
\U section~174.\fi

\M197. \P\?\X197:Cases for \\{loop\_scrap}\X\S\?\6
\\{squash}\?\1(\?\\{pp}\?,\39\?1\?,\39\?\\{beginning}\?,\39-\?2\?\2)\?\par
\U section~174.\fi

\M198. \P\?\X198:Cases for \\{math}\X\S\?\6
\&{if}\1\1\8\\{cat}\?\1(\?\\{pp}\?+\?1\?\to\?\\{pp}\?+\?2\?\2)=\1(\?\\{and%
\_scrap}\?,\39\?\\{math}\?\2)\?\2\8\&{then}\6
\\{squash}\?\1(\?\\{pp}\?,\39\?3\?,\39\?\\{math}\?,\39-\?3\?\2)\?;\6
\4\&{elsif}\1\8\\{cat}\?\1(\?\\{pp}\?+\?1\?\to\?\\{pp}\?+\?2\?\2)=\1(\?%
\\{colon}\?,\39\?\\{beginning}\?\2)\?\2\8\&{then}\6
\\{app}\?\1(\?\\{backup}\?\2)\?;\5
\\{app}\?\1(\?\\{pdollar}\?\2)\?;\5
\\{app1}\?\1(\?\\{pp}\?\2)\?;\5
\\{app}\?\1(\?\\{sspace}\?\2)\?;\5
\\{app}\?\1(\?\\{pdollar}\?\2)\?;\5
\\{app1}\?\1(\?\\{pp}\?+\?1\?\2)\?;\5
\\{app}\?\1(\?\\{sspace}\?\2)\?;\5
\\{app1}\?\1(\?\\{pp}\?+\?2\?\2)\?;\6
\\{reduce}\?\1(\?\\{pp}\?,\39\?3\?,\39\?\\{beginning}\?,\39-\?3\?\2)\?;\6
\4\&{elsif}\1\8\\{cat}\?\1(\?\\{pp}\?+\?1\?\to\?\\{pp}\?+\?2\?\2)=\1(\?%
\\{colon}\?,\39\?\\{decl}\?\2)\?\2\8\&{then}\6
\\{app}\?\1(\?\\{force}\?\2)\?;\5
\\{app}\?\1(\?\\{backup}\?\2)\?;\5
\\{app}\?\1(\?\\{pdollar}\?\2)\?;\5
\\{app1}\?\1(\?\\{pp}\?\2)\?;\5
\\{app}\?\1(\?\\{sspace}\?\2)\?;\5
\\{app}\?\1(\?\\{pdollar}\?\2)\?;\5
\\{app1}\?\1(\?\\{pp}\?+\?1\?\2)\?;\5
\\{app}\?\1(\?\\{sspace}\?\2)\?;\5
\\{app}\?\1(\?\\{cancel}\?\2)\?;\5
\\{app1}\?\1(\?\\{pp}\?+\?2\?\2)\?;\6
\\{reduce}\?\1(\?\\{pp}\?,\39\?3\?,\39\?\\{decl}\?,\39-\?3\?\2)\?;\6
\4\&{elsif}\1\8\\{cat}\?\1(\?\\{pp}\?+\?1\?\to\?\\{pp}\?+\?3\?\2)=\1(\?%
\\{colon}\?,\39\?\\{math}\?,\39\?\\{semi}\?\2)\?\2\8\&{then}\6
\\{app}\?\1(\?\\{pdollar}\?\2)\?;\5
\\{app3}\?\1(\?\\{pp}\?\2)\?;\5
\\{app}\?\1(\?\\{pdollar}\?\2)\?;\5
\\{app1}\?\1(\?\\{pp}\?+\?3\?\2)\?;\6
\\{reduce}\?\1(\?\\{pp}\?,\39\?4\?,\39\?\\{decl}\?,\39-\?3\?\2)\?;\6
\4\&{elsif}\1\8\\{cat}\?\1(\?\\{pp}\?+\?1\?\to\?\\{pp}\?+\?2\?\2)=\1(\?\\{else%
\_scrap}\?,\39\?\\{math}\?\2)\?\2\8\&{then}\6
\\{app1}\?\1(\?\\{pp}\?\2)\?;\5
\\{app}\?\1(\?\H{\\}\?\2)\?;\5
\\{app}\?\1(\?\H{V}\?\2)\?;\5
\\{app1}\?\1(\?\\{pp}\?+\?2\?\2)\?;\6
\\{reduce}\?\1(\?\\{pp}\?,\39\?3\?,\39\?\\{math}\?,\39-\?3\?\2)\?;\6
\4\&{elsif}\1\8\\{cat}\?\1(\?\\{pp}\?+\?1\?\to\?\\{pp}\?+\?3\?\2)=\1(\?%
\\{follows}\?,\39\?\\{math}\?,\39\?\\{semi}\?\2)\?\2\8\&{then}\6
\\{app}\?\1(\?\\{pdollar}\?\2)\?;\5
\\{app1}\?\1(\?\\{pp}\?\2)\?;\5
\\{app}\?\1(\?\\{pdollar}\?\2)\?;\5
\\{app1}\?\1(\?\\{pp}\?+\?1\?\2)\?;\5
\\{app}\?\1(\?\\{pdollar}\?\2)\?;\5
\\{app1}\?\1(\?\\{pp}\?+\?2\?\2)\?;\5
\\{app}\?\1(\?\\{pdollar}\?\2)\?;\5
\\{app1}\?\1(\?\\{pp}\?+\?3\?\2)\?;\6
\\{reduce}\?\1(\?\\{pp}\?,\39\?4\?,\39\?\\{decl}\?,\39-\?3\?\2)\?;\6
\4\&{elsif}\1\8\\{cat}\?\1(\?\\{pp}\?+\?1\?\to\?\\{pp}\?+\?2\?\2)=\1(\?%
\\{follows}\?,\39\?\\{record\_scrap}\?\2)\?\2\8\&{then}\6
\\{app}\?\1(\?\\{pdollar}\?\2)\?;\5
\\{app1}\?\1(\?\\{pp}\?\2)\?;\5
\\{app}\?\1(\?\\{pdollar}\?\2)\?;\5
\\{app1}\?\1(\?\\{pp}\?+\?1\?\2)\?;\5
\\{app}\?\1(\?\\{indent}\?\2)\?;\5
\\{app}\?\1(\?\\{force}\?\2)\?;\5
\\{app1}\?\1(\?\\{pp}\?+\?2\?\2)\?;\5
\\{app}\?\1(\?\\{indent}\?\2)\?;\5
\\{app}\?\1(\?\\{force}\?\2)\?;\6
\\{reduce}\?\1(\?\\{pp}\?,\39\?3\?,\39\?\\{decl}\?,\39-\?3\?\2)\?;\6
\4\&{elsif}\1\8\\{cat}\?\1(\?\\{pp}\?+\?1\?\2)=\?\\{loop\_scrap}\2\8\&{then}\6
\\{app}\?\1(\?\\{pdollar}\?\2)\?;\5
\\{app1}\?\1(\?\\{pp}\?\2)\?;\5
\\{app}\?\1(\?\\{pdollar}\?\2)\?;\5
\\{app}\?\1(\?\\{sspace}\?\2)\?;\5
\\{app1}\?\1(\?\\{pp}\?+\?1\?\2)\?;\6
\\{reduce}\?\1(\?\\{pp}\?,\39\?2\?,\39\?\\{beginning}\?,\39-\?2\?\2)\?;\6
\4\&{elsif}\1\8\\{cat}\?\1(\?\\{pp}\?+\?1\?\2)=\?\\{math}\2\8\&{then}\6
\\{squash}\?\1(\?\\{pp}\?,\39\?2\?,\39\?\\{math}\?,\39-\?3\?\2)\?;\6
\4\&{elsif}\1\8\\{cat}\?\1(\?\\{pp}\?+\?1\?\2)=\?\\{semi}\2\8\&{then}\6
\\{app}\?\1(\?\\{pdollar}\?\2)\?;\5
\\{app1}\?\1(\?\\{pp}\?\2)\?;\5
\\{app}\?\1(\?\\{pdollar}\?\2)\?;\5
\\{app1}\?\1(\?\\{pp}\?+\?1\?\2)\?;\6
\\{reduce}\?\1(\?\\{pp}\?,\39\?2\?,\39\?\\{stmt}\?,\39-\?3\?\2)\?;\6
\4\&{elsif}\1\8\\{cat}\?\1(\?\\{pp}\?+\?1\?\2)=\?\\{stmt}\2\8\&{then}\6
\\{app}\?\1(\?\\{pdollar}\?\2)\?;\5
\\{app1}\?\1(\?\\{pp}\?\2)\?;\5
\\{app}\?\1(\?\\{pdollar}\?\2)\?;\5
\\{app1}\?\1(\?\\{pp}\?+\?1\?\2)\?;\6
\\{reduce}\?\1(\?\\{pp}\?,\39\?2\?,\39\?\\{stmt}\?,\39-\?3\?\2)\?;\6
\4\&{elsif}\1\8\\{cat}\?\1(\?\\{pp}\?+\?1\?\2)=\?\\{separate\_scrap}\2\8%
\&{then}\6
\\{app}\?\1(\?\\{pdollar}\?\2)\?;\5
\\{app1}\?\1(\?\\{pp}\?\2)\?;\5
\\{app}\?\1(\?\\{pdollar}\?\2)\?;\5
\\{app1}\?\1(\?\\{pp}\?+\?1\?\2)\?;\6
\\{reduce}\?\1(\?\\{pp}\?,\39\?2\?,\39\?\\{separate\_scrap}\?,\39-\?1\?\2)\?;\6
\4\&{elsif}\1\8\\{cat}\?\1(\?\\{pp}\?+\?1\?\to\?\\{pp}\?+\?2\?\2)=\1(\?\\{when%
\_scrap}\?,\39\?\\{math}\?\2)\?\2\8\&{then}\6
\\{squash}\?\1(\?\\{pp}\?,\39\?3\?,\39\?\\{math}\?,\39-\?3\?\2)\?;\6
\4\&{elsif}\1\8\\{cat}\?\1(\?\\{pp}\?+\?1\?\2)=\?\\{with\_scrap}\2\8\&{then}\6
\\{app}\?\1(\?\\{pdollar}\?\2)\?;\5
\\{app1}\?\1(\?\\{pp}\?\2)\?;\5
\\{app}\?\1(\?\\{pspace}\?\2)\?;\5
\\{app}\?\1(\?\\{pdollar}\?\2)\?;\5
\\{app1}\?\1(\?\\{pp}\?+\?1\?\2)\?;\6
\\{reduce}\?\1(\?\\{pp}\?,\39\?2\?,\39\?\\{with\_scrap}\?,\39-\?1\?\2)\?;\2\6
\&{end}\8\2\&{if}\1\par
\U section~174.\fi

\M199. \P\?\X199:Cases for \\{mod\_scrap}\X\S\?\6
\&{if}\1\1\8\\{cat}\?\1(\?\\{pp}\?+\?1\?\2)=\?\\{semi}\2\8\&{then}\6
\\{app2}\?\1(\?\\{pp}\?\2)\?;\5
\\{app}\?\1(\?\\{force}\?\2)\?;\6
\\{reduce}\?\1(\?\\{pp}\?,\39\?2\?,\39\?\\{stmt}\?,\39-\?3\?\2)\?;\6
\4\&{else}\6
\\{squash}\?\1(\?\\{pp}\?,\39\?1\?,\39\?\\{math}\?,\39-\?3\?\2)\?;\2\6
\&{end}\8\2\&{if}\1\par
\U section~174.\fi

\M200. \P\?\X200:Cases for \\{new\_scrap}\X\S\?\6
\\{squash}\?\1(\?\\{pp}\?,\39\?1\?,\39\?\\{math}\?,\39-\?3\?\2)\?\par
\U section~174.\fi

\M201. \P\?\X201:Cases for \\{not\_scrap}\X\S\?\6
\&{if}\1\1\8\\{cat}\?\1(\?\\{pp}\?+\?1\?\2)=\?\\{in\_scrap}\2\8\&{then}\6
\\{app}\?\1(\?\H{\\}\?\2)\?;\5
\\{app}\?\1(\?\H{n}\?\2)\?;\5
\\{app}\?\1(\?\H{o}\?\2)\?;\5
\\{app}\?\1(\?\H{t}\?\2)\?;\5
\\{app}\?\1(\?\H{i}\?\2)\?;\5
\\{app}\?\1(\?\H{n}\?\2)\?;\6
\\{reduce}\?\1(\?\\{pp}\?,\39\?2\?,\39\?\\{math}\?,\39-\?3\?\2)\?;\6
\4\&{else}\6
\\{squash}\?\1(\?\\{pp}\?,\39\?1\?,\39\?\\{math}\?,\39-\?3\?\2)\?;\2\6
\&{end}\8\2\&{if}\1\par
\U section~174.\fi

\M202. \P\?\X202:Cases for \\{open\_scrap}\X\S\?\6
\&{if}\1\1\8\\{cat}\?\1(\?\\{pp}\?+\?1\?\2)=\?\\{decl}\2\8\&{then}\6
\\{app1}\?\1(\?\\{pp}\?\2)\?;\5
\\{app}\?\1(\?\\{pdollar}\?\2)\?;\5
\\{app1}\?\1(\?\\{pp}\?+\?1\?\2)\?;\5
\\{app}\?\1(\?\\{pdollar}\?\2)\?;\5
\\{app}\?\1(\?\H{\\}\?\2)\?;\5
\\{app}\?\1(\?\H{,}\?\2)\?;\5
\\{app}\?\1(\?\\{opt}\?\2)\?;\5
\\{app}\?\1(\?\H{5}\?\2)\?;\6
\\{reduce}\?\1(\?\\{pp}\?,\39\?2\?,\39\?\\{open\_scrap}\?,\39\?0\?\2)\?;\6
\4\&{elsif}\1\8\\{cat}\?\1(\?\\{pp}\?+\?1\?\to\?\\{pp}\?+\?2\?\2)=\1(\?\\{math}%
\?,\39\?\\{close\_scrap}\?\2)\?\2\8\&{then}\6
\\{app}\?\1(\?\\{indent}\?\2)\?;\5
\\{app2}\?\1(\?\\{pp}\?\2)\?;\5
\\{app}\?\1(\?\\{outdent}\?\2)\?;\5
\\{app1}\?\1(\?\\{pp}\?+\?2\?\2)\?;\6
\\{reduce}\?\1(\?\\{pp}\?,\39\?3\?,\39\?\\{math}\?,\39-\?3\?\2)\?;\6
\4\&{elsif}\1\8\\{cat}\?\1(\?\\{pp}\?+\?1\?\to\?\\{pp}\?+\?4\?\2)=\1(\?\\{math}%
\?,\39\?\\{colon}\?,\39\?\\{math}\?,\39\?\\{close\_scrap}\?\2)\?\2\8\&{then}\6
\\{app1}\?\1(\?\\{pp}\?\2)\?;\5
\\{app}\?\1(\?\\{indent}\?\2)\?;\5
\\{app3}\?\1(\?\\{pp}\?+\?1\?\2)\?;\5
\\{app}\?\1(\?\\{outdent}\?\2)\?;\5
\\{app1}\?\1(\?\\{pp}\?+\?4\?\2)\?;\6
\\{reduce}\?\1(\?\\{pp}\?,\39\?5\?,\39\?\\{math}\?,\39-\?3\?\2)\?;\2\6
\&{end}\8\2\&{if}\1\par
\U section~174.\fi

\M203. \P\?\X203:Cases for \\{period}\X\S\?\6
\\{squash}\?\1(\?\\{pp}\?,\39\?1\?,\39\?\\{math}\?,\39-\?3\?\2)\?\par
\U section~174.\fi

\M204. \P\?\X204:Cases for \\{private\_scrap}\X\S\?\6
\&{if}\1\1\8\\{cat}\?\1(\?\\{pp}\?+\?1\?\2)=\?\\{semi}\2\8\&{then}\6
\\{squash}\?\1(\?\\{pp}\?,\39\?1\?,\39\?\\{math}\?,\39-\?2\?\2)\?;\6
\4\&{else}\6
\\{app}\?\1(\?\\{force}\?\2)\?;\5
\\{app}\?\1(\?\\{backup}\?\2)\?;\5
\\{app1}\?\1(\?\\{pp}\?\2)\?;\5
\\{app}\?\1(\?\\{force}\?\2)\?;\6
\\{reduce}\?\1(\?\\{pp}\?,\39\?1\?,\39\?\\{decl}\?,\39-\?3\?\2)\?;\2\6
\&{end}\8\2\&{if}\1\par
\U section~174.\fi

\M205. \P\?\X205:Cases for \\{proc}\X\S\?\6
\&{if}\1\1\8\\{cat}\?\1(\?\\{pp}\?+\?1\?\to\?\\{pp}\?+\?3\?\2)=\1(\?\\{math}\?,%
\39\?\\{follows}\?,\39\?\\{new\_scrap}\?\2)\?\2\8\&{then}\6
\\{squash}\?\1(\?\\{pp}\?+\?2\?,\39\?2\?,\39\?\\{math}\?,\39-\?3\?\2)\?;\6
\4\&{elsif}\1\8\\{cat}\?\1(\?\\{pp}\?+\?1\?\to\?\\{pp}\?+\?3\?\2)=\1(\?\\{math}%
\?,\39\?\\{follows}\?,\39\?\\{separate\_scrap}\?\2)\?\2\8\&{then}\6
\\{squash}\?\1(\?\\{pp}\?+\?2\?,\39\?2\?,\39\?\\{math}\?,\39-\?3\?\2)\?;\6
\4\&{elsif}\1\8\\{cat}\?\1(\?\\{pp}\?+\?1\?\to\?\\{pp}\?+\?2\?\2)=\1(\?\\{math}%
\?,\39\?\\{follows}\?\2)\?\2\8\&{then}\6
\\{app1}\?\1(\?\\{pp}\?\2)\?;\5
\\{app}\?\1(\?\\{indent}\?\2)\?;\5
\\{app}\?\1(\?\\{pdollar}\?\2)\?;\5
\\{app1}\?\1(\?\\{pp}\?+\?1\?\2)\?;\5
\\{app}\?\1(\?\\{pdollar}\?\2)\?;\5
\\{app}\?\1(\?\\{outdent}\?\2)\?;\5
\\{app1}\?\1(\?\\{pp}\?+\?2\?\2)\?;\5
\\{app}\?\1(\?\\{cancel}\?\2)\?;\5
\\{app}\?\1(\?\\{indent}\?\2)\?;\6
\\{reduce}\?\1(\?\\{pp}\?,\39\?3\?,\39\?\\{decl}\?,\39-\?3\?\2)\?;\6
\4\&{elsif}\1\8\\{cat}\?\1(\?\\{pp}\?+\?1\?\to\?\\{pp}\?+\?2\?\2)=\1(\?\\{math}%
\?,\39\?\\{semi}\?\2)\?\2\8\&{then}\6
\\{app1}\?\1(\?\\{pp}\?\2)\?;\5
\\{app}\?\1(\?\\{indent}\?\2)\?;\5
\\{app}\?\1(\?\\{pdollar}\?\2)\?;\5
\\{app1}\?\1(\?\\{pp}\?+\?1\?\2)\?;\5
\\{app}\?\1(\?\\{pdollar}\?\2)\?;\5
\\{app}\?\1(\?\\{outdent}\?\2)\?;\5
\\{app1}\?\1(\?\\{pp}\?+\?2\?\2)\?;\6
\\{reduce}\?\1(\?\\{pp}\?,\39\?3\?,\39\?\\{decl}\?,\39-\?3\?\2)\?;\2\6
\&{end}\8\2\&{if}\1\par
\U section~174.\fi

\M206. \P\?\X206:Cases for \\{return\_scrap}\X\S\?\6
\&{if}\1\1\8\\{cat}\?\1(\?\\{pp}\?+\?1\?\2)=\?\\{semi}\2\8\&{then}\6
\\{squash}\?\1(\?\\{pp}\?,\39\?2\?,\39\?\\{stmt}\?,\39-\?3\?\2)\?;\6
\4\&{else}\6
\\{app}\?\1(\?\\{pspace}\?\2)\?;\5
\\{app1}\?\1(\?\\{pp}\?\2)\?;\5
\\{app}\?\1(\?\\{sspace}\?\2)\?;\6
\\{reduce}\?\1(\?\\{pp}\?,\39\?1\?,\39\?\\{math}\?,\39-\?3\?\2)\?;\2\6
\&{end}\8\2\&{if}\1\par
\U section~174.\fi

\M207. \P\?\X207:Cases for \\{separate\_scrap}\X\S\?\6
\&{if}\1\1\8\\{cat}\?\1(\?\\{pp}\?+\?1\?\2)=\?\\{math}\2\8\&{then}\6
\\{app1}\?\1(\?\\{pp}\?\2)\?;\5
\\{app}\?\1(\?\\{pdollar}\?\2)\?;\5
\\{app}\?\1(\?\\{sspace}\?\2)\?;\5
\\{app1}\?\1(\?\\{pp}\?+\?1\?\2)\?;\5
\\{app}\?\1(\?\\{pdollar}\?\2)\?;\6
\\{reduce}\?\1(\?\\{pp}\?,\39\?2\?,\39\?\\{decl}\?,\39-\?3\?\2)\?;\2\6
\&{end}\8\2\&{if}\1\par
\U section~174.\fi

\M208. \P\?\X208:Cases for \\{stmt}\X\S\?\6
\&{if}\1\1\8\\{cat}\?\1(\?\\{pp}\?+\?1\?\to\?\\{pp}\?+\?2\?\2)=\1(\?\\{else%
\_scrap}\?,\39\?\\{stmt}\?\2)\?\2\8\&{then}\6
\\{app1}\?\1(\?\\{pp}\?\2)\?;\5
\\{app}\?\1(\?\\{force}\?\2)\?;\5
\\{app}\?\1(\?\\{backup}\?\2)\?;\5
\\{app1}\?\1(\?\\{pp}\?+\?1\?\2)\?;\5
\\{app}\?\1(\?\\{force}\?\2)\?;\5
\\{app1}\?\1(\?\\{pp}\?+\?2\?\2)\?;\6
\\{reduce}\?\1(\?\\{pp}\?,\39\?3\?,\39\?\\{stmt}\?,\39-\?3\?\2)\?;\6
\4\&{elsif}\1\8\\{cat}\?\1(\?\\{pp}\?+\?1\?\2)=\?\\{stmt}\2\8\&{then}\6
\\{app1}\?\1(\?\\{pp}\?\2)\?;\5
\\{app}\?\1(\?\\{break\_space}\?\2)\?;\5
\\{app1}\?\1(\?\\{pp}\?+\?1\?\2)\?;\6
\\{reduce}\?\1(\?\\{pp}\?,\39\?2\?,\39\?\\{stmt}\?,\39-\?3\?\2)\?;\6
\4\&{elsif}\1\8\\{cat}\?\1(\?\\{pp}\?+\?1\?\2)=\?\\{terminator}\2\8\&{then}\6
\\{app1}\?\1(\?\\{pp}\?\2)\?;\5
\\{app}\?\1(\?\\{cancel}\?\2)\?;\5
\\{app}\?\1(\?\\{outdent}\?\2)\?;\5
\\{app}\?\1(\?\\{force}\?\2)\?;\5
\\{app1}\?\1(\?\\{pp}\?+\?1\?\2)\?;\5
\\{app}\?\1(\?\\{force}\?\2)\?;\6
\\{reduce}\?\1(\?\\{pp}\?,\39\?2\?,\39\?\\{stmt}\?,\39-\?3\?\2)\?;\2\6
\&{end}\8\2\&{if}\1\par
\U section~174.\fi

\M209. \P\?\X209:Cases for \\{when\_scrap}\X\S\?\6
\&{if}\1\1\8\\{cat}\?\1(\?\\{pp}\?+\?1\?\to\?\\{pp}\?+\?2\?\2)=\1(\?\\{math}\?,%
\39\?\\{arrow}\?\2)\?\2\8\&{then}\6
\\{app1}\?\1(\?\\{pp}\?\2)\?;\5
\\{app}\?\1(\?\\{pdollar}\?\2)\?;\5
\\{app1}\?\1(\?\\{pp}\?+\?1\?\2)\?;\5
\\{app}\?\1(\?\\{pdollar}\?\2)\?;\5
\\{app1}\?\1(\?\\{pp}\?+\?2\?\2)\?;\6
\\{reduce}\?\1(\?\\{pp}\?,\39\?3\?,\39\?\\{clause}\?,\39-\?1\?\2)\?;\2\6
\&{end}\8\2\&{if}\1\par
\U section~174.\fi

\M210. \P\?\X210:Cases for \\{with\_scrap}\X\S\?\6
\&{if}\1\1\8\\{cat}\?\1(\?\\{pp}\?+\?1\?\2)=\?\\{colon}\2\8\&{then}\6
\\{squash}\?\1(\?\\{pp}\?,\39\?2\?,\39\?\\{with\_scrap}\?,\39-\?1\?\2)\?;\6
\4\&{elsif}\1\8\\{cat}\?\1(\?\\{pp}\?+\?1\?\to\?\\{pp}\?+\?2\?\2)=\1(\?\\{math}%
\?,\39\?\\{semi}\?\2)\?\2\8\&{then}\6
\\{app1}\?\1(\?\\{pp}\?\2)\?;\5
\\{app}\?\1(\?\\{pdollar}\?\2)\?;\5
\\{app1}\?\1(\?\\{pp}\?+\?1\?\2)\?;\5
\\{app}\?\1(\?\\{pdollar}\?\2)\?;\5
\\{app1}\?\1(\?\\{pp}\?+\?2\?\2)\?;\6
\\{reduce}\?\1(\?\\{pp}\?,\39\?3\?,\39\?\\{decl}\?,\39-\?3\?\2)\?;\6
\4\&{elsif}\1\8\\{cat}\?\1(\?\\{pp}\?+\?1\?\2)=\?\\{proc}\2\8\&{then}\6
\\{squash}\?\1(\?\\{pp}\?,\39\?2\?,\39\?\\{gproc}\?,\39\?0\?\2)\?;\6
\4\&{elsif}\1\8\\{cat}\?\1(\?\\{pp}\?+\?1\?\to\?\\{pp}\?+\?2\?\2)=\1(\?%
\\{record\_scrap}\?,\39\?\\{colon}\?\2)\?\2\8\&{then}\6
\\{app}\?\1(\?\\{pdollar}\?\2)\?;\5
\\{app1}\?\1(\?\\{pp}\?\2)\?;\5
\\{app}\?\1(\?\\{indent}\?\2)\?;\5
\\{app}\?\1(\?\\{force}\?\2)\?;\5
\\{app}\?\1(\?\\{indent}\?\2)\?;\5
\\{app1}\?\1(\?\\{pp}\?+\?1\?\2)\?;\5
\\{app}\?\1(\?\\{pdollar}\?\2)\?;\6
\\{reduce}\?\1(\?\\{pp}\?,\39\?2\?,\39\?\\{math}\?,\39-\?1\?\2)\?;\6
\4\&{elsif}\1\8\\{cat}\?\1(\?\\{pp}\?+\?1\?\2)=\?\\{record\_scrap}\2\8\&{then}\6
\\{app1}\?\1(\?\\{pp}\?\2)\?;\5
\\{app}\?\1(\?\\{indent}\?\2)\?;\5
\\{app}\?\1(\?\\{force}\?\2)\?;\5
\\{app}\?\1(\?\\{indent}\?\2)\?;\5
\\{app1}\?\1(\?\\{pp}\?+\?1\?\2)\?;\6
\\{reduce}\?\1(\?\\{pp}\?,\39\?2\?,\39\?\\{decl}\?,\39-\?3\?\2)\?;\2\6
\&{end}\8\2\&{if}\1\par
\U section~174.\fi

\M211. The `\\{freeze\_text}' macro is used to give official status to a token
list.
Before saying \\{freeze\_text}, items are appended to the current token list,
and we know that the eventual number of this token list will be the current
value of \\{text\_ptr}. But no list of that number really exists as yet,
because no ending point for the current list has been
stored in the \\{tok\_start} array. After saying \\{freeze\_text}, the
old current token list becomes legitimate, and its number is the current
value of \\{text\_ptr}\?-\?1 since \\{text\_ptr} has been increased. The new
current token list is empty and ready to be appended to.
Note that \\{freeze\_text} does not check to see that \\{text\_ptr} hasn't
gotten
too large, since it is assumed that this test was done beforehand.

\Y\P\4\D\\{freeze\_text}\?\S\?\\{incr}\?\1(\?\\{text\_ptr}\?\2)\?;\5
\\{tok\_start}\?\1(\?\\{text\_ptr}\?\2)\K\?\\{tok\_ptr}\par
\fi

\M212. The `\\{reduce}' macro used in our code for productions actually calls
on
a procedure named `\\{red}', which makes the appropriate changes to the
scrap list.

\Y\P\?\4\X212:Procedures in \\{parsing}\X\S\?\6
\&{procedure}\8\1\\{red}\?(\?\|j\?:\?\\{sixteen\_bits};\?\,\35\?\|k\?:\?%
\\{eight\_bits};\?\,\35\?\|c\?:\?\\{eight\_bits};\?\,\35\1\?\|d\?:\?\p{INTEGER}%
\?\2)\?\2\8\&{is}\8\1\6
\4\&{begin}\6
\\{cat}\?\1(\?\|j\?\2)\K\?\|c;\5
\\{trans}\?\1(\?\|j\?\2)\K\?\\{text\_ptr};\5
\\{freeze\_text};\6
\&{if}\1\1\8\|k\?>\?1\2\8\&{then}\6
\1\&{for}\8\|i\?\in\?\|j\?+\?\|k\?\to\?\\{lo\_ptr}\8\&{loop}\6
\\{cat}\?\1(\?\|i\?-\?\|k\?+\?1\?\2)\K\?\\{cat}\?\1(\?\|i\?\2)\?;\5
\\{trans}\?\1(\?\|i\?-\?\|k\?+\?1\?\2)\K\?\\{trans}\?\1(\?\|i\?\2)\?;\2\6
\&{end}\8\&{loop};\6
\\{lo\_ptr}\?\K\?\\{lo\_ptr}\?-\?\|k\?+\?1;\2\6
\&{end}\8\2\&{if}\1;\6
\X213:Change \\{pp} to $\max(\\{scrap\_base},\\{pp}\?+\?\|d)$\X;\2\6
\&{end}\8\\{red};\par
\A sections~214, 218, 222, 238, and~240.
\U section~9\*.\fi

\M213. \P\?\X213:Change \\{pp} to $\max(\\{scrap\_base},\\{pp}\?+\?\|d)$\X\S\?\6
\&{if}\1\1\8\\{pp}\?+\?\|d\?\G\?\\{scrap\_base}\2\8\&{then}\6
\\{pp}\?\K\?\\{pp}\?+\?\|d;\6
\4\&{else}\6
\\{pp}\?\K\?\\{scrap\_base};\2\6
\&{end}\8\2\&{if}\1\par
\U sections~212 and~214.\fi

\M214. Similarly, the `\\{squash}' macro invokes a procedure called `\\{sq}'.
This
procedure takes advantage of the simplification that occurs when \|k\?=\?1.

\Y\P\?\4\X212:Procedures in \\{parsing}\X\mathrel{+}\S\?\6
\&{procedure}\8\1\\{sq}\?(\?\|j\?:\?\\{sixteen\_bits};\?\,\35\?\|k\?:\?\\{eight%
\_bits};\?\,\35\?\|c\?:\?\\{eight\_bits};\?\,\35\1\?\|d\?:\?\p{INTEGER}\?\2)\?%
\2\8\&{is}\8\1\6
\4\&{begin}\6
\&{if}\1\1\8\|k\?=\?1\2\8\&{then}\6
\\{cat}\?\1(\?\|j\?\2)\K\?\|c;\6
\X213:Change \\{pp} to $\max(\\{scrap\_base},\\{pp}\?+\?\|d)$\X;\6
\4\&{else}\6
\1\&{for}\8\|i\?\in\?\|j\?\to\?\|j\?+\?\|k\?-\?1\8\&{loop}\6
\\{app1}\?\1(\?\|i\?\2)\?;\2\6
\&{end}\8\&{loop};\6
\\{red}\?\1(\?\|j\?,\39\?\|k\?,\39\?\|c\?,\39\?\|d\?\2)\?;\2\6
\&{end}\8\2\&{if}\1;\2\6
\&{end}\8\\{sq};\par
\fi

\M215. Here now is the code that applies productions as long as possible.

\Y\P\?\4\X215:Reduce the scraps using the productions until no more rules apply%
\X\S\?\6
\1\&{loop}\6
\X216:Make sure the entries \\{cat}\?(\?\\{pp}\?\to\?\\{pp}\?+\?4\?)\? are
defined\X;\6
\&{if}\1 \?\1(\?\\{tok\_ptr}\?+\?11\?>\?\\{max\_toks}\?\2)\V\1(\?\\{text\_ptr}%
\?+\?4\?>\?\\{max\_texts}\?\2)\? \8\&{then} \6
\&{stat}\1\6
\&{if}\1\1\8\\{tok\_ptr}\?>\?\\{max\_tok\_ptr}\2\8\&{then}\6
\\{max\_tok\_ptr}\?\K\?\\{tok\_ptr};\2\6
\&{end}\8\2\&{if}\1;\6
\&{if}\1\1\8\\{text\_ptr}\?>\?\\{max\_txt\_ptr}\2\8\&{then}\6
\\{max\_txt\_ptr}\?\K\?\\{text\_ptr};\2\6
\&{end}\8\2\&{if}\1;\2\6
\&{tats}\6
\\{overflow}\?\1(\?\.{"token/text"}\?\2)\?;\2\6
\&{end}\8\2\&{if}\1;\6
\?\&{exit}\8\&{when}\?\8\\{pp}\?>\?\\{lo\_ptr};\6
\X174:Match a production at \\{pp}, or increase \\{pp} if there is no match\X;%
\2\6
\&{end}\8\&{loop}\par
\U section~219\*.\fi

\M216. If we get to the end of the scrap list, category codes equal to zero are
stored, since zero does not match anything in a production.

\Y\P\?\4\X216:Make sure the entries \\{cat}\?(\?\\{pp}\?\to\?\\{pp}\?+\?4\?)\?
are defined\X\S\?\6
\&{if}\1\1\8\\{lo\_ptr}\?<\?\\{pp}\?+\?4\2\8\&{then}\6
\1\&{loop}\6
\&{if}\1\1\8\\{hi\_ptr}\?\L\?\\{scrap\_ptr}\2\8\&{then}\6
\\{incr}\?\1(\?\\{lo\_ptr}\?\2)\?;\6
\\{cat}\?\1(\?\\{lo\_ptr}\?\2)\K\?\\{cat}\?\1(\?\\{hi\_ptr}\?\2)\?;\5
\\{trans}\?\1(\?\\{lo\_ptr}\?\2)\K\?\\{trans}\?\1(\?\\{hi\_ptr}\?\2)\?;\6
\\{incr}\?\1(\?\\{hi\_ptr}\?\2)\?;\2\6
\&{end}\8\2\&{if}\1;\6
\?\&{exit}\8\&{when}\8\1(\?\\{hi\_ptr}\?>\?\\{scrap\_ptr}\?\2)\V\1(\?\\{lo%
\_ptr}\?=\?\\{pp}\?+\?4\?\2)\?;\2\6
\&{end}\8\&{loop};\6
\1\&{for}\8\|i\?\in\?\\{lo\_ptr}\?+\?1\?\to\?\\{pp}\?+\?4\8\&{loop}\6
\\{cat}\?\1(\?\|i\?\2)\K\?0;\2\6
\&{end}\8\&{loop};\2\6
\&{end}\8\2\&{if}\1\par
\U section~215.\fi

\M217. The \\{translate} function assumes that scraps have been stored in
positions \\{scrap\_base} through \\{scrap\_ptr} of \\{cat} and \\{trans}. It
appends a \\{comment\_scrap} and begins to apply productions as much as
possible. The result is a token list containing the translation of
the given sequence of scraps.

After calling \\{translate}, we will have \\{text\_ptr}\?+\?3\?\L\?\\{max%
\_texts} and
\\{tok\_ptr}\?+\?6\?\L\?\\{max\_toks}, so it will be possible to create up to
three token
lists with up to six tokens without checking for overflow. Before calling
\\{translate}, we should have \\{text\_ptr}\?<\?\\{max\_texts} and \\{scrap%
\_ptr}\?<\?\\{max\_scraps},
since \\{translate} might add a new text and a new scrap before it checks
for overflow.

\Y\P\?\4\X217:Specification of procedures visible from \\{parsing}\X\S\?\6
\&{function}\8\1\\{translate}\?\8\&{return}\?\8\\{text\_pointer}\2;\par
\A sections~221, 237, and~239.
\U section~9\*.\fi

\M218. \P\?\X212:Procedures in \\{parsing}\X\mathrel{+}\S\?\6
\X175:Declaration of subprocedure for \\{translate}\X\6
\&{function}\8\1\\{translate}\?\8\&{return}\?\8\\{text\_pointer}\?\8\&{is}\8%
\&{separate}\?\2;\par
\fi

\M219\*.
\Y\P\O{parsing-translate.adb}\?\6
\?\&{separate}\?\8\1(\?\\{parsing}\?\2)\?\6
\&{function}\8\1\\{translate}\?\8\&{return}\?\8\\{text\_pointer}\2\8\&{is}\8%
\C{converts a sequence of scraps}\1\6
\|k\?:\?\\{long\_buffer\_index};\C{index into \\{buffer}}\6
\4\&{begin}\6
\\{pp}\?\K\?\\{scrap\_base};\5
\\{lo\_ptr}\?\K\?\\{pp}\?-\?1;\5
\\{hi\_ptr}\?\K\?\\{pp};\5
\X215:Reduce the scraps using the productions until no more rules apply\X;\6
\&{if}\1\1\?\8\1(\?\\{lo\_ptr}\?=\?\\{scrap\_base}\?\2)\W\1(\?\\{cat}\?\1(\?%
\\{lo\_ptr}\?\2)\I\?\\{math}\?\2)\?\2\8\&{then}\6
\?\&{return}\?\8\\{trans}\?\1(\?\\{lo\_ptr}\?\2)\?;\6
\4\&{else}\6
\X220:Combine the irreducible scraps that remain\X;\2\6
\&{end}\8\2\&{if}\1;\2\6
\&{end}\8\\{translate};\par
\fi

\M220. If the initial sequence of scraps does not reduce to a single scrap,
we concatenate the translations of all remaining scraps, separated by
blank spaces, with dollar signs surrounding the translations of \\{math}
scraps.

\Y\P\?\4\X220:Combine the irreducible scraps that remain\X\S\?\6
\1\&{for}\8\|j\?\in\?\\{scrap\_base}\?\to\?\\{lo\_ptr}\8\&{loop}\6
\&{if}\1\1\8\|j\?\I\?\\{scrap\_base}\2\8\&{then}\6
\\{app}\?\1(\?\H{\ }\?\2)\?;\2\6
\&{end}\8\2\&{if}\1;\6
\&{if}\1\1\8\\{cat}\?\1(\?\|j\?\2)=\?\\{math}\?\V\?\\{cat}\?\1(\?\|j\?\2)=\?%
\\{open\_scrap}\?\V\?\\{cat}\?\1(\?\|j\?\2)=\?\\{close\_scrap}\2\8\&{then}\6
\\{app}\?\1(\?\\{pdollar}\?\2)\?;\2\6
\&{end}\8\2\&{if}\1;\6
\\{app1}\?\1(\?\|j\?\2)\?;\6
\&{if}\1\1\8\\{cat}\?\1(\?\|j\?\2)=\?\\{math}\?\V\?\\{cat}\?\1(\?\|j\?\2)=\?%
\\{open\_scrap}\?\V\?\\{cat}\?\1(\?\|j\?\2)=\?\\{close\_scrap}\2\8\&{then}\6
\\{app}\?\1(\?\\{pdollar}\?\2)\?;\2\6
\&{end}\8\2\&{if}\1;\6
\&{if}\1\1\8\\{tok\_ptr}\?+\?6\?>\?\\{max\_toks}\2\8\&{then}\6
\\{overflow}\?\1(\?\.{"token"}\?\2)\?;\2\6
\&{end}\8\2\&{if}\1;\2\6
\&{end}\8\&{loop};\6
\\{freeze\_text};\5
\?\&{return}\?\8\\{text\_ptr}\?-\?1\par
\U section~219\*.\fi

\N221.  Initializing the scraps.
If we are going to use the powerful production mechanism just developed, we
must get the scraps set up in the first place, given a \ADA\ text. A table
of the initial scraps corresponding to \ADA\ tokens appeared above in the
section on parsing; our goal now is to implement that table. We shall do this
by implementing a subroutine called \\{ADA\_parse} that is analogous to the
\\{ADA\_xref} routine used during phase one.

Like \\{ADA\_xref}, the \\{ADA\_parse} procedure starts with the current
value of \\{next\_control} and it uses the operation \\{next\_control}\?\K\?%
\\{get\_next}
repeatedly to read \ADA\ text until encountering the next `\v' or
`\.\{', or until \\{next\_control}\?\G\?\\{format}. The scraps corresponding to
what
it reads are appended into the \\{cat} and \\{trans} arrays, and \\{scrap\_ptr}
is advanced.

Like \\{prod}, this procedure has to split into pieces so that each
part is short enough to be handled by \ADA\ compilers that discriminate
against long subroutines. This time there is one split-off routine,
called \\{easy\_cases}.

After studying \\{ADA\_parse}, we will look at the sub-procedure
\\{app\_comment}, that is used in some of its branches.

\Y\P\?\4\X217:Specification of procedures visible from \\{parsing}\X\mathrel{+}%
\S\?\6
\&{procedure}\8\1\\{app\_comment}\2;\par
\fi

\M222. \P\?\X212:Procedures in \\{parsing}\X\mathrel{+}\S\?\6
\X236:Declaration of the \\{app\_comment} procedure\X\6
\X226:Declaration of the \\{easy\_cases} procedure\X\6
\&{procedure}\8\1\\{ADA\_parse}\?\8\&{is}\8\&{separate}\?\2;\par
\fi

\M223\*.
\Y\P\O{parsing-ada\_parse.adb}\?\6
\?\&{separate}\?\8\1(\?\\{parsing}\?\2)\?\6
\&{procedure}\8\1\\{ADA\_parse}\2\8\&{is}\8\C{creates scraps from \ADA\ tokens}%
\1\6
\|j\?:\?\\{long\_buffer\_index};\C{index into \\{buffer}}\6
\|p\?:\?\\{name\_pointer};\C{identifier designator}\6
\4\&{begin}\6
\1\&{while}\8\\{next\_control}\?<\?\\{format}\8\&{loop}\6
\X225:Append the scrap appropriate to \\{next\_control}\X;\6
\\{next\_control}\?\K\?\\{get\_next};\6
\&{if}\1\1\?\8\1(\?\\{next\_control}\?=\?\H{|}\?\2)\V\1(\?\\{next\_control}\?=%
\?\H{\{}\?\2)\?\2\8\&{then}\6
\&{return};\2\6
\&{end}\8\2\&{if}\1;\2\6
\&{end}\8\&{loop};\2\6
\&{end}\8\\{ADA\_parse};\par
\fi

\M224. The macros defined here are helpful abbreviations for the operations
needed when generating the scraps. A scrap of category \|c whose
translation has three tokens $t_1$, $t_2$, $t_3$ is generated by
\\{sc3}$(t_1)(t_2)(t_3)(c)$, etc.

\Y\P\4\D\\{s0}\?\1(\?\#\?\2)\S\?\\{incr}\?\1(\?\\{scrap\_ptr}\?\2)\?;\5
\\{cat}\?\1(\?\\{scrap\_ptr}\?\2)\K\?\#;\5
\\{trans}\?\1(\?\\{scrap\_ptr}\?\2)\K\?\\{text\_ptr};\5
\\{freeze\_text}\par
\P\4\D\\{sc1}\?\1(\?\#\?\2)\S\?\\{app}\?\1(\?\#\?\2)\?;\5
\\{s0}\par
\P\4\D\\{sc2}\?\1(\?\#\?\2)\S\?\\{app}\?\1(\?\#\?\2)\?;\5
\\{sc1}\par
\P\4\D\\{sc3}\?\1(\?\#\?\2)\S\?\\{app}\?\1(\?\#\?\2)\?;\5
\\{sc2}\par
\P\4\D\\{sc4}\?\1(\?\#\?\2)\S\?\\{app}\?\1(\?\#\?\2)\?;\5
\\{sc3}\par
\P\4\D\\{sc5}\?\1(\?\#\?\2)\S\?\\{app}\?\1(\?\#\?\2)\?;\5
\\{sc4}\par
\P\4\D\\{sc6}\?\1(\?\#\?\2)\S\?\\{app}\?\1(\?\#\?\2)\?;\5
\\{sc5}\par
\P\4\D\\{sc0}\?\1(\?\#\?\2)\S\?\\{incr}\?\1(\?\\{scrap\_ptr}\?\2)\?;\5
\\{cat}\?\1(\?\\{scrap\_ptr}\?\2)\K\?\#;\5
\\{trans}\?\1(\?\\{scrap\_ptr}\?\2)\K\?0\par
\P\4\D\\{comment\_scrap}\?\1(\?\#\?\2)\S\?\\{app}\?\1(\?\#\?\2)\?;\5
\\{app\_comment}\par
\fi

\M225. \P\?\X225:Append the scrap appropriate to \\{next\_control}\X\S\?\6
\X229:Make sure that there is room for at least four more scraps, six more
tokens, and four more texts\X;\6
\4\BL\\{reswitch}\EL\8\&{case}\1\8\\{next\_control}\1\8\&{is}\8\6
\4\&{when}\8\\{str}\?\mathrel{\.|}\?\\{verbatim}\?\mathrel{\.|}\?\\{output%
\_tok}\?\mathrel{\.|}\?\\{character\_tok}\KR\8\X231:Append \(a \(string scrap%
\X;\6
\4\&{when}\8\\{ascii\_tok}\KR\8\X233:Append \(an ASCII-token scrap\X;\6
\4\&{when}\8\\{identifier}\KR\8\X235:Append \(an identifier scrap\X;\6
\4\&{when}\8\\{TeX\_string}\KR\8\X234:Append \(a \TeX\ string scrap\X;\6
\4\&{when}\8\&{others}\8\KR\8\\{easy\_cases};\2\6
\4\2\&{end}\8\&{case}\par
\U section~223\*.\fi

\M226. The \\{easy\_cases} each result in straightforward scraps.

\Y\P\?\4\X226:Declaration of the \\{easy\_cases} procedure\X\S\?\6
\&{procedure}\8\1\\{easy\_cases}\2\8\&{is}\8\C{a subprocedure of \\{ADA%
\_parse}}\1\6
\4\&{begin}\6
\&{case}\1\8\\{next\_control}\1\8\&{is}\8\6
\4\&{when}\8\\{set\_element\_sign}\KR\8\\{sc3}\?\1(\?\H{\\}\?\2)\1(\?\H{i}\?\2)%
\1(\?\H{n}\?\2)\1(\?\\{in\_scrap}\?\2)\?;\6
\4\&{when}\8\\{double\_dot}\KR\8\\{sc3}\?\1(\?\H{\\}\?\2)\1(\?\H{t}\?\2)\1(\?%
\H{o}\?\2)\1(\?\\{math}\?\2)\?;\6
\4\&{when}\8\H{\#}\?\mathrel{\.|}\?\H{\$}\?\mathrel{\.|}\?\H{\%}\?\mathrel{\.|}%
\?\H{\^}\?\mathrel{\.|}\?\H{\_}\KR\8\\{sc4}\?\1(\?\\{pdollar}\?\2)\1(\?\H{\\}\?%
\2)\1(\?\\{next\_control}\?\2)\1(\?\\{pdollar}\?\2)\1(\?\\{math}\?\2)\?;\6
\4\&{when}\8\H{\&}\KR\8\\{sc6}\?\1(\?\\{math\_bin}\?\2)\1(\?\H{\\}\?\2)\1(\?%
\H{.}\?\2)\1(\?\H{\\}\?\2)\1(\?\H{\&}\?\2)\1(\?\H{\}}\?\2)\1(\?\\{math}\?\2)\?;%
\6
\4\&{when}\8\H{!}\KR\8\\{sc5}\?\1(\?\\{math\_rel}\?\2)\1(\?\H{\\}\?\2)\1(\?%
\H{.}\?\2)\1(\?\H{|}\?\2)\1(\?\H{\}}\?\2)\1(\?\\{math}\?\2)\?;\6
\4\&{when}\8\\{ignore}\?\mathrel{\.|}\?\H{|}\?\mathrel{\.|}\?\\{xref\_roman}\?%
\to\?\\{xref\_reserved}\KR\8\&{null};\6
\4\&{when}\8\H{(}\KR\8\\{sc1}\?\1(\?\\{next\_control}\?\2)\1(\?\\{open\_scrap}%
\?\2)\?;\6
\4\&{when}\8\H{)}\KR\8\\{sc1}\?\1(\?\\{next\_control}\?\2)\1(\?\\{close\_scrap}%
\?\2)\?;\6
\4\&{when}\8\H{*}\KR\8\\{sc4}\?\1(\?\H{\\}\?\2)\1(\?\H{a}\?\2)\1(\?\H{s}\?\2)%
\1(\?\H{t}\?\2)\1(\?\\{math}\?\2)\?;\6
\4\&{when}\8\H{,}\KR\8\\{sc3}\?\1(\?\H{,}\?\2)\1(\?\\{opt}\?\2)\1(\?\H{9}\?\2)%
\1(\?\\{math}\?\2)\?;\6
\4\&{when}\8\H{0}\?\to\?\H{9}\KR\6
\&{if}\1\1\?\8\1(\?\\{buffer}\?\1(\?\\{loc}\?\2)=\?\H{\#}\?\2)\V\1(\?\\{buffer}%
\?\1(\?\\{loc}\?+\?1\?\2)=\?\H{\#}\?\2)\?\2\8\&{then}\6
\X227:Get a based number\X;\6
\4\&{else}\6
\\{sc3}\?\1(\?\\{pdollar}\?\2)\1(\?\\{next\_control}\?\2)\1(\?\\{pdollar}\?\2)%
\1(\?\\{digit\_scrap}\?\2)\?;\2\6
\&{end}\8\2\&{if}\1;\6
\4\&{when}\8\H{.}\KR\8\\{sc1}\?\1(\?\H{.}\?\2)\1(\?\\{period}\?\2)\?;\6
\4\&{when}\8\H{;}\KR\8\\{sc1}\?\1(\?\H{;}\?\2)\1(\?\\{semi}\?\2)\?;\6
\4\&{when}\8\H{:}\KR\8\\{sc1}\?\1(\?\H{:}\?\2)\1(\?\\{colon}\?\2)\?;\6
\hbox{\4}\X230:Cases involving nonstandard ASCII characters\X\6
\4\&{when}\8\\{exponent}\KR\8\\{sc3}\?\1(\?\H{\\}\?\2)\1(\?\H{E}\?\2)\1(\?\H{%
\{}\?\2)\1(\?\\{exp}\?\2)\?;\6
\4\&{when}\8\\{begin\_comment}\KR\8\\{sc2}\?\1(\?\H{\\}\?\2)\1(\?\H{B}\?\2)\1(%
\?\\{math}\?\2)\?;\6
\4\&{when}\8\\{end\_comment}\KR\8\\{sc2}\?\1(\?\H{\\}\?\2)\1(\?\H{T}\?\2)\1(\?%
\\{math}\?\2)\?;\6
\4\&{when}\8\\{force\_line}\KR\8\\{sc2}\?\1(\?\H{\\}\?\2)\1(\?\H{)}\?\2)\1(\?%
\\{math}\?\2)\?;\6
\4\&{when}\8\\{thin\_space}\KR\8\\{sc2}\?\1(\?\H{\\}\?\2)\1(\?\H{,}\?\2)\1(\?%
\\{math}\?\2)\?;\6
\4\&{when}\8\\{math\_break}\KR\8\\{sc2}\?\1(\?\\{opt}\?\2)\1(\?\H{0}\?\2)\1(\?%
\\{math}\?\2)\?;\6
\4\&{when}\8\\{line\_break}\KR\8\\{comment\_scrap}\?\1(\?\\{force}\?\2)\?;\6
\4\&{when}\8\\{big\_line\_break}\KR\8\\{comment\_scrap}\?\1(\?\\{big\_force}\?%
\2)\?;\6
\4\&{when}\8\\{no\_line\_break}\KR\8\\{app}\?\1(\?\\{big\_cancel}\?\2)\?;\5
\\{app}\?\1(\?\H{\\}\?\2)\?;\5
\\{app}\?\1(\?\H{8}\?\2)\?;\5
\\{comment\_scrap}\?\1(\?\\{big\_cancel}\?\2)\?;\6
\4\&{when}\8\\{pseudo\_semi}\KR\8\\{sc0}\?\1(\?\\{semi}\?\2)\?;\6
\4\&{when}\8\\{join}\KR\8\\{sc2}\?\1(\?\H{\\}\?\2)\1(\?\H{J}\?\2)\1(\?\\{math}%
\?\2)\?;\6
\4\&{when}\8\&{others}\8\KR\8\\{sc1}\?\1(\?\\{next\_control}\?\2)\1(\?\\{math}%
\?\2)\?;\2\6
\4\2\&{end}\8\&{case};\2\6
\&{end}\8\\{easy\_cases};\par
\U section~222.\fi

\M227. \P\?\X227:Get a based number\X\S\?\6
\X228:Append \(\\{digit\_scrap} scraps until a `\#' occurs\X;\6
\\{sc1}\?\1(\?\\{next\_control}\?\2)\1(\?\\{base}\?\2)\?;\5
\\{next\_control}\?\K\?\\{get\_next};\6
\X228:Append \(\\{digit\_scrap} scraps until a `\#' occurs\X;\6
\\{sc1}\?\1(\?\\{next\_control}\?\2)\1(\?\\{base}\?\2)\?\par
\U section~226.\fi

\M228. \P\?\X228:Append \(\\{digit\_scrap} scraps until a `\#' occurs\X\S\?\6
\1\&{loop}\6
\&{if}\1\1\8\\{next\_control}\?=\?\H{\_}\2\8\&{then}\6
\\{sc2}\?\1(\?\H{\\}\?\2)\1(\?\H{\_}\?\2)\1(\?\\{digit\_scrap}\?\2)\?;\6
\4\&{else}\6
\\{sc1}\?\1(\?\\{next\_control}\?\2)\1(\?\\{digit\_scrap}\?\2)\?;\2\6
\&{end}\8\2\&{if}\1;\6
\\{next\_control}\?\K\?\\{get\_next};\6
\?\&{exit}\8\&{when}\?\8\\{next\_control}\?=\?\H{\#};\2\6
\&{end}\8\&{loop}\par
\U section~227(2).\fi

\M229. \P\?\X229:Make sure that there is room for at least four more scraps,
six more tokens, and four more texts\X\S\?\6
\&{if}\1 \?\1(\?\\{scrap\_ptr}\?+\?4\?>\?\\{max\_scraps}\?\2)\V\1(\?\\{tok%
\_ptr}\?+\?6\?>\?\\{max\_toks}\?\2)\V\1(\?\\{text\_ptr}\?+\?4\?>\?\\{max%
\_texts}\?\2)\? \8\&{then} \6
\&{stat}\1\6
\&{if}\1\1\8\\{scrap\_ptr}\?>\?\\{max\_scr\_ptr}\2\8\&{then}\6
\\{max\_scr\_ptr}\?\K\?\\{scrap\_ptr};\2\6
\&{end}\8\2\&{if}\1;\6
\&{if}\1\1\8\\{tok\_ptr}\?>\?\\{max\_tok\_ptr}\2\8\&{then}\6
\\{max\_tok\_ptr}\?\K\?\\{tok\_ptr};\2\6
\&{end}\8\2\&{if}\1;\6
\&{if}\1\1\8\\{text\_ptr}\?>\?\\{max\_txt\_ptr}\2\8\&{then}\6
\\{max\_txt\_ptr}\?\K\?\\{text\_ptr};\2\6
\&{end}\8\2\&{if}\1;\2\6
\&{tats}\6
\\{overflow}\?\1(\?\.{"scrap/token/text"}\?\2)\?;\2\6
\&{end}\8\2\&{if}\1\par
\U section~225.\fi

\M230. Some nonstandard ASCII characters may have entered \.{AWEAVE} by means
of
standard ones. They are converted to \TeX\ control sequences so that it is
possible to keep \.{AWEAVE} from stepping beyond standard ASCII.

\Y\P\?\4\X230:Cases involving nonstandard ASCII characters\X\S\?\6
\4\&{when}\8\\{not\_equal}\KR\8\\{sc2}\?\1(\?\H{\\}\?\2)\1(\?\H{I}\?\2)\1(\?%
\\{math}\?\2)\?;\6
\4\&{when}\8\\{less\_or\_equal}\KR\8\\{sc2}\?\1(\?\H{\\}\?\2)\1(\?\H{L}\?\2)\1(%
\?\\{math}\?\2)\?;\6
\4\&{when}\8\\{greater\_or\_equal}\KR\8\\{sc2}\?\1(\?\H{\\}\?\2)\1(\?\H{G}\?\2)%
\1(\?\\{math}\?\2)\?;\6
\4\&{when}\8\\{equivalence\_sign}\KR\8\\{sc2}\?\1(\?\H{\\}\?\2)\1(\?\H{S}\?\2)%
\1(\?\\{math}\?\2)\?;\6
\4\&{when}\8\\{and\_sign}\KR\8\\{sc2}\?\1(\?\H{\\}\?\2)\1(\?\H{W}\?\2)\1(\?%
\\{and\_scrap}\?\2)\?;\6
\4\&{when}\8\\{or\_sign}\KR\8\\{sc2}\?\1(\?\H{\\}\?\2)\1(\?\H{V}\?\2)\1(\?%
\\{else\_scrap}\?\2)\?;\6
\4\&{when}\8\\{not\_sign}\KR\8\\{sc2}\?\1(\?\H{\\}\?\2)\1(\?\H{R}\?\2)\1(\?%
\\{not\_scrap}\?\2)\?;\6
\4\&{when}\8\\{left\_arrow}\KR\8\\{sc2}\?\1(\?\H{\\}\?\2)\1(\?\H{K}\?\2)\1(\?%
\\{math}\?\2)\?;\6
\4\&{when}\8\\{right\_arrow}\KR\8\\{sc3}\?\1(\?\H{\\}\?\2)\1(\?\H{K}\?\2)\1(\?%
\H{R}\?\2)\1(\?\\{arrow}\?\2)\?;\6
\4\&{when}\8\\{double\_star}\KR\8\\{sc4}\?\1(\?\H{\\}\?\2)\1(\?\H{d}\?\2)\1(\?%
\H{s}\?\2)\1(\?\H{t}\?\2)\1(\?\\{math}\?\2)\?;\6
\4\&{when}\8\\{left\_label\_bracket}\KR\8\\{sc3}\?\1(\?\H{\\}\?\2)\1(\?\H{B}\?%
\2)\1(\?\H{L}\?\2)\1(\?\\{begin\_label\_scrap}\?\2)\?;\6
\4\&{when}\8\\{right\_label\_bracket}\KR\8\\{sc3}\?\1(\?\H{\\}\?\2)\1(\?\H{E}\?%
\2)\1(\?\H{L}\?\2)\1(\?\\{end\_label\_scrap}\?\2)\?;\6
\4\&{when}\8\\{box\_sign}\KR\8\\{sc4}\?\1(\?\H{\\}\?\2)\1(\?\H{B}\?\2)\1(\?%
\H{O}\?\2)\1(\?\H{X}\?\2)\1(\?\\{math}\?\2)\?;\par
\U section~226.\fi

\M231. The following code must use \\{app\_tok} instead of \\{app} in order to
protect against overflow. Note that \\{tok\_ptr}\?+\?1\?\L\?\\{max\_toks} after
\\{app\_tok}
has been used, so another \\{app} is legitimate before testing again.

Many of the special characters in a string must be prefixed by `\.\\' so that
\TeX\ will print them properly.

\Y\P\?\4\X231:Append \(a \(string scrap\X\S\?\6
\\{app}\?\1(\?\\{pdollar}\?\2)\?;\5
\\{app}\?\1(\?\H{\\}\?\2)\?;\6
\&{case}\1\8\\{next\_control}\1\8\&{is}\8\6
\4\&{when}\8\\{verbatim}\KR\8\\{app}\?\1(\?\H{=}\?\2)\?;\6
\4\&{when}\8\\{output\_tok}\KR\8\\{app}\?\1(\?\H{O}\?\2)\?;\6
\4\&{when}\8\&{others}\8\KR\8\\{app}\?\1(\?\H{.}\?\2)\?;\2\6
\4\2\&{end}\8\&{case};\6
\\{app}\?\1(\?\H{\{}\?\2)\?;\5
\|j\?\K\?\\{id\_first};\6
\1\&{while}\8\|j\?<\?\\{id\_loc}\8\&{loop}\6
\X232:Append the character appropriate to \\{buffer}\?(\?\|j\?)\?\X;\6
\\{incr}\?\1(\?\|j\?\2)\?;\2\6
\&{end}\8\&{loop};\6
\&{if}\1\1\8\\{next\_control}\?=\?\\{output\_tok}\2\8\&{then}\6
\\{sc3}\?\1(\?\H{\}}\?\2)\1(\?\\{pdollar}\?\2)\1(\?\\{force}\?\2)\1(\?\\{math}%
\?\2)\?;\6
\4\&{else}\6
\\{sc2}\?\1(\?\H{\}}\?\2)\1(\?\\{pdollar}\?\2)\1(\?\\{math}\?\2)\?;\2\6
\&{end}\8\2\&{if}\1\par
\U section~225.\fi

\M232. \P\?\X232:Append the character appropriate to \\{buffer}\?(\?\|j\?)\?\X%
\S\?\6
\&{case}\1\8\\{buffer}\?\1(\?\|j\?\2)\?\1\8\&{is}\8\6
\4\&{when}\8\H{\ }\?\mathrel{\.|}\?\H{\\}\?\mathrel{\.|}\?\H{\%}\?\mathrel{\.|}%
\?\H{\$}\?\mathrel{\.|}\?\H{\^}\?\mathrel{\.|}\?\H{\'}\?\mathrel{\.|}\?\H{\`}\?%
\mathrel{\.|}\?\H{\{}\?\mathrel{\.|}\?\H{\#}\?\mathrel{\.|}\?\H{\}}\?\mathrel{%
\.|}\?\H{\~}\?\mathrel{\.|}\?\H{\&}\?\mathrel{\.|}\?\H{\_}\KR\6
\\{app}\?\1(\?\H{\\}\?\2)\?;\6
\4\&{when}\8\H{@}\KR\6
\&{if}\1\1\8\\{buffer}\?\1(\?\|j\?+\?1\?\2)=\?\H{@}\2\8\&{then}\6
\\{incr}\?\1(\?\|j\?\2)\?;\6
\4\&{else}\6
\\{err\_print}\?\1(\?\.{"!\ Double\ @\ should\ be\ used\ in\ strings\ and\
character\ literals"}\?\2)\?;\2\6
\&{end}\8\2\&{if}\1;\6
\4\&{when}\8\&{others}\8\KR\6
\&{null};\2\6
\4\2\&{end}\8\&{case};\6
\\{app\_tok}\?\1(\?\\{buffer}\?\1(\?\|j\?\2)\2)\?\par
\U sections~231 and~233.\fi

\M233. \P\?\X233:Append \(an ASCII-token scrap\X\S\?\6
\\{app}\?\1(\?\\{pdollar}\?\2)\?;\5
\\{app}\?\1(\?\H{\\}\?\2)\?;\5
\\{app}\?\1(\?\H{H}\?\2)\?;\5
\\{app}\?\1(\?\H{\{}\?\2)\?;\5
\|j\?\K\?\\{id\_first}\?+\?1;\5
\X232:Append the character appropriate to \\{buffer}\?(\?\|j\?)\?\X;\6
\\{sc2}\?\1(\?\H{\}}\?\2)\1(\?\\{pdollar}\?\2)\1(\?\\{math}\?\2)\?\par
\U section~225.\fi

\M234. \P\?\X234:Append \(a \TeX\ string scrap\X\S\?\6
\\{app}\?\1(\?\\{pdollar}\?\2)\?;\5
\\{app}\?\1(\?\H{\\}\?\2)\?;\5
\\{app}\?\1(\?\H{h}\?\2)\?;\5
\\{app}\?\1(\?\H{b}\?\2)\?;\5
\\{app}\?\1(\?\H{o}\?\2)\?;\5
\\{app}\?\1(\?\H{x}\?\2)\?;\5
\\{app}\?\1(\?\H{\{}\?\2)\?;\6
\1\&{for}\8\|l\?\in\?\\{id\_first}\?\to\?\\{id\_loc}\?-\?1\8\&{loop}\6
\\{app\_tok}\?\1(\?\\{buffer}\?\1(\?\|l\?\2)\2)\?;\2\6
\&{end}\8\&{loop};\6
\\{sc2}\?\1(\?\H{\}}\?\2)\1(\?\\{pdollar}\?\2)\1(\?\\{math}\?\2)\?\par
\U section~225.\fi

\M235. \P\?\X235:Append \(an identifier scrap\X\S\?\6
\|p\?\K\?\\{id\_lookup}\?\1(\?\\{normal}\?\2)\?;\6
\&{case}\1\8\\{ilk}\?\1(\?\|p\?\2)\?\1\8\&{is}\8\6
\4\&{when}\8\\{normal}\?\mathrel{\.|}\?\\{predefined}\KR\8\\{sc5}\?\1(\?%
\\{pspace}\?\2)\1(\?\\{pdollar}\?\2)\1(\?\\{id\_flag}\?+\?\|p\?\2)\1(\?%
\\{pdollar}\?\2)\1(\?\\{pspace}\?\2)\1(\?\\{math}\?\2)\?;\C{not a reserved
word}\6
\4\&{when}\8\\{accept\_like}\KR\8\\{sc2}\?\1(\?\\{res\_flag}\?+\?\|p\?\2)\1(\?%
\\{sspace}\?\2)\1(\?\\{accept\_scrap}\?\2)\?;\6
\4\&{when}\8\\{all\_like}\KR\8\\{sc5}\?\1(\?\\{pdollar}\?\2)\1(\?\\{pspace}\?%
\2)\1(\?\\{res\_flag}\?+\?\|p\?\2)\1(\?\\{pdollar}\?\2)\1(\?\\{pspace}\?\2)\1(%
\?\\{math}\?\2)\?;\6
\4\&{when}\8\\{and\_like}\KR\8\\{sc2}\?\1(\?\H{\\}\?\2)\1(\?\H{W}\?\2)\1(\?%
\\{and\_scrap}\?\2)\?;\6
\4\&{when}\8\\{array\_like}\KR\8\\{sc2}\?\1(\?\\{res\_flag}\?+\?\|p\?\2)\1(\?%
\\{sspace}\?\2)\1(\?\\{open\_scrap}\?\2)\?;\6
\4\&{when}\8\\{at\_like}\KR\8\\{sc3}\?\1(\?\\{sspace}\?\2)\1(\?\\{res\_flag}\?+%
\?\|p\?\2)\1(\?\\{sspace}\?\2)\1(\?\\{colon}\?\2)\?;\6
\4\&{when}\8\\{begin\_like}\KR\8\\{sc1}\?\1(\?\\{res\_flag}\?+\?\|p\?\2)\1(\?%
\\{beginning}\?\2)\?;\6
\4\&{when}\8\\{body\_like}\KR\8\\{sc4}\?\1(\?\\{pdollar}\?\2)\1(\?\\{res\_flag}%
\?+\?\|p\?\2)\1(\?\\{pdollar}\?\2)\1(\?\\{pspace}\?\2)\1(\?\\{math}\?\2)\?;\6
\4\&{when}\8\\{case\_like}\KR\8\\{sc1}\?\1(\?\\{res\_flag}\?+\?\|p\?\2)\1(\?%
\\{case\_scrap}\?\2)\?;\6
\4\&{when}\8\\{declare\_like}\KR\8\\{sc4}\?\1(\?\\{force}\?\2)\1(\?\\{res%
\_flag}\?+\?\|p\?\2)\1(\?\\{indent}\?\2)\1(\?\\{force}\?\2)\1(\?\\{decl}\?\2)%
\?;\6
\4\&{when}\8\\{goto\_like}\KR\8\\{sc5}\?\1(\?\\{pspace}\?\2)\1(\?\\{pdollar}\?%
\2)\1(\?\\{res\_flag}\?+\?\|p\?\2)\1(\?\\{pdollar}\?\2)\1(\?\\{sspace}\?\2)\1(%
\?\\{math}\?\2)\?;\6
\4\&{when}\8\\{is\_like}\KR\8\\{sc3}\?\1(\?\\{sspace}\?\2)\1(\?\\{res\_flag}\?+%
\?\|p\?\2)\1(\?\\{sspace}\?\2)\1(\?\\{follows}\?\2)\?;\6
\4\&{when}\8\\{else\_like}\KR\8\\{sc1}\?\1(\?\\{res\_flag}\?+\?\|p\?\2)\1(\?%
\\{else\_scrap}\?\2)\?;\6
\4\&{when}\8\\{elsif\_like}\KR\8\\{sc2}\?\1(\?\\{backup}\?\2)\1(\?\\{res\_flag}%
\?+\?\|p\?\2)\1(\?\\{if\_scrap}\?\2)\?;\6
\4\&{when}\8\\{end\_like}\KR\8\\{sc1}\?\1(\?\\{res\_flag}\?+\?\|p\?\2)\1(\?%
\\{ending}\?\2)\?;\6
\4\&{when}\8\\{exception\_like}\KR\8\\{sc1}\?\1(\?\\{res\_flag}\?+\?\|p\?\2)\1(%
\?\\{exception\_scrap}\?\2)\?;\6
\4\&{when}\8\\{proc\_like}\KR\8\\{sc2}\?\1(\?\\{res\_flag}\?+\?\|p\?\2)\1(\?%
\\{sspace}\?\2)\1(\?\\{proc}\?\2)\?;\6
\4\&{when}\8\\{new\_like}\KR\8\\{sc3}\?\1(\?\\{pspace}\?\2)\1(\?\\{res\_flag}%
\?+\?\|p\?\2)\1(\?\\{pspace}\?\2)\1(\?\\{new\_scrap}\?\2)\?;\6
\4\&{when}\8\\{generic\_like}\KR\8\\{sc1}\?\1(\?\\{res\_flag}\?+\?\|p\?\2)\1(\?%
\\{generic\_scrap}\?\2)\?;\6
\4\&{when}\8\\{if\_like}\KR\8\\{sc2}\?\1(\?\\{res\_flag}\?+\?\|p\?\2)\1(\?%
\\{indent}\?\2)\1(\?\\{if\_scrap}\?\2)\?;\6
\4\&{when}\8\\{in\_like}\KR\8\\{sc3}\?\1(\?\H{\\}\?\2)\1(\?\H{i}\?\2)\1(\?\H{n}%
\?\2)\1(\?\\{in\_scrap}\?\2)\?;\6
\4\&{when}\8\\{loop\_like}\KR\8\\{sc1}\?\1(\?\\{res\_flag}\?+\?\|p\?\2)\1(\?%
\\{loop\_scrap}\?\2)\?;\6
\4\&{when}\8\\{mod\_like}\KR\8\\{sc3}\?\1(\?\\{math\_bin}\?\2)\1(\?\\{res%
\_flag}\?+\?\|p\?\2)\1(\?\H{\}}\?\2)\1(\?\\{math}\?\2)\?;\6
\4\&{when}\8\\{not\_like}\KR\8\\{sc2}\?\1(\?\H{\\}\?\2)\1(\?\H{R}\?\2)\1(\?%
\\{not\_scrap}\?\2)\?;\6
\4\&{when}\8\\{of\_like}\KR\8\\{sc3}\?\1(\?\\{sspace}\?\2)\1(\?\\{res\_flag}\?+%
\?\|p\?\2)\1(\?\\{sspace}\?\2)\1(\?\\{close\_scrap}\?\2)\?;\6
\4\&{when}\8\\{or\_like}\KR\8\\{sc1}\?\1(\?\\{res\_flag}\?+\?\|p\?\2)\1(\?%
\\{else\_scrap}\?\2)\?;\6
\4\&{when}\8\\{pragma\_like}\KR\8\\{sc2}\?\1(\?\\{res\_flag}\?+\?\|p\?\2)\1(\?%
\\{sspace}\?\2)\1(\?\\{proc}\?\2)\?;\6
\4\&{when}\8\\{private\_like}\KR\8\\{sc1}\?\1(\?\\{res\_flag}\?+\?\|p\?\2)\1(\?%
\\{private\_scrap}\?\2)\?;\6
\4\&{when}\8\\{range\_like}\KR\8\\{sc5}\?\1(\?\\{sspace}\?\2)\1(\?\\{pdollar}\?%
\2)\1(\?\\{res\_flag}\?+\?\|p\?\2)\1(\?\\{pdollar}\?\2)\1(\?\\{sspace}\?\2)\1(%
\?\\{math}\?\2)\?;\6
\4\&{when}\8\\{record\_like}\KR\8\\{sc1}\?\1(\?\\{res\_flag}\?+\?\|p\?\2)\1(\?%
\\{record\_scrap}\?\2)\?;\6
\4\&{when}\8\\{return\_like}\KR\8\\{sc1}\?\1(\?\\{res\_flag}\?+\?\|p\?\2)\1(\?%
\\{return\_scrap}\?\2)\?;\6
\4\&{when}\8\\{separate\_like}\KR\8\\{sc1}\?\1(\?\\{res\_flag}\?+\?\|p\?\2)\1(%
\?\\{separate\_scrap}\?\2)\?;\6
\4\&{when}\8\\{type\_like}\KR\8\\{sc4}\?\1(\?\\{pdollar}\?\2)\1(\?\\{res\_flag}%
\?+\?\|p\?\2)\1(\?\\{pdollar}\?\2)\1(\?\\{sspace}\?\2)\1(\?\\{math}\?\2)\?;\6
\4\&{when}\8\\{then\_like}\KR\8\\{sc2}\?\1(\?\\{sspace}\?\2)\1(\?\\{res\_flag}%
\?+\?\|p\?\2)\1(\?\\{then\_scrap}\?\2)\?;\6
\4\&{when}\8\\{with\_like}\KR\8\\{sc3}\?\1(\?\\{pspace}\?\2)\1(\?\\{res\_flag}%
\?+\?\|p\?\2)\1(\?\\{sspace}\?\2)\1(\?\\{with\_scrap}\?\2)\?;\6
\4\&{when}\8\\{when\_like}\KR\8\\{sc2}\?\1(\?\\{res\_flag}\?+\?\|p\?\2)\1(\?%
\\{sspace}\?\2)\1(\?\\{when\_scrap}\?\2)\?;\6
\4\&{when}\8\&{others}\8\KR\8\&{null};\2\6
\4\2\&{end}\8\&{case}\par
\U section~225.\fi

\M236. A comment is incorporated into the previous scrap.
The \\{app\_comment} procedure takes care of placing a comment at the end of
the
current scrap list. When \\{app\_comment} is called, we assume that the current
token list is the translation of the comment involved.

\Y\P\?\4\X236:Declaration of the \\{app\_comment} procedure\X\S\?\6
\&{procedure}\8\1\\{app\_comment}\2\8\&{is}\8\C{append a comment to the scrap
list}\1\6
\4\&{begin}\6
\\{freeze\_text};\5
\\{app1}\?\1(\?\\{scrap\_ptr}\?\2)\?;\5
\\{app}\?\1(\?\\{text\_ptr}\?-\?1\?+\?\\{tok\_flag}\?\2)\?;\5
\\{trans}\?\1(\?\\{scrap\_ptr}\?\2)\K\?\\{text\_ptr};\5
\\{freeze\_text};\2\6
\&{end}\8\\{app\_comment};\par
\U section~222.\fi

\M237. When the `\v' that introduces \ADA\ text is sensed, a call on
\\{ADA\_translate} will return a pointer to the \TeX\ translation of
that text. If scraps exist in the \\{cat} and \\{trans} arrays, they are
unaffected by this translation process.

\Y\P\?\4\X217:Specification of procedures visible from \\{parsing}\X\mathrel{+}%
\S\?\6
\&{function}\8\1\\{ADA\_translate}\?\8\&{return}\?\8\\{text\_pointer}\2;\par
\fi

\M238. \P\?\X212:Procedures in \\{parsing}\X\mathrel{+}\S\?\6
\&{function}\8\1\\{ADA\_translate}\?\8\&{return}\?\8\\{text\_pointer}\2\8\&{is}%
\8\1\6
\|p\?:\?\\{text\_pointer};\C{points to the translation}\6
\\{save\_base}\?:\?\\{scrap\_index};\C{holds original value of \\{scrap\_base}}%
\6
\4\&{begin}\6
\\{save\_base}\?\K\?\\{scrap\_base};\5
\\{scrap\_base}\?\K\?\\{scrap\_ptr}\?+\?1;\5
\\{ADA\_parse};\C{get the scraps together}\6
\&{if}\1\1\8\\{next\_control}\?\I\?\H{|}\2\8\&{then}\6
\\{err\_print}\?\1(\?\.{"!\ Missing\ \`|\'\ after\ ADA\ text"}\?\2)\?;\2\6
\&{end}\8\2\&{if}\1;\6
\\{app\_tok}\?\1(\?\\{cancel}\?\2)\?;\5
\\{app\_comment};\C{place a \\{cancel} token as a final ``comment''}\6
\|p\?\K\?\\{translate};\C{make the translation}\6
\&{stat}\1\6
\&{if}\1\1\8\\{scrap\_ptr}\?>\?\\{max\_scr\_ptr}\2\8\&{then}\6
\\{max\_scr\_ptr}\?\K\?\\{scrap\_ptr};\2\6
\&{end}\8\2\&{if}\1;\2\6
\&{tats}\6
\\{scrap\_ptr}\?\K\?\\{scrap\_base}\?-\?1;\5
\\{scrap\_base}\?\K\?\\{save\_base};\C{scrap the scraps}\6
\?\&{return}\?\8\|p;\2\6
\&{end}\8\\{ADA\_translate};\par
\fi

\M239. The \\{outer\_parse} routine is to \\{ADA\_parse} as \\{outer\_xref}
is to \\{ADA\_xref}: It constructs a sequence of scraps for \ADA\ text
until \\{next\_control}\?\G\?\\{format}. Thus, it takes care of embedded
comments.

\Y\P\?\4\X217:Specification of procedures visible from \\{parsing}\X\mathrel{+}%
\S\?\6
\&{procedure}\8\1\\{outer\_parse}\2;\par
\fi

\M240. \P\?\X212:Procedures in \\{parsing}\X\mathrel{+}\S\?\6
\&{procedure}\8\1\\{outer\_parse}\2\8\&{is}\8\C{makes scraps from \ADA\ tokens
and comments}\1\6
\\{bal}\?:\?\\{eight\_bits};\C{brace level in comment}\6
\|p\?,\39\?\|q\?:\?\\{text\_pointer};\C{partial comments}\6
\4\&{begin}\6
\1\&{while}\8\\{next\_control}\?<\?\\{format}\8\&{loop}\6
\&{if}\1\1\8\\{next\_control}\?\I\?\H{\{}\2\8\&{then}\6
\\{ADA\_parse};\6
\4\&{else}\6
\X241:Make sure that there is room for at least seven more tokens, three more
texts, and one more scrap\X;\6
\\{app}\?\1(\?\H{\\}\?\2)\?;\5
\\{app}\?\1(\?\H{C}\?\2)\?;\5
\\{app}\?\1(\?\H{\{}\?\2)\?;\5
\\{bal}\?\K\?\\{copy\_comment}\?\1(\?1\?\2)\?;\5
\\{next\_control}\?\K\?\H{|};\6
\1\&{while}\8\\{bal}\?>\?0\8\&{loop}\6
\|p\?\K\?\\{text\_ptr};\5
\\{freeze\_text};\5
\|q\?\K\?\\{ADA\_translate};\C{at this point we have \\{tok\_ptr}\?+\?6\?\L\?%
\\{max\_toks}}\6
\\{app}\?\1(\?\\{tok\_flag}\?+\?\|p\?\2)\?;\5
\\{app}\?\1(\?\\{inner\_tok\_flag}\?+\?\|q\?\2)\?;\6
\&{if}\1\1\8\\{next\_control}\?=\?\H{|}\2\8\&{then}\6
\\{bal}\?\K\?\\{copy\_comment}\?\1(\?\\{bal}\?\2)\?;\6
\4\&{else}\6
\\{bal}\?\K\?0;\C{an error has been reported}\2\6
\&{end}\8\2\&{if}\1;\2\6
\&{end}\8\&{loop};\6
\\{app}\?\1(\?\\{force}\?\2)\?;\5
\\{app\_comment};\C{the full comment becomes a scrap}\2\6
\&{end}\8\2\&{if}\1;\2\6
\&{end}\8\&{loop};\2\6
\&{end}\8\\{outer\_parse};\par
\fi

\M241. \P\?\X241:Make sure that there is room for at least seven more tokens,
three more texts, and one more scrap\X\S\?\6
\&{if}\1 \?\1(\?\\{tok\_ptr}\?+\?7\?>\?\\{max\_toks}\?\2)\V\1(\?\\{text\_ptr}%
\?+\?3\?>\?\\{max\_texts}\?\2)\V\1(\?\\{scrap\_ptr}\?\G\?\\{max\_scraps}\?\2)\?
\8\&{then} \6
\&{stat}\1\6
\&{if}\1\1\8\\{scrap\_ptr}\?>\?\\{max\_scr\_ptr}\2\8\&{then}\6
\\{max\_scr\_ptr}\?\K\?\\{scrap\_ptr};\2\6
\&{end}\8\2\&{if}\1;\6
\&{if}\1\1\8\\{tok\_ptr}\?>\?\\{max\_tok\_ptr}\2\8\&{then}\6
\\{max\_tok\_ptr}\?\K\?\\{tok\_ptr};\2\6
\&{end}\8\2\&{if}\1;\6
\&{if}\1\1\8\\{text\_ptr}\?>\?\\{max\_txt\_ptr}\2\8\&{then}\6
\\{max\_txt\_ptr}\?\K\?\\{text\_ptr};\2\6
\&{end}\8\2\&{if}\1;\2\6
\&{tats}\6
\\{overflow}\?\1(\?\.{"token/text/scrap"}\?\2)\?;\2\6
\&{end}\8\2\&{if}\1\par
\U section~240.\fi

\N242.  Output of tokens.
So far our programs have only built up multi-layered token lists in
\.{AWEAVE}'s internal memory; we have to figure out how to get them into
the desired final form. The job of converting token lists to characters in
the \TeX\ output file is not difficult, although it is an implicitly
recursive process. Three main considerations had to be kept in mind when
this part of \.{AWEAVE} was designed:  (a) There are two modes of output,
\\{outer} mode that translates tokens like \\{force} into line-breaking
control sequences, and \\{inner} mode that ignores them except that blank
spaces take the place of line breaks. (b) The \\{cancel} instruction applies
to adjacent token or tokens that are output, and this cuts across levels
of recursion since `\\{cancel}' occurs at the beginning or end of a token
list on one level. (c) The \TeX\ output file will be semi-readable if line
breaks are inserted after the result of tokens like \\{break\_space} and
\\{force}.  (d) The final line break should be suppressed, and there should
be no \\{force} token output immediately after `\.{\\Y\\P}'.

\fi

\M243. The output process uses a stack to keep track of what is going on at
different ``levels'' as the token lists are being written out. Entries on
this stack have three parts:

\yskip\hang \\{end\_field} is the \\{tok\_mem} location where the token list of
a
particular level will end;

\yskip\hang \\{tok\_field} is the \\{tok\_mem} location from which the next
token
on a particular level will be read;

\yskip\hang \\{mode\_field} is the current mode, either \\{inner} or \\{outer}.

\yskip\noindent The current values of these quantities are referred to
quite frequently, so they are stored in a separate place instead of in the
\\{stack} array. We call the current values \\{cur\_end}, \\{cur\_tok}, and
\\{cur\_mode}.

The global variable \\{stack\_ptr} tells how many levels of output are
currently in progress. The end of output occurs when an \\{end\_translation}
token is found, so the stack is never empty except when we first begin the
output process.

\Y\P\4\D\\{init\_stack}\?\S\?\\{stack\_ptr}\?\K\?0;\5
\\{cur\_mode}\?\K\?\\{outer}\C{do this to initialize the stack}\par
\Y\P\?\4\X14:Types and objects visible from \\{data\_structures}\X\mathrel{+}\S%
\?\6
\&{type}\8\\{output\_mode}\8\&{is}\?\8\1(\?\\{inner}\?,\39\?\\{outer}\?\2)\?;\6
\&{type}\8\\{output\_state}\8\&{is}\8\1\6
\&{record}\1\6
\\{end\_field}\?:\?\\{sixteen\_bits};\C{ending location of token list}\6
\\{tok\_field}\?:\?\\{sixteen\_bits};\C{present location within token list}\6
\\{mode\_field}\?:\?\\{output\_mode};\C{interpretation of control tokens}\2\6
\2\&{end}\8\&{record};\7
\\{cur\_state}\?:\?\\{output\_state};\C{\\{cur\_end}, \\{cur\_tok}, \\{cur%
\_mode}}\6
\\{stack}\?:\1\&{array}\8\1(\?\p{INTEGER}\8\&{range}\81\?\to\?\\{stack\_size}\?%
\2)\2\8\&{of}\?\8\\{output\_state};\C{info for non-current levels}\6
\\{stack\_ptr}\?:\?\p{INTEGER}\8\&{range}\80\?\to\?\\{stack\_size};\C{first
unused location in the output state stack}\6
\&{stat}\1\6
\\{max\_stack\_ptr}\?:\?\p{INTEGER}\8\&{range}\80\?\to\?\\{stack\_size}\?\K\?0;%
\C{largest value assumed by \\{stack\_ptr}}\2\6
\&{tats}\7
\\{cur\_end}\?:\?\\{sixteen\_bits}\8\&{renames}\8\\{cur\_state}\?.\?\\{end%
\_field};\C{current ending location in \\{tok\_mem}}\6
\\{cur\_tok}\?:\?\\{sixteen\_bits}\8\&{renames}\8\\{cur\_state}\?.\?\\{tok%
\_field};\C{location of next output token in \\{tok\_mem}}\6
\\{cur\_mode}\?:\?\\{output\_mode}\8\&{renames}\8\\{cur\_state}\?.\?\\{mode%
\_field};\C{current mode of interpretation}\par
\fi

\M244. To insert token-list \|p into the output, the \\{push\_level} subroutine
is called; it saves the old level of output and gets a new one going.
The value of \\{cur\_mode} is not changed.

\Y\P\?\4\X244:Specification of procedures visible from \\{second\_phase}\X\S\?\6
\&{procedure}\8\1\\{push\_level}\?(\1\?\|p\?:\?\\{text\_pointer}\?\2)\?\2;\par
\A sections~249, 278, and~303.
\U section~10\*.\fi

\M245. \P\?\X245:Procedures in \\{second\_phase}\X\S\?\6
\&{procedure}\8\1\\{push\_level}\?(\1\?\|p\?:\?\\{text\_pointer}\?\2)\?\2\8%
\&{is}\8\C{suspends the current level}\1\6
\4\&{begin}\6
\&{if}\1\1\8\\{stack\_ptr}\?=\?\\{stack\_size}\2\8\&{then}\6
\\{overflow}\?\1(\?\.{"stack"}\?\2)\?;\6
\4\&{else}\6
\&{if}\1\1\8\\{stack\_ptr}\?>\?0\2\8\&{then}\6
\\{stack}\?\1(\?\\{stack\_ptr}\?\2)\K\?\\{cur\_state};\C{save \\{cur\_end}$\,%
\ldots\,$\\{cur\_mode}}\2\6
\&{end}\8\2\&{if}\1;\6
\\{incr}\?\1(\?\\{stack\_ptr}\?\2)\?; \6
\&{stat}\1\6
\&{if}\1\1\8\\{stack\_ptr}\?>\?\\{max\_stack\_ptr}\2\8\&{then}\6
\\{max\_stack\_ptr}\?\K\?\\{stack\_ptr};\2\6
\&{end}\8\2\&{if}\1;\2\6
\&{tats}\6
\\{cur\_tok}\?\K\?\\{tok\_start}\?\1(\?\|p\?\2)\?;\5
\\{cur\_end}\?\K\?\\{tok\_start}\?\1(\?\|p\?+\?1\?\2)\?;\2\6
\&{end}\8\2\&{if}\1;\2\6
\&{end}\8\\{push\_level};\par
\A sections~248, 250, 251, 268, 279, and~304.
\U section~10\*.\fi

\M246. Conversely, the \\{pop\_level} routine restores the conditions that were
in
force when the current level was begun. This subroutine will never be
called when \\{stack\_ptr}\?=\?1. It is so simple, we declare it as a macro:

\Y\P\4\D\\{pop\_level}\?\S\?\\{decr}\?\1(\?\\{stack\_ptr}\?\2)\?;\5
\\{cur\_state}\?\K\?\\{stack}\?\1(\?\\{stack\_ptr}\?\2)\?\C{do this when \\{cur%
\_tok} reaches \\{cur\_end}}\par
\fi

\M247. The \\{get\_output} function returns the next byte of output that is not
a
reference to a token list. It returns the values \\{identifier} or \\{res%
\_word}
or \\{mod\_name} if the next token is to be an identifier (typeset in
italics), a reserved word (typeset in boldface) or a module name (typeset
by a complex routine that might generate additional levels of output).
In these cases \\{cur\_name} points to the identifier or module name in
question.

\Y\P\?\4\X14:Types and objects visible from \\{data\_structures}\X\mathrel{+}\S%
\?\6
\\{res\_word}\?:\?\&{constant}\8\\{eight\_bits}\?\K\bn{8}{201}\?;\C{returned by
\\{get\_output} for reserved words}\6
\\{mod\_name}\?:\?\&{constant}\8\\{eight\_bits}\?\K\bn{8}{200}\?;\C{returned by
\\{get\_output} for module names}\par
\fi

\M248. \P\?\X245:Procedures in \\{second\_phase}\X\mathrel{+}\S\?\6
\&{function}\8\1\\{get\_output}\?\8\&{return}\?\8\\{eight\_bits}\2\8\&{is}\8%
\C{returns the next token of output}\1\6
\|a\?:\?\\{sixteen\_bits};\C{current item read from \\{tok\_mem}}\6
\4\&{begin}\6
\4\BL\\{restart}\EL\8\1\&{while}\8\\{cur\_tok}\?=\?\\{cur\_end}\8\&{loop}\6
\\{pop\_level};\2\6
\&{end}\8\&{loop};\6
\|a\?\K\?\\{tok\_mem}\?\1(\?\\{cur\_tok}\?\2)\?;\5
\\{incr}\?\1(\?\\{cur\_tok}\?\2)\?;\6
\&{if}\1\1\8\|a\?\G\bn{8}{400}\?\2\8\&{then}\6
\\{cur\_name}\?\K\?\|a\?\mathbin{\&{mod}}\?\\{id\_flag};\6
\&{case}\1\8\|a\?/\?\\{id\_flag}\1\8\&{is}\8\6
\4\&{when}\82\KR\8\|a\?\K\?\\{res\_word};\C{\|a\?=\?\\{res\_flag}\?+\?\\{cur%
\_name}}\6
\4\&{when}\83\KR\8\|a\?\K\?\\{mod\_name};\C{\|a\?=\?\\{mod\_flag}\?+\?\\{cur%
\_name}}\6
\4\&{when}\84\KR\8\\{push\_level}\?\1(\?\\{cur\_name}\?\2)\?;\5
\&{goto}\8\\{restart};\C{\|a\?=\?\\{tok\_flag}\?+\?\\{cur\_name}}\6
\4\&{when}\85\KR\8\\{push\_level}\?\1(\?\\{cur\_name}\?\2)\?;\5
\\{cur\_mode}\?\K\?\\{inner};\5
\&{goto}\8\\{restart};\C{\|a\?=\?\\{inner\_tok\_flag}\?+\?\\{cur\_name}}\6
\4\&{when}\8\&{others}\8\KR\8\|a\?\K\?\\{identifier};\C{\|a\?=\?\\{id\_flag}\?+%
\?\\{cur\_name}}\2\6
\4\2\&{end}\8\&{case};\2\6
\&{end}\8\2\&{if}\1;\6
\?\&{return}\?\8\|a;\2\6
\&{end}\8\\{get\_output};\par
\fi

\M249. The real work associated with token output is done by \\{make\_output}.
This procedure appends an \\{end\_translation} token to the current token list,
and then it repeatedly calls \\{get\_output} and feeds characters to the output
buffer until reaching the \\{end\_translation} sentinel. It is possible for
\\{make\_output} to
be called recursively, since a module name may include embedded \ADA\
text; however, the depth of recursion never exceeds one level, since
module names cannot be inside of module names.

A procedure called \\{output\_ADA} does the scanning, translation, and
output of \ADA\ text within `\ab' brackets, and this procedure uses
\\{make\_output} to output the current token list. Thus, the recursive call
of \\{make\_output} actually occurs when \\{make\_output} calls \\{output\_ADA}
while outputting the name of a module.

\Y\P\?\4\X244:Specification of procedures visible from \\{second\_phase}\X%
\mathrel{+}\S\?\6
\&{procedure}\8\1\\{make\_output}\2;\par
\fi

\M250. \P\?\X245:Procedures in \\{second\_phase}\X\mathrel{+}\S\?\6
\&{procedure}\8\1\\{output\_ADA}\2\8\&{is}\8\C{outputs the current token list}%
\1\6
\\{save\_tok\_ptr}\?,\39\?\\{save\_text\_ptr}\?,\39\?\\{save\_next\_control}\?:%
\?\\{sixteen\_bits};\C{values to be restored}\6
\|p\?:\?\\{text\_pointer};\C{translation of the \ADA\ text}\6
\4\&{begin}\6
\\{save\_tok\_ptr}\?\K\?\\{tok\_ptr};\5
\\{save\_text\_ptr}\?\K\?\\{text\_ptr};\5
\\{save\_next\_control}\?\K\?\\{next\_control};\5
\\{next\_control}\?\K\?\H{|};\5
\|p\?\K\?\\{ADA\_translate};\5
\\{app}\?\1(\?\|p\?+\?\\{inner\_tok\_flag}\?\2)\?;\5
\\{make\_output};\C{output the list}\6
\&{stat}\1\6
\&{if}\1\1\8\\{text\_ptr}\?>\?\\{max\_txt\_ptr}\2\8\&{then}\6
\\{max\_txt\_ptr}\?\K\?\\{text\_ptr};\2\6
\&{end}\8\2\&{if}\1;\6
\&{if}\1\1\8\\{tok\_ptr}\?>\?\\{max\_tok\_ptr}\2\8\&{then}\6
\\{max\_tok\_ptr}\?\K\?\\{tok\_ptr};\2\6
\&{end}\8\2\&{if}\1;\2\6
\&{tats}\6
\\{text\_ptr}\?\K\?\\{save\_text\_ptr};\5
\\{tok\_ptr}\?\K\?\\{save\_tok\_ptr};\C{forget the tokens}\6
\\{next\_control}\?\K\?\\{save\_next\_control};\C{restore \\{next\_control} to
original state}\2\6
\&{end}\8\\{output\_ADA};\par
\fi

\M251. Here is \.{AWEAVE}'s major output handler.

\Y\P\?\4\X245:Procedures in \\{second\_phase}\X\mathrel{+}\S\?\6
\&{procedure}\8\1\\{make\_output}\2\8\&{is}\8\C{outputs the equivalents of
tokens}\1\6
\|a\?:\?\\{eight\_bits};\C{current output byte}\6
\|b\?:\?\\{eight\_bits};\C{next output byte}\6
\|k\?,\39\?\\{k\_limit}\?:\?\\{byte\_index};\C{indices into \\{byte\_mem}}\6
\|w\?:\?\\{w\_range};\C{row of \\{byte\_mem}}\6
\|j\?:\?\\{long\_buffer\_index};\C{index into \\{buffer}}\6
\|d\?:\?\p{INTEGER};\C{counts consecutive \H{\$} signs or spaces}\6
\\{string\_delimiter}\?:\?\\{ASCII\_code};\C{first and last character of
string being copied}\6
\\{save\_loc}\?,\39\?\\{save\_limit}\?:\?\\{long\_buffer\_index};\C{\\{loc} and
\\{limit} to be restored}\6
\\{cur\_mod\_name}\?:\?\\{name\_pointer};\C{name of module being output}\6
\\{save\_mode}\?:\?\\{output\_mode};\C{value of \\{cur\_mode} before a sequence
of breaks}\6
\\{ignore\_spaces}\?:\?\p{BOOLEAN};\C{should spaces be ignored?}\6
\4\&{begin}\6
\\{app}\?\1(\?\\{end\_translation}\?\2)\?;\C{append a sentinel}\6
\\{freeze\_text};\5
\\{push\_level}\?\1(\?\\{text\_ptr}\?-\?1\?\2)\?;\6
\1\&{loop}\6
\|a\?\K\?\\{get\_output};\5
\\{ignore\_spaces}\?\K\?\p{FALSE};\6
\4\BL\\{reswitch}\EL\8\&{case}\1\8\|a\1\8\&{is}\8\6
\4\&{when}\8\\{end\_translation}\KR\6
\&{return};\6
\4\&{when}\8\\{identifier}\?\mathrel{\.|}\?\\{res\_word}\KR\6
\X253:Output an identifier\X;\6
\4\&{when}\8\\{mod\_name}\KR\6
\X257:Output a module name\X;\6
\4\&{when}\8\\{math\_bin}\?\mathrel{\.|}\?\\{math\_op}\?\mathrel{\.|}\?\\{math%
\_rel}\KR\6
\X254:Output a \.{\\math} operator\X;\6
\4\&{when}\8\\{cancel}\KR\6
\1\&{loop}\6
\|a\?\K\?\\{get\_output};\6
\?\&{exit}\8\&{when}\8\1(\?\|a\?<\?\\{backup}\?\2)\V\1(\?\|a\?>\?\\{big\_force}%
\?\2)\?;\2\6
\&{end}\8\&{loop};\6
\\{ignore\_spaces}\?\K\?\p{FALSE};\5
\&{goto}\8\\{reswitch};\6
\4\&{when}\8\\{big\_cancel}\KR\6
\1\&{loop}\6
\|a\?\K\?\\{get\_output};\6
\?\&{exit}\8\&{when}\8\1(\1(\?\|a\?<\?\\{backup}\?\2)\W\1(\?\|a\?\I\?\H{\ }\?%
\2)\W\1(\?\|a\?\notin\?\\{pspace}\?\to\?\\{sspace}\?\2)\2)\V\1(\?\|a\?>\?\\{big%
\_force}\?\2)\?;\2\6
\&{end}\8\&{loop};\6
\\{ignore\_spaces}\?\K\?\p{FALSE};\5
\&{goto}\8\\{reswitch};\6
\4\&{when}\8\\{pdollar}\8\\{pspace}\?\to\?\\{sspace}\KR\6
\X252:Count the \\{pdollar} and space tokens; output the appropriate control
sequences and \&{goto}\8\\{reswitch}\X;\6
\4\&{when}\8\\{indent}\?\to\?\\{big\_force}\KR\6
\X255:Output a \(control, look ahead in case of line breaks, possibly \&{goto}%
\8\\{reswitch}\X;\6
\4\&{when}\8\&{others}\8\KR\6
\\{out1}\?\1(\?\|a\?\2)\?;\C{otherwise \|a is an ASCII character}\2\6
\4\2\&{end}\8\&{case};\2\6
\&{end}\8\&{loop};\2\6
\&{end}\8\\{make\_output};\par
\fi

\M252. Here we look ahead to see whether we have to output a begin-math control
sequence and/or a space control sequence

\Y\P\?\4\X252:Count the \\{pdollar} and space tokens; output the appropriate
control sequences and \&{goto}\8\\{reswitch}\X\S\?\6
\|d\?\K\?0;\5
\|b\?\K\?0;\6
\1\&{while}\?\8\1(\?\|a\?=\?\\{pdollar}\?\2)\V\1(\?\|a\?\in\?\\{pspace}\?\to\?%
\\{sspace}\?\2)\?\8\&{loop}\6
\&{if}\1\1\8\|a\?=\?\\{pdollar}\2\8\&{then}\6
\\{incr}\?\1(\?\|d\?\2)\?;\6
\4\&{else}\6
\|b\?\K\?\|b\?+\?\|a\?-\?\\{pspace}\?+\?1;\2\6
\&{end}\8\2\&{if}\1;\6
\|a\?\K\?\\{get\_output};\6
\&{if}\1\1\8\\{cur\_mode}\?=\?\\{inner}\2\8\&{then}\6
\1\&{while}\8\|a\?\in\?\\{big\_cancel}\?\to\?\\{backup}\8\&{loop}\6
\&{if}\1\1\8\|a\?=\?\\{opt}\2\8\&{then}\6
\|a\?\K\?\\{get\_output};\C{\\{opt} is followed by digit}\2\6
\&{end}\8\2\&{if}\1;\6
\|a\?\K\?\\{get\_output};\2\6
\&{end}\8\&{loop};\2\6
\&{end}\8\2\&{if}\1;\2\6
\&{end}\8\&{loop};\6
\&{if}\1\1\?\8\1(\?\|d\?\mathbin{\&{mod}}\?2\?\2)=\?1\2\8\&{then}\6
\\{out2}\?\1(\?\H{\\}\?\2)\1(\?\H{?}\?\2)\?;\2\6
\&{end}\8\2\&{if}\1;\6
\&{if}\1\1\?\8\R\?\\{ignore\_spaces}\?\W\?\|b\?>\?1\2\8\&{then}\6
\\{out2}\?\1(\?\H{\\}\?\2)\1(\?\H{8}\?\2)\?;\2\6
\&{end}\8\2\&{if}\1;\6
\\{ignore\_spaces}\?\K\?\p{FALSE};\5
\&{goto}\8\\{reswitch}\par
\U section~251.\fi

\M253. An identifier of length one does not have to be enclosed in braces, and
it
looks slightly better if set in a math-italic font instead of a (slightly
narrower) text-italic font. Thus we output `\.{\\\char'174a}' but
`\.{\\\\\{aa\}}'.

\Y\P\?\4\X253:Output an identifier\X\S\?\6
\\{out1}\?\1(\?\H{\\}\?\2)\?;\6
\&{if}\1\1\8\|a\?=\?\\{identifier}\2\8\&{then}\6
\&{if}\1\1\8\\{ilk}\?\1(\?\\{cur\_name}\?\2)=\?\\{predefined}\2\8\&{then}\6
\\{out1}\?\1(\?\H{p}\?\2)\?;\6
\4\&{elsif}\1\8\\{length}\?\1(\?\\{cur\_name}\?\2)=\?1\2\8\&{then}\6
\\{out1}\?\1(\?\H{|}\?\2)\?;\6
\4\&{else}\6
\\{out1}\?\1(\?\H{\\}\?\2)\?;\2\6
\&{end}\8\2\&{if}\1;\6
\4\&{else}\6
\\{out1}\?\1(\?\H{\&}\?\2)\?;\C{\|a\?=\?\\{res\_word}}\2\6
\&{end}\8\2\&{if}\1;\6
\&{if}\1\1\8\\{length}\?\1(\?\\{cur\_name}\?\2)=\?1\2\8\&{then}\6
\\{out1}\?\1(\?\\{byte\_mem}\?\1(\?\\{cur\_name}\?\mathbin{\&{mod}}\?\\{ww}\?,%
\39\?\\{byte\_start}\?\1(\?\\{cur\_name}\?\2)\2)\2)\?;\6
\4\&{else}\6
\\{out\_name}\?\1(\?\\{cur\_name}\?\2)\?;\2\6
\&{end}\8\2\&{if}\1\par
\U section~251.\fi

\M254. \P\?\X254:Output a \.{\\math} operator\X\S\?\6
\\{out5}\?\1(\?\H{\\}\?\2)\1(\?\H{m}\?\2)\1(\?\H{a}\?\2)\1(\?\H{t}\?\2)\1(\?%
\H{h}\?\2)\?;\6
\&{if}\1\1\8\|a\?=\?\\{math\_bin}\2\8\&{then}\6
\\{out3}\?\1(\?\H{b}\?\2)\1(\?\H{i}\?\2)\1(\?\H{n}\?\2)\?;\6
\4\&{elsif}\1\8\|a\?=\?\\{math\_rel}\2\8\&{then}\6
\\{out3}\?\1(\?\H{r}\?\2)\1(\?\H{e}\?\2)\1(\?\H{l}\?\2)\?;\6
\4\&{else}\6
\\{out2}\?\1(\?\H{o}\?\2)\1(\?\H{p}\?\2)\?;\2\6
\&{end}\8\2\&{if}\1;\6
\\{out1}\?\1(\?\H{\{}\?\2)\?\par
\U section~251.\fi

\M255. The current mode does not affect the behavior of \.{AWEAVE}'s output
routine
except when we are outputting control tokens.

\Y\P\?\4\X255:Output a \(control, look ahead in case of line breaks, possibly %
\&{goto}\8\\{reswitch}\X\S\?\6
\&{if}\1\1\8\|a\?<\?\\{break\_space}\2\8\&{then}\6
\&{if}\1\1\8\\{cur\_mode}\?=\?\\{outer}\2\8\&{then}\6
\\{out2}\?\1(\?\H{\\}\?\2)\1(\?\|a\?-\?\\{cancel}\?+\?\H{0}\?\2)\?;\6
\&{if}\1\1\8\|a\?=\?\\{opt}\2\8\&{then}\6
\\{out1}\?\1(\?\\{get\_output}\?\2)\?;\C{\\{opt} is followed by a digit}\2\6
\&{end}\8\2\&{if}\1;\6
\4\&{elsif}\1\8\|a\?=\?\\{opt}\2\8\&{then}\6
\|b\?\K\?\\{get\_output};\C{ignore digit following \\{opt}}\2\6
\&{end}\8\2\&{if}\1;\6
\4\&{else}\6
\X256:Look ahead for strongest line break, \&{goto}\8\\{reswitch}\X;\2\6
\&{end}\8\2\&{if}\1\par
\U section~251.\fi

\M256. If several of the tokens \\{break\_space}, \\{force}, \\{big\_force}
occur in a
row, possibly mixed with blank spaces (which are ignored),
the largest one is used. A line break also occurs in the output file,
except at the very end of the translation. The very first line break
is suppressed (i.e., a line break that follows `\.{\\Y\\P}').

\Y\P\?\4\X256:Look ahead for strongest line break, \&{goto}\8\\{reswitch}\X\S\?%
\6
\|b\?\K\?\|a;\5
\\{save\_mode}\?\K\?\\{cur\_mode};\6
\1\&{loop}\6
\|a\?\K\?\\{get\_output};\6
\&{if}\1\1\?\8\1(\?\|a\?=\?\\{cancel}\?\2)\V\1(\?\|a\?=\?\\{big\_cancel}\?\2)\?%
\2\8\&{then}\6
\&{goto}\8\\{reswitch};\C{\\{cancel} overrides everything}\2\6
\&{end}\8\2\&{if}\1;\6
\&{if}\1\1\?\8\1(\1(\?\|a\?\I\?\H{\ }\?\2)\W\1(\?\|a\?<\?\\{break\_space}\?\2)%
\W\1(\?\|a\?\notin\?\\{pspace}\?\to\?\\{sspace}\?\2)\2)\V\1(\?\|a\?>\?\\{big%
\_force}\?\2)\?\2\8\&{then}\6
\&{if}\1\1\8\\{save\_mode}\?=\?\\{outer}\2\8\&{then}\6
\&{if}\1\1\?\8\1(\?\\{out\_ptr}\?>\?3\?\2)\ANT\1(\?\\{out\_buf}\?\1(\?\\{out%
\_ptr}\?-\?3\?\to\?\\{out\_ptr}\?\2)=\1(\?\H{\\}\?,\39\?\H{Y}\?,\39\?\H{\\}\?,%
\39\?\H{P}\?\2)\2)\?\2\8\&{then}\6
\&{goto}\8\\{reswitch};\2\6
\&{end}\8\2\&{if}\1;\6
\\{out2}\?\1(\?\H{\\}\?\2)\1(\?\|b\?-\?\\{cancel}\?+\?\H{0}\?\2)\?;\5
\\{ignore\_spaces}\?\K\?\|b\?\in\?\\{force}\?\to\?\\{big\_force};\6
\&{if}\1\1\8\|a\?\I\?\\{end\_translation}\2\8\&{then}\6
\\{finish\_line};\2\6
\&{end}\8\2\&{if}\1;\6
\4\&{elsif}\1\?\8\1(\?\|a\?\I\?\\{end\_translation}\?\2)\W\1(\?\\{cur\_mode}\?=%
\?\\{inner}\?\2)\?\2\8\&{then}\6
\\{out1}\?\1(\?\H{\ }\?\2)\?;\2\6
\&{end}\8\2\&{if}\1;\6
\&{goto}\8\\{reswitch};\2\6
\&{end}\8\2\&{if}\1;\6
\&{if}\1\1\8\|a\?>\?\|b\2\8\&{then}\6
\|b\?\K\?\|a;\2\6
\&{end}\8\2\&{if}\1;\C{if \|a\?=\?\H{\ } we have \|a\?<\?\|b}\2\6
\&{end}\8\&{loop}\par
\U section~255.\fi

\M257. The remaining part of \\{make\_output} is somewhat more complicated.
When we
output a module name, we may need to enter the parsing and translation
routines, since the name may contain \ADA\ code embedded in
\ab\ constructions. This \ADA\ code is placed at the end of the active
input buffer and the translation process uses the end of the active
\\{tok\_mem} area.

\Y\P\?\4\X257:Output a module name\X\S\?\6
\\{out2}\?\1(\?\H{\\}\?\2)\1(\?\H{X}\?\2)\?;\5
\\{cur\_xref}\?\K\?\\{xref}\?\1(\?\\{cur\_name}\?\2)\?;\6
\&{if}\1\1\8\\{num}\?\1(\?\\{cur\_xref}\?\2)\G\?\\{def\_flag}\2\8\&{then}\6
\\{out\_mod}\?\1(\?\\{num}\?\1(\?\\{cur\_xref}\?\2)-\?\\{def\_flag}\?\2)\?;\6
\&{if}\1\1\8\\{phase\_three}\2\8\&{then}\6
\\{cur\_xref}\?\K\?\\{xlink}\?\1(\?\\{cur\_xref}\?\2)\?;\6
\1\&{while}\8\\{num}\?\1(\?\\{cur\_xref}\?\2)\G\?\\{def\_flag}\8\&{loop}\6
\\{out2}\?\1(\?\H{,}\?\2)\1(\?\H{\ }\?\2)\?;\5
\\{out\_mod}\?\1(\?\\{num}\?\1(\?\\{cur\_xref}\?\2)-\?\\{def\_flag}\?\2)\?;\5
\\{cur\_xref}\?\K\?\\{xlink}\?\1(\?\\{cur\_xref}\?\2)\?;\2\6
\&{end}\8\&{loop};\2\6
\&{end}\8\2\&{if}\1;\6
\4\&{else}\6
\\{out1}\?\1(\?\H{0}\?\2)\?;\C{output the module number, or zero if it was
undefined}\2\6
\&{end}\8\2\&{if}\1;\6
\\{out1}\?\1(\?\H{:}\?\2)\?;\5
\X258:Output the text of the module name\X;\6
\\{out2}\?\1(\?\H{\\}\?\2)\1(\?\H{X}\?\2)\?\par
\U section~251.\fi

\M258. \P\?\X258:Output the text of the module name\X\S\?\6
\|k\?\K\?\\{byte\_start}\?\1(\?\\{cur\_name}\?\2)\?;\5
\|w\?\K\?\\{cur\_name}\?\mathbin{\&{mod}}\?\\{ww};\5
\\{k\_limit}\?\K\?\\{byte\_start}\?\1(\?\\{cur\_name}\?+\?\\{ww}\?\2)\?;\5
\\{cur\_mod\_name}\?\K\?\\{cur\_name};\6
\1\&{while}\8\|k\?<\?\\{k\_limit}\8\&{loop}\6
\|b\?\K\?\\{byte\_mem}\?\1(\?\|w\?,\39\?\|k\?\2)\?;\5
\\{incr}\?\1(\?\|k\?\2)\?;\6
\&{if}\1\1\8\|b\?=\?\H{@}\2\8\&{then}\6
\X259:Skip next character, give error if not `\.{@}'\X;\2\6
\&{end}\8\2\&{if}\1;\6
\&{if}\1\1\8\|b\?\I\?\H{|}\2\8\&{then}\6
\\{out1}\?\1(\?\|b\?\2)\?;\6
\4\&{else}\6
\X260:Copy the \ADA\ text into \\{buffer}\?(\?\\{limit}\?+\?1\?\to\?\|j\?)\?\X;%
\6
\\{save\_loc}\?\K\?\\{loc};\5
\\{save\_limit}\?\K\?\\{limit};\5
\\{loc}\?\K\?\\{limit}\?+\?2;\5
\\{limit}\?\K\?\|j\?+\?1;\5
\\{buffer}\?\1(\?\\{limit}\?\2)\K\?\H{|};\5
\\{output\_ADA};\5
\\{loc}\?\K\?\\{save\_loc};\5
\\{limit}\?\K\?\\{save\_limit};\2\6
\&{end}\8\2\&{if}\1;\2\6
\&{end}\8\&{loop}\par
\U section~257.\fi

\M259. \P\?\X259:Skip next character, give error if not `\.{@}'\X\S\?\6
\&{if}\1\1\8\\{byte\_mem}\?\1(\?\|w\?,\39\?\|k\?\2)\I\?\H{@}\2\8\&{then}\6
\\{print\_nl}\?\1(\?\.{"!\ Illegal\ control\ code\ in\ section\ name:"}\?\2)\?;%
\5
\\{print\_nl}\?\1(\?\.{"<"}\?\2)\?;\5
\\{print\_id}\?\1(\?\\{cur\_mod\_name}\?\2)\?;\5
\\{print}\?\1(\?\.{">\ "}\?\2)\?;\5
\\{mark\_error};\2\6
\&{end}\8\2\&{if}\1;\6
\\{incr}\?\1(\?\|k\?\2)\?\par
\U section~258.\fi

\M260. The \ADA\ text enclosed in \ab\ should not contain `\v' characters,
except within strings. We put a `\v' at the front of the buffer, so that an
error message that displays the whole buffer will look a little bit sensible.
The variable \\{string\_delimiter} is zero outside of strings, otherwise it
equals the delimiter that began the string being copied.

\Y\P\?\4\X260:Copy the \ADA\ text into \\{buffer}\?(\?\\{limit}\?+\?1\?\to\?\|j%
\?)\?\X\S\?\6
\|j\?\K\?\\{limit}\?+\?1;\5
\\{buffer}\?\1(\?\|j\?\2)\K\?\H{|};\5
\\{string\_delimiter}\?\K\?0;\6
\1\&{loop}\6
\&{if}\1\1\8\|k\?\G\?\\{k\_limit}\2\8\&{then}\6
\\{print\_nl}\?\1(\?\.{"!\ ADA\ text\ in\ section\ name\ didn\'t\ end:"}\?\2)%
\?;\5
\\{print\_nl}\?\1(\?\.{"<"}\?\2)\?;\5
\\{print\_id}\?\1(\?\\{cur\_mod\_name}\?\2)\?;\5
\\{print}\?\1(\?\.{">\ "}\?\2)\?;\5
\\{mark\_error};\5
\&{exit};\2\6
\&{end}\8\2\&{if}\1;\6
\|b\?\K\?\\{byte\_mem}\?\1(\?\|w\?,\39\?\|k\?\2)\?;\5
\\{incr}\?\1(\?\|k\?\2)\?;\6
\&{if}\1\1\8\|b\?=\?\H{@}\2\8\&{then}\6
\X261:Copy a control code into the buffer\X;\6
\4\&{else}\6
\&{if}\1\1\8\|b\?=\?\H{"}\2\8\&{then}\6
\&{if}\1\1\8\\{string\_delimiter}\?=\?0\2\8\&{then}\6
\\{string\_delimiter}\?\K\?\|b;\6
\4\&{elsif}\1\8\\{string\_delimiter}\?=\?\|b\2\8\&{then}\6
\\{string\_delimiter}\?\K\?0;\2\6
\&{end}\8\2\&{if}\1;\6
\4\&{elsif}\1\?\8\1(\?\|b\?=\?\H{\'}\?\2)\W\1(\?\\{byte\_mem}\?\1(\?\|w\?,\39\?%
\|k\?+\?1\?\2)=\?\H{\'}\?\2)\W\1(\?\\{string\_delimiter}\?=\?0\?\2)\?\2\8%
\&{then}\6
\|j\?\K\?\|j\?+\?2;\5
\\{buffer}\?\1(\?\|j\?-\?1\?\to\?\|j\?\2)\K\1(\?\H{\'}\?,\39\?\\{byte\_mem}\?%
\1(\?\|w\?,\39\?\|k\?\2)\2)\?;\5
\|k\?\K\?\|k\?+\?3;\5
\|b\?\K\?\H{\'};\2\6
\&{end}\8\2\&{if}\1;\6
\&{if}\1\1\?\8\1(\?\|b\?\I\?\H{|}\?\2)\V\1(\?\\{string\_delimiter}\?\I\?0\?\2)%
\?\2\8\&{then}\6
\&{if}\1\1\8\|j\?>\?\\{long\_buf\_size}\?-\?3\2\8\&{then}\6
\\{overflow}\?\1(\?\.{"buffer"}\?\2)\?;\2\6
\&{end}\8\2\&{if}\1;\6
\\{incr}\?\1(\?\|j\?\2)\?;\5
\\{buffer}\?\1(\?\|j\?\2)\K\?\|b;\6
\4\&{else}\6
\&{exit};\2\6
\&{end}\8\2\&{if}\1;\2\6
\&{end}\8\2\&{if}\1;\2\6
\&{end}\8\&{loop}\par
\U section~258.\fi

\M261. \P\?\X261:Copy a control code into the buffer\X\S\?\6
\&{if}\1\1\8\|j\?>\?\\{long\_buf\_size}\?-\?4\2\8\&{then}\6
\\{overflow}\?\1(\?\.{"buffer"}\?\2)\?;\2\6
\&{end}\8\2\&{if}\1;\6
\\{buffer}\?\1(\?\|j\?+\?1\?\2)\K\?\H{@};\5
\\{buffer}\?\1(\?\|j\?+\?2\?\2)\K\?\\{byte\_mem}\?\1(\?\|w\?,\39\?\|k\?\2)\?;\5
\|j\?\K\?\|j\?+\?2;\5
\\{incr}\?\1(\?\|k\?\2)\?\par
\U section~260.\fi

\N262.  Phase two processing.
We have assembled enough pieces of the puzzle in order to be ready to specify
the processing in \.{AWEAVE}'s main pass over the source file. Phase two
is analogous to phase one, except that more work is involved because we must
actually output the \TeX\ material instead of merely looking at the
\.{AWEB} specifications.

\Y\P\?\4\X262:Phase II: Read all the text again and translate it to \TeX\ form%
\X\S\?\6
\\{reset\_input};\5
\\{print\_nl}\?\1(\?\.{"Writing\ the\ output\ file..."}\?\2)\?;\5
\\{module\_count}\?\K\?0;\5
\\{copy\_limbo};\5
\\{finish\_line};\5
\\{flush\_buffer}\?\1(\?0\?,\39\?\p{FALSE}\?\2)\?;\C{insert a blank line, it
looks nice}\6
\1\&{while}\?\8\R\?\\{input\_has\_ended}\8\&{loop}\6
\X264:Translate the \(current module\X;\2\6
\&{end}\8\&{loop}\par
\U section~304.\fi

\M263. The output file will contain the control sequence \.{\\Y} between
non-null
sections of a module, e.g., between the \TeX\ and definition parts if both
are nonempty. This puts a little white space between the parts when they are
printed. However, we don't want \.{\\Y} to occur between two definitions
within a single module. The variables \\{out\_line} or \\{out\_ptr} will
change if a section is non-null, so the following macros `\\{save\_position}'
and `\\{emit\_space\_if\_needed}' are able to handle the situation:

\Y\P\4\D\\{save\_position}\?\S\?\\{save\_line}\?\K\?\\{out\_line};\5
\\{save\_place}\?\K\?\\{out\_ptr}\par
\P\4\D\\{emit\_space\_if\_needed}\?\S\?\6
\&{if}\1\1\?\8\1(\?\\{save\_line}\?\I\?\\{out\_line}\?\2)\V\1(\?\\{save\_place}%
\?\I\?\\{out\_ptr}\?\2)\?\2\8\&{then}\6
\\{out2}\?\1(\?\H{\\}\?\2)\1(\?\H{Y}\?\2)\?;\2\6
\&{end}\8\2\&{if}\1\par
\Y\P\?\4\X14:Types and objects visible from \\{data\_structures}\X\mathrel{+}\S%
\?\6
\\{save\_line}\?:\?\p{INTEGER};\C{former value of \\{out\_line}}\6
\\{save\_place}\?:\?\\{sixteen\_bits};\C{former value of \\{out\_ptr}}\par
\fi

\M264. \P\?\X264:Translate the \(current module\X\S\?\6
\\{incr}\?\1(\?\\{module\_count}\?\2)\?;\6
\X265:Output the code for the beginning of a new module\X;\6
\\{save\_position};\6
\X266:Translate the \TeX\ part of the current module\X;\6
\X267:Translate the \(definition part of the current module\X;\6
\X272:Translate the \(\ADA\ part of the current module\X;\6
\X275:Show cross references to this module\X;\6
\X281:Output the code for the end of a module\X\par
\U section~262.\fi

\M265. Modules beginning with the \.{AWEB} control sequence `\.{@\ }' start in
the
output with the \TeX\ control sequence `\.{\\M}', followed by the module
number. Similarly, `\.{@*}' modules lead to the control sequence `\.{\\N}'.
If this is a changed module, we put \.{*} just before the module number.

\Y\P\?\4\X265:Output the code for the beginning of a new module\X\S\?\6
\\{out1}\?\1(\?\H{\\}\?\2)\?;\6
\&{if}\1\1\8\\{buffer}\?\1(\?\\{loc}\?-\?1\?\2)\I\?\H{*}\2\8\&{then}\6
\\{out1}\?\1(\?\H{M}\?\2)\?;\6
\4\&{else}\6
\\{out1}\?\1(\?\H{N}\?\2)\?;\5
\\{print}\?\1(\?\.{"*"}\?\2)\?;\5
\\{print\_int}\?\1(\?\\{module\_count}\?,\39\?1\?\2)\?;\5
\\{update\_terminal};\C{print a progress report}\2\6
\&{end}\8\2\&{if}\1;\6
\\{out\_mod}\?\1(\?\\{module\_count}\?\2)\?;\5
\\{out2}\?\1(\?\H{.}\?\2)\1(\?\H{\ }\?\2)\?\par
\U section~264.\fi

\M266. In the \TeX\ part of a module, we simply copy the source text, except
that
index entries are not copied and \ADA\ text within \ab\ is translated.

\Y\P\?\4\X266:Translate the \TeX\ part of the current module\X\S\?\6
\1\&{loop}\6
\\{next\_control}\?\K\?\\{copy\_TeX};\6
\&{case}\1\8\\{next\_control}\1\8\&{is}\8\6
\4\&{when}\8\H{|}\KR\6
\\{init\_stack};\5
\\{output\_ADA};\6
\4\&{when}\8\H{@}\KR\6
\\{out1}\?\1(\?\H{@}\?\2)\?;\6
\4\&{when}\8\\{TeX\_string}\?\mathrel{\.|}\?\\{xref\_roman}\?\to\?\\{xref%
\_reserved}\?\mathrel{\.|}\?\\{module\_name}\KR\6
\\{loc}\?\K\?\\{loc}\?-\?2;\5
\\{next\_control}\?\K\?\\{get\_next};\C{skip to \.{@>}}\6
\&{if}\1\1\8\\{next\_control}\?=\?\\{TeX\_string}\2\8\&{then}\6
\\{err\_print}\?\1(\?\.{"!\ TeX\ string\ should\ be\ in\ ADA\ text\ only"}\?\2)%
\?;\2\6
\&{end}\8\2\&{if}\1;\6
\4\&{when}\8\\{begin\_comment}\?\mathrel{\.|}\?\\{end\_comment}\?\mathrel{\.|}%
\?\\{thin\_space}\?\mathrel{\.|}\?\\{math\_break}\?\mathrel{\.|}\?\\{line%
\_break}\?\mathrel{\.|}\?\\{big\_line\_break}\?\mathrel{\.|}\?\\{no\_line%
\_break}\?\mathrel{\.|}\?\\{join}\?\mathrel{\.|}\?\\{pseudo\_semi}\KR\6
\\{err\_print}\?\1(\?\.{"!\ You\ can\'t\ do\ that\ in\ TeX\ text"}\?\2)\?;\6
\4\&{when}\8\&{others}\8\KR\6
\&{null};\2\6
\4\2\&{end}\8\&{case};\6
\?\&{exit}\8\&{when}\?\8\\{next\_control}\?\G\?\\{format};\2\6
\&{end}\8\&{loop}\par
\U section~264.\fi

\M267. When we get to the following code we have \\{next\_control}\?\G\?%
\\{format}, and
the token memory is in its initial empty state.

\Y\P\?\4\X267:Translate the \(definition part of the current module\X\S\?\6
\&{if}\1\1\8\\{next\_control}\?\L\?\\{definition}\2\8\&{then}\C{definition part
non-empty}\6
\\{emit\_space\_if\_needed};\5
\\{save\_position};\2\6
\&{end}\8\2\&{if}\1;\6
\1\&{while}\8\\{next\_control}\?\L\?\\{definition}\8\&{loop}\C{\\{format} or %
\\{definition}}\6
\\{init\_stack};\6
\&{if}\1\1\8\\{next\_control}\?=\?\\{definition}\2\8\&{then}\6
\X269:Start a macro definition\X;\6
\4\&{else}\6
\X270:Start a format definition\X;\2\6
\&{end}\8\2\&{if}\1;\6
\\{outer\_parse};\5
\\{finish\_ADA};\2\6
\&{end}\8\&{loop}\par
\U section~264.\fi

\M268. The \\{finish\_ADA} procedure outputs the translation of the current
scraps, preceded by the control sequence `\.{\\P}' and followed by the
control sequence `\.{\\par}'. It also restores the token and scrap
memories to their initial empty state.

A \\{force} token is appended to the current scraps before translation
takes place, so that the translation will normally end with \.{\\6} or
\.{\\7} (the \TeX\ macros for \\{force} and \\{big\_force}). This \.{\\6} or
\.{\\7} is replaced by the concluding \.{\\par} or by \.{\\Y\\par}.

\Y\P\?\4\X245:Procedures in \\{second\_phase}\X\mathrel{+}\S\?\6
\&{procedure}\8\1\\{finish\_ADA}\2\8\&{is}\8\C{finishes a definition or a \ADA\
part}\1\6
\|p\?:\?\\{text\_pointer};\C{translation of the scraps}\6
\4\&{begin}\6
\\{out2}\?\1(\?\H{\\}\?\2)\1(\?\H{P}\?\2)\?;\5
\\{app\_tok}\?\1(\?\\{force}\?\2)\?;\5
\\{app\_comment};\5
\|p\?\K\?\\{translate};\5
\\{app}\?\1(\?\|p\?+\?\\{tok\_flag}\?\2)\?;\5
\\{make\_output};\C{output the list}\6
\&{if}\1\1\?\8\1(\?\\{out\_ptr}\?>\?1\?\2)\ANT\1(\?\\{out\_buf}\?\1(\?\\{out%
\_ptr}\?-\?1\?\2)=\?\H{\\}\?\2)\?\2\8\&{then}\6
\&{if}\1\1\8\\{out\_buf}\?\1(\?\\{out\_ptr}\?\2)=\?\H{6}\2\8\&{then}\6
\\{out\_ptr}\?\K\?\\{out\_ptr}\?-\?2;\6
\4\&{elsif}\1\8\\{out\_buf}\?\1(\?\\{out\_ptr}\?\2)=\?\H{7}\2\8\&{then}\6
\\{out\_buf}\?\1(\?\\{out\_ptr}\?\2)\K\?\H{Y};\2\6
\&{end}\8\2\&{if}\1;\2\6
\&{end}\8\2\&{if}\1;\6
\\{out4}\?\1(\?\H{\\}\?\2)\1(\?\H{p}\?\2)\1(\?\H{a}\?\2)\1(\?\H{r}\?\2)\?;\5
\\{finish\_line}; \6
\&{stat}\1\6
\&{if}\1\1\8\\{text\_ptr}\?>\?\\{max\_txt\_ptr}\2\8\&{then}\6
\\{max\_txt\_ptr}\?\K\?\\{text\_ptr};\2\6
\&{end}\8\2\&{if}\1;\6
\&{if}\1\1\8\\{tok\_ptr}\?>\?\\{max\_tok\_ptr}\2\8\&{then}\6
\\{max\_tok\_ptr}\?\K\?\\{tok\_ptr};\2\6
\&{end}\8\2\&{if}\1;\6
\&{if}\1\1\8\\{scrap\_ptr}\?>\?\\{max\_scr\_ptr}\2\8\&{then}\6
\\{max\_scr\_ptr}\?\K\?\\{scrap\_ptr};\2\6
\&{end}\8\2\&{if}\1;\2\6
\&{tats}\6
\\{tok\_ptr}\?\K\?1;\5
\\{text\_ptr}\?\K\?1;\5
\\{scrap\_ptr}\?\K\?0;\C{forget the tokens and the scraps}\2\6
\&{end}\8\\{finish\_ADA};\par
\fi

\M269. \P\?\X269:Start a macro definition\X\S\?\6
\\{sc5}\?\1(\?\\{pdollar}\?\2)\1(\?\\{backup}\?\2)\1(\?\H{\\}\?\2)\1(\?\H{D}\?%
\2)\1(\?\\{pdollar}\?\2)\1(\?\\{math}\?\2)\?;\C{this will produce `\&{define
}'}\6
\\{next\_control}\?\K\?\\{get\_next};\6
\&{if}\1\1\8\\{next\_control}\?\I\?\\{identifier}\2\8\&{then}\6
\\{err\_print}\?\1(\?\.{"!\ Improper\ macro\ definition"}\?\2)\?;\6
\4\&{else}\6
\\{sc3}\?\1(\?\\{pdollar}\?\2)\1(\?\\{id\_flag}\?+\?\\{id\_lookup}\?\1(\?%
\\{normal}\?\2)\2)\1(\?\\{pdollar}\?\2)\1(\?\\{math}\?\2)\?;\2\6
\&{end}\8\2\&{if}\1;\6
\\{next\_control}\?\K\?\\{get\_next}\par
\U section~267.\fi

\M270. \P\?\X270:Start a format definition\X\S\?\6
\\{sc5}\?\1(\?\\{pdollar}\?\2)\1(\?\\{backup}\?\2)\1(\?\H{\\}\?\2)\1(\?\H{F}\?%
\2)\1(\?\\{pdollar}\?\2)\1(\?\\{math}\?\2)\?;\C{this will produce `\&{format
}'}\6
\\{next\_control}\?\K\?\\{get\_next};\6
\&{if}\1\1\8\\{next\_control}\?=\?\\{identifier}\2\8\&{then}\6
\\{sc3}\?\1(\?\\{pdollar}\?\2)\1(\?\\{id\_flag}\?+\?\\{id\_lookup}\?\1(\?%
\\{normal}\?\2)\2)\1(\?\\{pdollar}\?\2)\1(\?\\{math}\?\2)\?;\5
\\{next\_control}\?\K\?\\{get\_next};\6
\&{if}\1\1\8\\{next\_control}\?=\?\\{equivalence\_sign}\2\8\&{then}\6
\\{sc2}\?\1(\?\H{\\}\?\2)\1(\?\H{S}\?\2)\1(\?\\{math}\?\2)\?;\C{output an
equivalence sign}\6
\\{next\_control}\?\K\?\\{get\_next};\6
\&{if}\1\1\8\\{next\_control}\?=\?\\{identifier}\2\8\&{then}\6
\\{sc3}\?\1(\?\\{pdollar}\?\2)\1(\?\\{id\_flag}\?+\?\\{id\_lookup}\?\1(\?%
\\{normal}\?\2)\2)\1(\?\\{pdollar}\?\2)\1(\?\\{math}\?\2)\?;\5
\\{sc0}\?\1(\?\\{semi}\?\2)\?;\C{insert an invisible semicolon}\6
\\{next\_control}\?\K\?\\{get\_next};\2\6
\&{end}\8\2\&{if}\1;\2\6
\&{end}\8\2\&{if}\1;\2\6
\&{end}\8\2\&{if}\1;\6
\&{if}\1\1\8\\{scrap\_ptr}\?\I\?5\2\8\&{then}\6
\\{err\_print}\?\1(\?\.{"!\ Improper\ format\ definition"}\?\2)\?;\2\6
\&{end}\8\2\&{if}\1\par
\U section~267.\fi

\M271. Finally, when the \TeX\ and definition parts have been treated, we have
\\{next\_control}\?\G\?\\{begin\_ada}. We will make the global variable \\{this%
\_module}
point to the current module name, if it has a name.

\Y\P\?\4\X14:Types and objects visible from \\{data\_structures}\X\mathrel{+}\S%
\?\6
\\{this\_module}\?:\?\\{name\_pointer};\C{the current module name, or zero}\par
\fi

\M272. \P\?\X272:Translate the \(\ADA\ part of the current module\X\S\?\6
\\{this\_module}\?\K\?0;\6
\&{if}\1\1\8\\{next\_control}\?\L\?\\{module\_name}\2\8\&{then}\6
\\{emit\_space\_if\_needed};\5
\\{init\_stack};\6
\&{if}\1\1\8\\{next\_control}\?=\?\\{begin\_ada}\2\8\&{then}\6
\\{next\_control}\?\K\?\\{get\_next};\6
\4\&{else}\6
\\{this\_module}\?\K\?\\{cur\_module};\6
\X273:Check that \?=\? or \?\S\? follows this module name, and emit the scraps
to start the module definition\X;\2\6
\&{end}\8\2\&{if}\1;\6
\1\&{while}\8\\{next\_control}\?\L\?\\{module\_name}\8\&{loop}\6
\\{outer\_parse};\6
\X274:Emit the scrap for a module name if present\X;\2\6
\&{end}\8\&{loop};\6
\\{finish\_ADA};\2\6
\&{end}\8\2\&{if}\1\par
\U section~264.\fi

\M273. \P\?\X273:Check that \?=\? or \?\S\? follows this module name, and emit
the scraps to start the module definition\X\S\?\6
\1\&{loop}\6
\\{next\_control}\?\K\?\\{get\_next};\6
\?\&{exit}\8\&{when}\?\8\\{next\_control}\?\I\?\H{+};\C{allow optional `%
\.{+=}'}\2\6
\&{end}\8\&{loop};\6
\&{if}\1\1\?\8\1(\?\\{next\_control}\?\I\?\H{=}\?\2)\W\1(\?\\{next\_control}\?%
\I\?\\{equivalence\_sign}\?\2)\?\2\8\&{then}\6
\\{err\_print}\?\1(\?\.{"!\ You\ need\ an\ =\ sign\ after\ the\ section\ name"}%
\?\2)\?;\6
\4\&{else}\6
\\{next\_control}\?\K\?\\{get\_next};\2\6
\&{end}\8\2\&{if}\1;\6
\&{if}\1\1\?\8\1(\?\\{out\_ptr}\?>\?1\?\2)\ANT\1(\?\\{out\_buf}\?\1(\?\\{out%
\_ptr}\?-\?1\?\to\?\\{out\_ptr}\?\2)=\1(\?\H{\\}\?,\39\?\H{Y}\?\2)\2)\?\2\8%
\&{then}\6
\\{app}\?\1(\?\\{backup}\?\2)\?;\C{the module name will be flush left}\2\6
\&{end}\8\2\&{if}\1;\6
\\{sc1}\?\1(\?\\{mod\_flag}\?+\?\\{this\_module}\?\2)\1(\?\\{mod\_scrap}\?\2)%
\?;\5
\\{cur\_xref}\?\K\?\\{xref}\?\1(\?\\{this\_module}\?\2)\?;\6
\&{if}\1\1\8\\{num}\?\1(\?\\{cur\_xref}\?\2)\I\?\\{module\_count}\?+\?\\{def%
\_flag}\2\8\&{then}\6
\\{sc3}\?\1(\?\\{math\_rel}\?\2)\1(\?\H{+}\?\2)\1(\?\H{\}}\?\2)\1(\?\\{math}\?%
\2)\?;\C{module name is multiply defined}\6
\\{this\_module}\?\K\?0;\C{so we won't give cross-reference info here}\2\6
\&{end}\8\2\&{if}\1;\6
\\{sc2}\?\1(\?\H{\\}\?\2)\1(\?\H{S}\?\2)\1(\?\\{math}\?\2)\?;\C{output an
equivalence sign}\6
\\{sc1}\?\1(\?\\{force}\?\2)\1(\?\\{semi}\?\2)\?\C{this forces a line break
unless `\.{@+}' follows}\par
\U section~272.\fi

\M274. \P\?\X274:Emit the scrap for a module name if present\X\S\?\6
\&{if}\1\1\8\\{next\_control}\?<\?\\{module\_name}\2\8\&{then}\6
\\{err\_print}\?\1(\?\.{"!\ You\ can\'t\ do\ that\ in\ ADA\ text"}\?\2)\?;\6
\\{next\_control}\?\K\?\\{get\_next};\6
\4\&{elsif}\1\8\\{next\_control}\?=\?\\{module\_name}\2\8\&{then}\6
\\{sc1}\?\1(\?\\{mod\_flag}\?+\?\\{cur\_module}\?\2)\1(\?\\{mod\_scrap}\?\2)\?;%
\5
\\{next\_control}\?\K\?\\{get\_next};\2\6
\&{end}\8\2\&{if}\1\par
\U section~272.\fi

\M275. Cross references relating to a named module are given after the module
ends.

\Y\P\?\4\X275:Show cross references to this module\X\S\?\6
\&{if}\1\1\8\\{this\_module}\?>\?0\2\8\&{then}\6
\X277:Rearrange the list pointed to by \\{cur\_xref}\X;\6
\\{footnote}\?\1(\?\\{def\_flag}\?\2)\?;\5
\\{footnote}\?\1(\?0\?\2)\?;\2\6
\&{end}\8\2\&{if}\1\par
\U section~264.\fi

\M276. To rearrange the order of the linked list of cross references, we need
four more variables that point to cross reference entries.  We'll end up
with a list pointed to by \\{cur\_xref}.

\Y\P\?\4\X14:Types and objects visible from \\{data\_structures}\X\mathrel{+}\S%
\?\6
\\{next\_xref}\?,\39\?\\{this\_xref}\?,\39\?\\{first\_xref}\?,\39\?\\{mid%
\_xref}\?:\?\\{xref\_number};\C{pointer variables for rearranging a list}\par
\fi

\M277. We want to rearrange the cross reference list so that all the entries
with
\\{def\_flag} come first, in ascending order; then come all the other
entries, in ascending order.  There may be no entries in either one or both
of these categories.

\Y\P\?\4\X277:Rearrange the list pointed to by \\{cur\_xref}\X\S\?\6
\\{first\_xref}\?\K\?\\{xref}\?\1(\?\\{this\_module}\?\2)\?;\5
\\{this\_xref}\?\K\?\\{xlink}\?\1(\?\\{first\_xref}\?\2)\?;\C{bypass current
module number}\6
\&{if}\1\1\8\\{num}\?\1(\?\\{this\_xref}\?\2)>\?\\{def\_flag}\2\8\&{then}\6
\\{mid\_xref}\?\K\?\\{this\_xref};\5
\\{cur\_xref}\?\K\?0;\C{this value doesn't matter}\6
\1\&{loop}\6
\\{next\_xref}\?\K\?\\{xlink}\?\1(\?\\{this\_xref}\?\2)\?;\5
\\{xlink}\?\1(\?\\{this\_xref}\?\2)\K\?\\{cur\_xref};\5
\\{cur\_xref}\?\K\?\\{this\_xref};\5
\\{this\_xref}\?\K\?\\{next\_xref};\6
\?\&{exit}\8\&{when}\?\8\\{num}\?\1(\?\\{this\_xref}\?\2)\L\?\\{def\_flag};\2\6
\&{end}\8\&{loop};\6
\\{xlink}\?\1(\?\\{first\_xref}\?\2)\K\?\\{cur\_xref};\6
\4\&{else}\6
\\{mid\_xref}\?\K\?0;\C{first list null}\2\6
\&{end}\8\2\&{if}\1;\6
\\{cur\_xref}\?\K\?0;\6
\1\&{while}\8\\{this\_xref}\?\I\?0\8\&{loop}\6
\\{next\_xref}\?\K\?\\{xlink}\?\1(\?\\{this\_xref}\?\2)\?;\5
\\{xlink}\?\1(\?\\{this\_xref}\?\2)\K\?\\{cur\_xref};\5
\\{cur\_xref}\?\K\?\\{this\_xref};\5
\\{this\_xref}\?\K\?\\{next\_xref};\2\6
\&{end}\8\&{loop};\6
\&{if}\1\1\8\\{mid\_xref}\?>\?0\2\8\&{then}\6
\\{xlink}\?\1(\?\\{mid\_xref}\?\2)\K\?\\{cur\_xref};\6
\4\&{else}\6
\\{xlink}\?\1(\?\\{first\_xref}\?\2)\K\?\\{cur\_xref};\2\6
\&{end}\8\2\&{if}\1;\6
\\{cur\_xref}\?\K\?\\{xlink}\?\1(\?\\{first\_xref}\?\2)\6
\?\par
\U section~275.\fi

\M278. The \\{footnote} procedure gives cross reference information about
multiply defined module names (if the \\{flag} parameter is \\{def\_flag}), or
about
the uses of a module name (if the \\{flag} parameter is zero). It assumes that
\\{cur\_xref} points to the first cross-reference entry of interest, and it
leaves \\{cur\_xref} pointing to the first element not printed.  Typical
outputs:
`\.{\\A\ section 101.}'; `\.{\\U\ sections 370 and 1009.}';
`\.{\\U\ section 227(2)}';
`\.{\\A\ sections 8, 27\\*, and 64.}'.

\Y\P\?\4\X244:Specification of procedures visible from \\{second\_phase}\X%
\mathrel{+}\S\?\6
\&{procedure}\8\1\\{footnote}\?(\1\?\\{flag}\?:\?\\{sixteen\_bits}\?\2)\?\2;\par
\fi

\M279. \P\?\X245:Procedures in \\{second\_phase}\X\mathrel{+}\S\?\6
\&{procedure}\8\1\\{footnote}\?(\1\?\\{flag}\?:\?\\{sixteen\_bits}\?\2)\?\2\8%
\&{is}\8\C{outputs module cross-references}\1\6
\\{count}\?:\?\p{INTEGER};\C{counts the number of occurences of a
cross-reference}\6
\4\&{begin}\6
\&{if}\1\1\8\\{num}\?\1(\?\\{cur\_xref}\?\2)\L\?\\{flag}\2\8\&{then}\6
\&{return};\2\6
\&{end}\8\2\&{if}\1;\6
\\{finish\_line};\5
\\{out1}\?\1(\?\H{\\}\?\2)\?;\6
\&{if}\1\1\8\\{flag}\?=\?0\2\8\&{then}\6
\\{out1}\?\1(\?\H{U}\?\2)\?;\6
\4\&{else}\6
\\{out1}\?\1(\?\H{A}\?\2)\?;\2\6
\&{end}\8\2\&{if}\1;\6
\\{out4}\?\1(\?\H{\ }\?\2)\1(\?\H{s}\?\2)\1(\?\H{e}\?\2)\1(\?\H{c}\?\2)\?;\5
\\{out4}\?\1(\?\H{t}\?\2)\1(\?\H{i}\?\2)\1(\?\H{o}\?\2)\1(\?\H{n}\?\2)\?;\5
\X280:Output all the module numbers on the reference list \\{cur\_xref}\X;\6
\\{out1}\?\1(\?\H{.}\?\2)\?;\2\6
\&{end}\8\\{footnote};\par
\fi

\M280. The following code distinguishes three cases, according as the number
of cross references is one, two, or more than two.

\Y\P\4\D\\{count\_xref}\?\S\?\\{count}\?\K\?1;\6
\1\&{loop}\6
\?\&{exit}\8\&{when}\8\1(\?\\{num}\?\1(\?\\{xlink}\?\1(\?\\{cur\_xref}\?\2)\2)%
\L\?\\{flag}\?\2)\V\30\1(\?\\{num}\?\1(\?\\{xlink}\?\1(\?\\{cur\_xref}\?\2)\2)%
\I\?\\{num}\?\1(\?\\{cur\_xref}\?\2)\2)\?;\5
\\{cur\_xref}\?\K\?\\{xlink}\?\1(\?\\{cur\_xref}\?\2)\?;\5
\\{incr}\?\1(\?\\{count}\?\2)\?;\2\6
\&{end}\8\&{loop}\par
\Y\P\?\4\X280:Output all the module numbers on the reference list \\{cur\_xref}%
\X\S\?\6
\\{count\_xref};\5
\\{first\_xref}\?\K\?\\{xlink}\?\1(\?\\{cur\_xref}\?\2)\?;\C{\\{first\_xref} is
the first entry of the second xref}\6
\&{if}\1\1\8\\{num}\?\1(\?\\{first\_xref}\?\2)>\?\\{flag}\2\8\&{then}\6
\\{out1}\?\1(\?\H{s}\?\2)\?;\C{plural}\2\6
\&{end}\8\2\&{if}\1;\6
\\{out1}\?\1(\?\H{\~}\?\2)\?;\6
\1\&{loop}\6
\\{out\_mod}\?\1(\?\\{num}\?\1(\?\\{cur\_xref}\?\2)-\?\\{flag}\?,\39\?\\{count}%
\?\2)\?;\5
\\{cur\_xref}\?\K\?\\{xlink}\?\1(\?\\{cur\_xref}\?\2)\?;\C{point to the next
cross reference to output}\6
\?\&{exit}\8\&{when}\?\8\\{num}\?\1(\?\\{cur\_xref}\?\2)\L\?\\{flag};\5
\\{next\_xref}\?\K\?\\{cur\_xref};\C{save \\{cur\_xref} before counting}\6
\\{count\_xref};\6
\&{if}\1\1\?\8\1(\?\\{num}\?\1(\?\\{xlink}\?\1(\?\\{cur\_xref}\?\2)\2)>\?%
\\{flag}\?\2)\V\1(\?\\{next\_xref}\?\I\?\\{first\_xref}\?\2)\?\2\8\&{then}\6
\\{out1}\?\1(\?\H{,}\?\2)\?;\C{not the last of two}\2\6
\&{end}\8\2\&{if}\1;\6
\\{out1}\?\1(\?\H{\ }\?\2)\?;\6
\&{if}\1\1\8\\{num}\?\1(\?\\{xlink}\?\1(\?\\{cur\_xref}\?\2)\2)\L\?\\{flag}\2\8%
\&{then}\6
\\{out4}\?\1(\?\H{a}\?\2)\1(\?\H{n}\?\2)\1(\?\H{d}\?\2)\1(\?\H{\~}\?\2)\?;%
\C{the last}\2\6
\&{end}\8\2\&{if}\1;\2\6
\&{end}\8\&{loop}\par
\U section~279.\fi

\M281. \P\?\X281:Output the code for the end of a module\X\S\?\6
\\{out3}\?\1(\?\H{\\}\?\2)\1(\?\H{f}\?\2)\1(\?\H{i}\?\2)\?;\5
\\{finish\_line};\5
\\{flush\_buffer}\?\1(\?0\?,\39\?\p{FALSE}\?\2)\?\C{insert a blank line, it
looks nice}\par
\U section~264.\fi

\N282.  Phase three processing.
We are nearly finished! \.{AWEAVE}'s only remaining task is to write out the
index, after sorting the identifiers and index entries.

\Y\P\?\4\X282:Phase III: Output the cross-reference index\X\S\?\6
\\{phase\_three}\?\K\?\p{TRUE};\5
\\{print\_nl}\?\1(\?\.{"Writing\ the\ index..."}\?\2)\?;\6
\&{if}\1\1\8\\{change\_exists}\2\8\&{then}\6
\\{finish\_line};\6
\X284:Tell about changed modules\X;\2\6
\&{end}\8\2\&{if}\1;\6
\\{finish\_line};\5
\\{out4}\?\1(\?\H{\\}\?\2)\1(\?\H{i}\?\2)\1(\?\H{n}\?\2)\1(\?\H{x}\?\2)\?;\5
\\{finish\_line};\6
\X286:Do the first pass of sorting\X;\6
\X292:Sort and output the index\X;\6
\\{out4}\?\1(\?\H{\\}\?\2)\1(\?\H{f}\?\2)\1(\?\H{i}\?\2)\1(\?\H{n}\?\2)\?;\5
\\{finish\_line};\6
\X300:Output all the module names\X;\6
\\{out4}\?\1(\?\H{\\}\?\2)\1(\?\H{c}\?\2)\1(\?\H{o}\?\2)\1(\?\H{n}\?\2)\?;\5
\\{finish\_line};\5
\\{print}\?\1(\?\.{"Done."}\?\2)\?\par
\U section~306.\fi

\M283. Just before the index comes a list of all the changed modules, including
the index module itself.

\Y\P\?\4\X14:Types and objects visible from \\{data\_structures}\X\mathrel{+}\S%
\?\6
\\{k\_module}\?:\?\\{module\_index};\C{runs through the modules}\par
\fi

\M284. \P\?\X284:Tell about changed modules\X\S\?\6
\\{k\_module}\?\K\?1;\6
\1\&{while}\?\8\R\?\\{changed\_module}\?\1(\?\\{k\_module}\?\2)\?\8\&{loop}\6
\\{incr}\?\1(\?\\{k\_module}\?\2)\?;\2\6
\&{end}\8\&{loop};\6
\\{out4}\?\1(\?\H{\\}\?\2)\1(\?\H{c}\?\2)\1(\?\H{h}\?\2)\1(\?\H{\ }\?\2)\?;\5
\\{out\_mod}\?\1(\?\\{k\_module}\?\2)\?;\6
\1\&{loop}\6
\1\&{loop}\6
\\{incr}\?\1(\?\\{k\_module}\?\2)\?;\6
\?\&{exit}\8\&{when}\?\8\\{changed\_module}\?\1(\?\\{k\_module}\?\2)\?;\2\6
\&{end}\8\&{loop};\6
\\{out2}\?\1(\?\H{,}\?\2)\1(\?\H{\ }\?\2)\?;\5
\\{out\_mod}\?\1(\?\\{k\_module}\?\2)\?;\6
\?\&{exit}\8\&{when}\?\8\\{k\_module}\?=\?\\{module\_count};\2\6
\&{end}\8\&{loop};\6
\\{out1}\?\1(\?\H{.}\?\2)\?\par
\U section~282.\fi

\M285. A left-to-right radix sorting method is used, since this makes it easy
to
adjust the collating sequence and since the running time will be at worst
proportional to the total length of all entries in the index. We put the
identifiers into 102 different lists based on their first characters.
(Uppercase letters are put into the same list as the corresponding lowercase
letters, since we want to have `$t<\\{TeX}<\&{then}$'.) The
list for character \|c begins at location \\{bucket}\?(\?\|c\?)\? and continues
through
the \\{blink} array.

\Y\P\?\4\X14:Types and objects visible from \\{data\_structures}\X\mathrel{+}\S%
\?\6
\\{bucket}\?:\1\&{array}\8\1(\?\\{ASCII\_code}\?\2)\2\8\&{of}\?\8\\{name%
\_pointer};\6
\\{next\_name}\?:\?\\{name\_pointer};\C{successor of \\{cur\_name} when
sorting}\6
\\{blink}\?:\?\\{name\_info};\C{links in the buckets}\6
\|c\?:\?\\{ASCII\_code};\C{index into bucket}\par
\fi

\M286. To begin the sorting, we go through all the hash lists and put each
entry
having a nonempty cross-reference list into the proper bucket.

\Y\P\?\4\X286:Do the first pass of sorting\X\S\?\6
\1\&{for}\8\|d\?\in\?\p{INTEGER}\8\&{range}\80\?\to\?127\8\&{loop}\6
\\{bucket}\?\1(\?\|d\?\2)\K\?0;\2\6
\&{end}\8\&{loop};\6
\1\&{for}\8\|h\?\in\?\p{INTEGER}\8\&{range}\80\?\to\?\\{hash\_size}\?-\?1\8%
\&{loop}\6
\\{next\_name}\?\K\?\\{hash}\?\1(\?\|h\?\2)\?;\6
\1\&{while}\8\\{next\_name}\?\I\?0\8\&{loop}\6
\\{cur\_name}\?\K\?\\{next\_name};\5
\\{next\_name}\?\K\?\\{link}\?\1(\?\\{cur\_name}\?\2)\?;\6
\&{if}\1\1\8\\{xref}\?\1(\?\\{cur\_name}\?\2)\I\?0\2\8\&{then}\6
\|c\?\K\?\\{byte\_mem}\?\1(\?\\{cur\_name}\?\mathbin{\&{mod}}\?\\{ww}\?,\39\?%
\\{byte\_start}\?\1(\?\\{cur\_name}\?\2)\2)\?;\6
\&{if}\1\1\?\8\1(\?\|c\?\L\?\H{Z}\?\2)\W\1(\?\|c\?\G\?\H{A}\?\2)\?\2\8\&{then}\6
\|c\?\K\?\|c\?+\bn{8}{40}\?;\2\6
\&{end}\8\2\&{if}\1;\6
\\{blink}\?\1(\?\\{cur\_name}\?\2)\K\?\\{bucket}\?\1(\?\|c\?\2)\?;\5
\\{bucket}\?\1(\?\|c\?\2)\K\?\\{cur\_name};\2\6
\&{end}\8\2\&{if}\1;\2\6
\&{end}\8\&{loop};\2\6
\&{end}\8\&{loop}\par
\U section~282.\fi

\M287. During the sorting phase we shall use the \\{cat} and \\{trans} arrays
from
\.{AWEAVE}'s parsing algorithm and rename them \\{depth} and \\{head}. They now
represent a stack of identifier lists for all the index entries that have
not yet been output. The variable \\{sort\_ptr} tells how many such lists are
present; the lists are output in reverse order (first \\{sort\_ptr}, then
\\{sort\_ptr}\?-\?1, etc.). The \|jth list starts at \\{head}\?(\?\|j\?)\?, and
if the first
\|k characters of all entries on this list are known to be equal we have
\\{depth}\?(\?\|j\?)=\?\|k.

\Y\P\?\4\X14:Types and objects visible from \\{data\_structures}\X\mathrel{+}\S%
\?\6
\\{depth}\?:\?\\{cat\_type}\8\&{renames}\8\\{cat};\C{reclaims memory that is no
longer needed for parsing}\6
\\{head}\?:\?\\{trans\_type}\8\&{renames}\8\\{trans};\C{ditto}\6
\\{sort\_ptr}\?:\?\\{scrap\_index}\8\&{renames}\8\\{scrap\_ptr};\C{ditto}\6
\\{max\_sorts}\?:\?\&{constant}\?\8\K\?\\{max\_scraps};\C{ditto}\7
\\{cur\_depth}\?:\?\\{eight\_bits};\C{depth of current buckets}\6
\\{cur\_byte}\?:\?\\{byte\_index};\C{index into \\{byte\_mem}}\6
\\{cur\_bank}\?:\?\\{w\_range};\C{row of \\{byte\_mem}}\6
\\{cur\_val}\?:\?\\{sixteen\_bits};\C{current cross reference number}\6
\&{stat}\1\6
\\{max\_sort\_ptr}\?:\?\p{INTEGER}\8\&{range}\80\?\to\?\\{max\_sorts}\?\K\?0;\2%
\6
\&{tats}\C{largest value of \\{sort\_ptr}}\par
\fi

\M288. The desired alphabetic order is specified by the \\{collate} array;
namely,
\\{collate}\?(\?0\?)<\?\\{collate}\?(\?1\?)<\?\hbox{$\cdots$}\?<\?\\{collate}%
\?(\?100\?)\?.

\Y\P\?\4\X14:Types and objects visible from \\{data\_structures}\X\mathrel{+}\S%
\?\6
\\{collate}\?:\1\&{array}\8\1(\?0\?\to\?100\?\2)\2\8\&{of}\?\8\\{ASCII\_code};%
\C{collation order}\par
\fi

\M289. We use the order $\hbox{null}<\.\ <\hbox{other characters}<\.\_<
\.A=\.a<\cdots<\.Z=\.z<\.0<\cdots<\.9.$

\Y\P\?\4\X19:Set initial values\X\mathrel{+}\S\?\6
\\{collate}\?\1(\?0\?\2)\K\?0;\5
\\{collate}\?\1(\?1\?\2)\K\?\H{\ };\6
\1\&{for}\8\|d\?\in\?\p{INTEGER}\8\&{range}\81\?\to\?\H{\ }\?-\?1\8\&{loop}\6
\\{collate}\?\1(\?\|d\?+\?1\?\2)\K\?\|d;\2\6
\&{end}\8\&{loop};\6
\1\&{for}\8\|d\?\in\?\p{INTEGER}\8\&{range}\8\H{\ }\?+\?1\?\to\?\H{0}\?-\?1\8%
\&{loop}\6
\\{collate}\?\1(\?\|d\?\2)\K\?\|d;\2\6
\&{end}\8\&{loop};\6
\1\&{for}\8\|d\?\in\?\p{INTEGER}\8\&{range}\8\H{9}\?+\?1\?\to\?\H{A}\?-\?1\8%
\&{loop}\6
\\{collate}\?\1(\?\|d\?-\?10\?\2)\K\?\|d;\2\6
\&{end}\8\&{loop};\6
\1\&{for}\8\|d\?\in\?\p{INTEGER}\8\&{range}\8\H{Z}\?+\?1\?\to\?\H{\_}\?-\?1\8%
\&{loop}\6
\\{collate}\?\1(\?\|d\?-\?36\?\2)\K\?\|d;\2\6
\&{end}\8\&{loop};\6
\\{collate}\?\1(\?\H{\_}\?-\?36\?\2)\K\?\H{\_}\?+\?1;\6
\1\&{for}\8\|d\?\in\?\p{INTEGER}\8\&{range}\8\H{z}\?+\?1\?\to\?126\8\&{loop}\6
\\{collate}\?\1(\?\|d\?-\?63\?\2)\K\?\|d;\2\6
\&{end}\8\&{loop};\6
\\{collate}\?\1(\?64\?\2)\K\?\H{\_};\6
\1\&{for}\8\|d\?\in\?\p{INTEGER}\8\&{range}\8\H{a}\?\to\?\H{z}\8\&{loop}\6
\\{collate}\?\1(\?\|d\?-\?\H{a}\?+\?65\?\2)\K\?\|d;\2\6
\&{end}\8\&{loop};\6
\1\&{for}\8\|d\?\in\?\p{INTEGER}\8\&{range}\8\H{0}\?\to\?\H{9}\8\&{loop}\6
\\{collate}\?\1(\?\|d\?-\?\H{0}\?+\?91\?\2)\K\?\|d;\2\6
\&{end}\8\&{loop};\par
\fi

\M290. Procedure \\{unbucket} goes through the buckets and adds nonempty lists
to the stack, using the collating sequence specified in the \\{collate} array.
The parameter to \\{unbucket} tells the current depth in the buckets.
Any two sequences that agree in their first 255 character positions are
regarded as identical.

\Y\P\?\4\X14:Types and objects visible from \\{data\_structures}\X\mathrel{+}\S%
\?\6
\\{infinity}\?:\?\&{constant}\8\\{eight\_bits}\?\K\?255;\C{$\infty$
(approximately)}\par
\fi

\M291. \P\?\X291:Procedures in \\{third\_phase}\X\S\?\6
\&{procedure}\8\1\\{unbucket}\?(\1\?\|d\?:\?\\{eight\_bits}\?\2)\?\2\8\&{is}\8%
\C{empties buckets having depth \|d}\1\6
\4\&{begin}\6
\1\&{for}\8\|f\?\in\?\&{reverse}\80\?\to\?100\8\&{loop}\6
\&{if}\1\1\8\\{bucket}\?\1(\?\\{collate}\?\1(\?\|f\?\2)\2)>\?0\2\8\&{then}\6
\&{if}\1\1\8\\{sort\_ptr}\?>\?\\{max\_sorts}\2\8\&{then}\6
\\{overflow}\?\1(\?\.{"sorting"}\?\2)\?;\2\6
\&{end}\8\2\&{if}\1;\6
\\{incr}\?\1(\?\\{sort\_ptr}\?\2)\?; \6
\&{stat}\1\6
\&{if}\1\1\8\\{sort\_ptr}\?>\?\\{max\_sort\_ptr}\2\8\&{then}\6
\\{max\_sort\_ptr}\?\K\?\\{sort\_ptr};\2\6
\&{end}\8\2\&{if}\1;\2\6
\&{tats}\6
\&{if}\1\1\8\|f\?=\?0\2\8\&{then}\6
\\{depth}\?\1(\?\\{sort\_ptr}\?\2)\K\?\\{infinity};\6
\4\&{else}\6
\\{depth}\?\1(\?\\{sort\_ptr}\?\2)\K\?\|d;\2\6
\&{end}\8\2\&{if}\1;\6
\\{head}\?\1(\?\\{sort\_ptr}\?\2)\K\?\\{bucket}\?\1(\?\\{collate}\?\1(\?\|f\?%
\2)\2)\?;\5
\\{bucket}\?\1(\?\\{collate}\?\1(\?\|f\?\2)\2)\K\?0;\2\6
\&{end}\8\2\&{if}\1;\2\6
\&{end}\8\&{loop};\2\6
\&{end}\8\\{unbucket};\par
\A sections~299 and~306.
\U section~11\*.\fi

\M292. \P\?\X292:Sort and output the index\X\S\?\6
\\{sort\_ptr}\?\K\?0;\5
\\{unbucket}\?\1(\?1\?\2)\?;\6
\1\&{while}\8\\{sort\_ptr}\?>\?0\8\&{loop}\6
\\{cur\_depth}\?\K\?\\{depth}\?\1(\?\\{sort\_ptr}\?\2)\?;\6
\&{if}\1\1\?\8\1(\?\\{blink}\?\1(\?\\{head}\?\1(\?\\{sort\_ptr}\?\2)\2)=\?0\?%
\2)\V\1(\?\\{cur\_depth}\?=\?\\{infinity}\?\2)\?\2\8\&{then}\6
\X294:Output index entries for the list at \\{sort\_ptr}\X;\6
\4\&{else}\6
\X293:Split the list at \\{sort\_ptr} into further lists\X;\2\6
\&{end}\8\2\&{if}\1;\2\6
\&{end}\8\&{loop}\par
\U section~282.\fi

\M293. \P\?\X293:Split the list at \\{sort\_ptr} into further lists\X\S\?\6
\\{next\_name}\?\K\?\\{head}\?\1(\?\\{sort\_ptr}\?\2)\?;\6
\1\&{loop}\6
\\{cur\_name}\?\K\?\\{next\_name};\5
\\{next\_name}\?\K\?\\{blink}\?\1(\?\\{cur\_name}\?\2)\?;\5
\\{cur\_byte}\?\K\?\\{byte\_start}\?\1(\?\\{cur\_name}\?\2)+\?\\{cur\_depth};\5
\\{cur\_bank}\?\K\?\\{cur\_name}\?\mathbin{\&{mod}}\?\\{ww};\6
\&{if}\1\1\8\\{cur\_byte}\?=\?\\{byte\_start}\?\1(\?\\{cur\_name}\?+\?\\{ww}\?%
\2)\?\2\8\&{then}\6
\|c\?\K\?0;\C{we hit the end of the name}\6
\4\&{else}\6
\|c\?\K\?\\{byte\_mem}\?\1(\?\\{cur\_bank}\?,\39\?\\{cur\_byte}\?\2)\?;\6
\&{if}\1\1\?\8\1(\?\|c\?\L\?\H{Z}\?\2)\W\1(\?\|c\?\G\?\H{A}\?\2)\?\2\8\&{then}\6
\|c\?\K\?\|c\?+\bn{8}{40}\?;\2\6
\&{end}\8\2\&{if}\1;\2\6
\&{end}\8\2\&{if}\1;\6
\\{blink}\?\1(\?\\{cur\_name}\?\2)\K\?\\{bucket}\?\1(\?\|c\?\2)\?;\5
\\{bucket}\?\1(\?\|c\?\2)\K\?\\{cur\_name};\6
\?\&{exit}\8\&{when}\?\8\\{next\_name}\?=\?0;\2\6
\&{end}\8\&{loop};\6
\\{decr}\?\1(\?\\{sort\_ptr}\?\2)\?;\5
\\{unbucket}\?\1(\?\\{cur\_depth}\?+\?1\?\2)\?\par
\U section~292.\fi

\M294. \P\?\X294:Output index entries for the list at \\{sort\_ptr}\X\S\?\6
\\{cur\_name}\?\K\?\\{head}\?\1(\?\\{sort\_ptr}\?\2)\?;\6
\1\&{loop}\6
\\{out2}\?\1(\?\H{\\}\?\2)\1(\?\H{:}\?\2)\?;\6
\X295:Output the name at \\{cur\_name}\X;\6
\X296:Output the cross-references at \\{cur\_name}\X;\6
\\{cur\_name}\?\K\?\\{blink}\?\1(\?\\{cur\_name}\?\2)\?;\6
\?\&{exit}\8\&{when}\?\8\\{cur\_name}\?=\?0;\2\6
\&{end}\8\&{loop};\6
\\{decr}\?\1(\?\\{sort\_ptr}\?\2)\?\par
\U section~292.\fi

\M295. \P\?\X295:Output the name at \\{cur\_name}\X\S\?\6
\&{case}\1\8\\{ilk}\?\1(\?\\{cur\_name}\?\2)\?\1\8\&{is}\8\6
\4\&{when}\8\\{normal}\KR\6
\&{if}\1\1\8\\{length}\?\1(\?\\{cur\_name}\?\2)=\?1\2\8\&{then}\6
\\{out2}\?\1(\?\H{\\}\?\2)\1(\?\H{|}\?\2)\?;\6
\4\&{else}\6
\\{out2}\?\1(\?\H{\\}\?\2)\1(\?\H{\\}\?\2)\?;\2\6
\&{end}\8\2\&{if}\1;\6
\4\&{when}\8\\{roman}\KR\8\&{null};\6
\4\&{when}\8\\{wildcard}\KR\8\\{out2}\?\1(\?\H{\\}\?\2)\1(\?\H{9}\?\2)\?;\6
\4\&{when}\8\\{italic}\KR\8\\{out2}\?\1(\?\H{\\}\?\2)\1(\?\H{\\}\?\2)\?;\6
\4\&{when}\8\\{slanted}\KR\8\\{out2}\?\1(\?\H{\\}\?\2)\1(\?\H{s}\?\2)\?;\6
\4\&{when}\8\\{predefined}\KR\8\\{out2}\?\1(\?\H{\\}\?\2)\1(\?\H{p}\?\2)\?;\6
\4\&{when}\8\\{typewriter}\KR\8\\{out2}\?\1(\?\H{\\}\?\2)\1(\?\H{.}\?\2)\?;\6
\4\&{when}\8\&{others}\8\KR\8\\{out2}\?\1(\?\H{\\}\?\2)\1(\?\H{\&}\?\2)\?;\C{%
\\{bold\_face} or reserved word}\2\6
\4\2\&{end}\8\&{case};\6
\\{out\_name}\?\1(\?\\{cur\_name}\?\2)\?\par
\U section~294.\fi

\M296. Section numbers that are to be underlined are enclosed in
`\.{\\[}$\,\ldots\,$\.]', wheras those to be output in italics are
enclosed in `\.{\\\\\{}$\,\ldots\,$\.\}'. Here we have to consider a
special case: if section numbers $s_1,\ldots,s_n$ have to be output
and there are indices $1\leq i<j\leq n$ such that
$s_i+1=s_{i+1},s_{i+1}+1=s_{i+2},\ldots,s_{j-1}+1=s_j$ and $j=n$ or
$s_j+1\not=s_{j+1}$, we output `$s_i$\.{--}$s_j$'.

\Y\P\?\4\X296:Output the cross-references at \\{cur\_name}\X\S\?\6
\X297:Invert the cross-reference list at \\{cur\_name}, making \\{cur\_xref}
the head\X;\6
\1\&{while}\8\\{cur\_xref}\?\I\?0\8\&{loop}\6
\\{out2}\?\1(\?\H{,}\?\2)\1(\?\H{\ }\?\2)\?;\5
\\{cur\_val}\?\K\?\\{num}\?\1(\?\\{cur\_xref}\?\2)\?;\5
\\{this\_xref}\?\K\?\\{cur\_xref};\6
\&{if}\1\1\8\\{cur\_val}\?<\?\\{def\_flag}\2\8\&{then}\6
\\{next\_xref}\?\K\?\\{xlink}\?\1(\?\\{this\_xref}\?\2)\?;\6
\1\&{while}\?\8\1(\?\\{next\_xref}\?\I\?0\?\2)\ANT\1(\1(\?\\{num}\?\1(\?\\{next%
\_xref}\?\2)<\?\\{def\_flag}\?\2)\W\30\1(\?\\{num}\?\1(\?\\{this\_xref}\?\2)=%
\1(\?\\{num}\?\1(\?\\{next\_xref}\?\2)-\?1\?\2)\2)\2)\?\8\&{loop}\6
\\{this\_xref}\?\K\?\\{next\_xref};\5
\\{next\_xref}\?\K\?\\{xlink}\?\1(\?\\{this\_xref}\?\2)\?;\2\6
\&{end}\8\&{loop};\6
\\{out\_mod}\?\1(\?\\{cur\_val}\?\2)\?;\6
\&{if}\1\1\8\\{cur\_xref}\?\I\?\\{this\_xref}\2\8\&{then}\6
\\{out2}\?\1(\?\H{-}\?\2)\1(\?\H{-}\?\2)\?;\5
\\{out\_mod}\?\1(\?\\{num}\?\1(\?\\{this\_xref}\?\2)\2)\?;\2\6
\&{end}\8\2\&{if}\1;\6
\4\&{elsif}\1\8\\{cur\_val}\?<\?\\{body\_flag}\2\8\&{then}\6
\\{out2}\?\1(\?\H{\\}\?\2)\1(\?\H{[}\?\2)\?;\5
\\{out\_mod}\?\1(\?\\{cur\_val}\?-\?\\{def\_flag}\?\2)\?;\5
\\{out1}\?\1(\?\H{]}\?\2)\?;\6
\4\&{else}\6
\\{out3}\?\1(\?\H{\\}\?\2)\1(\?\H{\\}\?\2)\1(\?\H{\{}\?\2)\?;\5
\\{out\_mod}\?\1(\?\\{cur\_val}\?-\?\\{body\_flag}\?\2)\?;\5
\\{out1}\?\1(\?\H{\}}\?\2)\?;\2\6
\&{end}\8\2\&{if}\1;\6
\\{cur\_xref}\?\K\?\\{xlink}\?\1(\?\\{this\_xref}\?\2)\?;\2\6
\&{end}\8\&{loop};\6
\\{out1}\?\1(\?\H{.}\?\2)\?;\5
\\{finish\_line}\par
\U section~294.\fi

\M297. List inversion is best thought of as popping elements off one stack and
pushing them onto another. In this case \\{cur\_xref} will be the head of
the stack that we push things onto.

\Y\P\?\4\X297:Invert the cross-reference list at \\{cur\_name}, making \\{cur%
\_xref} the head\X\S\?\6
\\{this\_xref}\?\K\?\\{xref}\?\1(\?\\{cur\_name}\?\2)\?;\5
\\{cur\_xref}\?\K\?0;\6
\1\&{loop}\6
\\{next\_xref}\?\K\?\\{xlink}\?\1(\?\\{this\_xref}\?\2)\?;\5
\\{xlink}\?\1(\?\\{this\_xref}\?\2)\K\?\\{cur\_xref};\5
\\{cur\_xref}\?\K\?\\{this\_xref};\5
\\{this\_xref}\?\K\?\\{next\_xref};\6
\?\&{exit}\8\&{when}\?\8\\{this\_xref}\?=\?0;\2\6
\&{end}\8\&{loop}\par
\U section~296.\fi

\M298. The following recursive procedure walks through the tree of module names
and
prints them.

\Y\P\?\4\X298:Specification of procedures visible from \\{third\_phase}\X\S\?\6
\&{procedure}\8\1\\{mod\_print}\?(\1\?\|p\?:\?\\{name\_pointer}\?\2)\?\2;\par
\A section~305.
\U section~11\*.\fi

\M299. \P\?\X291:Procedures in \\{third\_phase}\X\mathrel{+}\S\?\6
\&{procedure}\8\1\\{mod\_print}\?(\1\?\|p\?:\?\\{name\_pointer}\?\2)\?\2\8%
\&{is}\8\C{print all module names in subtree \|p}\1\6
\4\&{begin}\6
\&{if}\1\1\8\|p\?>\?0\2\8\&{then}\6
\\{mod\_print}\?\1(\?\\{llink}\?\1(\?\|p\?\2)\2)\?;\6
\\{out2}\?\1(\?\H{\\}\?\2)\1(\?\H{:}\?\2)\?;\6
\\{tok\_ptr}\?\K\?1;\5
\\{text\_ptr}\?\K\?1;\5
\\{scrap\_ptr}\?\K\?0;\5
\\{init\_stack};\5
\\{app}\?\1(\?\|p\?+\?\\{mod\_flag}\?\2)\?;\5
\\{make\_output};\5
\\{footnote}\?\1(\?0\?\2)\?;\C{\\{cur\_xref} was set by \\{make\_output}}\6
\\{finish\_line};\6
\\{mod\_print}\?\1(\?\\{rlink}\?\1(\?\|p\?\2)\2)\?;\2\6
\&{end}\8\2\&{if}\1;\2\6
\&{end}\8\\{mod\_print};\par
\fi

\M300. \P\?\X300:Output all the module names\X\S\?\8\\{mod\_print}\?\1(\?%
\\{root}\?\2)\?\par
\U section~282.\fi

\N301.  The main program.
Let's put it all together now: \.{AWEAVE} starts and ends here.

The main procedure has been split into three sub-procedures in order to
keep certain \ADA\ compilers from overflowing their capacity.

\Y\P\?\4\X128:Specification of procedures visible from \\{first\_phase}\X%
\mathrel{+}\S\?\6
\&{procedure}\8\1\\{Phase\_I}\2;\par
\fi

\M302. \P\?\X129:Procedures in \\{first\_phase}\X\mathrel{+}\S\?\6
\&{procedure}\8\1\\{Phase\_I}\2\8\&{is}\8\1\6
\4\&{begin}\6
\X126:Phase I: Read all the user's text and store the cross references\X;\2\6
\&{end}\8\\{Phase\_I};\par
\fi

\M303. \P\?\X244:Specification of procedures visible from \\{second\_phase}\X%
\mathrel{+}\S\?\6
\&{procedure}\8\1\\{Phase\_II}\2;\par
\fi

\M304. \P\?\X245:Procedures in \\{second\_phase}\X\mathrel{+}\S\?\6
\&{procedure}\8\1\\{Phase\_II}\2\8\&{is}\8\1\6
\4\&{begin}\6
\X262:Phase II: Read all the text again and translate it to \TeX\ form\X;\2\6
\&{end}\8\\{Phase\_II};\par
\fi

\M305. \P\?\X298:Specification of procedures visible from \\{third\_phase}\X%
\mathrel{+}\S\?\6
\&{procedure}\8\1\\{Phase\_III}\2;\par
\fi

\M306. \P\?\X291:Procedures in \\{third\_phase}\X\mathrel{+}\S\?\6
\&{procedure}\8\1\\{Phase\_III}\2\8\&{is}\8\1\6
\4\&{begin}\6
\X282:Phase III: Output the cross-reference index\X;\2\6
\&{end}\8\\{Phase\_III};\par
\fi

\M307.
\Y\P\4\D\\{aweb\_input\_file\_name}\?\S\?\.{"input.web"}\?\6
\?\par
\P\4\D\\{change\_input\_file\_name}\?\S\?\.{"input.ch"}\?\6
\?\par
\P\4\D\\{TeX\_output\_file\_name}\?\S\?\.{"aweb\_output.tex"}\?\6
\?\par
\Y\P\?\4\X14:Types and objects visible from \\{data\_structures}\X\mathrel{+}\S%
\?\6
\\{opened}\?:\?\p{BOOLEAN}\?\K\?\p{FALSE};\par
\fi

\M308\*.
\Y\P\O{aweave.adb}\?\6
\?\&{with}\8\\{Ada}\?.\?\\{Command\_Line};\6
\&{use}\8\\{Ada}\?.\?\\{Command\_Line};\6
\&{with}\8\\{Ustrings};\6
\&{use}\8\\{Ustrings};\6
\&{with}\8\p{TEXT\_IO}\?,\39\?\\{int\_number\_io};\6
\&{with}\8\\{data\_structures}\?,\39\?\\{searching}\?,\39\?\\{input\_output}\?,%
\39\?\\{first\_phase}\?,\39\?\\{second\_phase}\?,\39\?\\{third\_phase};\6
\&{use}\8\\{data\_structures}\?,\39\?\\{searching}\?,\39\?\\{input\_output}\?,%
\39\?\\{first\_phase}\?,\39\?\\{second\_phase}\?,\39\?\\{third\_phase};\6
\&{with}\8\\{Ada}\?.\?\\{Strings}\?.\?\\{Unbounded};\6
\&{use}\8\\{Ada}\?.\?\\{Strings}\?.\?\\{Unbounded};\6
\&{procedure}\8\1\\{AWEAVE}\2\8\&{is}\8\1\6
\\{ac}\?:\?\p{INTEGER}\?\K\?0;\6
\\{ic}\?:\?\p{INTEGER}\?\K\?0;\6
\\{yr\_output\_file}\?:\?\\{Ustring};\6
\4\&{begin}\6
\\{print\_ln}\?\1(\?\.{"Trying\ to\ open\ input\ file."}\?\2)\?;\5
\\{ac}\?\K\?\\{Argument\_Count};\6
\&{if}\1\1\8\\{ac}\?<\?1\2\8\&{then}\6
\\{fatal\_error}\?\1(\?\.{"No\ file\ on\ command\ line"}\?\2)\?;\2\6
\&{end}\8\2\&{if}\1;\6
\\{open\_a\_file}\?\1(\?\\{aweb\_file}\?,\39\?\p{TEXT\_IO}\?.\?\p{IN\_FILE}\?,%
\39\?\\{Argument}\?\1(\?1\?\2),\39\?\\{opened}\?\2)\?;\6
\&{if}\1\1\?\8\R\?\\{opened}\2\8\&{then}\6
\\{fatal\_error}\?\1(\?\.{"AWEB\ file\ not\ open"}\?\2)\?;\2\6
\&{end}\8\2\&{if}\1;\6
\\{opened}\?\K\?\p{FALSE};\6
\&{if}\1\1\8\\{ac}\?\G\?2\2\8\&{then}\6
\\{open\_a\_file}\?\1(\?\\{change\_file}\?,\39\?\p{TEXT\_IO}\?.\?\p{IN\_FILE}%
\?,\39\?\\{Argument}\?\1(\?2\?\2),\39\?\\{opened}\?\2)\?;\2\6
\&{end}\8\2\&{if}\1;\6
\&{if}\1\1\?\8\R\?\\{opened}\2\8\&{then}\6
\\{new\_ln};\5
\\{print\_ln}\?\1(\?\.{"!\ No\ change\ file."}\?\2)\?;\2\6
\&{end}\8\2\&{if}\1;\6
\\{ic}\?\K\?\\{Index}\?\1(\?\|U\?\1(\?\\{Argument}\?\1(\?1\?\2)\2),\39\?\.{"."}%
\?\2)\?;\6
\&{if}\1\1\8\\{ic}\?=\?0\2\8\&{then}\6
\\{yr\_output\_file}\?\K\?\|U\?\1(\?\\{Argument}\?\1(\?1\?\2)\2)\mathbin{\.\&}%
\?\.{".tex"};\6
\4\&{else}\6
\\{yr\_output\_file}\?\K\?\|U\?\1(\?\\{Argument}\?\1(\?1\?\2)\2)\?;\5
\\{Overwrite}\?\1(\?\\{yr\_output\_file}\?,\39\?\\{ic}\?,\39\?\.{".tex"}\?\2)%
\?;\5
\\{Ada}\?.\?\\{Strings}\?.\?\\{Unbounded}\?.\?\\{head}\?\1(\?\\{yr\_output%
\_file}\?,\39\?\\{ic}\?+\?3\?\2)\?;\2\6
\&{end}\8\2\&{if}\1;\6
\?--\?\\{open\_a\_file}\?\1(\?\\{TeX\_file}\?,\39\?\p{TEXT\_IO}\?.\?\p{OUT%
\_FILE}\?,\39\?\\{TeX\_output\_file\_name}\?,\39\?\\{opened}\?\2)\?;\5
\\{open\_a\_file}\?\1(\?\\{TeX\_file}\?,\39\?\p{TEXT\_IO}\?.\?\p{OUT\_FILE}\?,%
\39\?\|S\?\1(\?\\{yr\_output\_file}\?\2),\39\?\\{opened}\?\2)\?;\6
\&{if}\1\1\?\8\1(\R\?\\{opened}\?\2)\ANT\1(\R\?\\{create\_TeX\_file}\?\1(\?\|S%
\?\1(\?\\{yr\_output\_file}\?\2)\2)\2)\?\2\8\&{then}\6
\\{fatal\_error}\?\1(\?\.{"TeX\ file\ not\ open"}\?\2)\?;\2\6
\&{end}\8\2\&{if}\1;\6
\\{initialize};\5
\\{print\_ln}\?\1(\?\\{banner}\?\2)\?;\C{print a ``banner line''}\6
\X67:Store all the reserved words\X;\6
\X69:Store all the predefined identifiers\X;\6
\\{Phase\_I};\5
\\{Phase\_II};\5
\\{Phase\_III};\6
\X98:Check that all changes have been read\X;\6
\&{raise}\8\\{end\_of\_AWEAVE}; \4\&{exception} \6
\4\&{when}\8\\{end\_of\_AWEAVE}\KR\6
\&{stat}\1\6
\X309:Print statistics about memory usage\X;\2\6
\&{tats}\6
\hbox{\4\4}\?\C{here files should be closed if the operating system requires
it}\6
\?\6
\X310:Print the job \\{history}\X;\6
\&{return};\5
\hbox{\2}\2\6
\&{end}\8\\{AWEAVE};\par
\fi

\M309. \P\?\X309:Print statistics about memory usage\X\S\?\6
\\{print\_nl}\?\1(\?\.{"Memory\ usage\ statistics:\ "}\?\2)\?;\5
\\{print\_int}\?\1(\?\\{name\_ptr}\?,\39\?1\?\2)\?;\5
\\{print}\?\1(\?\.{"\ names,\ "}\?\2)\?;\5
\\{print\_int}\?\1(\?\\{xref\_ptr}\?,\39\?1\?\2)\?;\5
\\{print}\?\1(\?\.{"\ cross\ references,\ "}\?\2)\?;\5
\\{print\_int}\?\1(\?\\{byte\_ptr}\?\1(\?0\?\2),\39\?1\?\2)\?;\6
\1\&{for}\8\|i\?\in\?1\?\to\?\\{ww}\?-\?1\8\&{loop}\6
\\{print}\?\1(\?\.{"+"}\?\2)\?;\5
\\{print\_int}\?\1(\?\\{byte\_ptr}\?\1(\?\|i\?\2),\39\?1\?\2)\?;\2\6
\&{end}\8\&{loop};\6
\\{print}\?\1(\?\.{"\ bytes;"}\?\2)\?;\5
\\{print\_nl}\?\1(\?\.{"parsing\ required\ "}\?\2)\?;\5
\\{print\_int}\?\1(\?\\{max\_scr\_ptr}\?,\39\?1\?\2)\?;\5
\\{print}\?\1(\?\.{"\ scraps,\ "}\?\2)\?;\5
\\{print\_int}\?\1(\?\\{max\_txt\_ptr}\?,\39\?1\?\2)\?;\5
\\{print}\?\1(\?\.{"\ texts,\ "}\?\2)\?;\5
\\{print\_int}\?\1(\?\\{max\_tok\_ptr}\?,\39\?1\?\2)\?;\5
\\{print}\?\1(\?\.{"\ tokens,\ "}\?\2)\?;\5
\\{print\_int}\?\1(\?\\{max\_stack\_ptr}\?,\39\?1\?\2)\?;\5
\\{print}\?\1(\?\.{"\ levels;"}\?\2)\?;\5
\\{print\_nl}\?\1(\?\.{"sorting\ required\ "}\?\2)\?;\5
\\{print\_int}\?\1(\?\\{max\_sort\_ptr}\?,\39\?1\?\2)\?;\5
\\{print}\?\1(\?\.{"\ levels."}\?\2)\?\par
\U section~308\*.\fi

\M310. Some implementations may wish to pass the \\{history} value to the
operating system so that it can be used to govern whether or not other
programs are started. Here we simply report the history to the user.

\Y\P\?\4\X310:Print the job \\{history}\X\S\?\6
\&{case}\1\8\\{history}\1\8\&{is}\8\6
\4\&{when}\8\\{spotless}\KR\8\\{print\_nl}\?\1(\?\.{"(No\ errors\ were\
found.)"}\?\2)\?;\6
\4\&{when}\8\\{harmless\_message}\KR\8\\{print\_nl}\?\1(\?\.{"(Did\ you\ see\
the\ warning\ message\ above?)"}\?\2)\?;\6
\4\&{when}\8\\{error\_message}\KR\8\\{print\_nl}\?\1(\?\.{"(Pardon\ me,\ but\ I%
\ think\ I\ spotted\ something\ wrong.)"}\?\2)\?;\6
\4\&{when}\8\\{fatal\_message}\KR\8\\{print\_nl}\?\1(\?\.{"(That\ was\ a\ fatal%
\ error,\ my\ friend.)"}\?\2)\?;\2\6
\4\2\&{end}\8\&{case};\C{there are no other cases}\6
\\{new\_ln}\par
\U section~308\*.\fi

\N311.  System-dependent changes.
This module should be replaced, if necessary, by changes to the program
that are necessary to make \.{AWEAVE} work at a particular installation.
It is usually best to design your change file so that all changes to
previous modules preserve the module numbering; then everybody's version
will be consistent with the printed program. More extensive changes,
which introduce new modules, can be inserted here; then only the index
itself will get a new module number.

\fi

\N312\*.  Index.
If you have read and understood the code for Phase III above, you know what
is in this index and how it got here. All modules in which an identifier is
used are listed with that identifier, except that reserved words are
indexed only when they appear in format definitions, and the appearances
of identifiers in module names are not indexed. Underlined entries
correspond to where the identifier was declared. Error messages, control
sequences put into the output, and a few
other things like ``recursion'' are indexed here too.
\fi


\ch 3\*, 4\*, 5\*, 6\*, 7\*, 8\*, 9\*, 10\*, 11\*, 176\*, 219\*, 223\*, 308\*,
312\*.
\inx
\:\.{(...)}, 152.
\:\.{**}, 20, 112, 169.
\:\.{+=}, 273.
\:\.{--}, 296.
\:\.{..}, 112, 169.
\:\.{/=}, 20, 112, 169.
\:\.{:=}, 20, 112, 169.
\:\.{<<}, 20, 112, 169.
\:\.{<=}, 20, 112, 169.
\:\.{<>}, 20, 112, 169.
\:\.{==}, 20, 112, 169.
\:\.{=>}, 20, 112, 169.
\:\.{>=}, 20, 112, 169.
\:\.{>>}, 20, 112, 169.
\:\.{\$}, 169.
\:\.{\%}, 169.
\:\.{\&}, 169, 226.
\:\.{\\'}, 232.
\:\.{\\)}, 226.
\:\.{\\*}, 152.
\:\.{\\,}, 202, 226.
\:\.{\\.}, 168--169, 231, 295.
\:\.{\\:}, 294, 299.
\:\.{\\=}, 169, 231.
\:\.{\\?}, 168, 252.
\:\.{\\[}, 296.
\:\.{\\\ }, 232.
\:\.{\\\#}, 169, 226, 232.
\:\.{\\\$}, 169, 226, 232.
\:\.{\\\%}, 169, 226, 232.
\:\.{\\\&}, 168--169, 185, 232, 253, 295.
\:\.{\\\\}, 168--169, 232, 253, 295--296.
\:\.{\\\^}, 169, 226, 232.
\:\.{\\\{}, 232.
\:\.{\\\}}, 232.
\:\.{\\\~}, 232.
\:\.{\\`}, 232.
\:\.{\\|}, 253, 295.
\:\.{\\\_}, 155, 169, 226, 228, 232.
\:\.{\\A}, 279.
\:\.{\\ANT}, 178.
\:\.{\\ast}, 226.
\:\.{\\B}, 226.
\:\.{\\BL}, 169, 230.
\:\.{\\bn}, 180.
\:\.{\\BOX}, 169, 230.
\:\.{\\C}, 168, 240.
\:\.{\\ch}, 284.
\:\.{\\con}, 282.
\:\.{\\D}, 269.
\:\.{\\dst}, 169, 230.
\:\.{\\E}, 169, 226.
\:\.{\\EL}, 169, 230.
\:\.{\\F}, 270.
\:\.{\\fi}, 281.
\:\.{\\fin}, 282.
\:\.{\\G}, 169, 230.
\:\.{\\H}, 233.
\:\.{\\hbox}, 234.
\:\.{\\I}, 169, 230.
\:\.{\\in}, 226, 235.
\:\.{\\input awebmac}, 144.
\:\.{\\inx}, 282.
\:\.{\\J}, 226.
\:\.{\\K}, 169, 230.
\:\.{\\KR}, 169, 230.
\:\.{\\L}, 169, 230.
\:\.{\\M}, 265.
\:\.{\\mathbin}, 168, 254.
\:\.{\\mathop}, 168, 254.
\:\.{\\mathrel}, 168, 254.
\:\.{\\N}, 265.
\:\.{\\notin}, 201.
\:\.{\\O}, 169, 231.
\:\.{\\ORE}, 188.
\:\.{\\ORW}, 188.
\:\.{\\P}, 256, 268.
\:\.{\\p}, 168--169, 253, 295.
\:\.{\\par}, 268.
\:\.{\\R}, 230, 235.
\:\.{\\S}, 169, 230, 270, 273.
\:\.{\\s}, 295.
\:\.{\\T}, 226.
\:\.{\\to}, 226.
\:\.{\\U}, 279.
\:\.{\\V}, 198, 230.
\:\.{\\W}, 230, 235.
\:\.{\\X}, 168, 257.
\:\.{\\Y}, 256, 263, 268, 273.
\:\.{\\1}, 168, 255--256.
\:\.{\\2}, 168, 255--256.
\:\.{\\3}, 168, 255--256.
\:\.{\\4}, 168, 255--256.
\:\.{\\5}, 168, 255--256.
\:\.{\\6}, 168, 255--256, 268.
\:\.{\\7}, 168, 255--256, 268.
\:\.{\\8}, 168--169, 226, 252.
\:\.{\\9}, 295.
\:\.{\^}, 169.
\:\.{\AT!!}, 51, 101, 109, 133.
\:\.{\AT!'}, 101, 169.
\:\.{\AT!*}, 81, 99, 101, 265.
\:\.{\AT!+}, 99, 101, 169, 273.
\:\.{\AT!,}, 99, 101, 169.
\:\.{\AT!.}, 45, 99, 101, 132.
\:\.{\AT!/}, 99, 101, 169.
\:\.{\AT!:}, 45, 99, 101, 132.
\:\.{\AT!;}, 99, 101, 169.
\:\.{\AT!<}, 99, 101, 169.
\:\.{\AT!=}, 101, 169.
\:\.{\AT!>}, 101.
\:\.{\AT!?}, 101, 109.
\:\.{\AT!\ }, 81, 99, 101, 265.
\:\.{\AT!\#}, 99, 101, 169.
\:\.{\AT!\&}, 99, 101, 169.
\:\.{\AT!\\}, 101, 169.
\:\.{\AT!\^}, 45, 99, 101, 132.
\:\.{\AT!\{}, 101, 169.
\:\.{\AT!\}}, 101, 169.
\:\.{\AT!\~}, 101, 169.
\:\.{\AT!\AT!}, 101, 156.
\:\.{\AT!|}, 99, 101, 169.
\:\.{\AT!a}, 99, 101, 132.
\:\.{\AT!b}, 45, 99, 101, 132.
\:\.{\AT!d}, 51, 99, 101.
\:\.{\AT!f}, 51, 99, 101.
\:\.{\AT!i}, 45, 99, 101, 132.
\:\.{\AT!p}, 99, 101, 132.
\:\.{\AT!r}, 99, 101, 132.
\:\.{\AT!s}, 45, 99, 101.
\:\.{\AT!t}, 99, 101, 169.
\:\.{\AT!x}, 87--88.
\:\.{\AT!y}, 87, 90--91.
\:\.{\AT!z}, 87, 97.
\:\|{a}, \[152], \[248], \[251].
\:\.{abort}, 169.
\:\&{abort}, 45, \[67], 167.
\:\.{abs}, 169.
\:\&{abs}, 45, \[67], 167.
\:\\{ac}, 308\*.
\:\&{accept}, 45, \[67], 128, 167.
\:\.{accept}, 169.
\:\\{accept\_like}, \[45], 67, 129, 235.
\:\\{accept\_scrap}, \[167], 169--170, 176\*, 235.
\:\&{access}, 45, \[67], 167.
\:\.{access}, 169.
\:\\{Ada}, 308\*.
\:{ADA attributes}, 18.
\:\.{ADA text...didn't end}, 260.
\:\\{ADA\_parse}, 221, \[222], \\{223\*}, 226, 238--240.
\:\\{ADA\_translate}, \[237], \\{238}, 240, 250.
\:\\{ADA\_xref}, \[128], \\{129}, 130, 132, 221, 239.
\:\p{ADDRESS}, \[69].
\:\p{AFT}, \[69].
\:\.{all}, 169.
\:\&{all}, 45, \[67], 167.
\:\\{all\_like}, \[45], 67--68, 235.
\:\.{Ambiguous prefix}, 80.
\:\.{and}, 20, 169.
\:\&{and}, 45, \[67], 167.
\:\\{and\_like}, \[45], 67, 235.
\:\\{and\_scrap}, \[167], 169--170, 176\*, 198, 230, 235.
\:\\{and\_sign}, \[20], 230.
\:\\{app}, \[173], 177--178, 180--186, 188--194, 196, 198--199, 201--202,
204--210, 220, 224, 226, 231--234, 236, 240, 250--251, 268, 273, 299.
\:\\{app\_comment}, \[221], 224, \\{236}, 238, 240, 268.
\:\\{app\_tok}, \[162], 163--165, 173, 231--232, 234, 238, 268.
\:\\{append\_xref}, \[53], 54--55.
\:\\{app1}, \[173], 177, 180--187, 189--190, 192--194, 196, 198, 202, 204--210,
214, 220, 236.
\:\\{app2}, \[173], 186, 189, 191, 199, 202.
\:\\{app3}, \[173], 191, 193, 198, 202.
\:\\{Argument}, 308\*.
\:\\{Argument\_Count}, 308\*.
\:\.{array}, 169.
\:\&{array}, 45, \[67], 167.
\:\\{array\_like}, \[45], 67, 235.
\:\\{arrow}, \[167], 170, 176\*, 209, 230.
\:\p{ASCII}, \[69].
\:{ASCII code}, 16, 99.
\:\\{ASCII\_code}, \[16], 18, 20, 31--32, 42, 73, 100--101, 103, 107, 139, 148,
157, 163, 251, 285, 288.
\:\\{ascii\_tok}, \[99], 101, 109, 115, 225.
\:\&{at}, 45, \[67], 167.
\:\.{at}, 169.
\:\\{at\_like}, \[45], 67, 235.
\:\\{AWEAVE}, \\{308\*}.
\:\.{aweave.adb}, 308\*.
\:\.{AWEB file ended...}, 90.
\:\.{AWEB file not open}, 308\*.
\:\\{aweb\_file}, \[26], 37, 82, 84, 90, 93, 96, 98, 308\*.
\:\\{aweb\_input\_file\_name}, \[307].
\:\.{aweb\_output.tex}, 307.
\:\|{b}, \[107], \[141], \[163], \[251].
\:\\{backup}, \[168], 169--170, 184, 186, 189--190, 196, 198, 204, 208, 235,
251--252, 269--270, 273.
\:\\{bal}, 106, \[107], \[130], 162, \[163], \[240].
\:\\{banner}, \[1], 308\*.
\:\p{BASE}, \[69].
\:\\{base}, \[167], 169--170, 176\*, 180, 187, 227.
\:\.{begin}, 169.
\:\&{begin}, 45, \[67], 167.
\:\\{begin\_ada}, \[99], 101, 135, 271--272.
\:\\{begin\_comment}, \[99], 101, 226, 266.
\:\\{begin\_label\_scrap}, \[167], 170, 176\*, 230.
\:\\{begin\_like}, \[45], 67--68, 235.
\:\\{beginning}, 166, \[167], 169--170, 176\*--177, 186, 189--190, 197--198,
235.
\:\\{big\_cancel}, \[168], 169, 226, 251--252, 256.
\:\\{big\_force}, \[168], 169, 226, 251, 256, 268.
\:\\{big\_line\_break}, \[99], 101, 226, 266.
\:\\{blink}, \[285], 286, 292--294.
\:\&{body}, 45, \[67], 128, 167.
\:\.{body}, 169.
\:\\{body\_flag}, \[51], 54, 109, 115, 128--129, 132, 296.
\:\\{body\_like}, \[45], 67, 129, 235.
\:\\{bold\_face}, \[45], 295.
\:\p{BOOLEAN}, 27--28, 30, 32--34, 50, \[69], 82, 85, 109, 140--141, 251, 307.
\:\\{box\_sign}, \[20], 112, 230.
\:\\{break\_out}, 145--146, \[147], \\{148}.
\:\\{break\_space}, \[168], 170, 208, 242, 255--256.
\:\\{bucket}, \[285], 286, 291, 293.
\:\\{buf\_size}, \[14], 31--33.
\:\\{buffer}, \[31], 32--33, 36--37, 59, 61--62, 64--66, 85, 87, 89--91, 93,
95--98, 103, 105, 107--109, 111--115, 118--119, 121--124, 127, 143, 157--158,
160--161, 163--164, 219\*, 223\*, 226, 232, 234, 251, 258, 260--261, 265.
\:\\{buffer\_index}, \[31], 33, 36, 59, 61, 82, 84.
\:\\{buffer\_type}, \[31], 84.
\:\\{byte\_index}, \[42], 43, 61, 76, 80, 251, 287.
\:\\{byte\_mem}, 41, \[42], 43, 46, 49, 56, 61, 64--65, 76--78, 80, 155, 251,
253, 258--261, 286--287, 293.
\:\\{byte\_memory}, \[42].
\:\\{byte\_ptr}, \[43], 65, 77, 309.
\:\\{byte\_start}, 41, \[42], 43--44, 49, 54, 59, 64--65, 77--78, 109, 131,
155, 253, 258, 286, 293.
\:\|{c}, \[33], \[61], \[76], \[80], \[101], \[103], \[105], \[107], \[111], %
\[148], \[152], \[157], \[160], \[163], \[212], \[214], \[285].
\:\p{CALENDAR}, \[69].
\:\p{CALLABLE}, \[69].
\:\\{cancel}, \[168], 169--170, 186, 196, 198, 205, 208, 238, 242, 251,
255--256.
\:\\{carriage\_return}, \[20], 21, 32.
\:\&{case}, 45, \[67], 167.
\:\.{case}, 169.
\:\\{case\_like}, \[45], 67, 235.
\:\\{case\_scrap}, \[167], 169--170, 176\*, 189, 235.
\:\\{cat}, \[171], 174, 176\*--178, 180--189, 191--194, 196, 198--199,
201--202, 204--210, 212, 214, 216--217, 219\*--221, 224, 237, 287.
\:\\{cat\_type}, \[171], 287.
\:\.{Change file ended...}, 88, 90, 97.
\:\.{Change file entry did not match}, 98.
\:\\{change\_buffer}, \[84], 85--86, 89--90, 96, 98.
\:\\{change\_changing}, \[83], 90, 93, 97.
\:\\{change\_exists}, \[50], 126--127, 282.
\:\\{change\_file}, \[26], 37, 82, 84, 87--88, 90, 93, 97, 308\*.
\:\\{change\_input\_file\_name}, \[307].
\:\\{change\_limit}, \[84], 85--86, 89--90, 96, 98.
\:\\{changed\_module}, \[50], 95, 126--127, 152, 284.
\:\\{changing}, 37, \[82], 83--84, 86, 90, 93, 95, 98.
\:\p{CHARACTER}, 17, 19, \[69].
\:\\{character\_tok}, \[99], 109, 123, 225.
\:\\{check\_change}, \\{90}, 96.
\:\\{clause}, \[167], 170, 176\*, 209.
\:\p{CLOCK}, \[69].
\:\\{close\_scrap}, \[167], 169--170, 202, 220, 226, 235.
\:\p{COL}, \[69].
\:\\{collate}, \[288], 289--291.
\:\\{colon}, \[167], 169--170, 176\*, 185, 198, 202, 210, 226, 235.
\:\\{Command\_Line}, 308\*.
\:\\{comment}, \[169].
\:\\{comment\_scrap}, 217, \[224], 226.
\:\\{compare}, \[74], 76, 80.
\:\\{compress}, \[112].
\:\\{compress\_end}, \[189].
\:\\{confusion}, \[39].
\:\&{constant}, 45, \[67], 167.
\:\.{constant}, 169.
\:\p{CONSTRAINED}, \[69].
\:\p{CONSTRAINT\_ERROR}, \[69].
\:\\{continue}, 87.
\:\.{Control codes are forbidden...}, 121.
\:\.{Control text didn't end}, 121.
\:\\{control\_code}, \[100], \\{101}, 105, 109, 115, 161.
\:\p{CONTROLLED}, \[69].
\:\\{copy\_comment}, 156, \[162], \\{163}, 240.
\:\\{copy\_limbo}, \[156], \\{157}, 262.
\:\\{copy\_TeX}, 156, \[159], \\{160}, 266.
\:\\{count}, \[80], \[279], 280.
\:\p{COUNT}, \[69].
\:\\{count\_xref}, \[280].
\:\p{CREATE}, 30, \[69].
\:\\{create\_TeX\_file}, \[27], \\{30}, 308\*.
\:\\{cur\_bank}, \[287], 293.
\:\\{cur\_byte}, \[287], 293.
\:\\{cur\_depth}, \[287], 292--293.
\:\\{cur\_end}, \[243], 245--246, 248.
\:\\{cur\_mod\_name}, \[251], 258--260.
\:\\{cur\_mode}, \[243], 244--245, 248, 251--252, 255--256.
\:\\{cur\_module}, \[109], 116, 135, 272, 274.
\:\\{cur\_name}, \[66], 247--248, 253, 257--258, 285--286, 293--295, 297.
\:\\{cur\_state}, \[243], 245--246.
\:\\{cur\_tok}, \[243], 245--246, 248.
\:\\{cur\_val}, \[287], 296.
\:\\{cur\_xref}, \[136], 137, 257, 273, 276--280, 296--297, 299.
\:\p{CURRENT\_INPUT}, \[69].
\:\p{CURRENT\_OUTPUT}, \[69].
\:\|{d}, \[111], \[148], \[212], \[214], \[251], \[286], \[289], \[291].
\:\p{DATA}, \[70].
\:\p{DATA\_ERROR}, \[70].
\:\\{data\_structures}, \\{3\*}, 4\*--11\*, 308\*.
\:\.{data\_structures.adb}, 3\*.
\:\.{data\_structures.ads}, 3\*.
\:\p{DATE}, \[70].
\:\p{DAY}, \[70].
\:\p{DAY\_DURATION}, \[70].
\:\\{decl}, 166, \[167], 169--170, 176\*, 183--184, 186, 192--193, 198, 202,
204--205, 207, 210, 235.
\:\.{declare}, 169.
\:\&{declare}, 12, 45, \[67], 167.
\:\\{declare\_like}, \[45], 67, 235.
\:\\{decr}, \[13], 33, 108, 113, 118, 141, 148, 153, 161, 164--165, 246,
293--294.
\:\\{def\_flag}, \[51], 54--55, 109, 115, 128--129, 132--133, 135, 137, 151,
257, 273, 275, 277--278, 296.
\:\p{DEFAULT\_AFT}, \[70].
\:\p{DEFAULT\_BASE}, \[70].
\:\p{DEFAULT\_EXP}, \[70].
\:\p{DEFAULT\_FORE}, \[70].
\:\\{default\_name}, 27, \[28].
\:\p{DEFAULT\_SETTING}, \[70].
\:\p{DEFAULT\_WIDTH}, \[70].
\:\\{defined}, \[45], 63.
\:\\{definition}, \[99], 101, 133, 267.
\:\.{delay}, 169.
\:\&{delay}, 45, \[67], 167.
\:\p{DELETE}, \[70].
\:\p{DELTA}, \[70].
\:\.{delta}, 169.
\:\&{delta}, 45, \[67], 167.
\:\\{depth}, \[287], 291--292.
\:\p{DEVICE}, \[70].
\:\p{DEVICE\_ERROR}, \[70].
\:\\{dig}, \[150], 152--153.
\:\\{digit\_scrap}, \[167], 169--170, 176\*, 187, 191, 226, 228.
\:\p{DIGITS}, \[70].
\:\&{digits}, 45, \[67], 167.
\:\.{digits}, 169.
\:\p{DIRECT\_IO}, \[70].
\:\&{do}, 45, \[67], 167.
\:\.{do}, 169.
\:\\{do\_cases}, 174, \[175], \\{176\*}.
\:\.{Done.}, 282.
\:\.{Double \AT! required...}, 158.
\:\.{Double \AT! should be used...}, 232.
\:\\{double\_dot}, \[99], 112, 226.
\:\\{double\_star}, \[20], 112, 230.
\:\p{DURATION}, \[70].
\:\\{easy\_cases}, 221, 225, \\{226}.
\:\\{eight\_bits}, \[41], 60--61, 99--101, 104--107, 109--111, 125, 130,
140--141, 159--160, 162--163, 167--168, 171, 212, 214, 240, 247--248, 251, 287,
290--291.
\:\p{ELABORATE}, \[70].
\:\p{ELEMENT\_TYPE}, \[70].
\:\&{else}, 45, \[67], 167.
\:\.{else}, 169.
\:\\{else\_like}, \[45], 67, 235.
\:\\{else\_scrap}, \[167], 169--170, 176\*, 188, 198, 208, 230, 235.
\:\.{elsif}, 169.
\:\&{elsif}, 45, \[67], 167.
\:\\{elsif\_like}, \[45], 67, 235.
\:\p{EMAX}, \[70].
\:\\{emit\_space\_if\_needed}, \[263], 267, 272.
\:\.{end}, 169.
\:\&{end}, 12, 45, \[67], 167.
\:\\{end\_comment}, \[99], 101, 226, 266.
\:\p{END\_ERROR}, \[70].
\:\\{end\_field}, \[243].
\:\\{end\_label\_scrap}, \[167], 170, 181, 230.
\:\\{end\_like}, \[45], 67, 235.
\:\\{end\_of\_AWEAVE}, \[38], 308\*.
\:\p{END\_OF\_FILE}, 33, \[70].
\:\p{END\_OF\_LINE}, 33, \[70].
\:\p{END\_OF\_PAGE}, \[70].
\:\\{end\_translation}, \[168], 243, 249, 251, 256.
\:\\{ending}, \[167], 169--170, 176\*, 235.
\:\.{entry}, 169.
\:\&{entry}, 45, \[67], 128, 167.
\:\p{ENUM}, \[70].
\:\p{ENUMERATION\_IO}, \[70].
\:\p{EPSILON}, \[70].
\:\\{equal}, \[74], 76--78.
\:\\{equivalence\_sign}, \[20], 112, 134, 230, 270, 273.
\:\\{err\_print}, \[35], 76, 80, 83, 87--88, 90--91, 97--98, 101, 111, 114,
118--119, 121--122, 124, 158, 163--164, 232, 238, 266, 269--270, 273--274.
\:\\{error}, 33, \[35], \\{36}, 38, 91.
\:\\{error\_message}, \[15], 310.
\:\.{exception}, 169.
\:\&{exception}, 45, \[67], 167.
\:\\{exception\_like}, \[45], 67, 235.
\:\\{exception\_scrap}, \[167], 169--170, 176\*, 185, 235.
\:\&{exit}, 45, \[67], 167.
\:\.{exit}, 169.
\:\\{exp}, \[167], 169--170, 176\*, 226.
\:\p{EXP}, \[70].
\:\\{exponent}, \[109], 113, 226.
\:\\{extension}, \[74], 78.
\:\.{Extra \AT!>}, 101.
\:\|{f}, \[28], \[33], \[291].
\:\p{FALSE}, 28, 30, 32--34, \[70], 83--85, 90, 93, 98, 109, 111, 126--127,
140, 143, 148, 251--252, 262, 281, 307--308\*.
\:\\{fatal\_error}, \[38], 39--40, 308\*.
\:\\{fatal\_message}, \[15], 310.
\:\p{FIELD}, \[70].
\:\p{FILE}, 28, 30, \[70].
\:\p{FILE\_MODE}, 27--28, \[70].
\:\p{FILE\_TYPE}, 17, \[70].
\:\\{final\_limit}, \[33].
\:\p{FINE\_DELTA}, \[70].
\:\\{finish\_ADA}, 267, \\{268}, 272.
\:\\{finish\_line}, \[142], \\{143}, 144, 157, 160--161, 256, 262, 268, 279,
281--282, 296, 299.
\:\p{FIRST}, 17, \[70].
\:\p{FIRST\_BIT}, \[70].
\:\\{first\_phase}, \\{7\*}, 308\*.
\:\.{first\_phase.adb}, 7\*.
\:\.{first\_phase\_.ads}, 7\*.
\:\\{first\_xref}, \[276], 277, 280.
\:\p{FIXED\_IO}, \[70].
\:\\{flag}, 278, \[279], 280.
\:\p{FLOAT}, \[70].
\:\p{FLOAT\_IO}, \[70].
\:\\{flush\_buffer}, \[140], \\{141}, 143, 148--149, 262, 281.
\:\\{follows}, \[167], 169--170, 177, 183, 193, 198, 205, 235.
\:\\{footnote}, 275, \[278], \\{279}, 299.
\:\.{for}, 169.
\:\&{for}, 45, \[67], 128, 167.
\:\\{force}, \[168], 169--170, 172, 182--184, 186, 192, 194, 196, 198--199,
204, 208, 210, 226, 231, 235, 240, 242, 256, 268, 273.
\:\\{force\_line}, \[99], 101, 226.
\:\p{FORE}, \[70].
\:\p{FORM}, \[70].
\:\\{form\_feed}, \[20], 32.
\:\\{format}, \[99], 101, 128--130, 132--133, 221, 223\*, 239--240, 266--267.
\:\\{found}, 61, 64, 76, 173--174, 176\*--177.
\:\\{freeze\_text}, \[211], 212, 220, 224, 236, 240, 251.
\:\p{FROM}, \[70].
\:\&{function}, 45, \[67], 128, 167.
\:\.{function}, 169.
\:\.{generic}, 169.
\:\&{generic}, 45, \[67], 167.
\:\\{generic\_like}, \[45], 67, 235.
\:\\{generic\_scrap}, \[167], 169--170, 176\*, 192, 235.
\:\p{GET}, 33, \[70].
\:\p{GET\_LINE}, \[70].
\:\\{get\_line}, 82, \[94], \\{95}, 103, 105, 107, 111, 118, 142, 157, 160, 163.
\:\\{get\_next}, 109, \[110], \\{111}, 125, 129, 132--135, 221, 223\*,
227--228, 266, 269--270, 272--274.
\:\\{get\_output}, 247, \\{248}, 249, 251--252, 255--256.
\:\&{goto}, 45, \[67], 167.
\:\.{goto}, 169.
\:\\{goto\_like}, \[45], 67--68, 235.
\:\\{gproc}, \[167], 170, 176\*, 210.
\:\\{greater}, \[74], 76, 78, 80.
\:\\{greater\_or\_equal}, \[20], 112, 230.
\:\|{h}, \[61], \[286].
\:\\{harmless\_message}, \[15], 310.
\:\\{hash}, 43, \[59], 63, 286.
\:\\{hash\_index}, \[59], 61.
\:\\{hash\_size}, \[14], 59, 62, 286.
\:\\{head}, \[287], 291--294, 308\*.
\:\\{hi\_ptr}, \[171], 216, 219\*.
\:\\{history}, \[15], 310.
\:\\{history\_type}, \[15].
\:\.{Hmm... n of the preceding...}, 91.
\:\|{i}, \[22], \[61], \[120], \[129], \[212], \[214], \[216], \[309].
\:\\{ic}, 308\*.
\:\\{id}, \[66], 67--68.
\:\\{id\_first}, \[59], 61--62, 64--65, 67, 109, 113--114, 121--124, 231,
233--234.
\:\\{id\_flag}, \[172], 235, 248, 269--270.
\:\\{id\_loc}, \[59], 61--62, 64--66, 109, 113--114, 121--124, 231, 234.
\:\\{id\_lookup}, 59, \[60], \\{61}, 66, 109, 129, 132, 134, 235, 269--270.
\:\\{identifier}, \[109], 113, 128--129, 132, 134, 225, 247--248, 251, 253,
269--270.
\:\&{if}, 45, \[67], 167.
\:\.{if}, 169.
\:\\{if\_like}, \[45], 67, 235.
\:\\{if\_scrap}, \[167], 169--170, 174, 189, 235.
\:\\{ignore}, \[99], 101, 226.
\:\\{ignore\_spaces}, \[251], 252, 256.
\:\\{ilk}, 41, \[42], 45--46, 59--60, 63, 65, 129, 131--132, 134, 235, 253, 295.
\:\.{Illegal control code...}, 259.
\:\.{Illegal use of \AT!...}, 164.
\:\p{IMAGE}, \[70].
\:\.{Improper ASCII-code sequence...}, 124.
\:\.{Improper based number...}, 111.
\:\.{Improper format definition}, 270.
\:\.{Improper macro definition}, 269.
\:\&{in}, 45, \[67], 167.
\:\.{in}, 20, 169.
\:\p{IN\_FILE}, \[70], 308\*.
\:\\{in\_like}, \[45], 67, 235.
\:\\{in\_scrap}, \[167], 169--170, 174, 185, 201, 226, 235.
\:\.{Incompatible module names}, 76.
\:\\{incr}, \[13], 33, 53, 62, 64--65, 77--78, 80, 87--88, 90, 96--97, 103,
105, 107--108, 111--115, 118--119, 121--122, 124, 127, 141, 145, 153, 158,
161--164, 173--174, 176\*, 211, 216, 224, 231--232, 245, 248, 252, 258--261,
264, 280, 284, 291.
\:\\{indent}, \[168], 169--170, 177, 182--183, 192--194, 198, 202, 205, 210,
235, 251.
\:\p{INDEX}, \[70].
\:\\{Index}, 308\*.
\:\\{infinity}, \[290], 291--292.
\:\\{init\_stack}, \[243], 266--267, 272, 299.
\:\\{initialize}, \\{3\*}, 308\*.
\:\p{INLINE}, \[70].
\:\\{inner}, 242, \[243], 248, 252, 256.
\:\\{inner\_tok\_flag}, \[172], 240, 248, 250.
\:\p{INOUT\_FILE}, \[70].
\:\.{Input ended in mid-comment}, 163.
\:\.{Input ended in section name}, 118.
\:\.{Input line too long}, 33.
\:\.{input.ch}, 307.
\:\.{input.web}, 307.
\:\\{input\_has\_ended}, \[82], 90, 93, 96, 102--103, 105, 107, 111, 118, 126,
157, 160, 163, 262.
\:\\{input\_ln}, \[32], \\{33}, 87--88, 90, 96--97.
\:\\{input\_output}, 3\*, \\{4\*}, 5\*--11\*, 308\*.
\:\.{input\_output.adb}, 4\*.
\:\.{input\_output.ads}, 4\*.
\:\\{int\_number\_io}, 3\*--4\*, 6\*--11\*, 24, 308\*.
\:\p{INTEGER}, 16, 19, 31, 41--42, 45, 50--51, 56--57, 59, \[70], 73, 82, 90,
129, 139, 150--152, 171, 212, 214, 243, 251, 263, 279, 286--287, 289, 308\*.
\:\p{INTEGER\_IO}, \[70].
\:\p{IO\_EXCEPTIONS}, \[70].
\:\&{is}, 45, \[67], 167.
\:\.{is}, 169.
\:\\{is\_like}, \[45], 67, 235.
\:\p{IS\_OPEN}, 28, 30, 33, \[70], 93.
\:\\{italic}, \[45], 295.
\:\p{ITEM}, \[70].
\:\|{j}, \[76], \[80], \[111], \[141], \[149], \[212], \[214], \[220], \[223%
\*], \[251].
\:\\{join}, \[99], 101, 226, 266.
\:\\{jump\_out}, \[38].
\:\|{k}, \[37], \[49], \[61], \[76], \[80], \[111], \[141], \[143], \[148], %
\[152], \[155], \[212], \[214], \[219\*], \[251].
\:\\{k\_limit}, \[251], 258, 260.
\:\\{k\_module}, \[283], 284.
\:\|{l}, \[36], \[61], \[76], \[80], \[234].
\:\\{label}, \[167], 170, 174, 181, 196.
\:\p{LARGE}, \[71].
\:\p{LAST}, 17, \[71].
\:\p{LAST\_BIT}, \[71].
\:\p{LAYOUT\_ERROR}, \[71].
\:\\{left\_arrow}, \[20], 112, 230.
\:\\{left\_label\_bracket}, \[20], 112, 230.
\:\\{length}, \[43], 63, 253, 295.
\:\p{LENGTH}, \[71].
\:\\{length\_index}, \[139], 141, 148.
\:\\{less}, \[74], 76--78, 80.
\:\\{less\_or\_equal}, \[20], 112, 230.
\:\\{lexical\_scanning}, \\{6\*}, 7\*--10\*.
\:\.{lexical\_scanning.adb}, 6\*.
\:\.{lexical\_scannning.ads}, 6\*.
\:\\{lhs}, \[131], 132, 134.
\:\\{limit}, 32--33, 37, \[82], 85, 87--89, 91, 93--98, 103, 105, 107,
111--112, 114, 118, 121--122, 143, 157--158, 160--161, 163, 251, 258, 260.
\:\.{limited}, 169.
\:\&{limited}, 45, \[67], 167.
\:\\{line}, 37, \[82], 83, 87--88, 90, 93, 96--98.
\:\p{LINE}, \[71].
\:\.{Line had to be broken}, 149.
\:\\{line\_break}, \[99], 101, 226, 266.
\:\\{line\_feed}, \[20], 32.
\:\\{line\_length}, \[14], 139--140, 145.
\:\p{LINE\_LENGTH}, \[71].
\:\\{lines\_dont\_match}, \\{85}, 90.
\:\\{link}, 41, \[42], 43, 46, 63, 286.
\:\p{LIST}, \[71].
\:\\{llink}, \[46], 76--77, 80, 137, 299.
\:\\{lo\_ptr}, \[171], 212, 215--216, 219\*--220.
\:\\{loc}, 33, 37, \[82], 87, 91, 93--95, 97--98, 103, 105, 107--108, 111--115,
118--119, 121--124, 127, 132, 157--158, 160--161, 163--164, 226, 251, 258,
265--266.
\:\\{long\_buf\_size}, \[14], 31, 260--261.
\:\\{long\_buffer\_index}, \[31], 219\*, 223\*, 251.
\:\\{long\_name\_index}, \[73], 76, 80, 111.
\:\\{longest\_name}, \[14], 73, 118, 120.
\:\&{loop}, 45, \[67], 167.
\:\.{loop}, 169.
\:\\{loop\_like}, \[45], 67, 235.
\:\\{loop\_scrap}, 166, \[167], 169--170, 174, 189, 198, 235.
\:\p{LOW\_LEVEL\_IO}, \[71].
\:\p{LOWER\_CASE}, \[71].
\:\|{m}, \[54], \[77], \[152].
\:\p{MACHINE\_EMAX}, \[71].
\:\p{MACHINE\_EMIN}, \[71].
\:\p{MACHINE\_MANTISSA}, \[71].
\:\p{MACHINE\_OVERFLOWS}, \[71].
\:\p{MACHINE\_RADIX}, \[71].
\:\p{MACHINE\_ROUNDS}, \[71].
\:\\{make\_output}, \[249], 250, \\{251}, 257, 268, 299.
\:\p{MANTISSA}, \[71].
\:\\{mark\_error}, \[15], 36, 259--260.
\:\\{mark\_fatal}, \[15], 38.
\:\\{mark\_harmless}, \[15], 120, 137, 149.
\:\\{math}, \[167], 169--170, 173--174, 177--181, 183, 185, 187--189, 191,
193--195, 198--207, 209--210, 219\*--220, 226, 230--231, 233--235, 269--270,
273.
\:\\{math\_bin}, \[168], 169, 226, 235, 251, 254.
\:\\{math\_break}, \[99], 101, 226, 266.
\:\\{math\_op}, \[168], 251.
\:\\{math\_rel}, \[168], 169, 172, 226, 251, 254, 273.
\:\\{max\_bytes}, \[14], 42, 65, 77.
\:\p{MAX\_DIGITS}, \[71].
\:\p{MAX\_INT}, \[71].
\:\p{MAX\_MANTISSA}, \[71].
\:\\{max\_modules}, \[14], 50--51, 127.
\:\\{max\_names}, \[14], 42, 65, 77.
\:\\{max\_refs}, \[14], 51, 53.
\:\\{max\_scr\_ptr}, \[171], 229, 238, 241, 268, 309.
\:\\{max\_scraps}, \[14], 171, 217, 229, 241, 287.
\:\\{max\_sort\_ptr}, \[287], 291, 309.
\:\\{max\_sorts}, \[287], 291.
\:\\{max\_stack\_ptr}, \[243], 245, 309.
\:\\{max\_texts}, \[14], 56, 215, 217, 229, 241.
\:\\{max\_tok\_ptr}, \[57], 215, 229, 241, 250, 268, 309.
\:\\{max\_toks}, \[14], 57, 162, 215, 217, 220, 229, 231, 240--241.
\:\\{max\_txt\_ptr}, \[57], 215, 229, 241, 250, 268, 309.
\:\p{MEMORY\_SIZE}, \[71].
\:\\{mid\_xref}, \[276], 277.
\:\p{MIN\_INT}, \[71].
\:\.{Missing `|'...}, 238.
\:\&{mitpo}, \[12].
\:\.{mod}, 169.
\:\&{mod}, 45, \[67], 167.
\:\\{mod\_check}, \\{137}, 138.
\:\\{mod\_flag}, \[172], 248, 273--274, 299.
\:\\{mod\_like}, \[45], 67--68, 235.
\:\\{mod\_lookup}, 73--74, \[75], \\{76}, 116--117.
\:\\{mod\_name}, \[247], 248, 251.
\:\\{mod\_print}, \[298], \\{299}, 300.
\:\\{mod\_scrap}, \[167], 169--170, 174, 186, 273--274.
\:\\{mod\_text}, \[73], 76--80, 111, 116--120.
\:\\{mod\_xref\_switch}, \[51], 55, 135.
\:\p{MODE}, 28, 30, \[71].
\:\p{MODE\_ERROR}, \[71].
\:\\{mode\_field}, \[243].
\:\\{mode\_of\_file}, 27, \[28].
\:\\{module\_count}, \[50], 54--55, 95, 126--127, 262, 264--265, 273, 284.
\:\\{module\_index}, \[50], 283.
\:\\{module\_name}, \[99], 101, 109, 115, 132, 135, 266, 272, 274.
\:\p{MONTH}, \[71].
\:\p{MONTH\_NUMBER}, \[71].
\:\|{n}, \[54], \[90].
\:\p{NAME}, 28, 30, \[71].
\:\.{Name does not match}, 80.
\:\p{NAME\_ERROR}, 28, 30, \[71].
\:\\{name\_index}, \[42], 43, 80.
\:\\{name\_info}, \[42], 46, 285.
\:\\{name\_pointer}, \[43], 48--49, 53--55, 60--61, 66, 75--76, 79--80, 109,
129, 131, 137, 154--155, 171, 223\*, 251, 271, 285, 298--299.
\:\\{name\_ptr}, \[43], 49, 61, 63, 65, 77, 309.
\:\p{NATURAL}, \[71].
\:\.{Never defined: <section name>}, 137.
\:\.{Never used: <section name>}, 137.
\:\.{new}, 169.
\:\&{new}, 45, \[67], 167.
\:\\{new\_like}, \[45], 67, 235.
\:\p{NEW\_LINE}, 24, \[71], 141.
\:\\{new\_ln}, \[24], 37, 91, 149, 308\*, 310.
\:\\{new\_mod\_xref}, \[53], \\{55}, 135.
\:\\{new\_module}, \[99], 101, 104--105, 111, 160.
\:\p{NEW\_PAGE}, \[71].
\:\\{new\_scrap}, \[167], 169--170, 174, 205, 235.
\:\\{new\_xref}, \[53], \\{54}, 129, 131--132, 134.
\:\\{next\_control}, \[125], 128--130, 132--135, 221, 223\*, 225--228, 231,
238--240, 250, 266--267, 269--274.
\:\\{next\_name}, \[285], 286, 293.
\:\\{next\_xref}, \[276], 277, 280, 296--297.
\:\.{No change file}, 308\*.
\:\\{no\_line\_break}, \[99], 101, 226, 266.
\:\\{no\_underline}, \[99], 101, 115, 132.
\:\\{normal}, \[45], 60, 63, 128--129, 131--132, 134, 235, 269--270, 295.
\:\.{not}, 20, 169.
\:\&{not}, 45, \[67], 167.
\:\\{not\_equal}, \[20], 112, 230.
\:\\{not\_like}, \[45], 67, 235.
\:\\{not\_scrap}, \[167], 169--170, 174, 230, 235.
\:\\{not\_sign}, \[20], 230.
\:\&{null}, 45, \[67], 167.
\:\.{null}, 169.
\:\\{null\_array}, \[43], 44, 59.
\:\\{num}, \[51], 52--55, 137, 257, 273, 277, 279--280, 296.
\:\p{NUM}, \[71].
\:\\{num\_field}, \[51].
\:\p{NUMBER\_BASE}, \[71].
\:\p{NUMERIC\_ERROR}, \[71].
\:\&{of}, 45, \[67], 167.
\:\.{of}, 169.
\:\\{of\_like}, \[45], 67, 235.
\:\p{OFF}, \[71].
\:\p{ON}, \[71].
\:\p{OPEN}, 28, \[71].
\:\\{open\_a\_file}, \[27], \\{28}, 308\*.
\:\\{open\_scrap}, 166, \[167], 169--170, 174, 202, 220, 226, 235.
\:\\{opened}, \[307], 308\*.
\:\\{opt}, 166, \[168], 169--170, 196, 202, 226, 252, 255.
\:\&{optim}, \[12].
\:\p{OPTIMIZE}, \[71].
\:\&{or}, 45, \[67], 167.
\:\.{or}, 20, 169.
\:\\{or\_like}, \[45], 67, 235.
\:\\{or\_sign}, \[20], 230.
\:\\{other\_line}, \[82], 83, 93, 98.
\:\&{others}, 45, \[67], 167.
\:\.{others}, 169.
\:\.{out}, 169.
\:\&{out}, 45, \[68], 167.
\:\\{out\_buf}, \[139], 141, 144--146, 148--149, 256, 268, 273.
\:\p{OUT\_FILE}, 30, \[71], 308\*.
\:\\{out\_line}, \[139], 141, 144, 149, 263.
\:\\{out\_mod}, \[151], \\{152}, 257, 265, 280, 284, 296.
\:\\{out\_name}, \[154], \\{155}, 253, 295.
\:\\{out\_ptr}, \[139], 141, 143--145, 148--149, 161, 256, 263, 268, 273.
\:\\{outdent}, \[168], 170, 177, 186, 189, 192--194, 202, 205, 208.
\:\\{outer}, 242, \[243], 255--256.
\:\\{outer\_parse}, \[239], \\{240}, 267, 272.
\:\\{outer\_xref}, \\{130}, 133, 135, 239.
\:\\{output}, \\{8\*}, 9\*--11\*.
\:\.{output.adb}, 8\*.
\:\.{output.ads}, 8\*.
\:\\{output\_ADA}, 249, \\{250}, 258, 266.
\:\\{output\_mode}, \[243], 251.
\:\\{output\_state}, \[243].
\:\\{output\_tok}, \[99], 101, 109, 115, 122, 225, 231.
\:\\{out1}, \[145], 152--153, 155, 158, 161, 251, 253--258, 265--266, 279--280,
284, 296.
\:\\{out2}, \[145], 152, 252, 254--257, 263, 265, 268, 284, 294--296, 299.
\:\\{out3}, \[145], 254, 281, 296.
\:\\{out4}, \[145], 268, 279--280, 282, 284.
\:\\{out5}, \[145], 254.
\:\\{overflow}, \[40], 53, 65, 77, 127, 162, 215, 220, 229, 241, 245, 260--261,
291.
\:\\{Overwrite}, 308\*.
\:\|{p}, \[49], \[54], \[55], \[61], \[76], \[80], \[129], \[137], \[155], %
\[223\*], \[238], \[240], \[245], \[268].
\:\p{PACK}, 42, \[71].
\:\.{package}, 169.
\:\&{package}, 45, \[68], 128, 167.
\:\p{PAGE}, \[71].
\:\p{PAGE\_LENGTH}, \[71].
\:\\{param}, \[99].
\:\\{parsing}, \\{9\*}, 10\*, 176\*, 219\*, 223\*.
\:\.{parsing-ada\_parse.adb}, 223\*.
\:\.{parsing-do\_cases.adb}, 176\*.
\:\.{parsing-translate.adb}, 219\*.
\:\.{parsing.adb}, 9\*.
\:\.{parsing.ads}, 9\*.
\:\\{pdollar}, 166, \[168], 169--170, 177, 181, 183, 189, 193--194, 198, 202,
205, 207, 209--210, 220, 226, 231, 233--235, 251--252, 269--270.
\:\\{per\_cent}, 140, \[141].
\:\\{period}, \[167], 169--170, 174, 187, 226.
\:\\{Phase\_I}, \[301], \\{302}, 308\*.
\:\\{Phase\_II}, \[303], \\{304}, 308\*.
\:\\{Phase\_III}, \[305], \\{306}, 308\*.
\:\\{phase\_one}, \[34], 35, 126.
\:\\{phase\_three}, \[34], 126, 257, 282.
\:\\{pid}, \[66], 69--72.
\:\\{pop\_level}, \[246], 248.
\:\p{POS}, 18, \[71].
\:\p{POSITION}, \[71].
\:\p{POSITIVE}, \[71].
\:\p{POSITIVE\_COUNT}, \[71].
\:\\{pp}, \[171], 173--174, 176\*--210, 213, 215--216, 219\*.
\:\&{pragma}, 45, \[68], 167.
\:\.{pragma}, 169.
\:\\{pragma\_like}, \[45], 68, 235.
\:\p{PRED}, \[71].
\:\\{predefined}, \[45], 66, 235, 253, 295.
\:\\{prefix}, \[74], 78.
\:\\{prefix\_lookup}, \[79], \\{80}, 116.
\:\\{prime\_the\_change\_buffer}, \\{86}, 93, 97.
\:\\{print}, \[24], 37, 49, 91, 120, 127, 137, 149, 259--260, 265, 282, 309.
\:\\{print\_id}, \[48], \\{49}, 137, 259--260.
\:\\{print\_int}, \[24], 37, 91, 127, 149, 265, 309.
\:\\{print\_ln}, \[24], 37, 149, 308\*.
\:\\{print\_nl}, \[24], 33, 35, 38, 120, 137, 149, 259--260, 262, 282, 309--310.
\:\p{PRIORITY}, \[71].
\:\.{private}, 169.
\:\&{private}, 45, \[68], 167.
\:\\{private\_like}, \[45], 68, 235.
\:\\{private\_scrap}, \[167], 169--170, 174, 235.
\:\\{proc}, \[167], 169--170, 174, 192, 210, 235.
\:\\{proc\_like}, \[45], 67--68, 129, 235.
\:\.{procedure}, 169.
\:\&{procedure}, 45, \[68], 128, 167.
\:\\{prod}, 221.
\:{productions, table of}, 170.
\:\p{PROGRAM\_ERROR}, \[71].
\:\\{pseudo\_semi}, \[99], 101, 226, 266.
\:\\{psid}, \[66].
\:\\{pspace}, \[168], 169--170, 198, 206, 235, 251--252, 256.
\:\\{push\_level}, \[244], \\{245}, 248, 251.
\:\p{PUT}, 24, \[71], 141, 144.
\:\p{PUT\_LINE}, \[71].
\:\|{q}, \[54], \[55], \[61], \[76], \[80], \[240].
\:\|{r}, \[55], \[80].
\:\.{raise}, 169.
\:\&{raise}, 45, \[68], 167.
\:\.{range}, 169.
\:\&{range}, 45, \[68], 167.
\:\p{RANGE}, \[71].
\:\\{range\_like}, \[45], 68, 235.
\:\p{READ}, \[71].
\:\p{RECEIVE\_CONTROL}, \[71].
\:\&{record}, 45, \[68], 167.
\:\.{record}, 169.
\:\\{record\_like}, \[45], 68, 235.
\:\\{record\_scrap}, \[167], 169--170, 189, 198, 210, 235.
\:{recursion}, 137, 249, 298.
\:\\{red}, 173, \\{212}, 214.
\:\\{reduce}, \[173], 177--178, 180--194, 196, 198--199, 201--202, 204--210,
212.
\:\.{rem}, 169.
\:\&{rem}, 45, \[68], 167.
\:\.{renames}, 169.
\:\&{renames}, 45, \[68], 167.
\:\\{res\_flag}, \[172], 235, 248.
\:\\{res\_word}, \[247], 248, 251, 253.
\:\\{reserved}, \[45], 54.
\:\p{RESET}, \[71], 93.
\:\\{reset\_input}, \[92], \\{93}, 126, 262.
\:\\{restart}, 95, 111, 115, 248.
\:\\{reswitch}, 225, 251--252, 256.
\:\&{return}, 45, \[68], 167.
\:\.{return}, 169.
\:\\{return\_like}, \[45], 67--68, 235.
\:\\{return\_scrap}, \[167], 169--170, 174, 235.
\:\.{reverse}, 169.
\:\&{reverse}, 45, \[68], 167.
\:\\{rhs}, \[131], 134.
\:\\{right\_arrow}, \[20], 112, 230.
\:\\{right\_label\_bracket}, \[20], 112, 230.
\:\\{rlink}, \[46], 76--77, 80, 137, 299.
\:\\{roman}, \[45], 128, 295.
\:\\{root}, \[46], 47, 76, 80, 138, 300.
\:\p{SAFE\_EMAX}, \[72].
\:\p{SAFE\_LARGE}, \[72].
\:\p{SAFE\_SMALL}, \[72].
\:\\{save\_base}, \[238].
\:\\{save\_limit}, \[251], 258.
\:\\{save\_line}, \[263].
\:\\{save\_loc}, \[251], 258.
\:\\{save\_mode}, \[251], 256.
\:\\{save\_next\_control}, \[250].
\:\\{save\_place}, \[263].
\:\\{save\_position}, \[263], 264, 267.
\:\\{save\_text\_ptr}, \[250].
\:\\{save\_tok\_ptr}, \[250].
\:\\{scanning\_based\_number}, \[109], 111.
\:\\{scrap\_base}, \[171], 213, 217, 219\*--220, 238.
\:\\{scrap\_index}, \[171], 238, 287.
\:\\{scrap\_ptr}, \[171], 216--217, 221, 224, 229, 236, 238, 241, 268, 270,
287, 299.
\:\\{sc0}, \[224], 226, 270.
\:\\{sc1}, \[224], 226--228, 235, 273--274.
\:\\{sc2}, \[224], 226, 228, 230--231, 233--235, 270, 273.
\:\\{sc3}, \[224], 226, 230--231, 235, 269--270, 273.
\:\\{sc4}, \[224], 226, 230, 235.
\:\\{sc5}, \[224], 226, 235, 269--270.
\:\\{sc6}, \[224], 226.
\:\\{searching}, \\{5\*}, 6\*--7\*, 9\*--10\*, 308\*.
\:\.{searching.adb}, 5\*.
\:\.{searching.ads}, 5\*.
\:\\{second\_phase}, \\{10\*}, 11\*, 308\*.
\:\.{second\_phase.adb}, 10\*.
\:\.{second\_phase.ads}, 10\*.
\:\p{SECONDS}, \[72].
\:\.{Section name didn't end}, 119.
\:\.{Section name too long}, 120.
\:\&{select}, 45, \[68], 167.
\:\.{select}, 169.
\:\\{semi}, \[167], 169--170, 177, 186, 189, 193, 198--199, 204--206, 210, 226,
270, 273.
\:\p{SEND\_CONTROL}, \[72].
\:\&{separate}, 45, \[68], 167.
\:\.{separate}, 169.
\:\\{separate\_like}, \[45], 68, 235.
\:\\{separate\_scrap}, \[167], 169--170, 174, 198, 205, 235.
\:\p{SEQUENTIAL\_IO}, \[72].
\:\p{SET}, \[72].
\:\p{SET\_COL}, \[72].
\:\\{set\_element\_sign}, \[20], 226.
\:\p{SET\_INDEX}, \[72].
\:\p{SET\_INPUT}, \[72].
\:\p{SET\_LINE}, \[72].
\:\p{SET\_LINE\_LENGTH}, \[72].
\:\p{SET\_OUTPUT}, \[72].
\:\p{SET\_PAGE\_LENGTH}, \[72].
\:\p{SHARED}, \[72].
\:\\{sid}, \[66].
\:\\{sixteen\_bits}, \[41], 42, 46, 51, 54, 57, 59, 75--76, 79--80, 131, 172,
212, 214, 243, 248, 250, 263, 278--279, 287.
\:\p{SIZE}, \[72].
\:\\{skip\_comment}, \[106], \\{107}, 130, 156, 162.
\:\\{skip\_limbo}, \[102], \\{103}, 126, 156.
\:\p{SKIP\_LINE}, 33, \[72].
\:\p{SKIP\_PAGE}, \[72].
\:\\{skip\_TeX}, \[104], \\{105}, 132, 156.
\:\\{slanted}, \[45], 295.
\:\p{SMALL}, \[72].
\:\.{Sorry, x capacity exceeded}, 40.
\:\\{sort\_ptr}, \[287], 291--294.
\:\p{SPACE}, \[72].
\:\p{SPACING}, \[72].
\:{special string characters}, 231.
\:\p{SPLIT}, \[72].
\:{split procedures}, 174, 221, 301.
\:\\{spotless}, \[15], 310.
\:\\{sq}, 173, \\{214}.
\:\\{squash}, \[173], 177, 179, 185, 187, 189, 195, 197--201, 203--206, 210,
214.
\:\\{sspace}, \[168], 169--170, 183, 185, 189, 194, 196, 198, 206--207, 235,
251--252, 256.
\:\\{stack}, \[243], 245--246.
\:\\{stack\_ptr}, \[243], 245--246.
\:\\{stack\_size}, \[14], 243, 245.
\:\p{STANDARD}, \[72].
\:\p{STANDARD\_INPUT}, \[72].
\:\p{STANDARD\_OUTPUT}, 24, \[72].
\:\&{stat}, \[12].
\:\p{STATUS\_ERROR}, 28, 30, \[72].
\:\\{stmt}, \[167], 170, 174, 177, 182--184, 186, 194, 196, 198--199, 206, 208.
\:\p{STORAGE\_ERROR}, \[72].
\:\p{STORAGE\_SIZE}, \[72].
\:\p{STORAGE\_UNIT}, \[72].
\:\\{str}, \[109], 114, 225.
\:\p{STRING}, 27--28, 30, \[72].
\:\.{String constant didn't end}, 114.
\:\\{string\_delimiter}, \[251], 260.
\:\\{Strings}, 308\*.
\:\.{subtype}, 169.
\:\&{subtype}, 45, \[68], 128, 167.
\:\p{SUCC}, \[72].
\:\\{success}, 27, \[28].
\:\p{SUPPRESS}, \[72].
\:\p{SYSTEM}, \[72].
\:{system dependencies}, 17, 21, 24--25, 28, 30, 32, 37, 301, 310--311.
\:\p{SYSTEM\_NAME}, \[72].
\:\\{s0}, \[224].
\:\|{t}, \[61].
\:\\{tab\_mark}, \[20], 37, 101, 103, 108, 111, 118--119, 143, 158, 161.
\:\&{task}, 45, \[68], 128, 167.
\:\.{task}, 169.
\:\p{TASKING\_ERROR}, \[72].
\:\&{tats}, \[12].
\:\\{temp\_ilk}, \[131], 132.
\:\\{temp\_line}, \[82], 83.
\:\.{terminate}, 169.
\:\&{terminate}, 45, \[68], 167.
\:\p{TERMINATED}, \[72].
\:\\{terminator}, \[167], 170, 186, 189, 208.
\:\.{TeX file not open}, 308\*.
\:\.{TeX string should be...}, 266.
\:\\{TeX\_file}, \[29], 30, 141, 144, 308\*.
\:\\{TeX\_file\_name}, 27, \[30].
\:\\{TeX\_output\_file\_name}, \[307], 308\*.
\:\\{TeX\_string}, \[99], 101, 109, 115, 225, 266.
\:\\{text\_char}, \[17], 18--19, 24, 33.
\:\\{text\_file}, \[17], 26--29, 32--33.
\:\p{TEXT\_IO}, 3\*--11\*, 17, 24, 27--28, 30, 33, \[72], 93, 141, 144, 308\*.
\:\\{text\_pointer}, \[56], 57, 171, 217--219\*, 237--238, 240, 244--245, 250,
268.
\:\\{text\_ptr}, \[57], 211--212, 215, 217, 220, 224, 229, 236, 240--241,
250--251, 268, 299.
\:\&{then}, 45, \[68], 167.
\:\.{then}, 169.
\:\\{then\_like}, \[45], 68, 235.
\:\\{then\_scrap}, \[167], 169--170, 178, 194, 235.
\:\\{thin\_space}, \[99], 101, 226, 266.
\:\\{third\_phase}, \\{11\*}, 308\*.
\:\.{third\_phase.adb}, 11\*.
\:\.{third\_phase.ads}, 11\*.
\:\.{This can't happen}, 39.
\:\\{this\_module}, \[271], 272--273, 275, 277.
\:\\{this\_xref}, \[276], 277, 296--297.
\:\p{TICK}, \[72].
\:\p{TIME}, \[72].
\:\p{TIME\_ERROR}, \[72].
\:\p{TIME\_OF}, \[72].
\:\p{TO}, \[72].
\:\\{tok\_field}, \[243].
\:\\{tok\_flag}, \[172], 173, 236, 240, 248, 268.
\:\\{tok\_mem}, \[57], 162, \[172], 173, 243, 248, 257.
\:\\{tok\_ptr}, \[57], 162, 164, 173, 211, 215, 217, 220, 229, 231, 240--241,
250, 268, 299.
\:\\{tok\_start}, 56, \[57], 58, 171, 211, 245.
\:\\{token\_index}, \[57].
\:\\{trans}, \[171], 173, 212, 216--217, 219\*, 221, 224, 236--237, 287.
\:\\{trans\_type}, \[171], 287.
\:\\{translate}, 174, \[217], \[218], \\{219\*}, 238, 268.
\:\p{TRUE}, 28, 30, 32--34, \[72], 82--83, 85, 90, 93, 95--96, 98, 111,
126--127, 140, 148--149, 282.
\:\&{type}, 45, \[68], 128, 167.
\:\.{type}, 169.
\:\\{type\_like}, \[45], 67--68, 129, 235.
\:\p{TYPE\_SET}, \[72].
\:\\{typewriter}, \[45], 128, 295.
\:\p{UNBOUNDED}, \[72].
\:\\{Unbounded}, 308\*.
\:\\{unbucket}, 290, \\{291}, 292--293.
\:\p{UNCHECKED\_CONVERSION}, \[72].
\:\p{UNCHECKED\_DEALLOCATION}, \[72].
\:\\{underline}, \[99], 101, 115, 132.
\:\.{Unknown control code}, 101.
\:\\{update\_terminal}, \[25], 36, 127, 265.
\:\p{UPPER\_CASE}, \[72].
\:\.{use}, 169.
\:\&{use}, 45, \[68], 167.
\:\p{USE\_ERROR}, 28, 30, \[72].
\:\\{Ustring}, 308\*.
\:\\{Ustrings}, 308\*.
\:\p{VAL}, 18, \[72].
\:\p{VALUE}, \[72].
\:\\{verbatim}, \[99], 101, 115, 122, 225, 231.
\:\.{Verbatim string didn't end}, 122.
\:\|{w}, \[49], \[61], \[76], \[80], \[155], \[251].
\:\\{w\_range}, \[42], 43, 49, 61, 76, 80, 155, 251, 287.
\:\&{when}, 45, \[68], 167.
\:\.{when}, 169.
\:\\{when\_like}, \[45], 68, 235.
\:\\{when\_scrap}, \[167], 169--170, 174, 188, 198, 235.
\:\.{Where is the match...}, 87, 91, 97.
\:\.{while}, 169.
\:\&{while}, 45, \[68], 167.
\:\p{WIDTH}, \[72].
\:\\{wildcard}, \[45], 128, 295.
\:\&{with}, 45, \[68], 167.
\:\.{with}, 169.
\:\\{with\_like}, \[45], 68, 235.
\:\\{with\_scrap}, \[167], 169--170, 174, 186, 198, 210, 235.
\:\p{WRITE}, \[72].
\:\\{ww}, 14, \[42], 43--44, 49, 54, 64--65, 77--78, 155, 253, 258, 286, 293,
309.
\:\\{xchr}, \[18], 19, 21--22, 37, 49, 120, 141, 149.
\:\\{xlink}, \[51], 54--55, 137, 257, 277, 280, 296--297.
\:\\{xlink\_field}, \[51].
\:\\{xmem}, \[51].
\:\\{xmem\_field}, \[51].
\:\.{xor}, 169.
\:\&{xor}, 45, \[68], 167.
\:\\{xord}, \[18], 22, 33.
\:\\{xref}, 41, \[42], 51--52, 54--55, 65, 137, 257, 273, 277, 286, 297.
\:\\{xref\_bold\_face}, \[99], 101.
\:\\{xref\_italic}, \[99], 101.
\:\\{xref\_number}, \[51], 54--55, 136, 276.
\:\\{xref\_predefined}, \[99], 101, 129, 132.
\:\\{xref\_ptr}, \[51], 53--55, 309.
\:\\{xref\_reserved}, \[99], 101, 129, 132, 226, 266.
\:\\{xref\_roman}, \[99], 101, 109, 115, 128, 132, 226, 266.
\:\\{xref\_slanted}, \[99], 101.
\:\\{xref\_switch}, \[51], 54, 109, 115--116, 128--129, 132--133.
\:\\{xref\_typewriter}, \[99], 101, 128.
\:\\{xref\_wildcard}, \[99], 101, 128.
\:\p{YEAR}, \[72].
\:\p{YEAR\_NUMBER}, \[72].
\:\.{You can't do that...}, 266, 274.
\:\.{You need an = sign...}, 273.
\:\\{yr\_output\_file}, 308\*.
\fin
\:\X231:Append \(a \(string scrap\X
\U section~225.
\:\X234:Append \(a \TeX\ string scrap\X
\U section~225.
\:\X233:Append \(an ASCII-token scrap\X
\U section~225.
\:\X235:Append \(an identifier scrap\X
\U section~225.
\:\X228:Append \(\\{digit\_scrap} scraps until a `\#' occurs\X
\U section~227(2).
\:\X232:Append the character appropriate to \\{buffer}\?(\?\|j\?)\?\X
\U sections~231 and~233.
\:\X225:Append the scrap appropriate to \\{next\_control}\X
\U section~223\*.
\:\X177:Cases for \\{accept\_scrap}\X
\U section~176\*.
\:\X178:Cases for \\{and\_scrap}\X
\U section~176\*.
\:\X179:Cases for \\{arrow}\X
\U section~176\*.
\:\X180:Cases for \\{base}\X
\U section~176\*.
\:\X181:Cases for \\{begin\_label\_scrap}\X
\U section~176\*.
\:\X182:Cases for \\{beginning}\X
\U section~176\*.
\:\X183:Cases for \\{case\_scrap}\X
\U section~176\*.
\:\X184:Cases for \\{clause}\X
\U section~176\*.
\:\X185:Cases for \\{colon}\X
\U section~176\*.
\:\X186:Cases for \\{decl}\X
\U section~176\*.
\:\X187:Cases for \\{digit\_scrap}\X
\U section~176\*.
\:\X188:Cases for \\{else\_scrap}\X
\U section~176\*.
\:\X189:Cases for \\{ending}\X
\U section~176\*.
\:\X190:Cases for \\{exception\_scrap}\X
\U section~176\*.
\:\X191:Cases for \\{exp}\X
\U section~176\*.
\:\X192:Cases for \\{generic\_scrap}\X
\U section~176\*.
\:\X193:Cases for \\{gproc}\X
\U section~176\*.
\:\X194:Cases for \\{if\_scrap}\X
\U section~174.
\:\X195:Cases for \\{in\_scrap}\X
\U section~174.
\:\X196:Cases for \\{label}\X
\U section~174.
\:\X197:Cases for \\{loop\_scrap}\X
\U section~174.
\:\X198:Cases for \\{math}\X
\U section~174.
\:\X199:Cases for \\{mod\_scrap}\X
\U section~174.
\:\X200:Cases for \\{new\_scrap}\X
\U section~174.
\:\X201:Cases for \\{not\_scrap}\X
\U section~174.
\:\X202:Cases for \\{open\_scrap}\X
\U section~174.
\:\X203:Cases for \\{period}\X
\U section~174.
\:\X204:Cases for \\{private\_scrap}\X
\U section~174.
\:\X205:Cases for \\{proc}\X
\U section~174.
\:\X206:Cases for \\{return\_scrap}\X
\U section~174.
\:\X207:Cases for \\{separate\_scrap}\X
\U section~174.
\:\X208:Cases for \\{stmt}\X
\U section~174.
\:\X209:Cases for \\{when\_scrap}\X
\U section~174.
\:\X210:Cases for \\{with\_scrap}\X
\U section~174.
\:\X230:Cases involving nonstandard ASCII characters\X
\U section~226.
\:\X213:Change \\{pp} to $\max(\\{scrap\_base},\\{pp}\?+\?\|d)$\X
\U sections~212 and~214.
\:\X120:Check for overlong name\X
\U section~118.
\:\X98:Check that all changes have been read\X
\U section~308\*.
\:\X273:Check that \?=\? or \?\S\? follows this module name, and emit the
scraps to start the module definition\X
\U section~272.
\:\X165:Clear \|b and \?\&{return}\?\80\X
\U sections~163 and~164.
\:\X220:Combine the irreducible scraps that remain\X
\U section~219\*.
\:\X64:Compare name \|p with current identifier, \&{goto}\8\\{found} if equal\X
\U section~63.
\:\X112:Compress two-symbol combinations like `\.{:=}'\X
\U section~111.
\:\X62:Compute the hash code \|h\X
\U section~61.
\:\X63:Compute the name location \|p\X
\U section~61.
\:\X261:Copy a control code into the buffer\X
\U section~260.
\:\X164:Copy special things when \|c\?=\?\H{@}\?,\?\H{\\}\?,\?\H{\{}\?,\?\H{%
\}}; \&{return} at end\X
\U section~163.
\:\X260:Copy the \ADA\ text into \\{buffer}\?(\?\\{limit}\?+\?1\?\to\?\|j\?)\?\X
\U section~258.
\:\X161:Copy up to `\v' or control code, \?\&{return}\?\8\|c if finished\X
\U section~160.
\:\X158:Copy up to control code, \&{return} if finished\X
\U section~157.
\:\X252:Count the \\{pdollar} and space tokens; output the appropriate control
sequences and \&{goto}\8\\{reswitch}\X
\U section~251.
\:\X175:Declaration of subprocedure for \\{translate}\X
\U section~218.
\:\X236:Declaration of the \\{app\_comment} procedure\X
\U section~222.
\:\X226:Declaration of the \\{easy\_cases} procedure\X
\U section~222.
\:\X108:Do special things when \|c\?=\?\H{@}\?,\?\H{\\}\?,\?\H{\{}\?,\?\H{\}}; %
\&{return} at end\X
\U section~107.
\:\X286:Do the first pass of sorting\X
\U section~282.
\:\X274:Emit the scrap for a module name if present\X
\U section~272.
\:\X77:Enter a new module name into the tree\X
\U section~76.
\:\X65:Enter a new name into the table at position \|p\X
\U section~61.
\:\X227:Get a based number\X
\U section~226.
\:\X123:Get a character token\X
\U section~111.
\:\X114:Get a string\X
\U section~111.
\:\X113:Get an identifier\X
\U section~111.
\:\X115:Get control code and possible module name\X
\U section~111.
\:\X119:If end of name, \&{exit}\X
\U section~118.
\:\X91:If the current line starts with \.{@y}, report any discrepancies and %
\&{return}\X
\U section~90.
\:\X297:Invert the cross-reference list at \\{cur\_name}, making \\{cur\_xref}
the head\X
\U section~296.
\:\X256:Look ahead for strongest line break, \&{goto}\8\\{reswitch}\X
\U section~255.
\:\X229:Make sure that there is room for at least four more scraps, six more
tokens, and four more texts\X
\U section~225.
\:\X241:Make sure that there is room for at least seven more tokens, three more
texts, and one more scrap\X
\U section~240.
\:\X216:Make sure the entries \\{cat}\?(\?\\{pp}\?\to\?\\{pp}\?+\?4\?)\? are
defined\X
\U section~215.
\:\X174:Match a production at \\{pp}, or increase \\{pp} if there is no match\X
\U section~215.
\:\X89:Move \\{buffer} and \\{limit} to \\{change\_buffer} and \\{change%
\_limit}\X
\U sections~86 and~90.
\:\X255:Output a \(control, look ahead in case of line breaks, possibly %
\&{goto}\8\\{reswitch}\X
\U section~251.
\:\X254:Output a \.{\\math} operator\X
\U section~251.
\:\X257:Output a module name\X
\U section~251.
\:\X300:Output all the module names\X
\U section~282.
\:\X280:Output all the module numbers on the reference list \\{cur\_xref}\X
\U section~279.
\:\X253:Output an identifier\X
\U section~251.
\:\X294:Output index entries for the list at \\{sort\_ptr}\X
\U section~292.
\:\X265:Output the code for the beginning of a new module\X
\U section~264.
\:\X281:Output the code for the end of a module\X
\U section~264.
\:\X296:Output the cross-references at \\{cur\_name}\X
\U section~294.
\:\X153:Output the digits of \|a\X
\U section~152(2).
\:\X295:Output the name at \\{cur\_name}\X
\U section~294.
\:\X258:Output the text of the module name\X
\U section~257.
\:\X126:Phase I: Read all the user's text and store the cross references\X
\U section~302.
\:\X262:Phase II: Read all the text again and translate it to \TeX\ form\X
\U section~304.
\:\X282:Phase III: Output the cross-reference index\X
\U section~306.
\:\X37:Print error location based on input buffer\X
\U section~36.
\:\X138:Print error messages about unused or undefined module names\X
\U section~126.
\:\X309:Print statistics about memory usage\X
\U section~308\*.
\:\X310:Print the job \\{history}\X
\U section~308\*.
\:\X149:Print warning message, break the line, \&{return}\X
\U section~148.
\:\X49, 54, 55:Procedures in \\{data\_structures}\X
\U section~3\*.
\:\X129, 130, 137, 302:Procedures in \\{first\_phase}\X
\U section~7\*.
\:\X28, 30, 33, 36:Procedures in \\{input\_output}\X
\U section~4\*.
\:\X85, 86, 90, 93, 95, 101, 103, 105, 107, 111:Procedures in \\{lexical%
\_scanning}\X
\U section~6\*.
\:\X141, 143, 148, 152, 155, 157, 160, 163:Procedures in \\{output}\X
\U section~8\*.
\:\X212, 214, 218, 222, 238, 240:Procedures in \\{parsing}\X
\U section~9\*.
\:\X61, 76, 80:Procedures in \\{searching}\X
\U section~5\*.
\:\X245, 248, 250, 251, 268, 279, 304:Procedures in \\{second\_phase}\X
\U section~10\*.
\:\X291, 299, 306:Procedures in \\{third\_phase}\X
\U section~11\*.
\:\X134:Process a format definition\X
\U section~133.
\:\X118:Put module name into \\{mod\_text}\?(\?1\?\to\?\|k\?)\?\X
\U section~116.
\:\X96:Read from \\{aweb\_file} and maybe turn on \\{changing}\X
\U section~95.
\:\X97:Read from \\{change\_file} and maybe turn off \\{changing}\X
\U section~95.
\:\X277:Rearrange the list pointed to by \\{cur\_xref}\X
\U section~275.
\:\X215:Reduce the scraps using the productions until no more rules apply\X
\U section~219\*.
\:\X122:Scan a verbatim or output string\X
\U section~115.
\:\X124:Scan an ASCII token\X
\U section~115.
\:\X116:Scan the module name and make \\{cur\_module} point to it\X
\U section~115.
\:\X121:Scan to the next \.{@>}\X
\U section~115.
\:\X19, 21, 22, 44, 47, 52, 58, 117, 144, 146, 289:Set initial values\X
\U section~3\*.
\:\X78:Set variable \|c to the result of comparing the given name to name \|p\X
\U sections~76 and~80.
\:\X275:Show cross references to this module\X
\U section~264.
\:\X259:Skip next character, give error if not `\.{@}'\X
\U section~258.
\:\X87:Skip over comment lines in the change file; \&{return} if end of file\X
\U section~86.
\:\X88:Skip to the next nonblank line; \&{return} if end of file\X
\U section~86.
\:\X292:Sort and output the index\X
\U section~282.
\:\X48, 53:Specification of procedures visible from \\{data\_structures}\X
\U section~3\*.
\:\X128, 301:Specification of procedures visible from \\{first\_phase}\X
\U section~7\*.
\:\X27, 32, 35:Specification of procedures visible from \\{input\_output}\X
\U section~4\*.
\:\X92, 94, 100, 102, 104, 106, 110:Specification of procedures visible from %
\\{lexical\_scanning}\X
\U section~6\*.
\:\X140, 142, 147, 151, 154, 156, 159, 162:Specification of procedures visible
from \\{output}\X
\U section~8\*.
\:\X217, 221, 237, 239:Specification of procedures visible from \\{parsing}\X
\U section~9\*.
\:\X60, 75, 79:Specification of procedures visible from \\{searching}\X
\U section~5\*.
\:\X244, 249, 278, 303:Specification of procedures visible from \\{second%
\_phase}\X
\U section~10\*.
\:\X298, 305:Specification of procedures visible from \\{third\_phase}\X
\U section~11\*.
\:\X293:Split the list at \\{sort\_ptr} into further lists\X
\U section~292.
\:\X270:Start a format definition\X
\U section~267.
\:\X269:Start a macro definition\X
\U section~267.
\:\X69, 70, 71, 72:Store all the predefined identifiers\X
\U section~308\*.
\:\X67, 68:Store all the reserved words\X
\U section~308\*.
\:\X127:Store cross reference data for the current module\X
\U section~126.
\:\X135:Store cross references in the \(\ADA\ part of a module\X
\U section~127.
\:\X133:Store cross references in the \(definition part of a module\X
\U section~127.
\:\X132:Store cross references in the \TeX\ part of a module\X
\U section~127.
\:\X284:Tell about changed modules\X
\U section~282.
\:\X272:Translate the \(\ADA\ part of the current module\X
\U section~264.
\:\X264:Translate the \(current module\X
\U section~262.
\:\X267:Translate the \(definition part of the current module\X
\U section~264.
\:\X266:Translate the \TeX\ part of the current module\X
\U section~264.
\:\X14, 15, 16, 17, 18, 20, 26, 29, 31, 34, 38, 41, 42, 43, 45, 46, 50, 51, 56,
57, 59, 66, 73, 74, 82, 84, 99, 109, 125, 131, 136, 139, 150, 167, 168, 171,
172, 243, 247, 263, 271, 276, 283, 285, 287, 288, 290, 307:Types and objects
visible from \\{data\_structures}\X
\U section~3\*.
\con
